\documentclass[10pt,a4paper,twoside,bg=print,twocolumn,openany,nodeprecatedcode]{dndbook}
\usepackage[utf8]{inputenc}
\usepackage{tabulary}
\usepackage[normalem]{ulem}
\usepackage{color,soul}
\usepackage{float}
\usepackage{caption}
\usepackage{wrapfig}
\usepackage{tocloft}
\usepackage[toc]{multitoc}
\usepackage[hidelinks]{hyperref}
\raggedbottom
\extrafloats{2000}
\cftsetindents{section}{0em}{0em}
\cftsetindents{subsection}{0em}{0em}
\cftsetindents{subsubsection}{0.2em}{0em}
\setcounter{tocdepth}{1}
\renewcommand*{\multicolumntoc}{3}
\setlist[description]{nolistsep,listparindent=3pt,topsep=6pt}
\date{}
\begin{document}
\frontmatter
\tableofcontents
\mainmatter
\chapter{Foreword}
\DndQuote%
{}
{''We should call this book \textit{Oops! All Lairs.}''}
{Matthew Colville}
We ended up making two books. Before \textit{Flee, Mortals! The MCDM Monster Book} even launched on Kickstarter, we knew that hitting our lairs stretch goal might mean another book's worth of content. When \textit{Flee, Mortals!} reached the layout stage, we realized that our fateful prediction was correct. We could either cut fifty pages of content from the lairs and other parts of \textit{Flee, Mortals!}, or we could break out the lairs into their own book, add more content, and bring our dream book of dungeons and boss battles to life. Since you're reading this, you know which route we took.

The truth is that these lairs deserve to be in their own book. Each little slice of hell contains story hooks, background information, detailed hazards and NPCs, a villain, unique rewards, and all the stat blocks you need to run them. You can use this book without ever referencing \textit{Flee, Mortals!}

The real stars of this book are the villains. Each is abominable in their own way; every lair boss is a unique monster, from their personality to their stat block. When a group of heroes invades a villain's home, that evildoer fights like hell to end the threat. We've given each villain the beefy, playtested stat block they need to stand up to even the most optimized parties.

Each lair was specially crafted for the villain and their allies. A lair is an extension and example of its villain's power and personality. \textbf{Lord Syuul} performs his wicked experiments in a laboratory made of flesh. \textbf{Kiona the Dread Lord} obsessively researches the \textit{Codex Mortis} in her extradimensional library. The once-proud \textbf{Ithu'rath} plots their grand return to the top in a crumbling castle.

The boss battles and lairs in this book have saved my game nights. If I don't have anything prepared or if the players zig when I expect them to zag, I can easily grab a story hook or tempt them with some unique treasure and get them into one of these places where evil (and fun!) lives. When I'm scheming for the perfect capstone encounter to end a killer adventure, I reach for this book.

The authors packed so much delicious evil into these bosses and their homes. I am incredibly grateful to them as well as all the artists, consultants, editors, playtesters, and the rest of the MCDM team who worked super hard to make more pages of content in one year than we ever have before. We, uh, learned a lot—a whole friggin' lot.

We hope you enjoy reading and running these lairs as much as we enjoyed making them. Go, brave adventurers. Only you have the heart to face evil where it lives.

Ex animo

James Intracaso

MCDM Lead Game Designer

\addcontentsline{toc}{chapter}{Introduction}
\markboth{INTRODUCTION}{INTRODUCTION}
\chapter*{Introduction}
Twenty-two lairs that you can drop right into your campaign have arrived. Each is a home to a unique villain and is filled with lackeys, hazards, traps, and a treasure hoard with unique rewards. Every lair has been optimized to present a fun challenge for five player characters of a specific level, indicated in the lair's description and on the Lairs by Level table.

The lairs are presented in ascending order of their challenge.

\section{Maps}
Each lair has a map that indicates all hidden information, like traps and secret doors. Digital versions of these maps (along with player versions that don't include the secrets) are included with your purchase of this book.

\subsection{Cardinal Directions}
Most lair maps include a compass rose, even if the location is on a plane other than the Mundane World. Most planes of the timescape (see \textit{"The Timescape"}) are circular planets as well as planes, so they have a north, south, east, and west. But for the rare planes without cardinal directions, the compass rose is still included to make interpreting the map easier.

\section{Area Descriptions}
Each lair area begins with bullet points that overview the room's most important features, giving you all the information you need to run that area upfront. Some of these are meant to be shared with players, while others are for the GM's eyes only, as follows:


\begin{itemize}
\item "If you see a description in quotation marks like this, read or paraphrase it to the players when their characters first enter the area—provided they can sense their surroundings properly. This text describes what the characters immediately notice."
\item (\textit{If the characters did X in area Y...}) "If a quotation starts with a parenthetical condition like this, only read or paraphrase it to players if those conditions apply."
\item \textbf{(Secret)} If a bullet starts with this parenthetical subheading, don't read it aloud! This description contains information for the GM, and it shouldn't be revealed to players until their characters investigate the area further. Hidden creatures, traps, and more could be presented this way.
\end{itemize}

\begin{DndTable}[header=Lairs by Level]{cXXXc}

Optimized Level & Lair Name & Lair Boss & Boss Ancestry & Page\\
2nd & \textit{Jagged Edge Hideaway} & \textbf{Queen Bargnot} & Goblin & 16\\
3rd & \textit{Hanging Tree} & \textbf{Dohma Raskovar} & Orc & 29\\
4th & \textit{Shtriga Nonna's Hut} & \textbf{Shtriga Nonna} & Hag & 45\\
5th & \textit{Burnock Mill} & \textbf{Baron Uthrak} & Human & 55\\
6th & \textit{Camp Firefield} & \textbf{Bloodlord Varrox} & Hobgoblin & 68\\
7th & \textit{Shifting Library} & \textbf{Kiona the Dread Lord} & Incorporeal Undead & 82\\
8th & \textit{Cloud Fang Keep} & \textbf{Lady Emer} & Medusa & 94\\
9th & \textit{Molten Enclave} & \textbf{Zenith Aastrika} & Fire Giant & 110\\
10th & \textit{Ruins of Wrathrock} & \textbf{Ithu'rath} & Olothec & 121\\
10th & \textit{The White Tower} & \textbf{Xorannox} & Overmind & 137\\
11th & \textit{Canyon of the Tower Crown} & Aurumvas of Meaningless Greed & Demon & 152\\
11th & \textit{Tomb of the Keeper} & \textbf{Ashyra} & Mummy & 166\\
12th & \textit{Terminal Excrescence} & \textbf{Lord Syuul} & Voiceless Talker & 176\\
13th & \textit{Eighth City Advocacy Services} & \textbf{Chancellor Lazivos} & Devil & 193\\
13th & \textit{Shadowkeep} & \textbf{Count Rhodar von Glauer} & Vampire & 208\\
14th & \textit{Durixaviinox's Rest} & \textbf{Durixaviinox} & Dragon & 224\\
15th & \textit{Coronal Hollow} & \textbf{Qazyldrath} & Dragon & 236\\
16th & \textit{Eyes of the Mountain} & \textbf{Yserthrax} & Dragon & 251\\
17th & \textit{Glass Cavern} & \textbf{Xaantikorijek} & Dragon & 263\\
18th & \textit{Mount Brazen} & \textbf{Forzaantilirys} & Dragon & 274\\
19th & \textit{Ash Queen's Reliquary} & \textbf{Atæshia} & Elemental & 287\\
20th & \textit{Boughs of Eternity} & \textbf{High Mage Vairae} & Lich & 302\\
\end{DndTable}

\section{Treasure Hoards}
The bulk of a lair's treasure is in a hoard listed at the end of the lair. Each treasure hoard includes a unique magic or psionic item and another unique reward, such as a recipe for a new poison.

Treasures not included in the hoard are sprinkled throughout the lair. These appear in the "Treasure" subheading of the area where they are found.

\section{Stat Blocks}
Unless otherwise noted, when a creature appears in a lair's area description, their stat block is found at the end of the lair's entry. For instance, the stat block of \textbf{Queen Bargnot}, the ruler of \textit{Jagged Edge Hideaway}, appears at the end of the Jagged Edge Hideaway entry along with the stat blocks of all the other creatures in the lair. The stat blocks appear in alphabetical order.

\section{The Timescape}
The places, people, and other proper names in the lore of this book come from the timescape, the MCDM multiverse. Most of the creatures detailed in this book are found in Orden, a world of high fantasy and just one of the timescape's many planes (also called manifolds). Orden is also called the Mundane World. If you like what you're reading, you can find more lore about the timescape on the MCDM Patreon at mcdm.gg/patreon.

You can ignore any creature's lore if you wish to use their stat block in another way in your campaign.

\section{New Rules and Styles}
The creatures in this book generally follow the core rules, but we've made a few tweaks. These new rules and presentation styles are designed to make combat encounters easier to run, more fun, and more memorable.

\subsection{Allies and Enemies}
The effects in this book sometimes target just allies or enemies. A creature is your \textbf{ally} if they're inclined to help you or fight alongside you—or if you believe they're inclined to do so. However, you are not your own ally, so if you target your allies with an effect, it doesn't affect you. A creature is your \textbf{enemy} if they're inclined to oppose you or fight against you—or if you either believe they're inclined to do so, or you plan to pick a fight with them regardless. If in doubt, it's up to the GM to decide whether a creature counts as an ally, an enemy, or neither.

\subsection{Challenge Rating}
Each creature's stat block lists their challenge rating in the top right corner. The more obvious placement of the challenge rating makes it easier to find stat blocks and build encounters to challenge your players.

\subsection{Conditions}
Some of the creatures in this book inflict—or are immune to—the following new conditions.

\subsubsection{Dazed}
A dazed creature can only do one of the following things on their turn: move, use an action, or use a bonus action. If a creature becomes dazed during their turn, their turn ends. The cure ailment power, \textit{lesser restoration} spell, and \textit{greater restoration} spell remove the dazed condition. At the GM's discretion, other powers, spells, or effects might also remove the dazed condition.

When a dazed creature is affected by a spell or effect that gives them an extra action on their turn (like the \textit{haste} spell or the fighter's Action Surge feature), they can still take this extra action, in addition to the movement, action, or bonus action allowed by the dazed condition.

Some creatures in this book have immunity to the dazed condition. At the GM's discretion, a creature published in the core rules or another supplement who has immunity to the paralyzed or stunned condition also has immunity to the dazed condition.

\subsubsection{Flanked}
If your game uses the optional flanking rules in the core rules, a creature who is immune to the Flanked condition can't be flanked regardless of the position of their enemies.

\subsection{Creature Roles}
Each creature has a role listed next to their challenge rating. Roles are descriptive and most don't follow special rules—they simply help you build encounters and use the creature effectively in combat.

\subsubsection{Ambusher}
Ambushers are creatures who hide well—not just before an encounter, but during it. They utilize surprise and stealth to gain the upper hand.

\subsubsection{Artillery}
Artillery creatures fight best from afar. Whether they wield arrows or magical rays, these creatures always try to keep a distance from their foes.

\subsubsection{Brute}
Brutes are hardy creatures who have lots of hit points and deal lots of damage. They might not be the most disciplined warriors, but they make up for it with sheer toughness and aggression.

\subsubsection{Companion}
Companion creatures are meant to accompany player characters on adventures. For more information, see the "Companion Creatures" section.

\subsubsection{Controller}
Controllers debuff, move, and obstruct their enemies. They often have crowd-control actions that apply a debilitating effect or target multiple creatures at once.

\subsubsection{Leader}
A leader is an action-oriented creature who fights alongside underlings. For more information, see the "Action-Oriented Creatures" section.

\subsubsection{Minion}
Minions are weak creatures who find strength in numbers. For more information, see the "Minions" section.

\subsubsection{Retainer}
Retainers are sapient beings meant to accompany player characters on adventures. For more information, see the "Retainers" section.

\subsubsection{Skirmisher}
Skirmishers are mobile warriors who use hit-and-run tactics in combat. Their traits allow them to make the most of their position.

\subsubsection{Soldier}
Soldiers are well-armored creatures who draw the attacks of their foes, freeing allies to move around the battlefield. These trained warriors typically have higher attack bonuses and AC.

\subsubsection{Solo}
A solo creature is an action-oriented creature who can take on the player characters on their own. For more information, see the "Action-Oriented Creatures" section.

\subsubsection{Support}
Support creatures aid their allies, providing buffs, healing, movement, or action options.

\subsection{Mundane and Supernatural}
This book uses the term \textbf{mundane} to refer to attacks, items, and other effects that aren't magical or psionic.

On the other hand, the term \textbf{supernatural} describes an effect or item that is either magical or psionic. For instance, a creature's stat block might say they're resistant to "bludgeoning, piercing, and slashing damage from mundane attacks." This means that when you deal them bludgeoning, piercing, or slashing damage using a spell, power, magic weapon, or psionic weapon, the damage is supernatural and thus the creature isn't resistant to it.

\subsection{Pronouns}
The core rules often use the singular pronoun "it" to refer to a single creature. This book instead uses the singular "they" to clearly delineate between creatures and objects. Objects still use "it" as a singular pronoun.

\subsection{Psionic Powers}
Some creatures have psionic powers, creating effects with sheer force of will instead of with magic. When a creature manifests a power, they create an effect that has similar rules to casting spells, with the following exceptions.

\subsubsection{Not Magic}
Powers don't create magical effects, so they're unaffected by features and spells like \textit{antimagic field}, \textit{counterspell}, and \textit{dispel magic}. The damage from a power attack isn't magical. If a rule affects only spells or magical effects, it doesn't affect powers. However, psionic powers are supernatural, so if a rule affects supernatural effects, it does affect powers (see \textit{"Mundane and Supernatural"}).

\subsubsection{Power Orders}
Every power has an \textbf{order} that is an expression of its level of strength. Powers of the 1st order are the weakest (equivalent to cantrips in strength), and 6th-order powers are the strongest.

If a stat block feature should be considered a power, its order is noted in parentheses.

\subsubsection{No Components}
Powers don't have material, somatic, or verbal components.

\subsubsection{Manifesting Multiple Powers}
Unlike player characters, if a creature's stat block allows them to manifest a power as a bonus action, they can still manifest another power of 2nd order or higher as an action on that turn. For example, a creature who manifests the jaunt power as a bonus action can also make a 2nd-order power attack as an action. Similarly, many Multiattack actions let a creature make multiple power attacks on a turn, even though each attack is a power of its own.

\subsubsection{Power Attacks as Opportunity Attacks}
When an enemy's movement provokes an opportunity attack from a creature whose stat block includes a melee power attack, the creature can use their reaction to make a melee power attack against the creature, rather than making an opportunity attack.

\subsubsection{New Powers}
A character playing a talent—the class found in the MCDM supplement The Talent and Psionics—can use the rules in that book to learn new powers from the creatures in \textit{Where Evil Lives} Powers that aren't in \textit{The Talent and Psionics} are indicated in a creature's stat block with an asterisk (*) and provided in the "New Psionic Powers" chapter of this book.

Converter's Note: I left these asterisks out because they were ugly and, in the contexts in which this will be used, unnecessary.

\subsection{Save Ends Effects}
Some creatures have a trait or action that, on a failed saving throw, imposes a debilitating effect for an extended duration. But sometimes a lucky foe can retry their saving throw and potentially end the effect early. In such situations, the stat block specifies "save ends at start of turn" or "save ends at end of turn."

When you see this phrase, it means any creature affected by a \textbf{save ends effect} like this can repeat the saving throw on each of their turns, ending the effect on themself early on a success. They can only make that saving throw either at the start or end of their turn, as specified in the stat block.

Here is an example of an attack with a save ends effect from the \textbf{goblin cursespitter} stat block:

\textbf{Toxic Touch (Cantrip).} \textsl{Melee or Ranged Spell Attack:} 4 to hit, reach 5 ft. or range 30 ft., one target. \textsl{Hit:} 7 (2d6) poison damage, and the target must succeed on a DC 12 Constitution saving throw or be poisoned for 1 minute (save ends at end of turn).

\begin{DndSidebar}[float=!b]{Behind the Design: Save Ends Effects}
Save ends effects are one way of saving space in stat blocks and making them less complex for GMs to run at the table. Much like how stat blocks just reference the poisoned condition instead of spelling its rules out each time, we didn't want to fill up the pages by repeating long sentences similar to this one: "A target poisoned in this way can repeat the saving throw at the end of each of their turns, ending the effect on themself on a success."



\end{DndSidebar}

\subsection{Spells}
Some creatures have magical actions, bonus actions, and reactions. When these features should be considered spells (for the purpose of \textit{counterspell}, \textit{dispel magic}, and similar effects), the spell level is noted in parentheses. If the spell requires concentration, this is also noted in parentheses. Unless otherwise noted, these spells have somatic and verbal components.

\subsubsection{Casting Multiple Spells}
Unlike player characters, if a creature's stat block allows them to cast a spell as a bonus action, they can still cast another spell of 1st level or higher as an action on that turn. For example, a creature who casts the \textit{misty step} spell as a bonus action can also make a 2nd-level spell attack as an action. Similarly, many Multiattack actions let a creature make multiple spell attacks on a turn, even though each attack is a spell of its own.

\subsubsection{Spell Attacks as Opportunity Attacks}
When an enemy's movement provokes an opportunity attack from a creature whose stat block includes a melee spell attack, the creature can use their reaction to make a melee spell attack against the creature, rather than making an opportunity attack.

\subsubsection{Utility Spells}
In addition to combat-focused spells, some creatures can cast spells that are primarily used outside of combat. To streamline the stat blocks, these spells aren't written out in the creature's combat-oriented "Actions" section. Instead, their names are listed in a separate "Utility Spells" section at the end of the stat block. After each spell name, the casting time is indicated in superscript, as shown in the Casting Times table.

Note that at least for the moment, there is no "Utility Spells" section in these blocks on this site. The casting times superscripts are included, however.

\subsubsection{Casting Times}

\begin{DndTable}{cXX}

Superscript & Casting Time & \\
A & 1 action\\
B & 1 bonus action\\
R & 1 reaction\\
+ & Longer than 1 action (see spell description)\\
\end{DndTable}

For example, a spellcaster's utility spells might appear as follows:

1/day each: expeditious retreat\textsuperscript{B}, mage armor\textsuperscript{A}, phantom steed\textsuperscript{+}

\subsection{Section Structure}
In this book, creature sections begin with broad descriptions for all the creatures in that section. While not all the information is common in-character knowledge, you can share these early pages with your players if you don't mind them having a little meta-knowledge.

The creature stat blocks aren't presented until the end of each section, allowing you to share flavorful information with your players and flip fewer pages when you run a band of similar creatures.

\subsection{Unusual Needs}
Stat blocks in the core rules often indicate when a creature has an unusual nature and doesn't require air, food, sleep, or water to live and function. To keep the stat blocks in this book succinct, we've omitted those traits from individual stat blocks. Instead, you can refer to the following list to determine each creature type's unusual needs (or lack thereof):


\begin{description}
\item \textbf{Construct.} Constructs don't require air, food, drink, or sleep.
\item \textbf{Elemental.} Elementals don't require air, food, drink, or sleep.
\item \textbf{Ooze.} Oozes don't require air or sleep.
\item \textbf{Undead.} Undead don't require air, food, drink, or sleep.
\item \textbf{Other Creature Types.} Unless otherwise specified, other creature types require air, food, drink, and sleep.
\end{description}

\section{Action-Oriented Creatures}
The solo and leader creatures presented in this book are designed to be bosses: enemies who can take on an entire party by themselves or with a handful of underlings. Rather than simply increasing the challenge rating (an approach that often leads to underwhelming encounters), this book introduces \textbf{action-oriented creatures.}

A powerful villain needs plenty of opportunities to act and move when it's not their turn. Thus, each actionoriented creature has at least one special bonus action and reaction, as well as a special section with villain actions that let them dominate the battlefield.

These actions make the boss creatures dynamic and formidable. Whether fought as an exciting solo challenge or alongside a few easy-to-run underlings, action-oriented creatures challenge the characters with dramatic and powerful actions in combat.

\subsection{Villain Actions}
Every action-oriented creature has three \textbf{villain actions} they can use after an enemy's turn. Villain actions are similar to legendary actions with the following exceptions:


\begin{itemize}
\item A creature can use only one villain action per round (as such, villain actions tend to be more powerful than legendary actions).
\item Each villain action can only be used once during a combat encounter.
\end{itemize}

Like legendary actions, a creature can't use villain actions if incapacitated or otherwise unable to take actions.

\subsubsection{Choosing Villain Actions}
Each trio of villain actions has a recommended round order. These abilities give the battle a logical flow and a cinematic arc:


\begin{itemize}
\item The first villain action is an opener, which shows the characters they're not battling a typical creature. Openers generally deal some damage, summon a lackey or three, buff the boss, debuff the characters, or move the creature into an advantageous position. They're just a taste of what's to come.
\item The second villain action provides crowd control. It typically fires after the heroes have had a chance to respond once or twice, get into position, and surround the villain. This second action helps the villain regain the upper hand. Like an opener, this action comes in many flavors, but it's even more powerful than an opener.
\item The third and final villain action is an ultimate move or "ult"—a showstopper the villain can use to deal a devastating blow to the characters before the end of the battle.
\end{itemize}

While every creature has a recommended order of actions, you can take villain actions in any order if it makes your fight more dramatic. You could push back a villain action if a creature is stunned or might stay alive for more than three rounds, or you could perform the recommended third action in round two after several surprise critical hits!

\subsection{Unique Creatures}
Many of the action-oriented creatures in this book are unique. They typically have a specific name and backstory, though you can ignore these if you wish to use the creature's stat block in another way in your campaign. For instance, \textbf{Queen Bargnot}'s stat block could be used for any action-oriented goblin in your game.

\section{Companion Creatures}
A \textbf{companion} is a wild ally who adventures with characters. Each companion has unique traits and actions that make them a great ally. But beware! These creatures can be difficult to control in the heat of battle and just might bite the hand that feeds. Don't fret too much, though! A companion gelatinous cube or owlbear is worth the risk of an occasional wild rampage.

\subsection{Caregiver}
Every companion has a player character \textbf{caregiver} who commands the creature. The caregiver's player controls the companion most of the time during the game, though the GM can step in to take control if the companion and caregiver are separated, or if the caregiver mistreats the companion in some way.

In combat, a companion shares a turn with their caregiver and acts during the caregiver's turn. A companion can move and use their own reaction and bonus action independently. But they can take only the Dash, Disengage, or Dodge actions unless their caregiver uses a bonus action to command the companion to take a different action, including any of the actions noted in the companion's stat block. A companion must be able to see or hear their caregiver to receive a command. A companion can also take other actions if their caregiver is incapacitated or if the companion enters a rampage (as discussed below).

\begin{DndSidebar}[float=!b]{One Companion Per Group}
Companions are a lot of fun, but having more than one companion to manage can slow things down at the table. A companion is another member of the party, with statistics and actions to track, and additional companions can easily make combat slow to a crawl. If every player wants to get in on the companion action, it's simpler for the characters to take turns as a creature's caregiver, rather than running around adventuring with a menagerie.



For characters wishing to share a companion, the companion accepts a new caregiver at the end of a short or long rest. Because a companion's proficiency bonus and hit points depend on their caregiver's level, those statistics might fluctuate if the characters in a party are different levels, reflecting that a more experienced caregiver is more adept at directing a companion.



\end{DndSidebar}

\subsubsection{Charmed Caregiver}
A caregiver who is charmed can still command their companion, but they can't command the companion to attack a creature who charmed them.

\subsubsection{Incapacitated or Absent Caregiver}
If a companion's caregiver is incapacitated or dies, the GM determines who controls the companion—typically, the caregiver's player can maintain control of their companion as they take a heroic stand against the enemy.

However, in cases where a caregiver and their companion are physically separated, such as when a companion is captured, the GM might wish to take control of the companion to keep the player in suspense about the creature's fate.

\subsubsection{New Caregiver}
At the GM's discretion, a companion can abandon a caregiver character and choose a different willing creature as a caregiver.

\subsection{Ferocity}
Companions are dangerous creatures. Though often more docile than their wild counterparts, they aren't fully domesticated. Each companion's \textbf{ferocity} is a measure of their tenacity and fury, and of how those things build in battle. As a companion's ferocity increases, they gain access to powerful new actions, but they also become more difficult for a caregiver to control.

If a companion isn't incapacitated at the start of their and their caregiver's turn, their ferocity increases by 1d4 + the number of enemies the companion can see or hear within 5 feet of them. For the purpose of increasing ferocity, a group of creatures who share a single stat block (such as a swarm of rats) counts as one creature. Ferocity increases round after round during combat, and there is no maximum to the level of ferocity a companion can gain.

\begin{DndSidebar}[float=!b]{Bag of Rats}
We can already tell that some crafty players are scheming to stuff a bag full of slightly groggy rats, then open that bag up in front of their companion to build up their ferocity during a fight. However, a too-easy target doesn't rile up a companion the way being threatened by an enemy combatant in a battle for survival does. As such, the GM makes the final decision as to what constitutes an enemy for the purpose of increasing a companion's ferocity. Likewise, in the same way a swarm of creatures is counted as a single creature for the purpose of increasing ferocity, the GM is free to determine that two or three weak creatures might count as only one creature for that purpose.



\end{DndSidebar}

\subsubsection{Rampage}
After rolling to increase ferocity at the start of their turn, if a companion has 10 ferocity or more and isn't incapacitated, they run the risk of entering a \textbf{rampage}. The companion's caregiver can make a Wisdom (Animal Handling) check (no action required) to try to stop the companion from entering a rampage. To make the check, the caregiver must not be incapacitated, and the companion must be able to see or hear the caregiver. The DC for the check equals 5 + the companion's ferocity. On a success, the companion acts normally on their turn. On a failure, or if the caregiver doesn't make the check, the companion enters a rampage.

When a companion enters a rampage, they immediately move up to their speed toward the nearest creature they can sense and attack that creature with their signature attack (see below), dealing extra damage equal to half their ferocity if the attack hits. If at least one ally and one enemy are nearest and equidistant to the companion, the caregiver's player rolls any die. On an odd number, the companion attacks an ally. On an even number, the companion attacks an enemy. The caregiver's player determines which specific equidistant ally or enemy the companion engages (and can choose their own character if they wish).

A companion who can't reach a creature to attack while in a rampage uses the Dash action to move as far as they can toward the nearest creature they can sense. If a companion can't sense any potential targets, they move as far as they can in a random direction determined by the GM, avoiding danger.

When a companion who has entered a rampage resolves their action or ends their turn, their ferocity drops to 0 and they're no longer in a rampage.

\subsubsection{Reducing Ferocity}
To prevent a companion from entering a dangerous rampage, a caregiver has several options at their disposal for reducing the creature's ferocity.

\subparagraph{Ferocity Actions}
Each companion has three actions in their stat block that cost ferocity to use. To use one of these ferocity actions, the caregiver's character level must be equal to or greater than the ferocity action's level, and the companion must spend the necessary amount of ferocity before they use the action. If the companion doesn't have enough ferocity to spend, they can't use the action.

Ferocity actions always use the companion's action, meaning they can't be used as part of an opportunity attack. Ferocity actions can't be used while a companion is in a rampage.

\subparagraph{End of Combat}
When a combat encounter involving a companion ends and the companion isn't dying, the companion regains hit points equal to their ferocity, and their ferocity drops to 0. The GM determines when a combat encounter ends, typically at the point when creatures stop acting in initiative order.

\subsection{Dying Companions}
When a companion is reduced to 0 hit points, they are dying and make death saving throws just as characters do. Thus, characters always have a chance to save their furry (or scaly, or slimy, or exoskeletony) friends' lives! If combat ends while a companion is dying, their ferocity drops to 0 but they don't regain hit points (see \textit{"End of Combat"} above).

\subsection{Statistics}
In addition to their ferocity actions, a companion's statistics vary from the statistics of their wild counterparts. This makes a companion easier to run, keeps their power in line with other companions, and helps ensure companions never outshine the characters. As somewhat more social versions of wild creatures, companions are often cleverer and more versatile than their untamed counterparts.

\subsubsection{Hit Dice}
Most creatures without character classes have their Hit Die type determined by their size (d4 for Tiny creatures, d6 for Small creatures, and so forth). However, companions are special and use a d8 for Hit Dice regardless of size.

\subsubsection{Language}
A companion shares a unique bond with their caregiver and can understand basic commands in one language chosen by the caregiver. However, the companion can't read, speak, or write any language, even if similar creatures normally can.

\subsubsection{Proficiency Bonus}
Because a companion's effectiveness and survivability depend on the training and expertise of their caregiver, a companion's proficiency bonus is equal to their caregiver's proficiency bonus. Additionally, some of a companion's statistics refer to their proficiency bonus, abbreviated as PB. Other statistics use a number of dice equal to a companion's proficiency bonus; these are expressed with PB in place of the number of dice. For example, if a companion has a +2 proficiency bonus, PBd6 means 2d6.

\subsubsection{Signature Attack}
Each companion has an action designated as their signature attack. A signature attack is always a melee attack, and it's typically the creature's best natural attack. A companion uses their signature attack when they enter a rampage.

Each companion also has special actions that they can use only by spending ferocity during their turn, with some of those actions making use of the companion's signature attack. See \textit{"Reducing Ferocity"} above for more information on ferocity actions.

\subsection{Companion Mounts}
Many companions are large enough to ride, especially by caregivers who are Small, when outfitted with an exotic saddle similar to those worn by aquatic or flying mounts. When a caregiver rides a companion into combat, not much actually changes. The caregiver and companion still each act on the same turn, and the caregiver must use their bonus action to direct the companion to take any action other than the Dash, Disengage, or Dodge actions. Under some circumstances, a companion might allow themself to be ridden by a creature other than their caregiver, though that other creature can't give the companion commands.

If a companion rampages while bearing a rider, that rider counts as being within 5 feet of the companion when determining which creature the companion attacks. If a companion attacks their rider, they have disadvantage on the attack roll.

\begin{DndSidebar}[float=!b]{Companion Barding}
At the GM's discretion, characters can purchase barding for a companion, as discussed in the core rules, with the following adjustments:




\begin{itemize}
\item Barding for Small companions weighs half as much as the equivalent armor for Humanoids. Barding for Medium companions weighs the same as Humanoid armor, while barding for Large companions weighs four times as much.
\item Companions are proficient in any barding they wear.
\item When a companion wears barding, they don't add their caregiver's proficiency bonus to their AC.
\item Companions who are shapechangers (such as the mimic companion) can't use that ability while wearing barding.
\item Companions who have corrosive bodies or are amorphous (such as the gelatinous cube companion) can't wear nonmagical barding.
\end{itemize}



\end{DndSidebar}

\subsection{Companion Encounter Balance}
Unless a caregiver has the beastheart class (found in Beastheart and Monstrous Companions), the GM should consider a companion as akin to a powerful combat-focused magic item when building encounters. A companion gives a party a significant power boost, not just by dishing out more damage and providing more hit points for enemies to target, but also by creating complications that can make a fight more challenging for the party's foes. GMs can adjust encounter difficulty by one step (from easy to medium, hard to deadly, and so forth) to properly challenge a group of characters with a companion, particularly if the characters' average level is 7th or lower.

\subsection{NPCs with Companions}
Companion creatures are designed to accompany player characters on adventures. If an NPC has a pet or creature servant, it's recommended you use the creature's normal stat block. For instance, if an NPC ranger has an owlbear they work with, that creature would use the normal \textbf{owlbear} stat block and be played by the GM as normal. But if the characters are charged with rescuing a fallen druid NPC's owlbear pet from the clutches of a villain, the GM can choose to instead make the pet an \textbf{owlbear companion}, allowing the creature to join their rescuers on their adventures.

\section{Demons}
Demons consume mortal souls to fuel their fiendish powers. Use the following rules whenever you run one of the demons in this book.

\subsection{Demons and Souls}
Demons feast not on food or water, but on souls. These fuel their bloodthirsty powers, and while starved for souls, a demon can scarcely think.

\subsubsection{Soul Count}
A demon's stat block states the number of souls a given demon has already consumed at the beginning of combat. This number is presented right under their hit points in a similar fashion: both as a die expression and as an average number.

\subsubsection{Soul Devourer}
Demons can gain more souls by slaying other creatures, as described in their Soul Devourer trait. Unlike hit points, there is no maximum limit to a demon's soul count.

\subsubsection{Soul-Fueled Power}
Most demons have one or more features that require souls to use.

\subparagraph{Passive Traits}
Passive traits automatically take effect after a demon's soul count crosses the threshold described in that trait. The trait remains active until the demon's soul count drops below the threshold, and the trait reactivates when the demon's soul count equals or exceeds it.

\subparagraph{Active Features}
During an encounter, a demon can burn souls to use or enhance certain abilities. When they do, their soul count decreases by the number indicated. This cost is sometimes noted in parentheses at the beginning of an ability, such as "(Costs 1 Soul)." Other times, the text of the ability itself describes it, such as saying "the demon can burn 1 soul" for an additional effect.

\subsubsection{Lethe}
When a demon's soul count drops to 0, they fall into a state known as lethe—a violent hunger wherein they can only lash out in a desperate search for sustenance. Demons who have fallen into lethe become single-minded and violent, seeking only to consume. See a demon's Lethe trait for more information.

\subsubsection{Soulsight}
Demons have a special sense called soulsight, allowing them to perceive each creature, other than Constructs and Undead, within a certain radius. Soulsight doesn't rely on sight but counts as "seeing" for features that require sight. Neither physical objects (including total cover) nor supernatural effects can impede soulsight unless otherwise stated. A demon's stat block lists the radius of their soulsight under "Senses."

\begin{DndSidebar}[float=!b]{Adjusting Lethality}
Soul devouring makes demons deadly opponents. To give characters a better chance of living through their encounter with a demon, you can reduce the DC of the saving throw—or you can remove the saving throw entirely, and only allow a demon to devour a soul when a creature dies within 60 feet of them through other means. On the other hand, you might like to scale up the deadliness with more powerful demons! If so, instead of always using a DC 11 saving throw, you can set the DC equal to 8 + the demon's proficiency bonus.



\end{DndSidebar}

\DndQuote%
{}
{''I sometimes sit all day among my treasures. I have piles of coins, one-of-a-kind statuettes, and magic items of great power. Do you know what I think when I look upon my vast hoard? I think, "This is nothing compared to the taste of a mortal soul."''}
{Aurumvas of Meaningless Greed}
\section{Retainers}
\textbf{Retainers} are sapient beings who adventure alongside the player characters. Each retainer is a less experienced adventurer who a player character can take under their wing. Retainers are never meant to achieve the same power level as the player characters.

Rules for retainers first appeared in \textit{Strongholds \& Followers}. This book contains updated rules for these followers.

\subsection{Mentor}
Every retainer has a player character \textbf{mentor}. A retainer's mentor gives them orders, and the mentor's player also controls the retainer. A retainer acts on the same initiative count as their mentor in combat, acting immediately before or after the mentor (player's choice). As a mentor gains experience and levels up, so does their retainer (see \textit{"Statistics"} below).

\subsection{Statistics}
Retainers are designed to be easy to run so their player (who is already managing a complex character) doesn't get overwhelmed with even more details. Each retainer has a simple stat block—though these are similar to other creature stat blocks, retainers follow a few different rules.

\subsubsection{Level}
A retainer's level equals their mentor's level. As a retainer levels up, their hit points increase (see \textit{"Hit Points and Hit Dice"}) and they gain combat features (see \textit{"Features"}). Additionally, a few of their statistics increase when their mentor's proficiency bonus does, including attack bonuses, skills, and save DC (see \textit{"Proficiency Bonus"}). Their other statistics typically remain the same regardless of their level.

\subsubsection{Armor Class}
Unlike player characters, a retainer's armor class isn't determined by a particular set of armor they don (like studded leather). Instead, each retainer's armor class is simplified to one of three armor types:


\begin{itemize}
\item Retainers with light armor have AC 13.
\item Retainers with medium armor have AC 15.
\item Retainers with heavy armor have AC 18.
\end{itemize}

\subsubsection{Hit Points and Hit Dice}
Retainers gain one Hit Die per level. Their hit point maximum is based on the size of their Hit Die, as shown on the Retainer Hit Points table.

Retainers have exceptionally high hit points—sometimes even higher than their mentor. This allows the party to enjoy the presence of a beloved retainer without worrying about losing them to an unlucky roll of the dice or decision by a single player. The GM often awards characters with retainers in lieu of other treasure. Such rewards shouldn't be lost easily.

\begin{DndTable}[header=Retainer Hit Points]{cXX}

Hit Die Size & Hit Point Maximum & \\
d6 & 6 times their level\\
d8 & 7 times their level\\
d10 & 8 times their level\\
d12 & 9 times their level\\
\end{DndTable}

\subsubsection{Proficiency Bonus}
Because a retainer's effectiveness and survivability depend on the training and expertise of their mentor, a retainer's proficiency bonus is equal to their mentor's proficiency bonus.

Some of a retainer's statistics refer to their proficiency bonus, abbreviated as PB. Other statistics use a number of dice equal to a retainer's proficiency bonus; these are expressed with PB in place of the number of dice. For example, if a retainer has a +3 proficiency bonus, PBd10 means 3d10.

A retainer adds their proficiency bonus to any saving throw they make.

\subsubsection{Features}
Each retainer has a \textbf{signature attack} they can make using the Attack action each round. Retainers can use their signature attack to make opportunity attacks, even if their signature attack isn't a melee weapon attack.

A retainer gains new features at 3rd, 5th, and 7th level. Some features can only be used a certain number of times per day (as noted in the retainer's stat block).

Additionally, if a retainer's signature attack is a weapon attack, they typically gain an extra attack at 7th level (as noted in their stat block), allowing them to make two signature attacks per round instead of one.

\subsubsection{Gear}
When a retainer joins the party, they typically carry clothes appropriate to their position, a suit of armor, a weapon, and an explorer's pack. If the retainer casts spells, they also carry a spellcasting implement. Any additional equipment must be provided to them by their mentor.

\textit{\textbf{Magic Items.}} Retainers can use magic items like anyone else. For example, a \textit{+1 weapon} increases their attack and damage rolls by 1, and \textit{+1 armor} of the appropriate type (light, medium, or heavy) increases their AC by 1.

One happy side effect of having retainers is that as a player character levels up and acquires better equipment, they can pass their obsolete items on to the retainer.

\subsection{Dying Retainers}
When a retainer is reduced to 0 hit points, they follow the same rules as player characters. If not killed instantly by massive damage, they fall unconscious, make death saving throws, and can be stabilized or healed.

\subsection{Optional Rules: Shared Attacks}
If a player wants to speed up their retainer's turn, they can make \textbf{shared attacks} at the GM's discretion.

When the mentor hits a creature with an attack, if their retainer can see or hear them, the retainer is inspired by their mentor's success. On the inspired retainer's next turn, for each successful attack their mentor just made, the retainer can automatically hit a creature of their choice within range with their signature attack, up to the number of attacks the retainer can make.

If the mentor hit with fewer attacks than the numberof attacks the retainer can make, the retainer can make additional attack rolls for those attacks to see if they hit.

\subsection{Retainer Encounter Balance}
When building encounters, the GM should consider a retainer as akin to a powerful combat-focused magic item. A retainer gives a party a significant power boost; they not only dish out more damage and provide more hit points for enemies to target, but they also create complications to challenge the party's foes. GMs can increase encounter difficulty by one step (from easy to medium, hard to deadly, and so forth) to properly challenge a group of characters with a retainer.

\section{Minions}
A \textbf{minion} is a weak foe, designed to allow GMs to create dramatic combat encounters with hordes of enemies without overwhelming the characters. In fact, an encounter with minions makes characters feel heroic, since they can take on a myriad of foes and live to tell the tale.

However, minions still make threatening foes. Killing a minion still requires penetrating their defenses, and characters can't just shrug off damage from minion attacks.

So how do minions make running a horde of enemies quick and easy for the GM?


\begin{itemize}
\item Minions are simple to run. Their stat blocks are small and uncomplicated.
\item Minions act quickly. They don't multiattack, roll for damage, or take unique bonus actions or reactions, so their turns aren't long.
\item Minions die fast. A character can kill several minions with a single weapon attack!
\item Minions have strength in numbers. Their attacks can be grouped together to make them deadlier and faster to use at the table.
\end{itemize}

\subsection{No Hit Dice}
Minions have hit points but no Hit Dice, simplifying their design. Minions can't spend Hit Dice to heal during a short rest because they have none.

\subsection{No Damage Rolls}
Minions don't roll for damage because their attacks deal a static amount of damage. They also can't score critical hits.

\subsection{Shared Turns}
Typically, all minions of the same stat block act on the same turn. Since they share a turn, the minions can each move into position then each use an action if they wish, instead of each moving and taking an action individually.

\subsection{Minion Trait}
Every minion has the Minion trait, which affects the creature in the following ways:


\begin{itemize}
\item If the minion takes any damage from an attack or as the result of a failed saving throw, their hit points are reduced to 0.
\item If the minion takes damage from another effect, they die if the damage equals or exceeds their hit point maximum; otherwise they take no damage.
\end{itemize}

\subsection{Overkill Attacks}
Powerful weapon attacks can kill more than one minion in a single maneuver called an \textbf{overkill attack.}

As already discussed, a weapon attack requires only 1 point of damage to reduce a minion to 0 hit points, regardless of their hit point maximum. However, when a weapon attack's damage does exceed the target minion's hit point maximum, the attack becomes an overkill attack and the damage dealt beyond the minion's hit point maximum becomes \textbf{overkill damage.}

Overkill damage can be applied to a second minion who has the same stat block as the target and is in overkill range (see below). Damage against the second minion is counted as if you made a weapon attack against them; since it only takes 1 point of weapon damage to reduce a minion to 0 hit points, any amount of overkill damage immediately knocks them out. But wait, it gets better—if the initial attack's overkill damage exceeds the second minion's hit point maximum, the leftover overkill damage can roll over to a third minion, and so on! In other words, for each time the overkill damage exceeds the new target's hit point maximum, the attacker can choose an additional minion to reduce to 0 hit points.

For example, when a weapon attack deals 18 damage to a minion with a hit point maximum of 5, the overkill damage is 13. If there are three additional minions of the same stat block in overkill range, they can all three be immediately reduced to 0 hit points, since the overkill damage exceeded the target's hit point maximum more than twice over.

Overkill attacks can't be made as part of an opportunity attack.

\subsection{Overkill Range}
Minions must be within a certain range to qualify for an overkill attack, determined by whether the attack is a melee or ranged attack. In addition to the examples below, the "Overkill Damage Illustrated" sidebar demonstrates how to calculate overkill damage.

\subsubsection{Melee Overkill Attacks}
When a creature hits a minion with a melee weapon attack, other minions within reach of the attack are in \textbf{overkill range} and can be chosen as additional targets for an overkill attack. The overkill attack can't target minions outside the weapon attack's reach.

\textit{Lady Ulnock the paladin battles a horde of goblin minions (each with 6 hit points). She hits a goblin minion with her longsword and uses Divine Smite, dealing 8 slashing damage and 11 radiant damage to the target for a total of 19 damage. Since Lady Ulnock dealt 13 points of overkill damage—more than the hit point maximum of two additional minions—she can choose up to three additional goblin minions within 5 feet of her (the reach of her longsword attack) and reduce them to 0 hit points. If there are no other goblin minions within 5 feet of Lady Ulnock, she can't damage additional minions with this attack.}

\subsubsection{Ranged Overkill Attacks}
When a creature hits a minion with a ranged weapon attack, other minions in a line originating from the creature in the direction of the target, to a distance equal to the weapon's short range, are in \textbf{overkill range} and can be chosen as additional targets for an overkill attack. The overkill attack can't target minions outside the line or beyond the weapon's short range.

\textit{Perigold Quickfingers the rogue is hidden and takes aim at a group of zombie minions (each with 6 hit points) with his light crossbow (which has a short range of 80 feet). He hits a zombie minion with his crossbow, dealing extra damage thanks to his Sneak Attack, for a total of 14 damage. Since Perigold dealt 8 points of overkill damage—more than the hit point maximum of one additional minion—he can choose up to two additional zombie minions in an 80-foot-long line extending from Perigold in the direction of the target, reducing them to 0 hit points. If there are no other zombie minions in the line, then Perigold can't damage other minions with this attack.}

\begin{DndSidebar}[float=!b]{Behind the Design: Minion Trait}
You might ask, why not just give minions 1 hit point and take no damage when they save for half, like in fourth edition? First, spells that use a creature's hit points to determine effectiveness—like \textit{color spray} and \textit{sleep}—would devastate all minions, even those meant to challenge high-level characters. These spells are still effective againstminions, just not devastating!



Second, spells and effects that deal damage without any attack roll or save—like \textit{magic missile} and \textit{spike growth}—would lay waste to minions with 1 hit point. This fits the fiction for minions with low challenge ratings, like goblins and zombies. But the balance of combat and fiction breaks down to near-silliness at higher levels when the same spells easily take down powerful devil minions.



Finally, high-level spells with a save for half damage—like \textit{fireball} or \textit{meteor swarm}—would feel wasted against minions with 1 hit point. Why use a higher-level spell when a lower-level one will do? Similarly, the fourth edition design could lead to a kobold minion illogically surviving a fireball spell while a "stronger" standard kobold next to them dies, despite both creatures succeeding on their saving throw. By contrast, under this book's minion rules, spellcasters still have a good reason to use high-level spells against minions.



\end{DndSidebar}

\DndQuote%
{}
{''Outnumbered? Switch to area-of-effect powers. Otherwise you're gonna have a bad day, which will probably last the rest of your life.''}
{The Sun, Senior Pyrokinetic, The Society}
\begin{DndSidebar}[float=!b]{Overkill Damage Illustrated}
Imagine a fighter is making a melee attack with a shortsword against a lackey, and two more lackeys are within overkill range of that attack. Each lackey is a minion with 6 hit points. In example A, the attack hits and deals 4 piercing damage. This kills only the targeted lackey. A successful attack against a minion always kills the target.



In example B, the attack is a critical hit and deals 13 piercing damage. This kills the targeted lackey with 7 overkill damage remaining, which spreads to the other minions within the fighter's overkill range. So a second lackey within 5 feet of the fighter takes 6 piercing damage and dies, leaving 1 overkill damage remaining. The third lackey takes that last point of overkill damage and dies.



\end{DndSidebar}

\subsection{Group Attacks}
Each minion has at least one group attack action that speeds up play. In a group attack, two to five minions of the same stat block who share a turn can all use their action to join the attack, provided the target is within the original attack's reach or range for each minion.


\begin{description}
\item Make a single attack roll for the group attack. It counts as one attack.
\item A group attack roll gains a +1 bonus to the attack roll for each minion who joins the attack. (For example, if four goblin minions make a group attack together, the attack roll has a +4 bonus.)
\item If the group attack hits, multiply the damage by the number of minions who joined that group attack. (For example, if four goblin minions hit with a group attack that deals 1 damage, their group attack deals 4 damage.)
\end{description}

The GM decides how many minions join a group attack. For instance, if five minions surround a target, the GM may decide to have all five attack at once to speed up combat, or may break up the attacks among smaller groups to increase the odds that some minions hit while others miss. A single minion can even use their group attack action on their own—they make the attack as a normal creature would, and simply don't benefit from the group bonuses described above.

\subsubsection{Advantage and Disadvantage}
A group attack is only made with advantage or disadvantage if all the minions joining the group attack have advantage or disadvantage on the attack roll. Otherwise, the attack is made without advantage or disadvantage.

\subsubsection{Cover and Concealment}
If a target has cover or concealment from some but not all minions, the GM should divide the minions into multiple groups based on the type of cover or concealment they have, then make a separate attack for each group.

\subsubsection{Target Response Effects}
If a group attack triggers a reaction or similar effect that would normally affect a single attacker, such as the \textit{fire shield} or \textit{hellish rebuke} spell, the target of the group attack picks one minion who joined the attack to be affected by the effect.

\subsubsection{Group Opportunity Attacks}
If a creature provokes an opportunity attack from more than one minion of the same stat block at a time and those minions have a melee group attack action, the minions can each use their reaction to join a group attack as an opportunity attack.

\subsection{Optional Rule: Group Saving Throws}
Though minions often make saving throws individually, there are times when rolling individual saving throws for each minion could slow down the fight, like when a cleric surrounded by eighteen shade minions uses Turn Undead.

When many minions with the same stat block need to make a saving throw against the same effect at the same time, you can make one saving throw for a group of up to five minions at a time. All minions in a group use the result of the saving throw.

For instance, if thirty-four goblin minions need to make a saving throw against a hypnotic pattern spell, the minions would make a total of seven saving throws against the spell: six for thirty minions divided into six groups of five, and one more save for the remaining group of four minions.

\subsection{Optional Rule: Tough Minions}
Minions of a higher challenge rating, such as abyssal ghouls or fire giants, make for powerful foes. Consequently, it could break the game's verisimilitude for an NPC commoner to kill such a minion with a single attack. To keep minions believable, you can use the following rule.

When a minion's challenge rating is at least 6 higher than an NPC's challenge rating, that NPC's actions and traits affect a minion as if they didn't have the Minion trait, reducing the minion's hit points like a normal creature instead of automatically dropping them to 0.

This rule shouldn't be applied to player characters—they're heroes who can always kill a minion in one hit.

\subsection{Special Traits}
Many minions have traits that give them strength in numbers but become less powerful as their allies are defeated. For example, an enemy who starts their turn within 5 feet of three or more goblin lackeys must succeed on a saving throw or take damage from the lackeys' Tiny Stabs trait.

\subsection{Challenge Ratings}
Minions have a challenge rating just like any other creature. However, their experience point value depends on their challenge rating, as shown on the Minion Encounter Building table. Aside from their damage output, the Minion trait, and their experience point value, a minion's statistics are on par with a standard creature of the same challenge rating.

When the rules reference a challenge rating, such as the \textit{polymorph} spell or the cleric's Destroy Undead feature, use the minion's listed challenge rating as normal. For instance, a 5th-level cleric can affect Undead creatures with a challenge rating of 1/2 or lower with their Destroy Undead feature—so this feature can affect rotting zombie minions (CR 1/4) but not shade minions (CR 1).

\begin{DndSidebar}[float=!b]{Minion Bands}
A combat encounter with more than five minions per character can become deadly if all minions act on the same turn—the characters cut down minion hordes on their turns but then suffer massive damage as a sea of minions pour down on them. While many groups like this challenge, you have the option to divide the minions into different bands that act on different initiative counts. If you do this, minions can only join group attacks with minions in the same band. At the start of a new round, two or more bands of minions can reorganize into one band, acting on the lowest initiative count of the bands that combined.



If you divide minions into bands, be sure to distinguish which minions belong together. For gridded combat, you could use a colored marker for each miniature. In a theater of the mind encounter, try group descriptors like "zombie dwarves" and "zombie elves."



\end{DndSidebar}

\chapter{Jagged Edge Hideaway}
Jagged Edge Hideaway

\section{Jagged Edge Hideaway Stat Blocks}
The following stat blocks appear in the lair.

\begin{DndSidebar}[float=!b]{Goblin Tactics}
Goblins benefit from fighting in environments with features to climb on and hide behind, like trees and pillars. The Crafty trait means melee-focused goblins can run into combat, attack, then regroup with their allies. It also allows goblin minions to run past enemy warriors and surround their spellcasting foes to make use of their Tiny Stabs trait. Goblins who fight at range climb for better sightlines and defense.



If clearly losing a battle, goblins typically don't stick around and fight to the last warrior. Instead, they flee to untamed wilderness and tight tunnels, utilizing their natural agility to run and hide from threats. Goblins who escape with their lives take time to lick their wounds then gather more allies and rush back to face the threat with fresh, overwhelming numbers.



\end{DndSidebar}

\textbf{Goblin Assassin}

\DndQuote%
{}
{''Strike first, strike hard, grab what you can, and then disappear into the shadows. If you can't find a shadow, don't worry. Morky will make one for you.''}
{Queen Bargnot}
\textbf{Goblin Cursespitter}

\textbf{Goblin Lackey}

\textbf{Goblin Sneak}

\textbf{Goblin Sniper}

\DndQuote%
{}
{''Look, you saved me—that's enough to get my knives for a day. But if you help me stop Bargnot, we're friends for life.''}
{Toblobb}
\textbf{Goblin Spinecleaver}

\textbf{Goblin Underboss}

\textbf{Goblin Warrior}

\textbf{Swarm of Skitterlings}

\begin{DndSidebar}[float=!b]{Skitterling}
A six-legged, winged rodent the size of a housecat, a skitterling moves their clawed feet as they fly, appearing to scurry through the air. Goblins train these pets to claw at the faces of enemies, as their feet secrete a toxin that causes temporary blindness.



\end{DndSidebar}

\DndQuote%
{}
{''Skitterlings make excellent little guardians, but they leave quite the mess. Make sure they're fed. If not, you'll become the next mess they make.''}
{Queen Bargnot}
\textbf{Queen Bargnot}

\textbf{Swarm of Spiders}

\textbf{War Spider}

\chapter{Hanging Tree}
Hanging Tree

\section{Hanging Tree Stat Blocks}
The following stat blocks appear in the lair.

\textbf{Dohma Raskovar}

\textbf{Human Brawler}

\textbf{Human Guard}

\textbf{Mimic}

\DndQuote%
{}
{''\textit{Just} the doorknob, right? Not the \textit{door}, just that one doorknob. Little shit almost chewed my hand off before I even knew what was happening. But what are you gonna do? You can't just . . . stop opening doors. Just need to expect the unexpected. Not the most useful lesson, I realize.''}
{The Sun, Senior Pyrokinetic, The Society}
\textbf{Mohler}

\textbf{Mohler Companion}

\begin{DndSidebar}[float=!b]{Mystic Connection: Mohler}
If you're playing a beastheart and have a mohler companion, you gain the following benefit at 9th level when you gain the beastheart's Mystic Connection feature:



\paragraph{Optional Feature}
Mystic Connection: Mohler|FleeMortals



\end{DndSidebar}

\textbf{Orc Blitzer}

\textbf{Orc Bloodrunner}

\textbf{Orc Conduit}

\textbf{Orc Fury}

\DndQuote%
{}
{''The orcs are out there, and we're in here. I checked the Chronicle for advice on sieges. "Be on the outside," it said. Wonderful, thanks.''}
{Dancer, Chronicler of the Chain of Acheron, Heroes 216–231}
\textbf{Orc Garroter}

\textbf{Orc Godcaller}

\textbf{Orc Rampart}

\DndQuote%
{}
{''Everyone is welcome at the Hanging Tree. It's only when I find someone poking around where they don't belong that things turn violent. Care to explain why you're down here before we draw steel? I'm happy to listen after instead, but it'll be so much harder to understand you when your lips are bleeding.''}
{Dohma Raskovar}
\textbf{Orc Terranova}

\textbf{Scyza}

\chapter{Shtriga Nonna's Hut}
Shtriga Nonna's Hut

\section{Hut Stat Blocks}
The following stat blocks appear in the lair.

\textbf{Chimera}

\textbf{Human Mercenary}

\DndQuote%
{}
{''Honestly, it is a lot of work. Their hide is tough, they have so little marbling, and there's not a lot of meat at all once you get all those scales off. But if you prepare enough of them then make a stew and let it simmer for hours, you'll kick yourself for having not tried kobold sooner, my loves.''}
{Shtriga Nonna}
\textbf{Kobold Decanus}

\textbf{Shtriga Nonna}

\textbf{Troll}

\textbf{Troll Whelp}

\DndQuote%
{}
{''Trolls do make the most delightful lackeys. My beauties require so little maintenance. Their wounds get better without any care from me, and the trollies work for food! Oh, they're not fussy when it comes to payment—as long as it's meat, dearies. If you're really clever, you can carve off a piece of one troll to feed another and save the rarer meats for yourself.''}
{Shtriga Nonna}
\chapter{Burnock Mill}
Burnock Mill

\section{Mill Stat Blocks}
The following stat blocks appear in the lair.

\textbf{Baron Uthrak}

\textbf{Human Apprentice Mage}

\begin{DndSidebar}[float=!b]{Changing the Apprentice Mage}
Human wizards often follow arcane traditions to master a particular element of magic. In this stat block, the apprentice mage makes Lightning Strike and Thunder Crack attacks, but the GM can reflect a different elemental mastery by changing both damage types to acid, cold, or fire damage. A mage can only join other mages in this group attack if their group attack deals the same damage type.



\end{DndSidebar}

\textbf{Human Brawler}

\textbf{Human Death Acolyte}

\textbf{Human Death Cultist}

\DndQuote%
{}
{''Laws? You only need to follow the laws if you're weaker than the people enforcing them. Oh, I've made mistakes in my time, but I've never followed any law some rich brat who inherited a crown made up to control people. Nobles don't even follow the rules they write for others. It's a joke.''}
{''What you don't understand is that out here, I am the nobles. I am the law, and if the bodies at the bottom of the river could talk, they'd tell you I can enforce it.''}
{Baron Wigar Uthrak}
\textbf{Human Guard}

\textbf{Human Knave}

\textbf{Human Mercenary}

\textbf{Human Raider}

\textbf{Human Scoundrel}

\textbf{Human Storm Wizard}

\begin{DndSidebar}[float=!b]{Human Elementalist Wizards}
Human wizards often follow arcane traditions to master a particular element of magic. As GM, you can introduce new elementalist wizards by adjusting the \textbf{human storm wizard} stat block, as in the following examples.



\subsubsection{Human Cryomancer}
The wizard's Arcane Staff action and Arcane Shield reaction deal cold damage instead of lightning and thunder damage. Additionally, instead of casting the \textit{lightning bolt} spell, the wizard casts \textit{ice storm}.



\subsubsection{Human Pyromancer}
The wizard's Arcane Staff action and Arcane Shield reaction deal fire damage instead of lightning and thunder damage. Additionally, instead of casting the gust of \textit{wind spell}, the wizard casts \textit{flaming sphere}, and instead of \textit{lightning bolt}, the wizard casts \textit{fireball}.



\end{DndSidebar}

\textbf{Human Trickshot}

\DndQuote%
{}
{''Formally known as Terrans because of their unique relationship to the Mundane World, humans have no direct connection to their creator-god.''}
{Remainer}
\textbf{Sea Hag Apprentice}

\begin{DndSidebar}[float=!b]{Content Warning: Drowning}
Before using a sea hag, check in with your group to see if anyone is upset by descriptions of drowning or suffocation. If so, you can keep the game fun by describing the hag's Dry Drown action as a dizzying hex that affects creatures who live on land, instead of describing it as drowning.



\end{DndSidebar}

\chapter{Camp Firefield}
Camp Firefield

\section{Camp Firefield Stat Blocks}
\begin{DndSidebar}[float=!b]{Slaughter Demon}
When evil hobgoblins who embrace their fiendish heritage need to wipe an enemy off the map, their war mages ritualistically beseech an archdevil for the service of a grack'tanar, known as a slaughter demon in the Common tongue. Once summoned, this towering, serpent-bodied, six-clawed demon slithers to war alongside the hobgoblins who summoned them.



Devils captured the grack'tanars eons ago. Broken, these demons wait for a call to war, hungry and frothing in the Seven Cities of Hell. Their archdevil captors reward loyal hobgoblins by allowing the mortals to hold a grack'tanar's reins for a time. These slaughter demons are eager to kill and please their captors so they might be sent out again, and they rarely turn on hobgoblins unless they fall into lethe.



\end{DndSidebar}

The following stat blocks appear in the lair.

\textbf{Bloodlord Varrox}

\begin{DndSidebar}[float=!b]{Bloodlord Varrox Rides Again!}
If the characters managed to dispatch Bloodlord Varrox, the hobgoblin could easily return with the help of his infernal patrons to take revenge on the heroes who killed him. When Bloodlord Varrox returns from death, he can be more powerful than his previous incarnation, choosing from or inspired by one of the following options:




\begin{itemize}
\item Varrox returns with 168 hit points and a Strength of 24 (+7). Because of his new Strength score, he makes Strength (Athletics) checks with a +10 bonus and makes Soulsword and Death Skulls attacks with a +10 bonus to attack rolls and a +7 bonus to damage rolls.
\item Varrox returns as a \textbf{devil adjudicator}.
\item Varrox returns as a \textbf{vampire}.
\end{itemize}



\end{DndSidebar}

\textbf{Devil Notary}

\textbf{Grilp}

\textbf{Hobgoblin Brandbearer}

\textbf{Hobgoblin Death Captain}

\textbf{Hobgoblin Firerunner}

\textbf{Hobgoblin Incendiarist}

\textbf{Hobgoblin Recruit}

\DndQuote%
{}
{''Right, hobgoblins. These guys are bad. Like, really really bad. The Society takes on dragons, demons, things from other worlds, but these guys are dangerous the way we're dangerous. We're lucky they spend most of their time fighting each other. If they ever get organized, this world is in trouble.''}
{The Sun, Senior Pyrokinetic, The Society}
\textbf{Hobgoblin Smokebinder}

\textbf{Hobgoblin Tactician}

\textbf{Hobgoblin Trooper}

\textbf{Hobgoblin War Mage}

\textbf{Slaughter Demon}

\chapter{Shifting Library}
Shifting Library

\section{Library Stat Blocks}
\textbf{Blood-Borne Ooze}

\textbf{Hellhound}

\textbf{Kiona the Dread Lord}

\textbf{Lightthief}

\begin{DndSidebar}[float=!b]{Lightthief}
A lightthief is a fiend who comes to the Mundane World to feast on soul energy. They stalk graveyards and tombs, where they know vulnerable mourners come in the dark to pay their respects. These bat-winged creatures resemble flayed humanoid hands with a bloodshot eye leering from the palm—though their victims have little chance to notice them before darkness falls.



\end{DndSidebar}

\textbf{Shade}

\begin{DndSidebar}[float=!b]{Making Incorporeal Undead}
Though most incorporeal undead were once Humanoid, you can easily make spectral undead from any other creature type. First, choose the stat block you want to use for the undead: shade, shadow, specter, or wraith. Then adjust the stat block's size and Intelligence score to match the original creature's statistics. That's all you need to do to make an incorporeal undead version of a creature.



\end{DndSidebar}

\textbf{Shadow}

\textbf{Specter}

\DndQuote%
{}
{''I dunno why it's always wraiths. Dealt with a lot of deathless, but wraiths are the worst. They're always pissed about something.''}
{Dancer, Chronicler of the Chain of Acheron, Heroes 216–231}
\textbf{Wraith}

\chapter{Cloud Fang Keep}
Cloud Fang Keep

\section{Keep Stat Blocks}
The following stat blocks appear in the lair.

\textbf{Basilisk}

\textbf{Blood-Borne Ooze}

\textbf{Crux of Frost}

\textbf{Gem Jelly}

\textbf{Griffon}

\textbf{Griffon Companion}

\begin{DndSidebar}[float=!b]{Mystic Connection: Griffon}
If you're playing a beastheart and have a griffon companion, you gain the following benefit at 9th level when you gain the beastheart's Mystic Connection feature:



\paragraph{Optional Feature}
Mystic Connection: Griffon|FleeMortals



\end{DndSidebar}

\textbf{Hobgoblin Tactician}

\textbf{Human Mercenary}

\textbf{Kobold Decanus}

\textbf{Lady Emer}

\textbf{Lizardfolk Hunter}

\textbf{Minotaur Devastator}

\begin{DndSidebar}[float=!b]{Medium Minotaur}
The minotaur devestator originally appeared in \textit{Flee, Mortals! The MCDM Monster Book}. In that book, the retainer is Large. However, Worjo, the minotaur devestator in Lady Emer's lair, is Medium. The stat block here reflects Worjo's size.



\end{DndSidebar}

\textbf{Orc Blacksmith}

\textbf{Owlbear}

\DndQuote%
{}
{''Owlbears make spectacular guardians, and they're always a good topic of conversation for guests. Everyone wants to know how I trained them, what their names are, and if they do tricks. Well, they don't, and they smell atrocious. You probably would too if your diet consisted mainly of monster hunters and thieves.''}
{Lady Emer}
\textbf{Wyvern}

\DndQuote%
{}
{''I always feel better, seeing a wyvern. Makes me feel like the wilderness still has power, strength. Like the forest still has a chance.''}
{Lady Demelza, Courser}
\chapter{Molten Enclave}
Molten Enclave

\section{Enclave Stat Blocks}
\textbf{Basalt Stone Giant}

\textbf{Fire Giant Lightbearer}

\textbf{Fire Giant Red Fist}

\textbf{Fire Giant Trooper}

\textbf{Hellhound}

\textbf{Zenith Aastrika}

\DndQuote%
{}
{''I didn't find enlightenment atop some peaceful mountain high in the clouds or while sitting with my eyes closed in a serene field of high grass. I earned it in the hottest fires the earth produces. Would you like to feel that enlightenment for yourself?''}
{Zenith Aastrika}
\chapter{Ruins of Wrathrock}
Wrathrock

\section{Wrathrock Stat Blocks}
The following stat blocks appear in the lair.

\textbf{Ghost Sycophant}

\textbf{Haunt}

\textbf{Human Death Acolyte}

\textbf{Human Death Cultist}

\textbf{Human Knave}

\textbf{Human Storm Wizard}

\textbf{Human Trickshot}

\textbf{Ithu'rath}

\subsection{Slime Servants}
If a creature spends 1 continuous minute under the effect of the slime from an olothec's Devolving Tentacle attack, that effect ends and their flesh devolves into a translucent primordial state. The creature loses their memories and sense of self, becoming a slime servant NPC who is loyal to the olothec. The following stat blocks are Ithu'rath's slime servants.

\textbf{Goblin Underboss Slime Servant}

\textbf{Human Brawler Slime Servant}

\textbf{Human Trickshot Slime Servant}

\textbf{Lizardfolk Bloodeye Slime Servant}

\textbf{Minotaur Lackey Slime Servant}

\textbf{Orc Rampart Slime Servant}

\chapter{The White Tower}
The White Tower

\section{Tower Stat Blocks}
\textbf{Bugbear Channeler}

\textbf{Bugbear Commander}

\textbf{Gnoll Bonesplitter}

\textbf{Hobgoblin Firerunner}

\textbf{Human Guard}

\textbf{Human Knave}

\textbf{Human Scoundrel}

\textbf{Human Storm Wizard}

\textbf{Human Trickshot}

\textbf{Orc Godcaller}

\textbf{Time Raider Freebooter}

\DndQuote%
{}
{''Betrayal is inevitable. I'm no fool. Everyone wants what I have—even people working for me. They want to wield the Grasp and aren't content to simply share the spoils of my brilliance. I make it look easy, but those peons don't realize that my hard work gives them the good life. They could never be me. That's why they fail.''}
{Xorannox}
\textbf{Time Raider Nemesis}

\textbf{Xorannox}

\chapter{Canyon of the Tower Crown}
Canyon of the Tower Crown

\section{Canyon Stat Blocks}
\textbf{Abyssal Hyena}

\textbf{Abyssal Hyena Companion}

\textbf{Aurumvas}

\begin{DndSidebar}[float=!b]{Teleporting Aurumvas}
Aurumvas's Greed Is Good bonus action allows him to teleport all over the battlefield. By the time the characters are ready to face him in a mortal showdown, each likely has at least a few expensive magic items or pieces of equipment worth at least 100 gp. That allows Aurumvas to teleport to the most vulnerable and unprotected foe each round and unleash his Greedy Hands attacks upon them.



\end{DndSidebar}

\textbf{Chimeron}

\textbf{Gnoll Bonesplitter}

\begin{DndSidebar}[float=!b]{Gnolls}
Gnolls aren't humanoids, but hyena-faced fiends—created by a dead demon lord, they carry his evil legacy wherever they roam. Most gnolls live for the thrill of the hunt and the taste of flesh, and they aren't choosy with their targets.



\end{DndSidebar}

\DndQuote%
{}
{''"Gnoll" is a shortened form of their endonym Gnol'sylrak, meaning "Children of As'sylrak." As'sylrak held the title of Archion of the Cult of the Black Star when he died, and he passed his immortal hate of all things with souls on to his children.''}
{Remainer}
\textbf{Pitling}

\textbf{Ruinant}

\textbf{Tormenauk}

\textbf{Wobalas}

\chapter{Tomb of the Keeper}
Tomb of the Keeper

\section{Tomb Stat Blocks}
\textbf{Ashyra}

\textbf{Haunt}

\textbf{Koptourok}

\begin{DndSidebar}[float=!b]{Koptourok}
Derived from an archaic duergar dialect, koptourok (cop-TOUR-ock) roughly translates to "dead tourist." These lost dead were raised in a world with a sky then spent their last moments in a dark, breathless place. They rise as rasping paper-skinned husks, like a cross between an unwrapped mummy and accordion bellows that forever expand and contract. They hunger only for what they've lost—breath.



\end{DndSidebar}

\textbf{Mummy}

\chapter{Terminal Excrescence}
Terminal Excrescence

\section{Terminal Excrescence Stat Blocks}
The following stat blocks appear in the lair.

\textbf{Amalgam}

\textbf{Hulking Brain}

\textbf{Lord Syuul}

\textbf{Mindkiller}

\textbf{Mindkiller Whelp}

\textbf{Orc Blacksmith}

\textbf{Orc Conduit}

\DndQuote%
{}
{''That thing—he called it the enucleator—it was awful. I was ripped apart and examined, but I was still alive. I could feel every piece of me separate. I have no idea how you put me back together.''}
{Gorar}
\textbf{Troll}

\textbf{Voiceless Talker}

\textbf{Voiceless Talker Artillerist}

\DndQuote%
{}
{''But where does your magic come from? The humors—your very lifeblood? Perhaps it is deep within your marrow. Bones are such curious things. We must dissect again and again—every piece smaller than the last—until we arrive at the truth. Your contributions to this experiment could change everything.''}
{Lord Syuul}
\begin{DndSidebar}[float=!b]{Psi-Tech}
While most psi-tech weapons can be used only by voiceless talkers, a character using the talent class (from the MCDM supplement \textit{The Talent and Psionics}) can use these weapons by gaining 1 strain each time they make an attack with one. A creature using a psi-tech ranged weapon uses Intelligence instead of Dexterity for attack and damage rolls.



\end{DndSidebar}

\chapter{Eighth City Advocacy Services}
Eighth City Advocacy Services

\section{ECAS Stat Blocks}
The following stat blocks appear in the lair.

\textbf{Devil Adjudicator}

\textbf{Devil Jurist}

\textbf{Devil Legate}

\textbf{Devil Magistrate}

\textbf{Devil Notary}

\textbf{Infernal Chancellor Lazivos}

\textbf{Night Hag}

\DndQuote%
{}
{''I don't know why I'm in this damned place ... Yes, yes. I know that was a pun. Who cares? Get me my heart, and I'll reward you. I'll be extra impressed if you bring me the eyes of the devil guarding it. I need a snack.''}
{Sasha Darkdream}
\begin{DndSidebar}[float=!b]{Night Hag in a Coven}
A night hag in a coven gains the following action option.



\paragraph{Optional Feature}
Phantoms Abound|FleeMortals



\end{DndSidebar}

\textbf{Specter}

\textbf{Wraith}

\chapter{Shadowkeep}
Shadowkeep

\section{Shadowkeep Stat Blocks}
The following stat blocks appear in the lair.

\textbf{Count Rhodar von Glauer}

\textbf{Devil Magistrate}

\begin{DndSidebar}[float=!b]{Excerpt from the Chronicle of the Chain of Acheron}
\DndQuote%
{}
{''"In those days, the Chain was in service to the Emperor Gaius VIII, third of the Five Crusading Emperors, and 27th Emperor of Caelia. Gaius continued the crusades started by his illustrious granduncle. Those were profitable days for the Chain, the Wheel, and the Gate. While the other Helltroopers waded into the deserts of Khoursir and Khemhara, we battled in the forests of Vaslor. We had it easy. I hate deserts.''}
{''"Thinking his vampiric power would save him, Count Rhodar von Glauer waited until he was fulsome in his strength before trying to extend his dark realm into Caelian territory. But we know how to deal with vampires. The black and gold standard of the von Glauer family reads, "I do the devil's work." We can attest to that. We lost many souls to that long siege, but in the end, we staked him. An old woman howled at us saying, "He is the land!" but we hear that stuff all the time. We sacked his castle, burned his forest and orchards, and delivered his staff to the emperor. That was harder than it sounds. Finger had to drop it in one of his bottomless sacks to stop the mists from coming. What the emperor will do with it, I don't know. I fear that staff may hold more than the power to make the mist.''}
{Dancer, Chronicler of the Chain of Acheron, Heroes 228}


\end{DndSidebar}

\textbf{Human Mercenary}

\textbf{Koptourok}

\textbf{Night Hag}

\begin{DndSidebar}[float=!b]{Night Hag in a Coven}
A night hag in a coven gains the following action option.



\paragraph{Optional Feature}
Phantoms Abound|FleeMortals



\end{DndSidebar}

\DndQuote%
{}
{''Hush little beet, now don't you cry,''}
{''Granny's gonna make you into a pie.''}
{''If that pie is tasting sour,''}
{''Granny's gonna add some boney powder.''}
{''If that powder is too dry,''}
{''Granny's gonna make somebody die.''}
{Putrid Agnes}
\textbf{Shade}

\textbf{Specter}

\DndQuote%
{}
{''While we drink blood, specters, wraiths, and other dead creatures without form feed on a different lifeforce. Some say it is energy drawn from the souls of the living these ghosts pass through. Other sages take a more tangible approach, claiming the spirits siphon the warmth needed to survive from the bodies of mortals. Whatever the case, these undead have no interest in feeding on vampires, and therefore make perfect companions.''}
{Count Rhodar von Glauer}
\textbf{Vampire}

\textbf{Vampire Spawn}

\textbf{Wight}

\textbf{Wraith}

\chapter{Durixaviinox's Rest}
Durixaviinox's Rest

\section{Durixaviinox's Rest Stat Blocks}
The following stat blocks appear in the lair.

\textbf{Air Spark}

\textbf{Crux of Frost}

\begin{DndSidebar}[float=!b]{Crux of Frost}
An entity of mirrored ice, the crux of frost is almost always covered by a frosted facade. They are commonly called frigid glintmasks due to their ability to imperfectly mimic creatures, inflicting supernatural terror on those they reflect. This fear makes the reflected creature more vulnerable to harm, their form and spirit withering to a brittle consistency under the crux's scrutiny.



\end{DndSidebar}

\textbf{Durixaviinox}

\textbf{Frost Giant Wind Sprinter}

\textbf{Giant Zombie}

\DndQuote%
{}
{''My rule was glorious, and the dragons, giants, and others who I protected lived prosperous lives. Those that remain—the dead that walk—don't remember the time that was. They are mindless drones. The lucky ones died ages ago and stayed dead.''}
{Durixaviinox}
\textbf{Specter}

\DndQuote%
{}
{''My breath carries the sorrow of thousands of dead dragons. Their despair will freeze the heart of my sister and make brittle anything else that stands in the way of my revenge for those lost.''}
{Durixaviinox}
\textbf{Troll}

\textbf{Wight Deathguard}

\textbf{Wraith}

\chapter{Coronal Hollow}
Coronal Hollow

\section{Coronal Hollow Stat Blocks}
The following stat blocks appear in the lair.

\textbf{Deep Dreamer}

\begin{DndSidebar}[float=!b]{Deep Dreamer}
A fine network of barbed black filaments supports the deep dreamer's fungal body, making them resemble a massive sinister dandelion. These plants drift through the air on gentle psionic currents, propelled by the spores they constantly release. As other creatures dream, the deep dreamer uses their filaments to feed on those creatures' psychic energy. A deep dreamer's influence can compel creatures to become lost in their dreams, never to wake up.



\end{DndSidebar}

\textbf{Empowered Bog Body}

\textbf{Empyrean Stag}

\begin{DndSidebar}[float=!b]{Empyrean Stag}
As Celestial wardens of the forest, empyrean stags appear only to defend forest denizens from great threats or to warn mortals about forthcoming calamities. Their white fur is dappled with blue, and their eyes shine with the darkness of night.



\end{DndSidebar}

\textbf{Essence of Mist}

\textbf{Lacuna}

\textbf{Lightbender Companion}

\begin{DndSidebar}[float=!b]{Mystic Connection: Lightbender}
If you're playing a beastheart and have a lightbender companion, you gain the following benefit at 9th level when you gain the beastheart's Mystic Connection feature:



\paragraph{Optional Feature}
Mystic Connection: Lightbender|FleeMortals



\end{DndSidebar}

\textbf{Qazyldrath}

\textbf{Rot Angel}

\textbf{Sunlight Nexus}

\textbf{Vampire}

\textbf{Vampire Spawn}

\DndQuote%
{}
{''Many wish to witness the day I swallow the sun with darkness, but none are as eager as vampires. They are more faithful to me than to their own sires, for I offer what even the gods deny—a world where vampires can move about freely.''}
{Qazyldrath}
\chapter{Eyes of the Mountain}
Eyes of the Mountain

\section{Eyes of the Mountain Stat Blocks}
The following stat blocks appear in the lair.

\textbf{Gibbering Mouther}

\textbf{Hulking Brain}

\textbf{Mindkiller}

\textbf{Voiceless Talker}

\begin{DndSidebar}[float=!b]{Draconic Aberrations}
Just like Yserthrax is a dragon with eldritch power, you can give the voiceless talkers and other aberrations in this dungeon draconic features. You can implement the following changes to any Aberration stat block in the lair to surprise characters who have faced voiceless talkers and their creations before and show that spending time around a powerful dragon like Yserthrax can magically warp nearby creatures:




\begin{description}
\item[Scales.]
The Aberration is covered in dragon scales. Their AC increases by 2.

\item[Wings.]
The Aberration has leathery dragon wings. If the Aberration already has a flying speed, their flight is powered by the wings instead of psionic might. If the Aberration doesn't have a flying speed, they gain a flying speed equal to their walking speed.

\item[Poison Immunity.]
The Aberration has immunity to poison damage and the poisoned condition.

\item[Draconic Senses.]
The Aberration is proficient in the Perception skill, and their proficiency bonus is doubled for any ability check they make that uses Perception. If the Aberration doesn't have blindsight, they gain blindsight out to a range of 30 feet.

\item[Breath Weapon.]
As an action, the Aberration breathes poison in a 30-foot cone. Each creature in that area must make a DC 13 plus PB Constitution saving throw. On a failed save, a creature takes PBd12 poison damage and is poisoned for 1 minute (save ends at end of turn). On a successful save, a creature takes half as much damage and isn't poisoned. The Aberration can't use this action again until they finish a long rest.

\item[Poison-Infused Weapons.]
When the Aberration hits a creature with a melee weapon attack, the attack deals an extra 5 (1d10) poison damage.

\end{description}



\end{DndSidebar}

\textbf{Voiceless Talker Artillerist}

\textbf{Yserthrax}

\chapter{Glass Cavern}
Glass Cavern

\section{Cavern Stat Blocks}
The following stat blocks appear in the lair.

\textbf{Fire Giant Lightbearer}

\textbf{Fire Giant Red Fist}

\textbf{Kingfissure Worm}

\textbf{Lightbender}

\begin{DndSidebar}[float=!b]{Lightbender}
Lightbenders prowl deserts, plains, forests—any sunbathed wilderness. Their adaptations make them skilled daylight predators. This monstrous creature's fur bends and refracts light from the surrounding environment, producing mirages that distract and confuse their prey— hence their name.



At a distance, a lightbender looks akin to a regular lion, but closer inspection reveals their glowing eyes, iridescent mane, and pair of lashing tails spiked with refractive crystals. The lightbender's pelt magically warps light around them to disguise their movement, allowing them to teleport while leaving behind a past visual imprint. Unsuspecting prey rarely realize they're staring at an afterimage of the lightbender until the fearsome creature pounces.



\end{DndSidebar}

\textbf{Xaantikorijek}

\chapter{Mount Brazen}
Mount Brazen

\section{Mount Brazen Stat Blocks}
The following stat blocks appear in the lair.

\textbf{Earth Spark}

\DndQuote%
{}
{''I do not care what you have come for. You are but an annoyance and have simply earned the right to burn before the rest of the world.''}
{Forzaantilirys}
\textbf{Fire Giant Lightbearer}

\textbf{Fire Giant Red Fist}

\textbf{Fire Giant Trooper}

\begin{DndSidebar}[float=!b]{Dragon Tactics}
Dragons seek lairs with lofty ceilings (or none at all) so they can use one of their greatest defenses: flight. A dramatic encounter with a dragon can quickly become a slog if the wyrm lands and never moves around the battlefield. Dragons are smart! They know to use their long reach and breath weapon to make life difficult for their enemies.



Dragons choose extremely defensible places for their lairs, which are part home and part treasure vault. The wyrms rely on loyal guardians and an environment that is inhospitable to mortal treasure hunters and would-be heroes. For instance, Forzaantilirys makes her home in Mount Brazen, an enormous volcano that is hotter than the Seven Cities of Hell. Her treasure sits on an obsidian island floating in a lake of lava in the heart of the mountain. Just walking through the tunnels to reach Forzaantilirys's hoard would bake most people to death.



\end{DndSidebar}

\textbf{Force of Blood}

\begin{DndSidebar}[float=!b]{Force of Blood}
Forces of blood, also called blood sovereigns, are warlords and diplomats who equally represent the birth of mortal life and the moment of its ending. They could be mistaken for a massive humanoid clad entirely in black-and-red toothed armor, but in truth, their "armor" is blackened, gnarled flesh punctuated with veins of boiling red.



\end{DndSidebar}

\textbf{Force of Iron}

\begin{DndSidebar}[float=!b]{Force of Iron}
A force of iron, also called an iron stalwart, appears as a four-armed centaur made of malleable metal with a head shaped like an empty knight's helm. Most are staunch, loyal, and uncompromising. Some forces of iron even achieve fame by facing down wickedness where they find it—but others become wickedness, using their strength to vanquish anything that even mildly displeases them.



\end{DndSidebar}

\DndQuote%
{}
{''Her breath melts stone,''}
{''And her scales ignite,''}
{''But it's the dragon's ire''}
{''That burns so bright.''}
{''When our bones are black,''}
{''And the world is ash,''}
{''She will save the Voice''}
{''For the very last.''}
{''When his eyes flash blue,''}
{''While the earth just burns,''}
{''It will become clear.''}
{''His memories return.''}
{''Then her claws will slash,''}
{''And her fire will rage,''}
{''As the Voice gasps last,''}
{''"Thus will end my age."''}
{''Then the Ashen Song''}
{''Takes his heart away,''}
{''As she whispers sweet,''}
{''"Vengeance came today."''}
{Ballad of the Ashen Song}
\textbf{Forzaantilirys}

\textbf{Haunt}

\chapter{Ash Queen's Reliquary}
Ash Queen's Reliquary

\section{Reliquary Stat Blocks}
The following stat blocks appear in the lair.

\textbf{Air Spark}

\textbf{Atæshia}

\textbf{Earth Spark}

\begin{DndSidebar}[float=!b]{Sapient Elementals}
The elementals in the Ash Queen's Reliquary are sapient beings who live and fight together as a band. They have a rich history, hold grudges, and experience emotions just like humans do. Other elementals, called motes, exist in the timescape as well. Motes think and behave the way most experienced gamers expect when they hear the word "elemental.



\end{DndSidebar}

\textbf{Essence of Mist}

\textbf{Essence of Storms}

\textbf{Essence of Tides}

\textbf{Force of Blood}

\textbf{Force of Earth}

\textbf{Force of Iron}

\textbf{Oracle of Storms}

\begin{DndSidebar}[float=!b]{Elemental Appearances}
Elementals in the Mundane World take the shape of other creatures formed of elements.




\begin{description}
\item[Air Spark.]
An air spark is a wisp of cloud that takes the shape of a winged serpent.

\item[Earth Spark.]
An earth spark is a rhinoceros formed from hunks of dirt and stone.

\item[Essence of Mist.]
An essence of mist has the misty silhouette of a cloaked humanoid. A brutal weapon forms in their grasp as they strike, rendering their victims breathless.

\item[Essence of Storms.]
An essence of storms is formed from streaks of colored cloud-stuff woven into the silhouette of a large eagle or falcon.

\item[Essence of Tides.]
An essence of tides looks like a shimmering blue manta ray that glides over land as easily as through water.

\item[Force of Blood.]
A force of blood could be mistaken for a massive humanoid clad entirely in black-and-red toothed armor, but in truth, their "armor" is blackened, gnarled flesh punctuated with veins of boiling red.

\item[Force of Earth.]
A force of earth resembles a nine-foot-tall great ape formed of rough stone.

\item[Force of Iron.]
A force of iron appears as a four-armed centaur made of malleable metal with a head shaped like an empty knight's helm.

\item[Oracle of Storms.]
An oracle of storms takes the shape of a storm giant made of thunderclouds.

\item[Sunlight Nexus.]
A being of pure radiance, a sunlight nexus is easily mistaken for a winged celestial made of stained glass.

\item[Water Spark.]
A water spark is an octopus formed out of water.

\end{description}



\end{DndSidebar}

\textbf{Sunlight Nexus}

\textbf{Water Spark}

\chapter{Boughs of Eternity}
Boughs of Eternity

\section{Tower Stat Blocks}
The following stat blocks appear in the lair.

\textbf{Camel}

\textbf{Devil Adjudicator}

\textbf{Devil Legate}

\textbf{Devil Magistrate}

\textbf{Devil Retainer}

\textbf{Dirt Mournling}

\textbf{High Mage Vairae}

\textbf{Specter}

\textbf{Swamp Dryad}

\textbf{Thornblood}

\textbf{Vampire}

\textbf{Wraith}

\chapter{New Psionic Powers}
A character playing a talent—the class found in the MCDM supplement The Talent and Psionics—can use the rules in that book to learn the following new powers from the creatures in this book. Powers that aren't found in \textit{The Talent and Psionics} are listed in this section and indicated in a creature's stat block with an asterisk (*).

Note that some stat block powers work differently than the version for player characters presented below, usually because a creature's innate psionic ability enhances their power.

Dream Spores

Liquefy

\DndQuote%
{}
{''Our power lies in the brain's soft tissue. Most barking ones lack the capacity, but those who can manifest require study. With a few modifications, I could turn their parlor tricks into power''}
{Lord Syuul}
Memory Thief

Phase

\begin{DndSidebar}[float=!b]{What Is a Power?}
A power is a specific psionic effect created by energy drawn out of the manifester's body. A power is an incredible tool that can solve problems, harm enemies, or protect and aid you and your allies, but it can also drain your life energy. Heroes who wield these gifts use only their thoughts and bodily energy to hurl objects through the air, send telepathic messages, manipulate time, and reshape the world around them.



The stronger the power, the more energy it can take from you. Every power has an order from 1 to 6. The higher a power's order, the more difficult it is to learn and use.



\end{DndSidebar}

Piscine Transformation

Vanish for One

\chapter{Magic Items}

\chapter*{Printable Statblocks}
\setlist[description]{nolistsep,listparindent=3pt,topsep=6pt,font={\bfseries\sffamily\small}}
\pagestyle{empty}

\begin{DndMonster}[float=t]{Abyssal Hyena}
\DndMonsterType{Medium fiend (Minion), chaotic evil}
\DndMonsterBasics[
armor-class = {11},
hit-points = {\DndDice{8}},
speed = {50 ft.}
]
\DndMonsterAbilityScores[
 str = 14,
 dex = 13,
 con = 12,
 int = 5,
 wis = 12,
 cha = 7
]
\DndMonsterDetails[
senses = {darkvision 60 ft., passive perception 11},
languages = {understands Abyssal but can't speak},
challenge = 1,
]
\DndMonsterAction{Death Snap}
When the hyena is reduced to 0 hit points, they can deal 1 piercing damage to a creature within 5 feet of them (no action required), provided they weren't incapacitated before dropping to 0 hit points.

\DndMonsterAction{Minion}
If the hyena takes damage from an attack or as the result of a failed saving throw, their hit points are reduced to 0. If the hyena takes damage from another effect, they die if the damage equals or exceeds their hit point maximum; otherwise they take no damage.

\DndMonsterSection{Actions}
\DndMonsterAction{Bite (Group Attack)}
\textsl{Melee Weapon Attack:} 4 to hit, reach 5 ft., one target. \textsl{Hit:} 1 piercing damage. If the target is a creature and two or more hyenas joined the attack, the target must succeed on a Strength saving throw or be knocked prone. The DC for this saving throw equals 10 plus the number of hyenas who joined the attack.

\end{DndMonster}



\begin{DndMonster}[float=t]{Abyssal Hyena Companion}
\DndMonsterType{Medium fiend (Companion), U}
\DndMonsterBasics[
armor-class = {13 plus PB (natural armor)},
hit-points = {\DndDice{7 + seven times caregiver's level (number of d8 Hit Dice equal to their caregiver's level)}},
speed = {50 ft.}
]
\DndMonsterAbilityScores[
 str = 16,
 dex = 14,
 con = 12,
 int = 5,
 wis = 12,
 cha = 8
]
\DndMonsterDetails[
saving-throws = {Str +3 plus PB, Dex +2 plus PB},
skills = {Perception +1 plus PB},
senses = {darkvision 60 ft., passive perception 11 plus PB},
]
\DndMonsterSection{Actions}
\DndMonsterAction{Signature Attack (Bite)}
\textsl{Melee Weapon Attack:} 3 plus PB to hit, reach 5 ft., one target. \textsl{Hit:} 1d6 plus PB piercing damage.

\DndMonsterAction{1st Level: Unbalancing Attack (2 Ferocity)}
The hyena makes a signature attack. On a hit, the target also has disadvantage on the next Strength or Dexterity saving throw they make before the start of their next turn.

\DndMonsterAction{3rd Level: Gobbling Bite (5 Ferocity)}
The hyena makes a signature attack with advantage. On a hit, the attack deals an extra 0d1 + PB necrotic damage, and the hyena regains hit points equal to the necrotic damage dealt.

\DndMonsterAction{5th Level: Cackle (8 Ferocity)}
Each enemy within 30 feet of the hyena who can hear them must succeed on a DC 10 plus PB Wisdom saving throw or take PBd4 psychic damage and be frightened of the hyena until the start of the hyena's next turn.

\DndMonsterSection{Reactions}
\DndMonsterAction{Blood Frenzy (Recharges after a Short or Long Rest)}
When the hyena's caregiver is hit with an attack and the hyena is within 5 feet of the attacker, the hyena can make a signature attack against the attacker.

\end{DndMonster}



\begin{DndMonster}[float=t]{Air Spark}
\DndMonsterType{Small elemental (air, Minion), A}
\DndMonsterBasics[
armor-class = {13},
hit-points = {\DndDice{22}},
speed = {20 ft., fly 50 ft. \textit{(hover)}}
]
\DndMonsterAbilityScores[
 str = 8,
 dex = 16,
 con = 10,
 int = 9,
 wis = 11,
 cha = 12
]
\DndMonsterDetails[
skills = {Perception +5},
senses = {darkvision 60 ft., passive perception 15},
languages = {Auran, Common},
challenge = 15,
]
\DndMonsterAction{Minion}
If the spark takes damage from an attack or as the result of a failed saving throw, their hit points are reduced to 0. If the spark takes damage from another effect, they die if the damage equals or exceeds their hit point maximum; otherwise they take no damage.

\DndMonsterSection{Actions}
\DndMonsterAction{Wind Sweep (Group Attack)}
\textsl{Melee Spell Attack:} 8 to hit, reach 5 ft., one target. \textsl{Hit:} 8 slashing damage, and if this attack was made by more than one air spark against a Large or smaller creature, the sparks can move that creature up to 5 feet horizontally for each spark that joined the attack.

\end{DndMonster}



\begin{DndMonster}[float*=b,width=\textwidth + 8pt]{Ashyra}
\begin{multicols}{2}
\DndMonsterType{Medium undead (Solo), lawful evil}
\DndMonsterBasics[
armor-class = {17 (natural armor)},
hit-points = {\DndDice{30d8 + 90}},
speed = {30 ft.}
]
\DndMonsterAbilityScores[
 str = 18,
 dex = 10,
 con = 17,
 int = 15,
 wis = 19,
 cha = 16
]
\DndMonsterDetails[
saving-throws = {Con +8, Int +7, Wis +9},
skills = {Arcana +7, History +7, Perception +9, Religion +7},
damage-resistances = {bludgeoning, piercing, slashing from mundane attacks},
damage-immunities = {necrotic, poison},
condition-immunities = {blinded, charmed, deafened, exhaustion, frightened, poisoned},
senses = {blindsight 120 ft., passive perception 19},
languages = {Common, telepathy 60 ft.},
challenge = 15,
]
\DndMonsterAction{Atonement (3/Day)}
If Ashyra fails a saving throw, she can choose to succeed instead. When she does, she must either end the effect of her Ancient Curse on one enemy within 120 feet of her or be dazed until the end of her next turn.

\DndMonsterAction{Mummy Dust}
Whenever Ashyra takes piercing or slashing damage, each creature within 5 feet of her takes 9 (2d8) poison damage.

\DndMonsterSection{Actions}
\DndMonsterAction{Multiattack}
Ashyra makes two attacks using Desiccating Slam, Soul Fire, or both, and she uses Ancient Curse.

\DndMonsterAction{Desiccating Slam}
\textsl{Melee Weapon Attack:} 9 to hit, reach 5 ft., one target. \textsl{Hit:} 14 (3d6 + 4) bludgeoning damage plus 10 (3d6) necrotic damage. The target's speed is reduced to 0, and they can't regain hit points until the start of Ashyra's next turn.

\DndMonsterAction{Soul Fire (3rd-Level Spell)}
\textsl{Ranged Spell Attack:} 9 to hit, range 60 ft., one target. \textsl{Hit:} 20 (3d10 + 4) radiant damage, and if the target is a creature, their hit point maximum decreases by a number equal to the damage taken. This reduction lasts until the target finishes a long rest or is targeted by a cure ailment power of 4th order or higher, a \textit{greater restoration} spell, or a similar supernatural effect. The target dies if this effect reduces their hit point maximum to 0.

\DndMonsterAction{Ancient Curse}
Ashyra inflicts weakness on one creature she can see within 120 feet of her. The target must succeed on a DC 17 Wisdom saving throw or be cursed. While the target is cursed in this way, their attacks deal half damage. The curse lasts until Ashyra chooses to end it or until removed by a cure ailment power, a \textit{remove curse} spell, or a similar supernatural effect.

\DndMonsterSection{Bonus Actions}
\DndMonsterAction{Misty Step (2nd-Level Spell)}
Ashyra teleports up to 30 feet to an unoccupied space she can see.

\DndMonsterSection{Reactions}
\DndMonsterAction{Weight of Ages}
When a creature Ashyra can see within 60 feet of her deals damage to her, Ashyra brings down the weight of ages on them. The creature must succeed on a DC 17 Strength saving throw or be restrained until the end of their next turn.

\DndMonsterSection{Legendary Actions}
Ashyra has three villain actions. She can take each action once during an encounter after an enemy's turn. She can take these actions in any order but can use only one per round.

\begin{DndMonsterLegendaryActions}
\DndMonsterLegendaryAction{Action 1: Sands of the Ancients}{A storm of stinging dirt whirls around Ashyra in a 15-foot-radius until the end of her next turn. For the duration, that area is heavily obscured, and when a creature starts their turn in or enters that area for the first time on a turn, they must succeed on a DC 17 Constitution saving throw or take 22 (4d10) necrotic damage.
}
\DndMonsterLegendaryAction{Action 2: Dust to Dust}{Ashyra bursts into a cloud of dust and makes one Desiccating Slam attack against each enemy within 60 feet of her, then she reforms in an unoccupied space within 60 feet of her original location.
}
\DndMonsterLegendaryAction{Action 3: Last Rites}{Ashyra commands the land to entomb her foes. Each enemy within 30 feet of her must make a DC 17 Strength saving throw. On a failed save, a target takes 21 (6d6) bludgeoning damage and is entombed in stone. On a successful save, a target takes half as much damage and is not entombed. An entombed creature is restrained, has total cover against attacks and other effects outside the stone, and takes 7 (2d6) necrotic damage at the end of each of their turns. The creature can be freed if the stone is destroyed (AC 17; 30 hit points; immunity to poison and psychic damage).
}
\end{DndMonsterLegendaryActions}
\end{multicols}
\end{DndMonster}



\begin{DndMonster}[float*=b,width=\textwidth + 8pt]{Atæshia}
\begin{multicols}{2}
\DndMonsterType{Medium elemental (air, earth, fire, Leader), N}
\DndMonsterBasics[
armor-class = {20 (natural armor)},
hit-points = {\DndDice{35d8 + 175}},
speed = {50 ft., fly 30 ft. \textit{(hover)}}
]
\DndMonsterAbilityScores[
 str = 19,
 dex = 21,
 con = 21,
 int = 22,
 wis = 23,
 cha = 23
]
\DndMonsterDetails[
saving-throws = {Con +12, Wis +13},
skills = {Arcana +13, Deception +13, History +13, Nature +13, Perception +13, Persuasion +13, Religion +13},
damage-resistances = {cold, necrotic, poison, psychic, radiant},
damage-immunities = {fire},
condition-immunities = {charmed, frightened, poisoned},
senses = {blindsight 120 ft., passive perception 23},
languages = {all},
challenge = 23,
]
\DndMonsterAction{Offer to Remembrance (3/Day)}
When Atæshia fails a saving throw, a willing creature within 60 feet of her who can see her can give their life to protect her. This creature dies instantly, and Atæshia succeeds on the saving throw.

\DndMonsterAction{Reborn From War}
When Atæshia is reduced to 0 hit points, she disintegrates into ash on a solemn wind. For the next 100 years, if a Celestial, a Elemental, or a Humanoid whose challenge rating or level is 5 or higher dies within 60 feet of where Atæshia died, Atæshia returns to life where she died with 25 hit points.

\DndMonsterSection{Actions}
\DndMonsterAction{Multiattack}
Atæshia makes three Pale Flame attacks. She can replace one attack with a use of Moment of Loss, if available.

\DndMonsterAction{Pale Flame}
\textsl{Melee or Ranged Spell Attack:} 13 to hit, reach 5 ft. or range 120 ft., one target. \textsl{Hit:} 20 (4d6 + 6) radiant damage plus 14 (4d6) fire damage. If this attack reduces a creature to 0 hit points, they die and all their remains burn to nothing.

\DndMonsterAction{Moment of Loss (Recharge 6)}
Atæshia transforms one creature she can see within 60 feet of her into the form of a living ash cloud. While transformed to ash, the target is immune to all damage, can't provide cover to other creatures, and can't take reactions. At the start of the target's next turn, they reform and must make a DC 21 Constitution saving throw. On a failed save, the target is reduced to 1 hit point. On a successful save, the target takes 44 (8d10) necrotic damage. This damage can't reduce the target's hit points below 1.

\DndMonsterAction{From Death, Life (1/Day)}
Atæshia uses her Kindled Sacrifice villain action, even if it is unavailable, to sacrifice herself.

\DndMonsterSection{Bonus Actions}
\DndMonsterAction{Flickering Flame}
Atæshia forms into a cloud of ash, teleports to an unoccupied space she can see within 60 feet of her, and reforms into her true form. Each creature within 5 feet of the space she left must succeed on a DC 21 Dexterity saving throw or take 5 (1d10) fire damage.

\DndMonsterSection{Reactions}
\DndMonsterAction{Convocation of Ash}
When another Elemental Atæshia can see within 60 feet of her is hit with an attack, the damage from the attack is halved and the target can move up to half their speed without provoking opportunity attacks.

\DndMonsterSection{Legendary Actions}
Atæshia has three villain actions. She can take each action once during an encounter after an enemy's turn. She can take these actions in any order but can use only one per round.

\begin{DndMonsterLegendaryActions}
\DndMonsterLegendaryAction{Action 1: Wave of Sorrow}{Atæshia unleashes a pulse of overwhelming sadness. Each enemy Atæshia can see within 120 feet of her must make a DC 21 Wisdom saving throw. On a failed save, a creature is dazed for 1 minute (save ends at end of turn). On a successful save, a creature can't take reactions until the end of their next turn.
}
\DndMonsterLegendaryAction{Action 2: Kindled Sacrifice}{Atæshia chooses a willing creature within 60 feet of her to die so their allies may live. The target dies, and up to two other dead creatures within 60 feet of the target return to life. These revived creatures regain the same number of hit points the sacrificed creature had, up to the revived creature's hit point maximum.
}
\DndMonsterLegendaryAction{Action 3: Life's Finale}{Atæshia releases a torrent of primordial anguish to unmake the marrow of creation around her. Each enemy within 60 feet of her must make a DC 21 Charisma saving throw. On a failed save, they take 55 (10d10) necrotic damage and can't regain hit points until the end of Atæshia's next turn. On a successful save, they take half as much damage, but suffer no other effect. Succeed or fail, if this damage reduces a creature to 0 hit points, they immediately fail one death saving throw.
}
\end{DndMonsterLegendaryActions}
\end{multicols}
\end{DndMonster}



\begin{DndMonster}[float*=b,width=\textwidth + 8pt]{Aurumvas}
\begin{multicols}{2}
\DndMonsterType{Huge fiend (demon, category 6, Leader), chaotic evil}
\DndMonsterBasics[
armor-class = {18 (natural armor)},
hit-points = {\DndDice{26d12 + 104}},
speed = {40 ft., fly 40 ft.}
]
\DndMonsterAbilityScores[
 str = 25,
 dex = 15,
 con = 18,
 int = 22,
 wis = 20,
 cha = 21
]
\DndMonsterDetails[
saving-throws = {Dex +7, Wis +10, Cha +10},
skills = {Arcana +11, Athletics +12, Deception +10, History +11, Insight +10, Perception +10, Persuasion +10},
damage-resistances = {necrotic; psychic; bludgeoning, piercing, slashing from mundane attacks},
condition-immunities = {charmed, dazed, exhaustion, frightened, paralyzed, stunned},
senses = {darkvision 120 ft., soulsight 30 ft., passive perception 20},
languages = {all, telepathy 120 ft.},
challenge = 14,
]
\DndMonsterAction{Abyssal Resistance (Costs 2 Souls)}
When Aurumvas fails a saving throw, he succeeds instead.

\DndMonsterAction{Lethe}
When Aurumvas's soul count is 0, he has advantage on attack rolls, disadvantage on saving throws, and his Intelligence score becomes 3 (-4). Additionally, Aurumvas must use his movement on each of his turns to move as close as possible to the nearest creature he can sense with his soulsight, and then if he is able, he must use his action to attack and attempt to kill that creature. Aurumvas can't act with any other purpose until he adds 1 to his soul count.

\DndMonsterAction{Soul Devourer}
When Aurumvas reduces a creature who isn't a Construct or an Undead to 0 hit points or deals damage to a dying creature, the creature must make a DC 11 Wisdom saving throw. On a failed save, Aurumvas consumes the creature's soul and adds 1 to his soul count. A creature whose soul is consumed in this way immediately dies, and they can't be restored to life by any means short of a \textit{wish} spell.

\DndMonsterAction{Soul Weapons}
While Aurumvas's soul count is 2 or higher, his weapon attacks are magical, and allies within 60 feet of him deal an extra 5 (1d10) psychic damage with weapon attacks.

\DndMonsterSection{Actions}
\DndMonsterAction{Multiattack}
Aurumvas makes three attacks using Greedy Hands, Covetous Bolt, or both.

\DndMonsterAction{Greedy Hands}
\textsl{Melee Weapon Attack:} 12 to hit, reach 10 ft., one creature. \textsl{Hit:} 14 (2d6 + 7) bludgeoning damage plus 10 (3d6) psychic damage. Aurumvas can burn 1 soul to gain temporary hit points equal to the psychic damage dealt.

\DndMonsterAction{Covetous Bolt}
\textsl{Ranged Spell Attack:} 11 to hit, range 120 ft., one target. \textsl{Hit:} 22 (5d8) force damage. If the target has spell slots, Aurumvas can burn 2 souls to force the target to expend their highest-level spell slot available with no effect.

\DndMonsterSection{Bonus Actions}
\DndMonsterAction{Greed Is Good}
Aurumvas chooses an object he can see within 30 feet of him that is worth at least 100 gp. He teleports to an unoccupied space within 5 feet of that object.

\DndMonsterSection{Reactions}
\DndMonsterAction{Absorb Soul}
If a demon with a soul count of 1 or higher who Aurumvas can see dies within 60 feet of him, Aurumvas gains 1 soul.

\DndMonsterSection{Legendary Actions}
Aurumvas has three villain actions. He can take each action once during an encounter after an enemy's turn. He can take these actions in any order but can use only one per round.

\begin{DndMonsterLegendaryActions}
\DndMonsterLegendaryAction{Action 1: Time Is Money}{Aurumvas warps time with his supernatural greed. Each ally he can see within 120 feet of him moves up to their speed. At the end of this round of combat, Aurumvas rearranges the initiative order in any way he chooses.
}
\DndMonsterLegendaryAction{Action 2: That's Ours}{Aurumvas chooses up to three magic items he can see within 30 feet of him that aren't artifacts or of legendary rarity and that don't require attunement. He teleports each item to an ally who he can see within 30 feet of him, placing it either into their hands or at their feet. If a targeted object is being worn or carried by an enemy, the enemy can make a DC 19 Charisma saving throw, preventing that object from being teleported on a success.
}
\DndMonsterLegendaryAction{Action 3: At Any Cost}{Aurumvas teleports Tiny explosive treasures from his vault to four different points he can see within 120 feet of him. Each creature within 10 feet of one or more of these points must make a DC 19 Dexterity saving throw, taking 33 (6d10) force damage on a failed save, or half as much damage on a successful one.
}
\end{DndMonsterLegendaryActions}
\end{multicols}
\end{DndMonster}



\begin{DndMonster}[float*=b,width=\textwidth + 8pt]{Baron Uthrak}
\begin{multicols}{2}
\DndMonsterType{Medium humanoid (human, Leader), lawful evil}
\DndMonsterBasics[
armor-class = {16 (half plate armor)},
hit-points = {\DndDice{17d8 + 68}},
speed = {30 ft.}
]
\DndMonsterAbilityScores[
 str = 19,
 dex = 12,
 con = 19,
 int = 14,
 wis = 14,
 cha = 16
]
\DndMonsterDetails[
saving-throws = {Str +7, Con +7},
skills = {Athletics +10, Deception +6, Insight +5, Intimidation +9, Perception +5},
languages = {Common},
challenge = 7,
]
\DndMonsterAction{Bloodstones (3 Uses)}
When Baron Uthrak fails a saving throw while holding his magic longsword Nine Lives, he can choose to succeed instead. If he does so, he takes 7 (2d6) necrotic damage as one of the bloodstones set in the sword's pommel fills with his blood.

\DndMonsterAction{Exploit Opening (3/Day)}
When Baron Uthrak makes an attack, he has advantage on the attack roll.

\DndMonsterSection{Actions}
\DndMonsterAction{Multiattack}
Baron Uthrak makes two attacks using Nine Lives, Whip, or both.

\DndMonsterAction{Nine Lives}
\textsl{Melee Weapon Attack:} 8 to hit, reach 5 ft., one target. \textsl{Hit:} 9 (1d8 + 5) slashing damage plus 3 (1d6) necrotic damage.

\DndMonsterAction{Whip}
\textsl{Melee Weapon Attack:} 7 to hit, reach 15 ft., one target. \textsl{Hit:} 6 (1d4 + 4) slashing damage, and the target must succeed on a DC 15 Strength saving throw or be pulled up to 15 feet toward Baron Uthrak.

\DndMonsterSection{Bonus Actions}
\DndMonsterAction{Kneel, Peasant!}
Baron Uthrak kicks a Large or smaller creature within 5 feet of him. The target must succeed on a DC 15 Strength saving throw or take 6 (1d4 + 4) bludgeoning damage and be knocked prone.

\DndMonsterSection{Reactions}
\DndMonsterAction{Riposte}
When a creature Baron Uthrak can see misses him with a melee attack, Baron Uthrak makes a melee attack against that creature.

\DndMonsterSection{Legendary Actions}
Baron Uthrak has three villain actions. He can take each action once during an encounter after an enemy's turn. He can take these actions in any order but can use only one per round.

\begin{DndMonsterLegendaryActions}
\DndMonsterLegendaryAction{Action 1: Shoot!}{Each ally within 60 feet of Baron Uthrak who can hear him can make a ranged weapon attack (no action required).
}
\DndMonsterLegendaryAction{Action 2: Form Up!}{Baron Uthrak and each ally within 60 feet of him who can hear him can move up to their speed without provoking opportunity attacks. Baron Uthrak and each ally enters a defensive stance that lasts until that creature moves or is restrained or incapacitated. When two or more allies in a defensive stance are within 5 feet of each other, attack rolls against them are made with disadvantage.
}
\DndMonsterLegendaryAction{Action 3: Lead from the Front}{Baron Uthrak moves up to twice his speed without provoking opportunity attacks. During or after this movement, he can make up to four Nine Lives or Whip attacks with advantage, each against a different target. After he attacks a target in this way, each ally within 5 feet of that target can make a melee attack with advantage against the same target (no action required).
}
\end{DndMonsterLegendaryActions}
\end{multicols}
\end{DndMonster}



\begin{DndMonster}[float=t]{Basalt Stone Giant}
\DndMonsterType{Huge giant (Controller), A}
\DndMonsterBasics[
armor-class = {16 (natural armor)},
hit-points = {\DndDice{12d12 + 72}},
speed = {40 ft.}
]
\DndMonsterAbilityScores[
 str = 24,
 dex = 15,
 con = 22,
 int = 13,
 wis = 15,
 cha = 12
]
\DndMonsterDetails[
saving-throws = {Con +9, Wis +5},
skills = {Athletics +10, Perception +5},
damage-resistances = {bludgeoning, piercing, slashing from mundane attacks that aren't adamantine},
senses = {darkvision 60 ft., passive perception 15},
languages = {Giant},
challenge = 8,
]
\DndMonsterAction{False Appearance}
While motionless, the giant is indistinguishable from a statue.

\DndMonsterAction{Stone Flesh}
When a creature attacks the giant with a melee weapon that isn't adamantine and rolls a 1 on the d20 for that attack, that weapon is destroyed if it is mundane. If that weapon is supernatural, it loses all supernatural properties for 1 hour.

\DndMonsterSection{Actions}
\DndMonsterAction{Multiattack}
The giant makes two Rune-Signed Blade attacks. They can replace one attack with a Forked Knives attack.

\DndMonsterAction{Rune-Signed Blade}
\textsl{Melee Weapon Attack:} 10 to hit, reach 10 ft., one target. \textsl{Hit:} 25 (4d8 + 7) slashing damage. If the target is a creature, each time they are hit by this attack, their movement speed is reduced by a cumulative 10 feet (to a minimum speed of 5 feet) until they regain at least 1 hit point.

\DndMonsterAction{Forked Knives}
\textsl{Melee or Ranged Weapon Attack:} 10 to hit, reach 10 ft. or range 60/240 ft., one target. \textsl{Hit:} 17 (3d6 + 7) piercing damage, and if the target is a Medium or smaller creature on the ground, they are knocked prone and restrained as they are pinned by the knife. The target or a creature who can reach the knife can use an action to make a DC 18 Strength (Athletics) check to try to dislodge the knife, freeing the pinned target and ending the restrained condition on a success.

\DndMonsterSection{Reactions}
\DndMonsterAction{Resonate Rune}
When the giant is hit by a melee attack, each creature within 10 feet of them must succeed on a DC 17 Strength saving throw or be pushed up to 10 feet away from the giant.

\end{DndMonster}



\begin{DndMonster}[float=t]{Basilisk}
\DndMonsterType{Medium monstrosity (Brute), U}
\DndMonsterBasics[
armor-class = {15 (natural armor)},
hit-points = {\DndDice{10d8 + 20}},
speed = {30 ft.}
]
\DndMonsterAbilityScores[
 str = 16,
 dex = 10,
 con = 15,
 int = 2,
 wis = 8,
 cha = 7
]
\DndMonsterDetails[
damage-immunities = {poison},
condition-immunities = {petrified, poisoned},
senses = {darkvision 60 ft., passive perception 9},
challenge = 3,
]
\DndMonsterAction{Alchemical Ingredients}
After the basilisk dies, a creature can make a DC 12 Wisdom (Medicine) check using an herbalism kit and the basilisk's gullet, destroying the gullet in the process. On a success, the creature creates three doses of a salve. One dose of this salve can be applied to a petrified creature as an action, and 1 minute after the salve is applied, the petrified condition ends for that creature.

\DndMonsterAction{Toxic Ichor}
A creature who deals piercing or slashing damage to the basilisk while within 5 feet of them takes 2 (1d4) poison damage.

\DndMonsterSection{Actions}
\DndMonsterAction{Multiattack}
The basilisk uses Petrifying Eye Beams and makes one Bite attack or uses Poison Fumes, if available.

\DndMonsterAction{Bite}
\textsl{Melee Weapon Attack:} 5 to hit, reach 5 ft., one target. \textsl{Hit:} 6 (1d6 + 3) piercing damage plus 5 (2d4) poison damage.

\DndMonsterAction{Petrifying Eye Beams}
The basilisk targets one or two creatures they can see within 30 feet of them. Each target must succeed on a DC 12 Constitution saving throw or be restrained as they magically begin to turn to stone. Until the end of the restrained creature's next turn, they or an ally within 5 feet of them can use an action to cut the encroaching stone from the restrained creature's body using a weapon that deals slashing damage. When they do, the restrained creature takes 13 (2d12) slashing damage, which can't be reduced in any way, and they are no longer restrained.

A creature restrained in this way must repeat the saving throw at the end of their turn. On a successful save, the effect ends. On a failed save, this creature is petrified until freed by a cure ailment power of 4th order or higher, a \textit{greater restoration} spell, or a similar supernatural effect.

\DndMonsterAction{Poison Fumes (Recharge 4\textendash 6)}
The basilisk exhales poisonous fumes in a 15-foot cone. Each creature in this area must make a DC 12 Constitution saving throw, taking 10 (4d4) poison damage on a failed save, or half as much damage on a successful one.

\DndMonsterSection{Reactions}
\DndMonsterAction{Poison Splash}
When the basilisk takes piercing or slashing damage, their wound spews poison. Each creature within 10 feet of the basilisk must succeed on a DC 12 Constitution saving throw or be poisoned for 1 minute (save ends at end of turn).

\end{DndMonster}



\begin{DndMonster}[float=t]{Blood-Borne Ooze}
\DndMonsterType{Tiny ooze (Ambusher), neutral evil}
\DndMonsterBasics[
armor-class = {13 (natural armor)},
hit-points = {\DndDice{4d4 + 4}},
speed = {20 ft.}
]
\DndMonsterAbilityScores[
 str = 8,
 dex = 15,
 con = 13,
 int = 1,
 wis = 10,
 cha = 2
]
\DndMonsterDetails[
skills = {Stealth +4},
condition-immunities = {blinded, charmed, deafened, exhaustion, flanked, prone},
senses = {blindsight 60 ft. (blind beyond this radius), passive perception 10},
challenge = 1/4,
]
\DndMonsterAction{Amorphous}
The ooze can move through a space as narrow as 1 inch wide without squeezing.

\DndMonsterAction{False Appearance}
While the ooze remains motionless, they are indistinguishable from a pool of blood.

\DndMonsterSection{Actions}
\DndMonsterAction{Pseudopod}
\textsl{Melee Weapon Attack:} 4 to hit, reach 5 ft., one target. \textsl{Hit:} 4 (1d4 + 2) bludgeoning damage plus 2 (1d4) necrotic damage. If the target is a creature who isn't a Construct or an Undead, they must succeed on a DC 11 Constitution saving throw or the ooze melds into the target's body.

While inside a creature, the ooze has total cover against attacks and other effects originating outside that host.

If the host creature takes 5 damage or more on a single turn from a source other than the ooze, the ooze must succeed on a DC 12 Constitution saving throw at the end of that turn or exit the host, entering the nearest unoccupied space of the ooze's choice. A cure ailment power, \textit{protection from poison} spell, or \textit{lesser restoration} spell cast on the host also forces the ooze out. By spending 5 feet of their movement, the ooze can voluntarily leave the host's body.

\DndMonsterAction{Crimson Feast (Inside Host Only)}
The ooze consumes their host creature's bodily fluids. The host must make a DC 11 Constitution saving throw. On a failed save, the host takes 5 (2d4) necrotic damage and is poisoned until the end of their next turn. On a successful save, the host takes half as much damage and isn't poisoned.

\end{DndMonster}



\begin{DndMonster}[float*=b,width=\textwidth + 8pt]{Bloodlord Varrox}
\begin{multicols}{2}
\DndMonsterType{Medium humanoid (hobgoblin, Leader), lawful evil}
\DndMonsterBasics[
armor-class = {18 (plate armor)},
hit-points = {\DndDice{14d8 + 56}},
speed = {30 ft.}
]
\DndMonsterAbilityScores[
 str = 18,
 dex = 11,
 con = 18,
 int = 11,
 wis = 14,
 cha = 15
]
\DndMonsterDetails[
saving-throws = {Wis +5, Cha +5},
skills = {Arcana +3, Athletics +7, History +3, Intimidation +5},
damage-resistances = {fire, necrotic},
senses = {darkvision 60 ft., passive perception 12},
languages = {Common, Goblin, Infernal},
challenge = 8,
]
\DndMonsterAction{Infernal Ichor}
When Varrox dies, his corpse unleashes a spray of burning orange ichor. Each creature within 5 feet of him takes 5 fire damage.

\DndMonsterAction{Infernal Sacrifice (3/Day)}
When Varrox fails a saving throw, he can call on Hell to succeed instead. When he does, each ally within 60 feet of him takes 5 necrotic damage.

\DndMonsterSection{Actions}
\DndMonsterAction{Multiattack}
Varrox makes three attacks using Soulsword, Death Skulls, or both.

\DndMonsterAction{Soulsword}
\textsl{Melee Weapon Attack:} 7 to hit, reach 5 ft., one target. \textsl{Hit:} 11 (2d6 + 4) slashing damage plus 10 (3d6) necrotic damage. If the target is a creature, they can't regain hit points until the start of Varrox's next turn.

\DndMonsterAction{Death Skulls}
\textsl{Ranged Weapon Attack:} 7 to hit, range 80/320 ft., one target. \textsl{Hit:} 8 (1d8 + 4) piercing damage plus 9 (2d8) fire damage, and Varrox can teleport to an unoccupied space he can see within 5 feet of the target.

\DndMonsterSection{Bonus Actions}
\DndMonsterAction{Advance and Attack}
Varrox barks orders at one ally within 60 feet of him who can hear him. That ally can use their reaction to move up to their speed and make a weapon attack.

\DndMonsterSection{Reactions}
\DndMonsterAction{An Army from Blood}
When a non-minion hobgoblin who Varrox can see within 60 feet of him takes damage, Varrox magically summons three hobgoblin recruits who appear in unoccupied spaces nearest to the damaged hobgoblin.

\DndMonsterSection{Legendary Actions}
Varrox has three villain actions. He can take each action once during an encounter after an enemy's turn. He can take these actions in any order but can use only one per round.

\begin{DndMonsterLegendaryActions}
\DndMonsterLegendaryAction{Action 1: Uncanny Leadership}{Varrox chooses up to six non-minion allies within 60 feet of him who can see or hear him. Varrox and each chosen creature gain 10 temporary hit points and can move up to their speed.
}
\DndMonsterLegendaryAction{Action 2: Hellish Transport}{Varrox and each ally within 60 feet of him can teleport up to 60 feet to an unoccupied space they can see.
}
\DndMonsterLegendaryAction{Action 3: I Am Fire and Death}{Each enemy within 20 feet of Varrox must make a DC 15 Constitution saving throw, taking 11 (2d10) fire damage plus 11 (2d10) necrotic damage on a failed save, or half as much damage on a successful one. After the explosion, a shroud of black flame covers Varrox for 1 minute. For the duration, each time another creature hits Varrox with an attack while within 5 feet of him or touches him, that creature takes 5 (1d10) fire damage plus 5 (1d10) necrotic damage.
}
\end{DndMonsterLegendaryActions}
\end{multicols}
\end{DndMonster}



\begin{DndMonster}[float=t]{Bugbear Channeler}
\DndMonsterType{Medium humanoid (bugbear, Controller), A}
\DndMonsterBasics[
armor-class = {14 (studded leather armor)},
hit-points = {\DndDice{12d8 + 48}},
speed = {30 ft.}
]
\DndMonsterAbilityScores[
 str = 10,
 dex = 15,
 con = 18,
 int = 15,
 wis = 16,
 cha = 18
]
\DndMonsterDetails[
saving-throws = {Wis +6, Cha +7},
skills = {Arcana +5, Perception +6, Stealth +5, Survival +6},
senses = {darkvision 60 ft., passive perception 16},
languages = {Common, Goblin},
challenge = 7,
]
\DndMonsterAction{Bugbear's Inspiration}
Allied Humanoids within 15 feet of the channeler who can see or hear them have advantage on Wisdom and Charisma saving throws.

\DndMonsterAction{Spellcasting (Utility)}
In addition to any other spells in this stat block, the channeler can cast the following spells, using Charisma as the spellcasting ability (spell save DC 15):
\begin{DndMonsterSpells}
  \DndInnateSpellLevel{light\textsuperscript{Casting Times}, \textit{message}\textsuperscript{Casting Times}}
  \DndInnateSpellLevel[1]{detect magic\textsuperscript{Casting Times}, \textit{disguise self}\textsuperscript{Casting Times}, \textit{identify}\textsuperscript{Casting Times}}
\end{DndMonsterSpells}
\DndMonsterSection{Actions}
\DndMonsterAction{Multiattack}
The channeler makes two Channeling Staff attacks.

\DndMonsterAction{Channeling Staff}
\textsl{Melee Spell Attack:} 7 to hit, reach 5 ft., one target. \textsl{Hit:} 8 (1d8 + 4) bludgeoning damage plus 9 (2d8) lightning damage. If the target is a Large or smaller creature, the channeler can move them up to 10 feet horizontally.

\DndMonsterAction{Vine Eruption (3/Day; 4th-Level Spell)}
The channeler attempts to use primal magic to overload a creature they can see within 60 feet of them. The target must make a DC 15 Constitution saving throw as thorny vines erupt from their body and attempt to envelop them. On a failed save, the target takes 27 (5d10) piercing damage and is restrained until the vines are destroyed. On a successful save, a target takes half as much damage and isn't restrained. A creature restrained in this way takes 5 (1d10) piercing damage at the start of each of their turns. The vines have AC 15, 10 hit points, and immunity to psychic damage.

\DndMonsterAction{Ball Lightning (1/Day; 7th-Level Spell; Concentration)}
The channeler creates a 20-foot-radius sphere of electricity centered on a point they can see within 120 feet of them. The sphere lasts for 1 minute. When an enemy starts their turn in that area, they must succeed on a DC 15 Dexterity saving throw or take 27 (6d8) lightning damage. As a bonus action, the channeler can move the sphere up to 20 feet.

\DndMonsterSection{Reactions}
\DndMonsterAction{Shared Shield (2/Day; 2nd-Level Spell)}
When a creature the channeler can see within 30 feet of them is hit with an attack, the channeler can create a barrier of magical force around that target. Until the start of the channeler's next turn, the target gains a +5 bonus to AC, including against the triggering attack.

\end{DndMonster}



\begin{DndMonster}[float=t]{Bugbear Commander}
\DndMonsterType{Medium humanoid (bugbear, Support), A}
\DndMonsterBasics[
armor-class = {18 (plate armor)},
hit-points = {\DndDice{15d8 + 30}},
speed = {30 ft.}
]
\DndMonsterAbilityScores[
 str = 18,
 dex = 15,
 con = 15,
 int = 17,
 wis = 14,
 cha = 15
]
\DndMonsterDetails[
saving-throws = {Wis +5, Cha +5},
skills = {Intimidation +5, Persuasion +5, Stealth +5, Survival +5},
senses = {darkvision 60 ft., passive perception 12},
languages = {Common, Goblin},
challenge = 5,
]
\DndMonsterAction{Commander's Inspiration}
Allied Humanoids within 30 feet of the commander who can see or hear them have advantage on Wisdom and Charisma saving throws and gain a +1 bonus to weapon attack and damage rolls.

\DndMonsterSection{Actions}
\DndMonsterAction{Multiattack}
The commander makes two Greatsword or two Longbow attacks.

\DndMonsterAction{Greatsword}
\textsl{Melee Weapon Attack:} 7 to hit, reach 5 ft., one target. \textsl{Hit:} 11 (2d6 + 4) slashing damage, and one ally within 30 feet of the target who can see the commander has advantage on their next melee attack roll made against the target before the start of the commander's next turn.

\DndMonsterAction{Longbow}
\textsl{Ranged Weapon Attack:} 5 to hit, range 150/600 ft., one target. \textsl{Hit:} 6 (1d8 + 2) piercing damage.

\DndMonsterAction{Bark Orders (Recharge 6)}
Each ally within 60 feet of the commander who can hear and understand them can use a reaction to move up to their speed and make a weapon attack.

\DndMonsterSection{Reactions}
\DndMonsterAction{Bolstering Pride}
When another creature within 60 feet of the commander who can hear and understand them makes a saving throw, the commander can give that creature advantage on the saving throw.

\end{DndMonster}



\begin{DndMonster}[float=t]{Camel}
\DndMonsterType{Large beast (Controller), U}
\DndMonsterBasics[
armor-class = {11},
hit-points = {\DndDice{3d10 + 12}},
speed = {50 ft.}
]
\DndMonsterAbilityScores[
 str = 17,
 dex = 13,
 con = 19,
 int = 4,
 wis = 12,
 cha = 5
]
\DndMonsterDetails[
saving-throws = {Str +5, Con +6},
challenge = 1/2,
]
\DndMonsterAction{Beast of Burden}
The camel is considered to be one size larger for the purpose of determining their carrying capacity.

\DndMonsterAction{Desert Dweller}
The camel ignores difficult terrain created by dirt or sand. They can tolerate temperatures as high as 120° Fahrenheit.

\DndMonsterAction{Metabolic Reserves}
The camel can go without water for up to 7 days, and without food for up to 90 days.

\DndMonsterAction{Stubborn}
The camel has advantage on saving throws made to avoid or end the charmed, frightened, or prone condition on themself.

\DndMonsterSection{Actions}
\DndMonsterAction{Kick}
\textsl{Melee Weapon Attack:} 5 to hit, reach 5 ft., one target. \textsl{Hit:} 8 (1d10 + 3) bludgeoning damage, and the target must succeed on a DC 13 Strength saving throw or be knocked prone.

\DndMonsterSection{Bonus Actions}
\DndMonsterAction{Bilious Spit (Recharge 5\textendash 6)}
The camel vomits bile onto a creature they can see within 10 feet of them. The target must succeed on a DC 14 Constitution saving throw or be poisoned for 1 minute (save ends at end of turn). If the target fails the save by 5 or more, they are also blinded while poisoned in this way.

\end{DndMonster}



\begin{DndMonster}[float=t]{Chimera}
\DndMonsterType{Large monstrosity (Skirmisher), N}
\DndMonsterBasics[
armor-class = {15 (natural armor)},
hit-points = {\DndDice{14d10 + 42}},
speed = {30 ft., fly 60 ft.}
]
\DndMonsterAbilityScores[
 str = 18,
 dex = 14,
 con = 16,
 int = 3,
 wis = 14,
 cha = 11
]
\DndMonsterDetails[
skills = {Athletics +7, Perception +5},
condition-immunities = {frightened},
senses = {darkvision 60 ft., passive perception 15},
languages = {understands any one language but can't speak},
challenge = 6,
]
\DndMonsterAction{Volant}
When the chimera reduces a creature to 0 hit points, the chimera can move up to their speed toward an enemy they can see.

\DndMonsterSection{Actions}
\DndMonsterAction{Multiattack}
The chimera makes three Bite attacks. They can replace one attack with a Lion's Toss attack or a use of Dragon's Eruption, if available.

\DndMonsterAction{Bite}
\textsl{Melee Weapon Attack:} 7 to hit, reach 5 ft., one target. \textsl{Hit:} 11 (2d6 + 4) piercing damage.

\DndMonsterAction{Lion's Toss}
\textsl{Melee Weapon Attack:} 7 to hit, reach 5 ft., one target. \textsl{Hit:} 9 (2d4 + 4) piercing damage, and the target must succeed on a DC 15 Strength saving throw or be moved up to 20 feet in any direction.

\DndMonsterAction{Dragon's Eruption (Recharge 6)}
The chimera spits a volcanic explosion at a point they can see within 120 feet of them. Each creature in a 10-foot-radius sphere centered on that point must make a DC 14 Dexterity saving throw, taking 27 (6d8) fire damage on a failed save, or half as much damage on a successful one.

\DndMonsterSection{Bonus Actions}
\DndMonsterAction{Roar}
Each enemy who can hear the chimera and is within 60 feet of them must succeed on a DC 13 Wisdom saving throw or be frightened of the chimera for 1 minute (save ends at end of turn). If a creature succeeds on a saving throw against this effect or if the effect ends for them, the creature is immune to the chimera's Roar for the next 24 hours.

\DndMonsterSection{Reactions}
\DndMonsterAction{Ram's Defiance}
When a creature within 30 feet of the chimera misses them with an attack, the chimera can move up to half their speed in a straight line toward the creature. If the chimera ends this movement within 5 feet of the creature, the creature must succeed on a DC 15 Strength saving throw or be knocked prone.

\end{DndMonster}



\begin{DndMonster}[float=t]{Chimeron}
\DndMonsterType{Huge fiend (demon, category 5, Brute), chaotic evil}
\DndMonsterBasics[
armor-class = {17 (natural armor)},
hit-points = {\DndDice{20d12 + 100}},
speed = {30 ft., fly 30 ft.}
]
\DndMonsterAbilityScores[
 str = 24,
 dex = 12,
 con = 20,
 int = 16,
 wis = 14,
 cha = 12
]
\DndMonsterDetails[
saving-throws = {Str +11, Dex +5, Cha +5},
skills = {Intimidation +9, Perception +6},
damage-resistances = {cold, fire, thunder},
senses = {darkvision 120 ft., soulsight 30 ft., passive perception 16},
languages = {Abyssal, Common, telepathy 120 ft.},
challenge = 12,
]
\DndMonsterAction{Lethe}
While the chimeron's soul count is 0, they have advantage on attack rolls, disadvantage on saving throws, and their Intelligence score becomes 3 (-4). Additionally, the chimeron must use their movement on each of their turns to move as close as possible to the nearest creature they can sense with their soulsight, and then if they are able, they must use their action to attack and attempt to kill that creature. The chimeron can't act with any other purpose until they add 1 to their soul count.

\DndMonsterAction{Soul Devourer}
When the chimeron reduces a creature who isn't a Construct or an Undead to 0 hit points or deals damage to a dying creature, the creature must make a DC 11 Wisdom saving throw. On a failed save, the chimeron consumes the creature's soul and adds 1 to the chimeron's soul count. A creature whose soul is consumed in this way immediately dies, and they can't be restored to life by any means short of a \textit{wish} spell.

\DndMonsterAction{Soul Resistance}
While the chimeron's soul count is is 2 or higher, they have advantage on saving throws against powers, spells, and other supernatural effects.

\DndMonsterSection{Actions}
\DndMonsterAction{Multiattack}
The chimeron makes up to three attacks using Powerful Bite, Scorching Bite, Booming Bite, or Chilling Bite. They can't use any attack more than once.

\DndMonsterAction{Powerful Bite}
\textsl{Melee Weapon Attack:} 11 to hit, reach 10 ft., one target. \textsl{Hit:} 25 (4d8 + 7) piercing damage.

\DndMonsterAction{Scorching Bite}
\textsl{Melee Weapon Attack:} 11 to hit, reach 10 ft., one target. \textsl{Hit:} 11 (1d8 + 7) piercing damage plus 9 (2d8) fire damage, and the target is set on fire until the target or a creature who can reach them uses an action to extinguish the flames. A creature who is on fire at the start of their turn takes 7 (2d6) fire damage. If a creature who is already on fire is set on fire again on a subsequent turn, the damage isn't cumulative.

\DndMonsterAction{Booming Bite (Costs 1 Soul)}
\textsl{Melee Weapon Attack:} 11 to hit, reach 10 ft., one target. \textsl{Hit:} 11 (1d8 + 7) piercing damage plus 9 (2d8) thunder damage, and the target is knocked prone.

\DndMonsterAction{Chilling Bite (Costs 1 Soul)}
\textsl{Melee Weapon Attack:} 11 to hit, reach 10 ft., one target. \textsl{Hit:} 11 (1d8 + 7) piercing damage plus 9 (2d8) cold damage, the target's speed is halved until the end of their next turn, and they can only make one attack during their next turn.

\DndMonsterSection{Reactions}
\DndMonsterAction{Pain Absorption (Costs 1 Soul)}
When the chimeron takes damage, they reduce the damage taken by 16 (3d10).

\end{DndMonster}



\begin{DndMonster}[float*=b,width=\textwidth + 8pt]{Count Rhodar von Glauer}
\begin{multicols}{2}
\DndMonsterType{Medium undead (Solo), neutral evil}
\DndMonsterBasics[
armor-class = {18 (natural armor)},
hit-points = {\DndDice{34d8 + 136}},
speed = {30 ft., fly 30 ft.}
]
\DndMonsterAbilityScores[
 str = 23,
 dex = 20,
 con = 18,
 int = 20,
 wis = 17,
 cha = 22
]
\DndMonsterDetails[
saving-throws = {Dex +11, Wis +9, Cha +12},
skills = {Deception +12, History +11, Intimidation +12, Perception +9, Persuasion +12, Stealth +11},
damage-resistances = {necrotic; bludgeoning, piercing, slashing from mundane attacks},
condition-immunities = {charmed, flanked, frightened, prone, stunned},
senses = {blindsight 60 ft., darkvision 120 ft., passive perception 19},
languages = {Common},
challenge = 19,
]
\DndMonsterAction{Resting Place}
When Rhodar drops to 0 hit points and isn't in his resting place, he teleports to his resting place. While in his resting place, Rhodar is stable. After spending 1 hour in his resting place with 0 hit points, he regains 1 hit point.

\DndMonsterAction{Spear Sacrifice (3/Day)}
Rhodar is surrounded by three floating spears of shimmering darkness, which he can summon to a free hand on his turn (no action required). When Rhodar fails an ability check or saving throw, he can choose to destroy one of these spears and succeed on the ability check or saving throw instead. Once Rhodar destroys all three spears, he can't use his Spear of the Damned attack or Spear Sacrifice until the next dusk, when he manifests any destroyed spears.

\DndMonsterAction{Spellcasting (Utility)}
In addition to any other spells in this stat block, Rhodar can cast the following spells, using Charisma as the spellcasting ability (spell save DC 20):
\begin{DndMonsterSpells}
  \DndInnateSpellLevel{charm person\textsuperscript{Casting Times}, \textit{detect thoughts}\textsuperscript{Casting Times}, \textit{disguise self}\textsuperscript{Casting Times}, \textit{sending}\textsuperscript{Casting Times}}
  \DndInnateSpellLevel[3]{clairvoyance\textsuperscript{Casting Times}, \textit{geas} (at 9th level) \textsuperscript{Casting Times}, \textit{legend lore}\textsuperscript{Casting Times}}
\end{DndMonsterSpells}
\DndMonsterSection{Actions}
\DndMonsterAction{Multiattack}
Rhodar makes two Sanguinus attacks, and he can make one Spear of the Damned attack for each Spear Sacrifice use he has remaining.

\DndMonsterAction{Sanguinus}
\textsl{Melee Weapon Attack:} 12 to hit, reach 5 ft., one target. \textsl{Hit:} 13 (2d6 + 6) slashing damage plus 14 (4d6) necrotic damage. If the target is a creature, their hit point maximum is reduced by a number equal to the necrotic damage taken, and Rhodar regains hit points equal to that number. This reduction lasts until reversed by a cure ailment power of 4th order or higher, a \textit{greater restoration} spell, or a similar supernatural effect.

The target dies if this effect reduces their hit point maximum to 0. A humanoid slain in this way rises the following midnight as a vampire spawn under Rhodar's control.

\DndMonsterAction{Spear of the Damned}
\textsl{Ranged Spell Attack:} 12 to hit, range 60 ft., one target. \textsl{Hit:} 10 (3d6) force damage, and if the target is a creature, they are knocked prone and restrained as a spear of darkness impales them. A creature can attempt to free themself or another creature within their reach by using an action or a bonus action to make a DC 20 Strength (Athletics) check. On a success, the creature is freed and the restrained condition ends for them. The creature is also freed if Rhodar summons that spear to attack another target or destroys it using Spear Sacrifice.

\DndMonsterSection{Bonus Actions}
\DndMonsterAction{Beguile}
Rhodar targets one creature he can see within 30 feet of him and who can see him. The target must make a DC 20 Wisdom saving throw. On a failed save, the creature uses their reaction, if available, to move up to their speed to a location chosen by Rhodar, then either makes a melee attack against a target of Rhodar's choice or falls prone (Rhodar's choice).

\DndMonsterAction{Inhuman Agility}
Rhodar moves up to his speed without provoking opportunity attacks.

\DndMonsterSection{Reactions}
\DndMonsterAction{Withering Stare}
When Rhodar is hit by an attack by a creature he can see within 60 feet of him, his eyes bore into the creature. The target must succeed on a DC 20 Wisdom saving throw or take 10 (3d6) necrotic damage.

\DndMonsterSection{Legendary Actions}
Rhodar has three villain actions. He can take each action once during an encounter after an enemy's turn. He can take these actions in any order but can use only one per round.

\begin{DndMonsterLegendaryActions}
\DndMonsterLegendaryAction{Action 1: Bloodstained}{Rhodar causes blood to explode in a 20-foot-radius sphere centered on a point he can see within 120 feet of him. Creatures who aren't Undead in that area must succeed on a DC 20 Dexterity saving throw or be bloodstained until the end of Rhodar's next turn. Rhodar has advantage on attack rolls against bloodstained creatures, and bloodstained creatures have disadvantage on saving throws against Rhodar's Beguile bonus action.
}
\DndMonsterLegendaryAction{Action 2: Fire Drake}{Rhodar, along with anything he is wearing or carrying, turns into a Large dragon made of fire. When he does, he releases any creature he was grappling, and if Rhodar was grappled or restrained, that effect ends for him. He then flies up to twice his speed to an unoccupied space. During this move, Rhodar ignores difficult terrain, doesn't provoke opportunity attacks, and can move through creatures and objects. Each creature Rhodar passes through must make a DC 20 Dexterity saving throw, taking 36 (8d8) fire damage on a failed save, or half as much damage on a successful one. At the end of this movement, Rhodar transforms back into his previous form.
}
\DndMonsterLegendaryAction{Action 3: Mist of Blades}{Rhodar fills the area within 60 feet of him with thick swirling mists, heavily obscuring that area for other creatures. He can teleport up to four times into any unoccupied space in the mist and make one Sanguinus attack after each teleport. The mist dissipates at the end of this action.
}
\end{DndMonsterLegendaryActions}
\end{multicols}
\end{DndMonster}



\begin{DndMonster}[float=t]{Crux of Frost}
\DndMonsterType{Medium elemental (water, Controller), A}
\DndMonsterBasics[
armor-class = {16 (natural armor)},
hit-points = {\DndDice{30d8}},
speed = {30 ft.}
]
\DndMonsterAbilityScores[
 str = 12,
 dex = 14,
 con = 10,
 int = 12,
 wis = 16,
 cha = 18
]
\DndMonsterDetails[
skills = {Deception +10, Insight +9, Perception +6},
damage-immunities = {cold, poison},
condition-immunities = {invisible, poisoned},
senses = {darkvision 60 ft., passive perception 16},
languages = {Aquan, Common, plus the languages known by a creature reflected by Frosted Reflection},
challenge = 8,
]
\DndMonsterAction{Scintillating}
The crux can't benefit from being invisible.

\DndMonsterSection{Actions}
\DndMonsterAction{Multiattack}
The crux makes two Rimeglass Touch attacks, and they can use Frosted Reflection or It Stares Back, if available.

\DndMonsterAction{Rimeglass Touch}
\textsl{Melee or Ranged Spell Attack:} 7 to hit, reach 5 ft. or range 60 ft., one target. \textsl{Hit:} 17 (3d8 + 4) cold damage, and the target can't see or hear creatures other than themself and the crux until the start of the crux's next turn.

\DndMonsterAction{Frosted Reflection}
The crux reflects the appearance of a specific Medium or Small creature they've seen within the last week. This appearance is illusory and imperfect, showing a frostbitten version of the creature. The illusion lasts until the crux takes thunder damage, uses this action again, or is incapacitated. For the duration, the crux additionally knows any languages that creature knows.

\DndMonsterAction{It Stares Back (Recharge 5\textendash 6)}
The crux imposes fragility on a creature within 60 feet of them who is reflected by the crux's Frosted Reflection. If the target can see the crux, the target must make a DC 15 Wisdom saving throw. On a failed save, the target is frightened of the crux for 1 minute (save ends at end of turn), or until the crux stops reflecting them. While frightened in this way, whenever the target takes damage, they take an extra 11 (2d10) psychic damage. On a successful save, the target isn't frightened, and they are immune to the It Stares Back of all cruxes for 24 hours.

\DndMonsterSection{Bonus Actions}
\DndMonsterAction{Convocation of Ice (1/Day)}
The crux imbues the power of ice in themself or another Elemental they can see within 30 feet of them, granting one of the following effects of that Elemental's choice:


\begin{description}
\item \textbf{Frigid Sheen.} The Elemental gains a mirror sheen that reflects damage for 1 minute. When the Elemental is hit with an attack by a creature they can see within 30 feet of them, the Elemental can use a reaction to reflect the attack. The Elemental is unaffected by the attack, and the attacker must make a DC 15 Dexterity saving throw. On a failed save, the attack's damage and effects are reflected back at the attacker as if the attack originated from the Elemental, turning the attacker into the target. On a successful save, the attacker takes half as much damage, but suffers no other effect.
\item \textbf{Frost Squall.} Frost swirls around the Elemental for 1 minute. For the duration, each non-Elemental creature who starts their turn within 15 feet of the Elemental has their speed reduced to 15 feet unless it is already lower until the end of their turn, and the first time they willingly move before the start of their next turn, they take 10 (3d6) cold damage.
\end{description}

\end{DndMonster}



\begin{DndMonster}[float=t]{Deep Dreamer}
\DndMonsterType{Large plant (fungus, Controller), N}
\DndMonsterBasics[
armor-class = {17 (natural armor)},
hit-points = {\DndDice{21d10 + 63}},
speed = {0 ft., fly 20 ft. \textit{(hover)}}
]
\DndMonsterAbilityScores[
 str = 17,
 dex = 10,
 con = 16,
 int = 20,
 wis = 12,
 cha = 16
]
\DndMonsterDetails[
saving-throws = {Int +10},
skills = {Deception +8, Insight +6, Persuasion +8},
damage-immunities = {psychic},
condition-immunities = {charmed, frightened, prone},
senses = {darkvision 120 ft., passive perception 11},
languages = {Deep Speech, Undercommon, telepathy 120 ft.},
challenge = 13,
]
\DndMonsterAction{Psionic Shroud}
The dreamer is immune to divination spells and to any effect that would sense their emotions or read their thoughts.

\DndMonsterSection{Actions}
\DndMonsterAction{Multiattack}
The dreamer makes two Tendrils attacks and uses Dream Spores, if available.

\DndMonsterAction{Tendrils}
\textsl{Melee Psionic Attack:} 10 to hit, reach 10 ft., one target. \textsl{Hit:} 12 (2d6 + 5) slashing damage plus 22 (4d10) psychic damage. If the target is charmed by the dreamer, the dreamer regains hit points equal to half the psychic damage dealt.

\DndMonsterAction{*Dream Spores (5th-Order Power; Recharge 5-6)}
The dreamer releases a cloud of psionic spores. Each creature within 60 feet of them must succeed on a DC 18 Intelligence saving throw or be charmed by the dreamer for 1 minute.

While charmed in this way, a creature must use their movement during their turn to move within 10 feet of the dreamer by the safest available route.

A creature charmed in this way can repeat the saving throw whenever they take damage, ending the effect on themself on a success.

\DndMonsterSection{Reactions}
\DndMonsterAction{Rouse Dreamer}
When a creature makes an attack against the dreamer, the dreamer chooses a creature charmed by the dreamer within the attack's range, and the charmed creature becomes the target of the attack instead.

\end{DndMonster}



\begin{DndMonster}[float=t]{Devil Adjudicator}
\DndMonsterType{Medium fiend (devil, Controller), lawful evil}
\DndMonsterBasics[
armor-class = {18 (natural armor)},
hit-points = {\DndDice{24d8 + 96}},
speed = {30 ft., fly 40 ft.}
]
\DndMonsterAbilityScores[
 str = 12,
 dex = 16,
 con = 18,
 int = 16,
 wis = 14,
 cha = 20
]
\DndMonsterDetails[
saving-throws = {Con +9, Wis +7, Cha +10},
skills = {Deception +10, Insight +7, Perception +7, Persuasion +10, Religion +8},
damage-resistances = {bludgeoning, piercing, slashing from mundane attacks},
damage-immunities = {fire},
condition-immunities = {charmed, frightened},
senses = {darkvision 120 ft., passive perception 17},
languages = {Common, Infernal},
challenge = 14,
]
\DndMonsterAction{True Name}
If a creature the adjudicator can hear within 60 feet of them speaks the adjudicator's true name aloud, the adjudicator loses their damage resistances, damage immunities, and Devilish Charm reaction for 24 hours.

\DndMonsterSection{Actions}
\DndMonsterAction{Multiattack}
The adjudicator makes two Infernal Injunction attacks and uses Adjudicator's Interdiction, if available.

\DndMonsterAction{Infernal Injunction}
\textsl{Melee or Ranged Spell Attack:} 10 to hit, reach 5 ft. or range 120 ft., one target. \textsl{Hit:} 27 (5d8 + 5) fire damage, and the target must make a DC 18 Wisdom saving throw. On a failed save, the adjudicator chooses for the target to become either charmed by or frightened of them until the end of the target's next turn.

\DndMonsterAction{Adjudicator's Interdiction (Recharge 5\textendash 6)}
The adjudicator places an infernal seal on one creature they can see within 120 feet of them. The target must succeed on a DC 18 Charisma saving throw or be interdicted until the adjudicator dies. While interdicted, a creature's speed is halved, they can't take reactions, and they can't regain hit points. A cure ailment power, a \textit{lesser restoration} or \textit{remove curse} spell, or a similar supernatural effect ends the effect early.

\DndMonsterAction{Bad Deal (1/Day)}
Three creatures the adjudicator can see within 60 feet of them must make a DC 18 Charisma saving throw. On a failed save, a target must choose to either take a -5 penalty to AC or a -5 penalty to ability checks and attack rolls. This penalty lasts for 10 minutes (save ends at end of turn).

\DndMonsterSection{Reactions}
\DndMonsterAction{Devilish Charm (2/Day)}
When the adjudicator is targeted by an attack, power, spell, or other supernatural effect by a creature they can see within 60 feet of them, the creature must make a DC 18 Charisma saving throw. On a failed save, the creature is charmed by the adjudicator until the start of the creature's next turn, and the adjudicator chooses a new target the adjudicator can see for the triggering effect. The new target must be within the triggering effect's range.

\end{DndMonster}


\clearpage

\begin{DndMonster}[float=t]{Devil Jurist}
\DndMonsterType{Medium fiend (devil, Artillery), lawful evil}
\DndMonsterBasics[
armor-class = {16 (studded leather armor)},
hit-points = {\DndDice{24d8 + 48}},
speed = {30 ft., fly 40 ft.}
]
\DndMonsterAbilityScores[
 str = 11,
 dex = 18,
 con = 15,
 int = 16,
 wis = 14,
 cha = 19
]
\DndMonsterDetails[
saving-throws = {Int +7, Wis +6, Cha +8},
skills = {Arcana +7, Deception +8, Perception +6, Persuasion +8},
damage-resistances = {bludgeoning, piercing, slashing from mundane attacks},
damage-immunities = {fire},
condition-immunities = {charmed},
senses = {darkvision 120 ft., passive perception 16},
languages = {Common, Infernal},
challenge = 10,
]
\DndMonsterAction{Hellfire}
Fire damage dealt by the jurist ignores damage resistance.

\DndMonsterAction{True Name}
If a creature the jurist can hear within 60 feet of them speaks the jurist's true name aloud, the jurist loses their damage resistances, damage immunities, and Devilish Charm reaction for 24 hours.

\DndMonsterSection{Actions}
\DndMonsterAction{Multiattack}
The jurist makes two Fire and Brim stone attacks.

\DndMonsterAction{Fire and Brimstone}
\textsl{Melee or Ranged Spell Attack:} 8 to hit, reach 5 ft. or range 120 ft., one target. \textsl{Hit:} 17 (3d8 + 4) fire damage, and the target must succeed on a DC 16 Constitution saving throw or gain one festering wound. A creature takes 4 (1d8) fire damage at the start of each of their turns for each festering wound they have. A creature who receives supernatural healing loses all festering wounds.

\DndMonsterAction{Inferno (Recharge 5\textendash 6)}
The jurist makes one Fire and Brimstone attack against each enemy they can see within 60 feet of them.

\DndMonsterSection{Bonus Actions}
\DndMonsterAction{Ashes to Ashes}
The jurist chooses a creature they can see within 120 feet of them who has one or more festering wounds. One of the target's wounds ignites and they take 9 (2d8) fire damage.

\DndMonsterSection{Reactions}
\DndMonsterAction{Devilish Charm (1/Day)}
When the jurist is targeted by an attack, power, spell, or other supernatural effect by a creature they can see within 60 feet of them, the creature must make a DC 16 Charisma saving throw. On a failed save, the creature is charmed by the jurist until the start of the creature's next turn, and the jurist chooses a new target the jurist can see for the triggering effect. The new target must be within the triggering effect's range.

\end{DndMonster}



\begin{DndMonster}[float=t]{Devil Legate}
\DndMonsterType{Medium fiend (devil, Soldier), lawful evil}
\DndMonsterBasics[
armor-class = {17 (half plate armor)},
hit-points = {\DndDice{15d8 + 60}},
speed = {30 ft.}
]
\DndMonsterAbilityScores[
 str = 18,
 dex = 15,
 con = 18,
 int = 11,
 wis = 14,
 cha = 18
]
\DndMonsterDetails[
saving-throws = {Str +7, Wis +5, Cha +7},
skills = {Athletics +7, Deception +7, Persuasion +7},
damage-resistances = {bludgeoning, piercing, slashing from mundane attacks},
damage-immunities = {fire},
condition-immunities = {charmed},
senses = {darkvision 120 ft., passive perception 12},
languages = {Common, Infernal},
challenge = 8,
]
\DndMonsterAction{Hellish Resistance}
While in Hell or within 60 feet of a devil with a challenge rating of 9 or higher, the legate has advantage on saving throws against powers, spells, and other supernatural effects.

\DndMonsterAction{True Name}
If a creature the legate can hear within 60 feet of them speaks the legate's true name aloud, the legate loses their damage resistances, damage immunities, and Devilish Charm reaction for 24 hours.

\DndMonsterSection{Actions}
\DndMonsterAction{Multiattack}
The legate makes two Infernal Pike attacks.

\DndMonsterAction{Infernal Pike}
\textsl{Melee Weapon Attack:} 7 to hit, reach 10 ft., one target. \textsl{Hit:} 9 (1d10 + 4) piercing damage plus 11 (2d10) fire damage, and the legate chooses one of the following effects:


\begin{description}
\item[Focused Hate.]
The target has disadvantage on attack rolls made against creatures other than the legate until the end of the legate's next turn.

\item[Hellbite.]
The target takes an extra 5 (2d4) fire damage.

\item[Skewer.]
The legate deals 5 (1d10) piercing damage plus 5 (1d10) fire damage to another creature within 5 feet of the target.

\end{description}

\DndMonsterSection{Reactions}
\DndMonsterAction{Devilish Charm (1/Day)}
When the legate is targeted by an attack, power, spell, or other supernatural effect by a creature they can see within 60 feet of them, the creature must make a DC 15 Charisma saving throw. On a failed save, the creature is charmed by the legate until the start of the creature's next turn, and the legate chooses a new target the legate can see for the triggering effect. The new target must be within the triggering effect's range.

\end{DndMonster}



\begin{DndMonster}[float=t]{Devil Magistrate}
\DndMonsterType{Medium fiend (devil, Skirmisher), lawful evil}
\DndMonsterBasics[
armor-class = {18 (studded leather armor)},
hit-points = {\DndDice{28d8 + 56}},
speed = {50 ft.}
]
\DndMonsterAbilityScores[
 str = 12,
 dex = 22,
 con = 14,
 int = 11,
 wis = 16,
 cha = 19
]
\DndMonsterDetails[
saving-throws = {Dex +10, Wis +7, Cha +8},
skills = {Acrobatics +10, Deception +8, Insight +7, Persuasion +8, Stealth +10},
damage-resistances = {bludgeoning, piercing, slashing from mundane attacks},
damage-immunities = {fire},
condition-immunities = {charmed},
senses = {darkvision 120 ft., passive perception 13},
languages = {Common, Infernal},
challenge = 12,
]
\DndMonsterAction{True Name}
If a creature the magistrate can hear within 60 feet of them speaks the magistrate's true name aloud, the magistrate loses their damage resistances, damage immunities, and Devilish Charm reaction for 24 hours.

\DndMonsterSection{Actions}
\DndMonsterAction{Multiattack}
The magistrate makes two Infernal Knife attacks and one Obsidian Kris attack.

\DndMonsterAction{Infernal Knife}
\textsl{Melee or Ranged Weapon Attack:} 10 to hit, reach 5 ft. or range 20/60 ft., one target. \textsl{Hit:} 15 (2d8 + 6) piercing damage plus 9 (2d8) fire damage, and the target is frightened of their allies until the end of their next turn.

\DndMonsterAction{Obsidian Kris}
\textsl{Melee Weapon Attack:} 10 to hit, reach 5 ft., one target. \textsl{Hit:} 15 (2d8 + 6) piercing damage, and if the target is frightened, they can't take reactions until the frightened condition ends for them.

\DndMonsterSection{Reactions}
\DndMonsterAction{Devilish Charm (2/Day)}
When the magistrate is targeted by an attack, power, spell, or other supernatural effect by a creature they can see within 60 feet of them, the creature must make a DC 16 Charisma saving throw. On a failed save, the creature is charmed by the magistrate until the start of the creature's next turn, and the magistrate chooses a new target the magistrate can see for the triggering effect. The new target must be within the triggering effect's range.

\end{DndMonster}



\begin{DndMonster}[float=t]{Devil Notary}
\DndMonsterType{Medium fiend (devil, Minion), lawful evil}
\DndMonsterBasics[
armor-class = {15 (\textit{chain shirt})},
hit-points = {\DndDice{13}},
speed = {30 ft.}
]
\DndMonsterAbilityScores[
 str = 11,
 dex = 15,
 con = 14,
 int = 12,
 wis = 14,
 cha = 16
]
\DndMonsterDetails[
damage-immunities = {fire},
senses = {darkvision 120 ft., passive perception 12},
languages = {Common, Infernal},
challenge = 6,
]
\DndMonsterAction{Importunity}
At the start of the notary's turn, they can grant an infernal blessing to a non-minion ally they can see within 60 feet of them (no action required). The next time that ally makes an attack roll or saving throw before the start of the notary's next turn, the ally can add a +1 bonus to the roll. This bonus increases by 1 (maximum bonus of +5) for each notary who grants an infernal blessing to the creature.

\DndMonsterAction{Minion}
If the notary takes damage from an attack or as the result of a failed saving throw, their hit points are reduced to 0. If the notary takes damage from another effect, they die if the damage equals or exceeds their hit point maximum; otherwise they take no damage.

\DndMonsterAction{True Name}
If a creature the notary can hear within 60 feet of them speaks the notary's true name aloud, the notary loses their damage immunities and Importunity trait for 24 hours.

\DndMonsterSection{Actions}
\DndMonsterAction{Brimstone (Group Attack)}
\textsl{Melee or Ranged Spell Attack:} 6 to hit, reach 5 ft. or range 30 ft., one target. \textsl{Hit:} 4 fire damage.

\end{DndMonster}



\begin{DndMonster}[float=t]{Devil Retainer}
\DndMonsterType{Medium fiend (devil, Retainer), A}
\DndMonsterBasics[
armor-class = {15 (medium armor)},
hit-points = {\DndDice{Eight times their level (number of d10 Hit Dice equal to their level)}},
speed = {30 ft.}
]
\DndMonsterAbilityScores[
 str = 10,
 dex = 16,
 con = 10,
 int = 10,
 wis = 10,
 cha = 14
]
\DndMonsterDetails[
saving-throws = {Str +PB, Dex +PB, Con +PB, Int +PB, Wis +PB, Cha +PB},
skills = {Deception +2 plus PB, Persuasion +2 plus PB},
damage-resistances = {fire},
condition-immunities = {charmed},
senses = {darkvision 120 ft., passive perception 10},
languages = {Common, Infernal},
]
\DndMonsterSection{Actions}
\DndMonsterAction{Signature Attack (Daggers)}
\textsl{Melee or Ranged Weapon Attack:} 3 plus PB to hit, reach 5 ft. or range 20/60 ft. \textsl{Hit:} 2d4 + PB piercing damage. Beginning at 7th level, the retainer can make this attack twice, instead of once, when they take the Attack action on their turn.

\DndMonsterAction{3rd Level: Misleading Strike (3/Day)}
As an action, the retainer makes a signature attack. On a hit, the target must use their reaction, if available, to make a weapon attack against a creature of the retainer's choice.

\DndMonsterAction{5th Level: Adjuration (3/Day)}
As a bonus action, the retainer grants an infernal blessing to an ally they can see within 60 feet of them. The next time that ally makes an attack roll before the start of the retainer's next turn, the ally adds a +PB bonus to the roll. If the attack hits, it deals an extra 0d1 + PB damage.

\DndMonsterAction{7th Level: Hellbitten Knives (1/Day)}
As a bonus action, the retainer imbues their daggers with hellfire. For 1 minute, their signature attack deals an extra 0d1 + PB fire damage and ignores damage resistances.

\end{DndMonster}



\begin{DndMonster}[float=t]{Dirt Mournling}
\DndMonsterType{Huge undead (Controller), U}
\DndMonsterBasics[
armor-class = {19 (natural armor)},
hit-points = {\DndDice{22d12 + 110}},
speed = {40 ft., burrow 40 ft., climb 40 ft.}
]
\DndMonsterAbilityScores[
 str = 22,
 dex = 17,
 con = 20,
 int = 10,
 wis = 12,
 cha = 10
]
\DndMonsterDetails[
saving-throws = {Str +12, Con +11, Cha +6},
skills = {Athletics +12, Perception +7},
damage-immunities = {cold, fire, poison},
condition-immunities = {charmed, exhaustion, frightened, poisoned},
senses = {darkvision 120 ft., passive perception 17},
languages = {understands the languages of their creator but can't speak},
challenge = 18,
]
\DndMonsterAction{Anguish Unleashed}
When the mournling dies, the sadness they contain is released into the world. Each creature within 30 feet of the mournling who can hear them must make a DC 19 Wisdom saving throw, taking 22 (4d10) psychic damage on a failed save, or half as much damage on a successful one.

\DndMonsterAction{Immutable Form}
The mournling is immune to any power, spell, or effect that would alter their form.

\DndMonsterAction{Puddle of Mud}
When the mournling leaves a space, the ground in that space becomes difficult terrain for their enemies for 10 minutes.

\DndMonsterSection{Actions}
\DndMonsterAction{Multiattack}
The mournling makes three attacks using Slam, Mudslide, or both, and they can use Dust Storm, if available.

\DndMonsterAction{Slam}
\textsl{Melee Weapon Attack:} 12 to hit, reach 10 ft., one target. \textsl{Hit:} 28 (4d10 + 6) bludgeoning damage, and the target is grappled (escape DC 20). While grappled in this way, the target is restrained.

\DndMonsterAction{Mudslide}
\textsl{Ranged Weapon Attack:} 12 to hit, range 30/60 ft., one target. \textsl{Hit:} 22 (3d10 + 6) bludgeoning damage. If the target is a Huge or smaller creature, the target and each Huge or smaller creature of the mournling's choice within 5 feet of the target are pulled up to 30 feet toward the mournling.

\DndMonsterAction{Dust Storm (1/Day)}
The mournling hurls dirt, gravel, and sand. Each creature within 60 feet of them must succeed on a DC 19 Constitution saving throw or be blinded for 1 minute (save ends at end of turn).

\DndMonsterSection{Bonus Actions}
\DndMonsterAction{Break and Reform}
The mournling teleports to an unoccupied space they can see within 30 feet of them.

\DndMonsterSection{Reactions}
\DndMonsterAction{Stuck Fast}
When the mournling is hit by a melee weapon attack, they capture the weapon that struck them in quickhardening mud, disarming the attacker. A creature who can reach the weapon can use an action or bonus action to make a DC 19 Strength (Athletics) check, freeing the weapon on a success. The mournling can't use this reaction on attacks made with natural weapons.

\end{DndMonster}



\begin{DndMonster}[float*=b,width=\textwidth + 8pt]{Dohma Raskovar}
\begin{multicols}{2}
\DndMonsterType{Medium humanoid (orc, Leader), lawful evil}
\DndMonsterBasics[
armor-class = {14 (\textit{chain shirt})},
hit-points = {\DndDice{11d8 + 44}},
speed = {35 ft.}
]
\DndMonsterAbilityScores[
 str = 17,
 dex = 12,
 con = 18,
 int = 16,
 wis = 15,
 cha = 16
]
\DndMonsterDetails[
saving-throws = {Str +6, Con +7, Wis +5},
skills = {Athletics +6, Deception +6, Insight +5, Intimidation +6, Perception +5},
senses = {darkvision 60 ft., passive perception 15},
languages = {Common, Orc},
challenge = 5,
]
\DndMonsterAction{Adrenaline Rush}
If Raskovar fails a saving throw, he can expend a use of his Dohma Rally without gaining its benefits (no action required) and succeed instead.

\DndMonsterAction{Relentless (1/Turn)}
When Raskovar isn't incapacitated and he is reduced to 0 hit points but not killed outright, he can make an attack against an enemy (no action required) before the hit point reduction is resolved. If the attack hits and its damage reduces the target to 0 hit points, Raskovar drops to 1 hit point instead of 0 hit points.

\DndMonsterSection{Actions}
\DndMonsterAction{Multiattack}
Raskovar makes three Repeating Crossbow attacks or two War Mace attacks.

\DndMonsterAction{War Mace}
\textsl{Melee Weapon Attack:} 6 to hit, reach 5 ft., one target. \textsl{Hit:} 8 (2d4 + 3) bludgeoning damage, and the target is dazed until the end of their next turn.

\DndMonsterAction{Repeating Crossbow}
\textsl{Ranged Weapon Attack:} 4 to hit, range 30/120 ft., one target. \textsl{Hit:} 7 (1d12 + 1) piercing damage, and the target's speed is reduced by 10 feet until the end of their next turn.

\DndMonsterAction{Reinforcements (1/Day)}
Raskovar takes the Dodge action, and ten orc blitzers appear in unoccupied spaces within 30 feet of him.

\DndMonsterSection{Bonus Actions}
\DndMonsterAction{Dohma Rally (3/Day)}
Raskovar regains 20 hit points, and he gains resistance to bludgeoning, piercing, and slashing damage until the end of his next turn.

\DndMonsterSection{Reactions}
\DndMonsterAction{On the House}
When a creature Raskovar can see within 30 feet of him fails an ability check or misses with an attack roll, that creature can reroll the check or attack, choosing either result.

\DndMonsterSection{Legendary Actions}
Raskovar has three villain actions. He can take each action once during an encounter after an enemy's turn. He can take these actions in any order but can use only one per round.

\begin{DndMonsterLegendaryActions}
\DndMonsterLegendaryAction{Action 1: Close In}{Each ally within 60 feet of Raskovar who can hear him can move up to their speed. After this movement, each enemy within 5 feet of at least one creature who moved must make a DC 15 Wisdom saving throw or be frightened of Raskovar for 1 minute (save ends at end of turn).
}
\DndMonsterLegendaryAction{Action 2: Lockdown}{Each ally within 60 feet of Raskovar who can hear him can make a Strength (Athletics) check with advantage to grapple an enemy within their reach (no action required). Then Raskovar can move up to his speed.
}
\DndMonsterLegendaryAction{Action 3: Do It Myself}{Raskovar moves up to his speed without provoking opportunity attacks. At the end of the movement, he can make up to five War Mace attacks with advantage.
}
\end{DndMonsterLegendaryActions}
\end{multicols}
\end{DndMonster}



\begin{DndMonster}[float*=b,width=\textwidth + 8pt]{Durixaviinox}
\begin{multicols}{2}
\DndMonsterType{Gargantuan dragon (Solo), lawful evil}
\DndMonsterBasics[
armor-class = {18 (natural armor)},
hit-points = {\DndDice{20d20 + 160}},
speed = {40 ft., burrow 40 ft., fly 80 ft., swim 40 ft.}
]
\DndMonsterAbilityScores[
 str = 26,
 dex = 10,
 con = 26,
 int = 18,
 wis = 14,
 cha = 18
]
\DndMonsterDetails[
saving-throws = {Con +14, Wis +8, Cha +10},
skills = {Intimidation +10, Perception +8, Persuasion +10},
damage-immunities = {cold, force},
condition-immunities = {charmed, dazed, frightened, paralyzed, petrified, stunned},
senses = {blindsight 120 ft., truesight 60 ft., passive perception 18},
languages = {Common, Draconic},
challenge = 20,
]
\DndMonsterAction{Frosted Resistance (3/Day)}
When Durixaviinox fails a saving throw, he can succeed instead. When he does, his speed is halved and he can't take the Disengage action until the end of his next turn.

\DndMonsterAction{Hoarfrost}
Cold damage dealt by Durixaviinox ignores damage resistance.

\DndMonsterAction{Ice Walk}
Durixaviinox can move across and climb icy surfaces without needing to make an ability check. Additionally, difficult terrain composed of ice or snow doesn't cost him extra movement.

\DndMonsterSection{Actions}
\DndMonsterAction{Multiattack}
Durixaviinox makes one Bite attack and two Claw attacks.

\DndMonsterAction{Bite}
\textsl{Melee Weapon Attack:} 14 to hit, reach 15 ft., one target. \textsl{Hit:} 19 (2d10 + 8) piercing damage plus 9 (2d8) cold damage, and the target can't take reactions until the end of their next turn.

\DndMonsterAction{Claw}
\textsl{Melee Weapon Attack:} 14 to hit, reach 10 ft., one target. \textsl{Hit:} 15 (2d6 + 8) slashing damage, and Durixaviinox can move the target up to 15 feet horizontally.

\DndMonsterAction{Numbing Breath (Recharge 5\textendash 6)}
Durixaviinox exhales a freezing blast in a 90-foot cone. Each creature in that area must make a DC 22 Constitution saving throw. On a failed save, a creature takes 42 (12d6) cold damage and is restrained until the end of their next turn. On a successful save, a creature takes half as much damage and isn't restrained.

\DndMonsterSection{Bonus Actions}
\DndMonsterAction{Icy Grip}
Durixaviinox encases a creature he can see within 60 feet of him in ice. The target must make a DC 18 Strength saving throw. On a failed save, a creature takes 22 (4d10) cold damage and is restrained. On a successful save, a target takes half as much damage and isn't restrained. A creature can use an action to break themself or another creature they can reach out of the ice, ending the restrained condition.

\DndMonsterSection{Reactions}
\DndMonsterAction{Withering Frost}
When a creature within 5 feet of Durixaviinox hits him with an attack, Durixaviinox releases a burst of cold energy. Each creature within 5 feet of him must succeed on a DC 18 Constitution saving throw or take 10 (3d6) cold damage.

\DndMonsterSection{Legendary Actions}
Durixaviinox has three villain actions. He can take each action once during an encounter after an enemy's turn. He can take these actions in any order but can use only one per round.

\begin{DndMonsterLegendaryActions}
\DndMonsterLegendaryAction{Action 1: Raging Blizzard}{Durixaviinox beats his wings, and frigid winds and hail swirl around him. Each creature within 120 feet of him must make a DC 18 Strength saving throw. On a failed save, a creature takes 21 (6d6) bludgeoning damage, is pushed 60 feet directly away from Durixaviinox, and is knocked prone. On a successful save, a creature takes half as much damage and isn't pushed or knocked prone.
}
\DndMonsterLegendaryAction{Action 2: Royal Defense}{Durixaviinox calls the faithful soul of a spectral guard to fight at his side. The guard uses the frost giant wind sprinter stat block, but they are an Undead instead of a Giant. The guard acts immediately after Durixaviinox and follows his verbal commands. While the guard is within 60 feet of Durixaviinox, Durixaviinox gains resistance to all damage. If Durixaviinox dies, the guard is destroyed.
}
\DndMonsterLegendaryAction{Action 3: Frost Cataclysm}{Durixaviinox releases a pulse of frigid energy followed by a shock wave of force in a 120-foot-radius sphere centered on him. Each enemy in that area must make a DC 18 Constitution saving throw, taking 45 (10d8) cold damage on a failed save, or half as much damage on a successful one. Immediately after, each enemy in that area must make a DC 18 Strength saving throw, taking 45 (10d8) force damage on a failed save, or half as much damage on a successful one.
}
\end{DndMonsterLegendaryActions}
\end{multicols}
\end{DndMonster}



\begin{DndMonster}[float=t]{Earth Spark}
\DndMonsterType{Large elemental (earth, Minion), A}
\DndMonsterBasics[
armor-class = {18 (natural armor)},
hit-points = {\DndDice{17}},
speed = {30 ft., burrow 20 ft.}
]
\DndMonsterAbilityScores[
 str = 16,
 dex = 8,
 con = 16,
 int = 9,
 wis = 10,
 cha = 12
]
\DndMonsterDetails[
damage-immunities = {poison},
condition-immunities = {poisoned},
senses = {darkvision 60 ft., passive perception 10},
languages = {Common, Terran},
challenge = 10,
]
\DndMonsterAction{Minion}
If the spark takes damage from an attack or as the result of a failed saving throw, their hit points are reduced to 0. If the spark takes damage from another effect, they die if the damage equals or exceeds their hit point maximum; otherwise they take no damage.

\DndMonsterAction{Trampling Rampage}
If an enemy who is touching the ground starts their turn within 5 feet of two or more earth sparks, the enemy must make a Strength saving throw with a DC equal to 12 plus the number of earth sparks within 5 feet of the enemy. On a failed save, the enemy falls prone, and they can't stand up until there are fewer than two earth sparks within 5 feet of them.

\DndMonsterSection{Actions}
\DndMonsterAction{Shatter Smash (Group Attack)}
\textsl{Melee Weapon Attack:} 7 to hit, reach 5 ft., one target. \textsl{Hit:} 5 bludgeoning damage.

\DndMonsterAction{Snort Rocks (Group Attack)}
\textsl{Ranged Weapon Attack:} 7 to hit, range 30/60 ft., one target. \textsl{Hit:} 5 bludgeoning damage. If the target is a flying creature, they must succeed on a Strength saving throw or fall prone. The DC for this saving throw equals 12 plus the number of sparks who joined the attack.

\end{DndMonster}



\begin{DndMonster}[float=t]{Empyrean Stag}
\DndMonsterType{Large celestial (Support), neutral good}
\DndMonsterBasics[
armor-class = {20 (natural armor)},
hit-points = {\DndDice{22d10 + 66}},
speed = {60 ft.}
]
\DndMonsterAbilityScores[
 str = 24,
 dex = 16,
 con = 16,
 int = 18,
 wis = 20,
 cha = 20
]
\DndMonsterDetails[
saving-throws = {Dex +8, Wis +10, Cha +10},
skills = {Athletics +12, Insight +10, Perception +10},
damage-resistances = {necrotic; bludgeoning, piercing, slashing from mundane attacks},
damage-immunities = {radiant},
condition-immunities = {charmed, exhaustion, frightened},
senses = {darkvision 120 ft., passive perception 20},
languages = {all, telepathy 120 ft.},
challenge = 13,
]
\DndMonsterAction{Supernatural Resistance}
The stag has advantage on saving throws against powers, spells, and other supernatural effects.

\DndMonsterSection{Actions}
\DndMonsterAction{Multiattack}
The stag makes three Ram or Lightburst attacks.

\DndMonsterAction{Ram}
\textsl{Melee Weapon Attack:} 12 to hit, reach 10 ft., one target. \textsl{Hit:} 25 (4d8 + 7) piercing damage.

\DndMonsterAction{Lightburst}
\textsl{Ranged Spell Attack:} 10 to hit, range 120 ft., one target. \textsl{Hit:} 18 (3d8 + 5) radiant damage, and the target is outlined in a golden light, giving attack rolls against them advantage until the start of the stag's next turn.

\DndMonsterAction{Starcall (1/Day)}
The stag calls down a magical shower of stars on an area that is 60 feet long and 5 feet wide within 120 feet of the stag. Each enemy in that area must make a DC 18 Dexterity saving throw, taking 28 (8d6) radiant damage on a failed save, or half as much damage on a successful one. Each other willing creature in that area regains 15 hit points.

The stars then create a 10-foot-high wall of opaque radiant light in that area. Each creature in the wall's space is pushed out of it by the shortest route. The creature chooses which side of the wall to end up on. The wall lasts for 1 minute or until the stag is incapacitated or dies. Any enemy who enters the wall of light for the first time on a turn must make a DC 18 Constitution saving throw, taking 14 (4d6) radiant damage on a failed save, or half as much damage on a successful one.

\DndMonsterSection{Bonus Actions}
\DndMonsterAction{Starlight Infusion}
The stag infuses radiant energy into a willing creature they can see within 60 feet of them. That creature deals an extra 5 radiant damage with weapon attacks until the start of the stag's next turn.

\DndMonsterSection{Reactions}
\DndMonsterAction{Nature's Ward}
When the stag or a creature they can see within 30 feet of them takes damage, the stag magically creates a protective barrier around that creature. The barrier reduces that damage to the protected creature by 20, then vanishes.

\end{DndMonster}



\begin{DndMonster}[float=t]{Essence of Mist}
\DndMonsterType{Medium elemental (air, water, Ambusher), A}
\DndMonsterBasics[
armor-class = {15},
hit-points = {\DndDice{22d8 + 66}},
speed = {40 ft., fly 20 ft. \textit{(hover)}}
]
\DndMonsterAbilityScores[
 str = 14,
 dex = 20,
 con = 17,
 int = 12,
 wis = 12,
 cha = 16
]
\DndMonsterDetails[
saving-throws = {Dex +10, Cha +8},
skills = {Deception +8, Sleight Of Hand +10, Stealth +10},
damage-resistances = {psychic},
condition-immunities = {blinded, charmed, frightened, grappled, prone, restrained},
senses = {blindsight 60 ft., darkvision 120 ft., passive perception 11},
languages = {Aquan, Auran, Common},
challenge = 13,
]
\DndMonsterAction{Flowing Form}
The essence can enter an enemy's space and stop there, and they can move through a space as narrow as 1 inch wide without squeezing.

\DndMonsterAction{Gasping Aura}
While within 10 feet of the essence, creatures who need to breathe can't take reactions or speak.

\DndMonsterSection{Actions}
\DndMonsterAction{Multiattack}
The essence makes two Breathless Rend attacks.

\DndMonsterAction{Breathless Rend}
\textsl{Melee Weapon Attack:} 10 to hit, reach 5 ft., one target. \textsl{Hit:} 14 (2d8 + 5) slashing damage plus 22 (4d10) psychic damage.

\DndMonsterSection{Bonus Actions}
\DndMonsterAction{Convocation of Breath (2/Day)}
The essence imbues the power of breath in themself or another Elemental they can see within 30 feet of them, granting one of the following effects of that Elemental's choice:


\begin{description}
\item \textbf{Hide in Breath.} The Elemental tries to hide inside the body of a creature within 10 feet of them who isn't a Construct, an Elemental, an Ooze, or an Undead. The target must succeed on a DC 18 Constitution saving throw or the Elemental retreats inside the target until the start of the Elemental's next turn. While inside the target, the Elemental can't take actions and has total cover against attacks and other effects outside the target.
\item \textbf{Toxic Breath.} The Elemental exudes a noxious aura 10 feet in every direction for 1 minute or until dispersed by a moderate or stronger wind. That area is lightly obscured, and each creature who needs to breathe who starts their turn in that area must succeed on a DC 16 Constitution saving throw or be poisoned until the end of their next turn.
\end{description}

\DndMonsterAction{Breathstealer's Shimmer (1/Day)}
The essence becomes invisible then takes the Hide action. This invisibility ends if the essence makes an attack or uses Convocation of Breath.

\end{DndMonster}



\begin{DndMonster}[float=t]{Essence of Storms}
\DndMonsterType{Medium elemental (air, Controller), A}
\DndMonsterBasics[
armor-class = {14 (natural armor)},
hit-points = {\DndDice{12d8 + 24}},
speed = {20 ft., fly 60 ft. \textit{(hover)}}
]
\DndMonsterAbilityScores[
 str = 12,
 dex = 17,
 con = 14,
 int = 9,
 wis = 11,
 cha = 13
]
\DndMonsterDetails[
skills = {Acrobatics +6, Perception +6},
damage-resistances = {lightning, poison, thunder},
condition-immunities = {dazed, grappled, paralyzed, poisoned, restrained},
senses = {darkvision 60 ft., passive perception 16},
languages = {Auran, Common},
challenge = 5,
]
\DndMonsterSection{Actions}
\DndMonsterAction{Multiattack}
The essence makes two Wind Talons attacks. They can replace one attack with a use of Bluster or Lightning Squall, if available.

\DndMonsterAction{Wind Talons}
\textsl{Melee Weapon Attack:} 6 to hit, reach 5 ft., one target. \textsl{Hit:} 10 (2d6 + 3) piercing damage.

\DndMonsterAction{Bluster}
The essence unleashes a 30-foot cone of wind. Each creature in that area must succeed on a DC 14 Strength saving throw or be moved up to 10 feet in any direction.

\DndMonsterAction{Lightning Squall (3/Day)}
The essence hurls a small lightning storm at a creature they can see within 30 feet of them. The target must make a DC 14 Constitution saving throw. On a failed save, the target takes 28 (8d6) lightning damage and is dazed for 1 minute (save ends at end of turn). On a successful save, they take half as much damage and aren't dazed.

\DndMonsterSection{Bonus Actions}
\DndMonsterAction{Convocation of Air (1/Day)}
The essence imbues the power of air in themself or another Elemental they can see within 30 feet of them, granting one of the following effects of that Elemental's choice:


\begin{description}
\item \textbf{Rising Whirlwind.} Each creature of the Elemental's choice within 30 feet of them must make a DC 14 Strength saving throw against a violent whirlwind. On a failed save, the target rises 10 feet upward and is restrained for 1 minute (save ends at end of turn). A target restrained in this way rises an additional 10 feet at the start of each of their turns. Another creature who can reach the target can use an action to pull them free of the effect, ending the restrained condition for that creature. The whirlwind ends if the Elemental dies or chooses to end it (no action required).
\item \textbf{Vortex Terrain.} A vortex swirls around the Elemental. For 1 minute, the area within 20 feet of the Elemental is difficult terrain for enemies.
\end{description}

\end{DndMonster}



\begin{DndMonster}[float=t]{Essence of Tides}
\DndMonsterType{Medium elemental (water, Skirmisher), A}
\DndMonsterBasics[
armor-class = {15 (natural armor)},
hit-points = {\DndDice{13d8 + 26}},
speed = {15 ft., swim 50 ft.}
]
\DndMonsterAbilityScores[
 str = 12,
 dex = 18,
 con = 14,
 int = 9,
 wis = 16,
 cha = 10
]
\DndMonsterDetails[
skills = {Acrobatics +7, Medicine +6},
damage-resistances = {acid, fire, poison},
condition-immunities = {grappled, poisoned, restrained},
senses = {darkvision 60 ft., passive perception 13},
languages = {Aquan, Common},
challenge = 5,
]
\DndMonsterAction{Water Glide}
While within 5 feet of a liquid or solid surface, the essence gains a flying speed equal to their swimming speed. While flying, the essence doesn't provoke opportunity attacks.

\DndMonsterSection{Actions}
\DndMonsterAction{Multiattack}
The essence makes two Water Wing attacks and one Lightning Tail Whip attack.

\DndMonsterAction{Water Wing}
\textsl{Melee Weapon Attack:} 7 to hit, reach 5 ft., one target. \textsl{Hit:} 7 (1d6 + 4) bludgeoning damage. If the same creature is hit twice with this attack on a turn, they are vulnerable to lightning damage until the end of the essence's turn.

\DndMonsterAction{Lightning Tail Whip}
\textsl{Melee Spell Attack:} 6 to hit, reach 15 ft., one target. \textsl{Hit:} 7 (2d6) lightning damage, and the target is pushed up to 10 feet away from the essence.

\DndMonsterSection{Bonus Actions}
\DndMonsterAction{Convocation of Water (1/Day)}
The essence imbues the power of water in themself or another Elemental they can see within 30 feet, granting one of the following effects of that Elemental's choice:


\begin{description}
\item \textbf{Scalding Steam.} The Elemental spouts blistering steam. Each non-Elemental creature within 15 feet of them must make a DC 15 Constitution saving throw. On a failed save, a target takes 18 (4d8) fire damage and is blinded until the end of their next turn. On a successful save, a target takes half as much damage and isn't blinded.
\item \textbf{Waters of Vitality.} The Elemental regains 25 hit points and can end one of the following conditions affecting them: blinded, charmed, dazed, deafened, paralyzed, poisoned, or stunned.
\end{description}

\end{DndMonster}



\begin{DndMonster}[float=t]{Fire Giant Lightbearer}
\DndMonsterType{Huge giant (Support), A}
\DndMonsterBasics[
armor-class = {18 (natural armor)},
hit-points = {\DndDice{11d12 + 66}},
speed = {40 ft.}
]
\DndMonsterAbilityScores[
 str = 22,
 dex = 16,
 con = 23,
 int = 10,
 wis = 21,
 cha = 13
]
\DndMonsterDetails[
saving-throws = {Str +10, Dex +7, Con +10},
skills = {Athletics +10, Perception +9},
damage-immunities = {fire},
languages = {Giant},
challenge = 10,
]
\DndMonsterAction{Healing Heat}
When the lightbearer targets a fire giant with an attack, spell, or other supernatural effect that deals fire damage, that target instead regains a number of hit points equal to the damage that would be dealt. The target can choose to be hit by this attack or fail their saving throw against this effect.

\DndMonsterAction{Molten Flesh}
The first time a creature other than a fire giant touches the lightbearer or hits them with a melee attack on a turn, that creature takes 7 (2d6) fire damage.

\DndMonsterSection{Actions}
\DndMonsterAction{Multiattack}
The lightbearer makes two attacks using Slam, Living Blaze, or both. They can replace one attack with a use of Flamelash.

\DndMonsterAction{Slam}
\textsl{Melee Weapon Attack:} 10 to hit, reach 10 ft., one target. \textsl{Hit:} 16 (3d6 + 6) bludgeoning damage plus 7 (2d6) fire damage, and if the target is a Huge or smaller creature, the target must choose between being knocked prone or pushed 10 feet away from the lightbearer.

\DndMonsterAction{Living Blaze}
\textsl{Ranged Spell Attack:} 10 to hit, range 180 ft., one target. \textsl{Hit:} 16 (3d6 + 6) fire damage, and the lightbearer can ricochet the flame toward another target within 5 feet of the first. The lightbearer makes another spell attack roll against the second target with disadvantage, dealing the same damage on a hit.

\DndMonsterAction{Flamelash}
The lightbearer extends a whiplike flame to strike a creature they can see within 30 feet of them. The target must make a DC 18 Dexterity saving throw. On a failed save, the target takes 20 (4d6 + 6) fire damage, and if the target is Large or smaller, the lightbearer moves the target up to 10 feet horizontally.

\DndMonsterSection{Bonus Actions}
\DndMonsterAction{Travel by Fire (3/Day)}
The lightbearer chooses two willing creatures they can see within 30 feet of them. Each target takes 14 (4d6) fire damage and teleports, swapping places with the other target.

\end{DndMonster}



\begin{DndMonster}[float=t]{Fire Giant Red Fist}
\DndMonsterType{Huge giant (Soldier), A}
\DndMonsterBasics[
armor-class = {17 (natural armor)},
hit-points = {\DndDice{13d12 + 78}},
speed = {40 ft.}
]
\DndMonsterAbilityScores[
 str = 25,
 dex = 14,
 con = 23,
 int = 10,
 wis = 14,
 cha = 13
]
\DndMonsterDetails[
saving-throws = {Str +11, Dex +6, Con +10},
skills = {Athletics +11, Perception +6},
damage-immunities = {fire},
languages = {Giant},
challenge = 9,
]
\DndMonsterAction{Heat and Pressure}
When an enemy moves out of the red fist's reach, the enemy must succeed on a DC 18 Constitution saving throw or gain a level of exhaustion. Creatures with resistance or immunity to fire damage are immune to this effect.

\DndMonsterAction{Molten Flesh}
The first time a creature other than a fire giant touches the red fist or hits them with a melee attack on a turn, that creature takes 7 (2d6) fire damage.

\DndMonsterSection{Actions}
\DndMonsterAction{Multiattack}
The red fist makes two Unarmed Strike attacks.

\DndMonsterAction{Unarmed Strike}
\textsl{Melee Weapon Attack:} 11 to hit, reach 10 ft., one target. \textsl{Hit:} 21 (4d6 + 7) bludgeoning damage plus 7 (2d6) fire damage. If the same creature is hit twice with this attack on a turn, they are blinded until the end of their next turn.

\DndMonsterAction{Hurl Flame}
\textsl{Ranged Spell Attack:} 10 to hit, range 180 ft., one target. \textsl{Hit:} 34 (8d6 + 6) fire damage.

\DndMonsterSection{Reactions}
\DndMonsterAction{Guardian Block (3/Day)}
When an enemy within 10 feet of the red fist hits a creature other than the red fist with an attack, the red fist halves the damage dealt by that attack. The red fist then makes an Unarmed Strike against the attacker.

\end{DndMonster}



\begin{DndMonster}[float=t]{Fire Giant Trooper}
\DndMonsterType{Huge giant (Minion), A}
\DndMonsterBasics[
armor-class = {18 (natural armor)},
hit-points = {\DndDice{19}},
speed = {40 ft.}
]
\DndMonsterAbilityScores[
 str = 22,
 dex = 14,
 con = 19,
 int = 10,
 wis = 14,
 cha = 13
]
\DndMonsterDetails[
damage-immunities = {fire},
languages = {Giant},
challenge = 12,
]
\DndMonsterAction{Minion}
If the trooper takes damage from an attack or as the result of a failed saving throw, their hit points are reduced to 0. If the trooper takes damage from another effect, they die if the damage equals or exceeds their hit point maximum; otherwise they take no damage.

\DndMonsterAction{Molten Flesh}
The first time a creature other than a fire giant touches the trooper or hits them with a melee attack on a turn, that creature takes 7 fire damage.

\DndMonsterSection{Actions}
\DndMonsterAction{Unarmed Strike (Group Attack)}
\textsl{Melee Weapon Attack:} 10 to hit, reach 10 ft., one target. \textsl{Hit:} 3 bludgeoning damage plus 3 fire damage, and the target must succeed on a Strength saving throw or fall prone. The DC for this saving throw equals 12 + the number of troopers who joined the group attack.

\DndMonsterAction{Rock (Group Attack)}
\textsl{Ranged Weapon Attack:} 10 to hit, range 60/240 ft., one target. \textsl{Hit:} 6 bludgeoning damage.

\end{DndMonster}



\begin{DndMonster}[float=t]{Force of Blood}
\DndMonsterType{Large elemental (earth, fire, water, Brute), A}
\DndMonsterBasics[
armor-class = {21 (natural armor)},
hit-points = {\DndDice{20d10 + 120}},
speed = {50 ft.}
]
\DndMonsterAbilityScores[
 str = 24,
 dex = 20,
 con = 22,
 int = 19,
 wis = 18,
 cha = 20
]
\DndMonsterDetails[
saving-throws = {Con +12, Wis +10, Cha +11},
skills = {Intimidation +11, Medicine +10, Persuasion +11},
damage-resistances = {fire, necrotic},
condition-immunities = {charmed, exhaustion, frightened},
senses = {darkvision 120 ft., passive perception 14},
languages = {Common, Dwarvish, Elvish, Giant, Primordial},
challenge = 19,
]
\DndMonsterAction{Visceral Sight}
The force of blood can see any creature within 120 feet of them who isn't a Construct or an Undead. This sight ignores effects that obscure sight.

\DndMonsterSection{Actions}
\DndMonsterAction{Multiattack}
The force of blood makes three Scarlet Longspear attacks. They can replace one attack with a Rend from Within attack.

\DndMonsterAction{Scarlet Longspear}
\textsl{Melee or Ranged Weapon Attack:} 13 to hit, reach 15 ft. or range 30/60 ft., one target. \textsl{Hit:} 18 (2d10 + 7) piercing damage. If the target is a creature, the force of blood regains 5 hit points.

\DndMonsterAction{Rend from Within}
\textsl{Melee Spell Attack:} 12 to hit, reach 5 ft., one creature. \textsl{Hit:} 13 (2d6 + 6) slashing damage plus 26 (4d12) necrotic damage, and the target spends 1 Hit Die with no effect and can't take reactions until the start of their next turn. Miss: 13 (2d12) necrotic damage.

\DndMonsterSection{Bonus Actions}
\DndMonsterAction{Convocation of Blood}
The force of blood imbues the power of blood in themself or another Elemental they can see within 30 feet of them, granting one of the following effects of that Elemental's choice:


\begin{description}
\item \textbf{Boiling Blood.} The Elemental boils the blood of a creature within 5 feet of them who isn't a Construct or an Undead. The target must make a DC 19 Constitution saving throw. On a failed save, the target must choose to either take 27 (5d10) fire damage or make an attack with a weapon that has a damage die of 1d4 or greater against a creature of the Elemental's choice. If the attack hits, it deals an extra 11 (2d10) fire damage.
\item \textbf{Essence Thief.} The Elemental steals the essence of two creatures the Elemental can see within 30 feet of them. Each target must make a DC 19 Constitution saving throw, taking 19 (3d12) necrotic damage on a failed save, or half as much damage on a successful one. The Elemental can then teleport to an unoccupied space they can see within 5 feet of one of the targets.
\end{description}

\end{DndMonster}



\begin{DndMonster}[float=t]{Force of Earth}
\DndMonsterType{Large elemental (earth, Soldier), A}
\DndMonsterBasics[
armor-class = {18 (natural armor)},
hit-points = {\DndDice{10d10 + 40}},
speed = {30 ft., burrow 20 ft.}
]
\DndMonsterAbilityScores[
 str = 19,
 dex = 9,
 con = 18,
 int = 10,
 wis = 14,
 cha = 9
]
\DndMonsterDetails[
skills = {Athletics +7, History +6},
damage-immunities = {poison},
condition-immunities = {petrified, poisoned},
senses = {darkvision 60 ft., tremorsense 30 ft., passive perception 12},
languages = {Common, Terran},
challenge = 5,
]
\DndMonsterAction{Earth Glide}
The force of earth can burrow through mundane, unworked earth and stone. While doing so, the force of earth doesn't disturb the material they move through.

\DndMonsterSection{Actions}
\DndMonsterAction{Multiattack}
The force of earth makes two Earthbind Strike attacks.

\DndMonsterAction{Earthbind Strike}
\textsl{Melee Weapon Attack:} 7 to hit, reach 5 ft., one target. \textsl{Hit:} 13 (2d8 + 4) bludgeoning damage, and the target's speed becomes 0 until the start of the force of earth's next turn.

\DndMonsterAction{Shatter Rock}
\textsl{Ranged Weapon Attack:} 7 to hit, range 60/120 ft., one target. \textsl{Hit:} 23 (3d12 + 4) bludgeoning damage, or 47 (6d12 + 8) bludgeoning damage if the target is an object or structure.

\DndMonsterSection{Bonus Actions}
\DndMonsterAction{Convocation of Earth (1/Day)}
The force of earth imbues the power of earth in themself or another Elemental they can see within 30 feet of them, granting one of the following effects of that Elemental's choice:


\begin{description}
\item \textbf{Earthen Aura.} Stones rapidly whir around the Elemental in a 10-foot-radius sphere for 1 minute or until they're killed or incapacitated. A creature who isn't an Elemental takes 5 slashing damage for every 5 feet they move into or within that area.
\item \textbf{Stone Armor.} Stone grows over the Elemental in a protective carapace, and they gain 20 temporary hit points. While the Elemental has these temporary hit points, they are immune to the dazed, paralyzed, petrified, and stunned conditions.
\end{description}

\end{DndMonster}



\begin{DndMonster}[float=t]{Force of Iron}
\DndMonsterType{Large elemental (earth, fire, Brute), A}
\DndMonsterBasics[
armor-class = {19 (natural armor) (22 in Shield Form)},
hit-points = {\DndDice{12d10 + 84}},
speed = {60 ft.}
]
\DndMonsterAbilityScores[
 str = 21,
 dex = 10,
 con = 25,
 int = 11,
 wis = 13,
 cha = 12
]
\DndMonsterDetails[
saving-throws = {Str +9},
skills = {Athletics +9, Perception +5},
damage-resistances = {fire},
damage-immunities = {poison},
condition-immunities = {exhaustion, frightened, petrified, poisoned},
senses = {darkvision 60 ft., tremorsense 30 ft., passive perception 15},
languages = {Common, Ignan, Terran},
challenge = 10,
]
\DndMonsterSection{Actions}
\DndMonsterAction{Multiattack}
The force of iron uses Form Stance then makes three Ferrous Blade attacks or three Iron Spike attacks.

\DndMonsterAction{Ferrous Greatblade}
\textsl{Melee Weapon Attack:} 9 to hit, reach 5 ft. (20 ft. in Polearm Form), one target. \textsl{Hit:} 19 (4d6 + 5) slashing damage.

\DndMonsterAction{Iron Spike}
\textsl{Ranged Weapon Attack:} 9 to hit, range 30/60 ft., one target. \textsl{Hit:} 14 (2d8 + 5) piercing damage.

\DndMonsterAction{Form Stance}
The force of iron shapes their body to gain one of the following benefits until the start of their next turn:

\DndMonsterAction{Polearm Form}
The force of iron's melee attacks have a reach of 20 feet instead of 5 feet.

\DndMonsterAction{Reckless Form}
The force of iron has advantage on all melee attack rolls, but attack rolls against the force of iron have advantage.

\DndMonsterAction{Shield Form}
The force of iron gains a +3 bonus to AC.

\DndMonsterAction{Fire Form (1/Day)}
The force of iron's attacks deal an extra 13 (2d12) fire damage on a hit.

\DndMonsterSection{Bonus Actions}
\DndMonsterAction{Convocation of Iron (1/Day)}
The force of iron imbues the power of iron in themself or another Elemental they can see within 30 feet of them, granting one of the following effects of that Elemental's choice:


\begin{description}
\item \textbf{Earth Reinforcements.} Five \textbf{earth spark}s appear in unoccupied spaces within 20 feet of the Elemental. The sparks understand the Elemental, follow their verbal commands, and act immediately after the Elemental in the initiative order.
\item \textbf{Iron Skin.} Metal covers the Elemental, granting them resistance to bludgeoning, piercing, and slashing damage for 1 minute or until their concentration is broken (as if concentration on a spell).
\end{description}

\end{DndMonster}



\begin{DndMonster}[float=t]{Frost Giant Wind Sprinter}
\DndMonsterType{Huge giant (Skirmisher), A}
\DndMonsterBasics[
armor-class = {15 (issenblau plating)},
hit-points = {\DndDice{11d12 + 44}},
speed = {60 ft., climb 30 ft.}
]
\DndMonsterAbilityScores[
 str = 23,
 dex = 9,
 con = 18,
 int = 9,
 wis = 10,
 cha = 10
]
\DndMonsterDetails[
saving-throws = {Str +9, Con +7},
skills = {Athletics +9, Intimidation +3, Perception +3},
damage-immunities = {cold},
languages = {Giant},
challenge = 8,
]
\DndMonsterAction{Crush Underfoot}
When the wind sprinter enters a Medium or smaller enemy's space, the enemy must choose to either fall prone or take 10 (3d6) bludgeoning damage. Each enemy can be affected by this trait only once per turn.

\DndMonsterAction{Oncoming Storm}
Wisps of stinging snow and icy winds surround the wind sprinter. When an enemy concentration on a power, a spell, or a similar effect starts their turn within 15 feet of the wind sprinter, the enemy must succeed on a DC 15 Constitution saving throw or lose concentration.

\DndMonsterSection{Actions}
\DndMonsterAction{Multiattack}
The wind sprinter makes two Ice Axe attacks.

\DndMonsterAction{Ice Axe}
\textsl{Melee Weapon Attack:} 9 to hit, reach 10 ft., one target. \textsl{Hit:} 22 (3d10 + 6) slashing damage.

\DndMonsterAction{Rock}
\textsl{Ranged Weapon Attack:} 9 to hit, range 60/240 ft., one target. \textsl{Hit:} 28 (4d10 + 6) bludgeoning damage.

\DndMonsterAction{Blizzard Surge (Recharge 5\textendash 6)}
The wind sprinter moves up to their speed in a straight line. At any point during this move, they can make up to three Ice Axe attacks against different creatures.

\DndMonsterSection{Reactions}
\DndMonsterAction{Begone, Smallfolk!}
After a Medium or smaller creature makes an opportunity attack against the wind sprinter, the wind sprinter can make an Ice Axe attack against that creature.

\end{DndMonster}


\clearpage

\begin{DndMonster}[float=t]{Gem Jelly}
\DndMonsterType{Large ooze (Brute), U}
\DndMonsterBasics[
armor-class = {10 (20 in crystalline form)},
hit-points = {\DndDice{16d10 + 64}},
speed = {25 ft., climb 25 ft. \textit{(0 in crystalline form)}}
]
\DndMonsterAbilityScores[
 str = 18,
 dex = 10,
 con = 18,
 int = 6,
 wis = 14,
 cha = 4
]
\DndMonsterDetails[
skills = {Stealth +4},
damage-resistances = {acid; bludgeoning, piercing, slashing from mundane attacks},
condition-immunities = {blinded, charmed, deafened, exhaustion, frightened, prone},
senses = {blindsight 120 ft., passive perception 12},
challenge = 9,
]
\DndMonsterAction{Amorphous}
The gem jelly can move through a space as narrow as 1 inch wide without squeezing.

\DndMonsterAction{False Appearance}
While the jelly remains motionless, they are indistinguishable from a natural crystal formation.

\DndMonsterSection{Actions}
\DndMonsterAction{Multiattack}
The gem jelly makes two Pseudopod attacks.

\DndMonsterAction{Pseudopod (True Form Only)}
\textsl{Melee Weapon Attack:} 8 to hit, reach 10 ft., one target. \textsl{Hit:} 18 (4d6 + 4) piercing damage.

\DndMonsterAction{Crystalline Shatter (Crystalline Form Only)}
The jelly shatters their hardened outer layer, returning to their true form and sending crystal shrapnel in all directions. Each creature within 10 feet of the gem jelly must make a DC 16 Dexterity saving throw, taking 36 (8d8) piercing damage on a failed save, or half as much damage on a successful one. Gem jellies are immune to this damage.

\DndMonsterSection{Bonus Actions}
\DndMonsterAction{Harden}
The gem jelly transforms into a psionically hardened crystalline form or back into their true form. While in crystalline form, the jelly's AC increases to 20, their speed becomes 0, they can't make Pseudopod attacks, they have disadvantage on Dexterity saving throws, and they gain vulnerability to thunder damage.

\DndMonsterSection{Reactions}
\DndMonsterAction{Reactive Shatter (Crystalline Form Only)}
When the jelly takes thunder damage, they use Crystalline Shatter.

\end{DndMonster}



\begin{DndMonster}[float=t]{Ghost Sycophant}
\DndMonsterType{Medium undead (undead, Retainer), A}
\DndMonsterBasics[
armor-class = {13 (light armor)},
hit-points = {\DndDice{Six times their level (number of d6 Hit Dice equal to their level)}},
speed = {30 ft., fly 30 ft. \textit{(hover)}}
]
\DndMonsterAbilityScores[
 str = 10,
 dex = 12,
 con = 10,
 int = 10,
 wis = 12,
 cha = 16
]
\DndMonsterDetails[
saving-throws = {Str +PB, Dex +PB, Con +PB, Int +PB, Wis +PB, Cha +PB},
skills = {Intimidation +3 plus PB, Stealth +1 plus PB},
damage-resistances = {acid, cold, fire, lightning, thunder},
damage-immunities = {necrotic},
senses = {darkvision 60 ft., passive perception 11},
languages = {Common, any two languages they knew in life},
]
\DndMonsterSection{Actions}
\DndMonsterAction{Signature Attack (Dark Essence)}
\textsl{Melee or Ranged Spell Attack:} 3 plus PB to hit, reach 5 ft. or range 60 ft., one target. \textsl{Hit:} 5 (1d10) necrotic damage. This spell's damage increases by 1d10 when the sycophant reaches 5th level (2d10), 11th level (3d10), and 17th level (4d10).

\DndMonsterAction{3rd Level: Call of the Dead (3/Day)}
When the sycophant hits a creature with a signature attack, the target must also succeed on a DC 10 plus PB Wisdom saving throw or be frightened of the sycophant for 1 minute (save ends at end of turn).

\DndMonsterAction{5th Level: Astral Presence (3/Day)}
As an action, the sycophant becomes invisible for 1 minute or until they use a bonus action to end this effect. While invisible, the sycophant can move through creatures and objects as if they were difficult terrain, and the sycophant's signature attacks deal an extra 0d1 + PB damage to creatures who can't see them. The sycophant takes 5 (1d10) force damage if they end their turn inside an object.

\DndMonsterAction{7th Level: Whispers of the Dead (3/Day)}
As an action, the sycophant chooses up to PB creatures they can see within 30 feet of them. Each target must succeed on a DC 10 plus PB Charisma saving throw or make an attack against their nearest ally (no action required). A frightened creature has disadvantage on the saving throw.

\end{DndMonster}



\begin{DndMonster}[float=t]{Giant Zombie}
\DndMonsterType{Huge undead (Brute), neutral evil}
\DndMonsterBasics[
armor-class = {14 (natural armor)},
hit-points = {\DndDice{8d12 + 32}},
speed = {30 ft.}
]
\DndMonsterAbilityScores[
 str = 19,
 dex = 8,
 con = 18,
 int = 3,
 wis = 5,
 cha = 4
]
\DndMonsterDetails[
damage-vulnerabilities = {radiant},
damage-immunities = {poison},
condition-immunities = {exhaustion, poisoned},
senses = {darkvision 60 ft., passive perception 7},
languages = {understands the languages they knew in life but can't speak},
challenge = 4,
]
\DndMonsterAction{Negative Nerves}
Whenever a creature deals 7 damage or less to the zombie, the zombie takes no damage instead.

\DndMonsterSection{Actions}
\DndMonsterAction{Multiattack}
The zombie makes two Fist attacks.

\DndMonsterAction{Fist}
\textsl{Melee Weapon Attack:} 6 to hit, reach 10 ft., one target. \textsl{Hit:} 11 (2d6 + 4) bludgeoning damage, and if the target is Large or smaller, they are grappled (escape DC 14). The zombie has two fists, each of which can grapple one target.

\DndMonsterAction{Toxic Exhalation (Recharge 5\textendash 6)}
The zombie spews toxic vomit in a 15-foot cone. Each creature in that area must make a DC 14 Constitution saving throw, taking 14 (4d6) poison damage on a failed save, or half as much damage on a successful one.

\DndMonsterSection{Reactions}
\DndMonsterAction{Shove Away}
When a Huge or smaller creature within 10 feet of the zombie hits them with a melee attack, the zombie pushes the creature up to 10 feet away.

\end{DndMonster}



\begin{DndMonster}[float=t]{Gibbering Mouther}
\DndMonsterType{Medium aberration (Controller), chaotic neutral}
\DndMonsterBasics[
armor-class = {11 (natural armor)},
hit-points = {\DndDice{8d8 + 24}},
speed = {10 ft., fly 10 ft. \textit{(hover)}, swim 10 ft.}
]
\DndMonsterAbilityScores[
 str = 12,
 dex = 12,
 con = 16,
 int = 7,
 wis = 10,
 cha = 5
]
\DndMonsterDetails[
saving-throws = {Con +5},
condition-immunities = {charmed, frightened, prone},
senses = {darkvision 60 ft., passive perception 10},
languages = {speaks all languages but doesn't understand any},
challenge = 2,
]
\DndMonsterAction{Immutable Form}
The gibbering mouther is immune to any power, spell, or effect that would alter their form.

\DndMonsterAction{Primordial Influence}
A dazed enemy who starts their turn within 20 feet of the gibbering mouther must make a DC 13 Wisdom saving throw. On a failed save, that creature's body is altered in an otherworldly way of the GM's choice—they might flicker in and out of reality, sprout miniature fingers from their fingers, or have their physical form altered in a similarly alien way. This alteration doesn't affect the creature's game statistics. A cure ailment power of 4th order or higher, a \textit{greater restoration} spell, or a similar supernatural effect reverses this alteration.

While altered in this way, a creature must repeat the saving throw whenever they finish a long rest. On a failed save, their previous alteration worsens or they experience another chaotic alteration of the GM's choice. If a creature fails this saving throw three times after their initial alteration, they transform into a gibbering mouther controlled by the GM, and only a \textit{wish} spell can restore the creature to their original form.

\DndMonsterAction{Viscous Vicinity}
The gibbering mouther envelops their surroundings in their shifting reality. The area within 20 feet of them is difficult terrain for other creatures.

\DndMonsterSection{Actions}
\DndMonsterAction{Multiattack}
The gibbering mouther makes two Reality Rend attacks. They can replace one attack with a use of Pull.

\DndMonsterAction{Reality Rend}
\textsl{Melee Psionic Attack:} 5 to hit, reach 5 ft., one target. \textsl{Hit:} 10 (2d6 + 3) psychic damage, and the target must succeed on a DC 13 Wisdom saving throw or be dazed until the start of the gibbering mouther's next turn.

\DndMonsterAction{Pull}
The gibbering mouther warps reality around up to three creatures they can see within 60 feet of them. Each target must succeed on a DC 13 Strength saving throw or be pulled up to 30 feet directly toward the gibbering mouther.

\end{DndMonster}



\begin{DndMonster}[float=t]{Gnoll Bonesplitter}
\DndMonsterType{Medium fiend (gnoll, Brute), chaotic evil}
\DndMonsterBasics[
armor-class = {13 (\textit{hide armor})},
hit-points = {\DndDice{12d8 + 36}},
speed = {30 ft.}
]
\DndMonsterAbilityScores[
 str = 18,
 dex = 12,
 con = 17,
 int = 10,
 wis = 12,
 cha = 12
]
\DndMonsterDetails[
saving-throws = {Con +5},
senses = {darkvision 60 ft., passive perception 11},
languages = {Abyssal, Gnoll},
challenge = 4,
]
\DndMonsterSection{Actions}
\DndMonsterAction{Multiattack}
The bonesplitter makes two attacks using Bite, Spiked Flail, or both. They can replace one attack with a use of Bloody Roar, if available.

\DndMonsterAction{Bite}
\textsl{Melee Weapon Attack:} 6 to hit, reach 5 ft., one creature. \textsl{Hit:} 7 (1d6 + 4) piercing damage, or 11 (2d6 + 4) piercing damage if the target is restrained.

\DndMonsterAction{Spiked Flail}
\textsl{Melee Weapon Attack:} 6 to hit, reach 5 ft., one target. \textsl{Hit:} 13 (2d8 + 4) piercing damage. If the target is a Medium or smaller creature, they are grappled (escape DC 14). Until this grapple ends, the target is restrained and the bonesplitter can't make Spiked Flail attacks.

\DndMonsterAction{Bloody Roar (Recharge 6)}
The bonesplitter roars, spitting a spray of blood. Each ally within 10 feet of them who can see them can make a weapon attack (no action required), provided that ally isn't incapacitated.

\DndMonsterSection{Bonus Actions}
\DndMonsterAction{Over Here!}
The bonesplitter yanks a creature they are grappling. The target must succeed on a DC 14 Strength saving throw or swap spaces with the bonesplitter.

\DndMonsterSection{Reactions}
\DndMonsterAction{Death Frenzy}
When an ally the bonesplitter can see within 30 feet of them is reduced to 0 hit points, the bonesplitter moves up to half their speed and makes a Bite attack.

\end{DndMonster}



\begin{DndMonster}[float=t]{Goblin Assassin}
\DndMonsterType{Small humanoid (goblin, Ambusher), A}
\DndMonsterBasics[
armor-class = {15 (studded leather armor)},
hit-points = {\DndDice{3d6 + 6}},
speed = {30 ft., climb 20 ft.}
]
\DndMonsterAbilityScores[
 str = 8,
 dex = 16,
 con = 14,
 int = 10,
 wis = 10,
 cha = 8
]
\DndMonsterDetails[
skills = {Stealth +7},
senses = {darkvision 60 ft., passive perception 10},
languages = {Common, Goblin},
challenge = 1/2,
]
\DndMonsterAction{Backstab}
When the assassin has advantage on their attack roll against a creature who isn't a Construct or an Undead, their attacks deal an extra 3 (1d6) damage and inflict a bleeding wound on the target that lasts until the bleeding creature regains at least 1 hit point. A bleeding creature loses 2 hit points for each bleeding wound they have at the start of their turn. Any creature who can reach the target can use an action to stanch all the target's wounds, ending the effect.

\DndMonsterAction{Crafty}
The assassin doesn't provoke opportunity attacks when they move out of an enemy's reach.

\DndMonsterSection{Actions}
\DndMonsterAction{Scimitar}
\textsl{Melee Weapon Attack:} 5 to hit, reach 5 ft., one target. \textsl{Hit:} 6 (1d6 + 3) slashing damage.

\DndMonsterAction{Dagger}
\textsl{Melee or Ranged Weapon Attack:} 5 to hit, reach 5 ft. or range 20/60 ft., one target. \textsl{Hit:} 5 (1d4 + 3) piercing damage.

\DndMonsterAction{Summon Shadows (1/Day)}
A 10-foot-radius sphere of magical darkness emanates from a point the assassin can see for 1 minute. The darkness spreads around corners. Except for the assassin, a creature with darkvision can't see through this darkness, and mundane light can't illuminate it. At the start of their turn, the assassin can move the darkness up to 30 feet to a point they can see (no action required). If the assassin takes damage, the effect ends.

\DndMonsterSection{Bonus Actions}
\DndMonsterAction{Sneak}
The assassin takes the Hide action.

\end{DndMonster}



\begin{DndMonster}[float=t]{Goblin Cursespitter}
\DndMonsterType{Small humanoid (goblin, Controller), A}
\DndMonsterBasics[
armor-class = {15 (\textit{leather armor})},
hit-points = {\DndDice{5d6 + 10}},
speed = {30 ft., climb 20 ft.}
]
\DndMonsterAbilityScores[
 str = 8,
 dex = 14,
 con = 14,
 int = 10,
 wis = 10,
 cha = 15
]
\DndMonsterDetails[
saving-throws = {Wis +2},
skills = {Stealth +4},
senses = {darkvision 60 ft., passive perception 10},
languages = {Common, Goblin},
challenge = 1,
]
\DndMonsterAction{Crafty}
The cursespitter doesn't provoke opportunity attacks when they move out of an enemy's reach.

\DndMonsterSection{Actions}
\DndMonsterAction{Toxic Touch (Cantrip)}
\textsl{Melee or Ranged Spell Attack:} 4 to hit, reach 5 ft. or range 30 ft., one target. \textsl{Hit:} 7 (2d6) poison damage, and the target must succeed on a DC 12 Constitution saving throw or be poisoned for 1 minute (save ends at end of turn).

\DndMonsterAction{Brittle Bone Hex (Cantrip)}
The cursespitter chooses one creature they can see within 60 feet of them. The target's bones are wracked with pain until the end of their next turn. The first time the target willingly moves or uses an action, a bonus action, or a reaction before then, they must succeed on a DC 12 Constitution saving throw or take 9 (2d8) necrotic damage.

\DndMonsterAction{To Me!}
The cursespitter chooses up to two willing creatures they can see within 30 feet of them. Each creature is teleported to an unoccupied space within 5 feet of the cursespitter.

\DndMonsterAction{Dizzying Hex (2/Day; 1st-Level Spell)}
The cursespitter chooses one creature they can see within 60 feet of them. The target must make a DC 12 Wisdom saving throw. On a failed save, the target falls prone and can't stand back up for 1 minute (save ends at end of turn).

\DndMonsterSection{Reactions}
\DndMonsterAction{Cowardly Commander}
When a creature the cursespitter can see hits them with an attack, the cursespitter chooses a willing ally within 5 feet of them. The attack hits the ally instead.

\end{DndMonster}



\begin{DndMonster}[float=t]{Goblin Lackey}
\DndMonsterType{Small humanoid (goblin, Minion), A}
\DndMonsterBasics[
armor-class = {14 (\textit{leather armor})},
hit-points = {\DndDice{6}},
speed = {30 ft., climb 20 ft.}
]
\DndMonsterAbilityScores[
 str = 8,
 dex = 13,
 con = 10,
 int = 10,
 wis = 8,
 cha = 8
]
\DndMonsterDetails[
senses = {darkvision 60 ft., passive perception 9},
languages = {Common, Goblin},
challenge = 1/4,
]
\DndMonsterAction{Crafty}
The lackey doesn't provoke opportunity attacks when they move out of an enemy's reach.

\DndMonsterAction{Minion}
If the lackey takes damage from an attack or as the result of a failed saving throw, their hit points are reduced to 0. If the lackey takes damage from another effect, they die if the damage equals or exceeds their hit point maximum; otherwise they take no damage.

\DndMonsterAction{Tiny Stabs}
If an enemy starts their turn within 5 feet of three or more lackeys who can see them, the enemy must succeed on a Dexterity saving throw or take 1 piercing damage for each lackey within 5 feet. The DC for this saving throw equals 10 + the number of lackeys within 5 feet of the enemy.

\DndMonsterSection{Actions}
\DndMonsterAction{Dagger (Group Attack)}
\textsl{Melee or Ranged Weapon Attack:} 3 to hit, range 20/60 ft., one target. \textsl{Hit:} 1 piercing damage.

\end{DndMonster}



\begin{DndMonster}[float=t]{Goblin Sneak}
\DndMonsterType{Small humanoid (goblin, Retainer), A}
\DndMonsterBasics[
armor-class = {15 (medium armor)},
hit-points = {\DndDice{Seven times their level (number of d8 Hit Dice equal to their level)}},
speed = {30 ft., climb 20 ft.}
]
\DndMonsterAbilityScores[
 str = 10,
 dex = 16,
 con = 10,
 int = 12,
 wis = 10,
 cha = 10
]
\DndMonsterDetails[
saving-throws = {Str +PB, Dex +PB, Con +PB, Int +PB, Wis +PB, Cha +PB},
skills = {Acrobatics +3 plus PB, Stealth +3 plus PB},
senses = {darkvision 60 ft., passive perception 10},
languages = {Common, Goblin},
]
\DndMonsterSection{Actions}
\DndMonsterAction{Signature Attack (Daggers)}
\textsl{Melee or Ranged Weapon Attack:} 3 plus PB to hit, reach 5 ft. or range 20/60 ft., one target. \textsl{Hit:} 2d4 plus PB piercing damage. Beginning at 7th level, the sneak can make this attack twice, instead of once, when they take the Attack action on their turn.

\DndMonsterAction{3rd Level: Weaving Knives (3/Day)}
As an action, the sneak moves up to their speed without provoking opportunity attacks. Before, during, or after the move, they can make two signature attacks.

\DndMonsterAction{5th Level: Sneak and Stab (3/Day)}
As a bonus action, the sneak takes the Hide action. If the sneak hits a creature they are hidden from with an attack on the same turn, the creature takes an extra PBd10 piercing damage, and the sneak can immediately take the Hide action (no action required).

\DndMonsterAction{7th Level: Poisoned Blade (1/Day)}
As a bonus action, the sneak covers a dagger in a special poison, which lasts for 1 hour or until the sneak hits a creature with a signature attack. A creature hit with the poisoned dagger must make a DC 10 plus PB Constitution saving throw. On a failed save, the target takes PBd12 poison damage and is poisoned for 1 minute (save ends at the end of turn). On a successful save, the target takes only half as much damage and isn't poisoned.

\end{DndMonster}



\begin{DndMonster}[float=t]{Goblin Sniper}
\DndMonsterType{Small humanoid (goblin, Artillery), A}
\DndMonsterBasics[
armor-class = {14 (\textit{leather armor})},
hit-points = {\DndDice{3d6 + 3}},
speed = {30 ft., climb 20 ft.}
]
\DndMonsterAbilityScores[
 str = 8,
 dex = 16,
 con = 12,
 int = 10,
 wis = 12,
 cha = 8
]
\DndMonsterDetails[
skills = {Perception +3, Stealth +5},
senses = {darkvision 60 ft., passive perception 13},
languages = {Common, Goblin},
challenge = 1/2,
]
\DndMonsterAction{Crafty}
The sniper doesn't provoke opportunity attacks when they move out of an enemy's reach.

\DndMonsterAction{Sniper}
If the sniper misses with a ranged weapon attack while they are hidden, they remain hidden. Additionally, if the sniper hits a target with a ranged weapon attack while they have advantage on the attack roll, the attack deals an extra 3 (1d6) damage.

\DndMonsterSection{Actions}
\DndMonsterAction{Dagger}
\textsl{Melee or Ranged Weapon Attack:} 5 to hit, reach 5 ft. or range 20/60 ft., one target. \textsl{Hit:} 5 (1d4 + 3) piercing damage.

\DndMonsterAction{Shortbow}
\textsl{Ranged Weapon Attack:} 5 to hit, range 80/320 ft., one target. \textsl{Hit:} 6 (1d6 + 3) piercing damage.

\DndMonsterSection{Bonus Actions}
\DndMonsterAction{Sneak}
The sniper takes the Hide action.

\end{DndMonster}



\begin{DndMonster}[float=t]{Goblin Spinecleaver}
\DndMonsterType{Small humanoid (goblin, Brute), A}
\DndMonsterBasics[
armor-class = {14 (\textit{hide armor})},
hit-points = {\DndDice{6d6 + 12}},
speed = {30 ft., climb 20 ft.}
]
\DndMonsterAbilityScores[
 str = 16,
 dex = 14,
 con = 14,
 int = 10,
 wis = 10,
 cha = 8
]
\DndMonsterDetails[
saving-throws = {Con +4},
skills = {Athletics +5},
senses = {darkvision 60 ft., passive perception 10},
languages = {Common, Goblin},
challenge = 1,
]
\DndMonsterAction{Crafty}
The spinecleaver doesn't provoke opportunity attacks when they move out of an enemy's reach.

\DndMonsterAction{Strong Grip}
Wielding a heavy weapon doesn't impose disadvantage on the spinecleaver's attack rolls.

\DndMonsterSection{Actions}
\DndMonsterAction{Greataxe}
\textsl{Melee Weapon Attack:} 5 to hit, reach 5 ft., one target. \textsl{Hit:} 9 (1d12 + 3) slashing damage.

\DndMonsterAction{Handaxe}
\textsl{Melee or Ranged Weapon Attack:} 5 to hit, range 20/60 ft., one target. \textsl{Hit:} 6 (1d6 + 3) slashing damage.

\DndMonsterSection{Reactions}
\DndMonsterAction{Tricksy Warrior}
When a creature within 5 feet of the spinecleaver misses them with an attack, the spinecleaver can make a melee attack against the creature with disadvantage.

\end{DndMonster}



\begin{DndMonster}[float=t]{Goblin Underboss}
\DndMonsterType{Small humanoid (goblin, Support), A}
\DndMonsterBasics[
armor-class = {17 (studded leather armor)},
hit-points = {\DndDice{8d6 + 8}},
speed = {30 ft., climb 20 ft.}
]
\DndMonsterAbilityScores[
 str = 10,
 dex = 16,
 con = 13,
 int = 12,
 wis = 12,
 cha = 10
]
\DndMonsterDetails[
saving-throws = {Dex +5, Wis +3},
skills = {Insight +3, Intimidation +2, Stealth +5},
senses = {darkvision 60 ft., passive perception 11},
languages = {Common, Goblin},
challenge = 2,
]
\DndMonsterAction{Crafty}
The underboss doesn't provoke opportunity attacks when they move out of an enemy's reach.

\DndMonsterSection{Actions}
\DndMonsterAction{Multiattack}
The underboss makes two Shortsword attacks or two Shortbow attacks. They can replace one attack with a use of Command.

\DndMonsterAction{Shortsword}
\textsl{Melee Weapon Attack:} 5 to hit, reach 5 ft., one target. \textsl{Hit:} 6 (1d6 + 3) piercing damage.

\DndMonsterAction{Shortbow}
\textsl{Ranged Weapon Attack:} 5 to hit, range 80/320 ft., one target. \textsl{Hit:} 6 (1d6 + 3) piercing damage.

\DndMonsterAction{Command}
The underboss chooses one ally they can see within 30 feet of them. If the target can hear the underboss, the target can use their reaction to move up to their speed or make one weapon attack.

\DndMonsterSection{Bonus Actions}
\DndMonsterAction{Get Reckless (Recharge 6)}
Each willing ally within 30 feet of the underboss who can hear them becomes reckless until the start of the underboss's next turn. While reckless, a creature has advantage on attack rolls, and attack rolls against the creature have advantage.

\DndMonsterSection{Reactions}
\DndMonsterAction{Cowardly Commander}
When a creature the underboss can see hits them with an attack, the underboss chooses a willing ally within 5 feet of them. The attack hits the ally instead.

\end{DndMonster}



\begin{DndMonster}[float=t]{Goblin Warrior}
\DndMonsterType{Small humanoid (goblin, Skirmisher), A}
\DndMonsterBasics[
armor-class = {15 (\textit{leather armor})},
hit-points = {\DndDice{2d6 + 2}},
speed = {30 ft., climb 20 ft.}
]
\DndMonsterAbilityScores[
 str = 8,
 dex = 14,
 con = 12,
 int = 10,
 wis = 10,
 cha = 8
]
\DndMonsterDetails[
skills = {Acrobatics +4, Stealth +4},
senses = {darkvision 60 ft., passive perception 10},
languages = {Common, Goblin},
challenge = 1/4,
]
\DndMonsterAction{Crafty}
The warrior doesn't provoke opportunity attacks when they move out of an enemy's reach.

\DndMonsterSection{Actions}
\DndMonsterAction{Shortsword}
\textsl{Melee Weapon Attack:} 4 to hit, reach 5 ft., one target. \textsl{Hit:} 5 (1d6 + 2) piercing damage.

\DndMonsterAction{Shortbow}
\textsl{Ranged Weapon Attack:} 4 to hit, range 80/320 ft., one target. \textsl{Hit:} 5 (1d6 + 2) piercing damage.

\DndMonsterSection{Reactions}
\DndMonsterAction{Fleet Foot}
When a creature within 5 feet of the warrior misses them with a melee attack, the warrior can move up to half their speed.

\end{DndMonster}



\begin{DndMonster}[float=t]{Griffon}
\DndMonsterType{Large beast (Controller), U}
\DndMonsterBasics[
armor-class = {12},
hit-points = {\DndDice{6d10 + 24}},
speed = {30 ft., fly 60 ft.}
]
\DndMonsterAbilityScores[
 str = 18,
 dex = 14,
 con = 18,
 int = 6,
 wis = 14,
 cha = 8
]
\DndMonsterDetails[
saving-throws = {Con +6},
skills = {Perception +4},
challenge = 2,
]
\DndMonsterAction{Steady}
The griffon has advantage on ability checks and saving throws against being knocked prone.

\DndMonsterAction{Swoop}
If the griffon flies more than 20 feet in a straight line toward their target and then hits them with a Claw attack on the same turn, the target is grappled (escape DC 14). While the target is grappled in this way, the griffon has advantage on Beak attacks against the target but can't make Claw attacks. Grappling a Medium or smaller creature imposes no penalty on the griffon's speed.

\DndMonsterSection{Actions}
\DndMonsterAction{Multiattack}
The griffon makes one Beak attack, and they can make one Claw attack, use Buffet, or use Piercing Cry.

\DndMonsterAction{Beak}
\textsl{Melee Weapon Attack:} 6 to hit, reach 5 ft., one target. \textsl{Hit:} 7 (1d6 + 4) piercing damage.

\DndMonsterAction{Claw}
\textsl{Melee Weapon Attack:} 6 to hit, reach 5 ft., one target. \textsl{Hit:} 9 (2d4 + 4) slashing damage.

\DndMonsterAction{Buffet}
The griffon beats their wings in mighty gusts of wind. Each creature not grappled by the griffon within a 30-foot cone originating from the griffon must make a DC 14 Strength saving throw. On a failed save, a creature takes 5 (2d4) bludgeoning damage, is pushed 5 feet away from the griffon, and is knocked prone.

\DndMonsterAction{Piercing Cry}
The griffon lets out an earsplitting screech. Each enemy within 60 feet of the griffon who can hear them must succeed on a DC 12 Wisdom saving throw or be frightened of the griffon for 1 minute (save ends at end of turn). If a creature's saving throw is successful or the effect ends for them, they are immune to the Piercing Cry of all griffons for the next 24 hours.

\end{DndMonster}



\begin{DndMonster}[float=t]{Griffon Companion}
\DndMonsterType{Large beast (Companion), U}
\DndMonsterBasics[
armor-class = {13 plus PB (natural armor)},
hit-points = {\DndDice{7 + seven times caregiver's level (number of d8 Hit Dice equal to their caregiver's level)}},
speed = {30 ft. (see learning to fly)}
]
\DndMonsterAbilityScores[
 str = 16,
 dex = 12,
 con = 15,
 int = 6,
 wis = 12,
 cha = 8
]
\DndMonsterDetails[
saving-throws = {Con +2 plus PB},
skills = {Perception +1 plus PB},
]
\DndMonsterAction{Learning to Fly}
The griffon has a flying speed of 10 × PB feet.

\DndMonsterSection{Actions}
\DndMonsterAction{Signature Attack (Beak)}
\textsl{Melee Weapon Attack:} 3 plus PB to hit, reach 5 ft., one target. \textsl{Hit:} 1d6 plus PB piercing damage.

\DndMonsterAction{1st Level: Violent Attack (2 Ferocity)}
The griffon makes a signature attack. On a hit, the attack deals an extra 0d1 + PB slashing damage, and the griffon can move the target 5 feet in any direction.

\DndMonsterAction{3rd Level: Grabbed from Above (5 Ferocity)}
The griffon makes a signature attack against a Medium or smaller creature. On a hit, the target is also grappled (escape DC 10 plus PB) if they weren't already. While grappled in this way, the target is restrained and the griffon can't make a signature attack against another target.

\DndMonsterAction{3rd Level: Bombs Away (5 Ferocity)}
The griffon flies up to their speed without provoking opportunity attacks. If the griffon has a creature grappled, the grapple doesn't halve the griffon's speed and the griffon can end the grapple at any point during this movement. If the griffon ends the grapple before the end of their turn and the target takes bludgeoning damage from that fall, the target takes an extra 0d1 + PB bludgeoning damage.

\DndMonsterAction{5th Level: Backdraft (8 Ferocity)}
The griffon flies up to their speed without provoking opportunity attacks. When the griffon ends this movement, each creature within 15 feet of them must make a DC 10 plus PB Strength saving throw. On a failed save, a creature takes PBd6 bludgeoning damage and is knocked prone. On a successful save, a creature takes half as much damage and isn't knocked prone.

\end{DndMonster}



\begin{DndMonster}[float=t]{Grilp}
\DndMonsterType{Tiny fiend (devil, Ambusher), lawful evil}
\DndMonsterBasics[
armor-class = {13 (natural armor)},
hit-points = {\DndDice{3d4 + 3}},
speed = {30 ft., fly 30 ft.}
]
\DndMonsterAbilityScores[
 str = 6,
 dex = 14,
 con = 12,
 int = 10,
 wis = 12,
 cha = 11
]
\DndMonsterDetails[
skills = {Stealth +6, Perception +3},
damage-immunities = {fire, necrotic},
senses = {darkvision 60 ft., passive perception 13},
languages = {Goblin, Infernal},
challenge = 1/4,
]
\DndMonsterAction{Shifting Camouflage}
The grilp has advantage on Dexterity (Stealth) checks and can hide from creatures who are observing them.

\DndMonsterSection{Actions}
\DndMonsterAction{Bite}
\textsl{Melee Weapon Attack:} 4 to hit, reach 5 ft., one target. \textsl{Hit:} 4 (1d4 + 2) piercing damage, and if the target is a creature, they must succeed on a DC 11 Constitution saving throw or gain vulnerability to fire and necrotic damage until the start of the grilp's next turn.

\DndMonsterAction{Necrotic Mote}
\textsl{Ranged Spell Attack:} 3 to hit, range 30 ft., one target. \textsl{Hit:} 3 (1d6) necrotic damage.

\end{DndMonster}



\begin{DndMonster}[float=t]{Haunt}
\DndMonsterType{Huge undead (Skirmisher), chaotic evil}
\DndMonsterBasics[
armor-class = {15 (natural armor)},
hit-points = {\DndDice{18d12 + 36}},
speed = {0 ft., fly 40 ft. \textit{(hover)}}
]
\DndMonsterAbilityScores[
 str = 10,
 dex = 16,
 con = 14,
 int = 8,
 wis = 16,
 cha = 18
]
\DndMonsterDetails[
saving-throws = {Wis +7, Cha +8},
skills = {Intimidation +8, Perception +7},
damage-resistances = {acid; fire; lightning; thunder; bludgeoning, piercing, slashing from mundane attacks},
damage-immunities = {cold, necrotic, poison},
condition-immunities = {charmed, dazed, exhaustion, flanked, frightened, grappled, paralyzed, petrified, poisoned, prone, restrained},
senses = {darkvision 120 ft., passive perception 17},
languages = {telepathy 60 ft.},
challenge = 9,
]
\DndMonsterAction{Incorporeal Cloud}
The haunt can occupy another creature's space and vice versa. In addition, the haunt can move through creatures and objects as if they were difficult terrain.

The haunt takes 5 (1d10) force damage if they end their turn inside an object.

\DndMonsterAction{Invisibility}
The haunt is invisible.

\DndMonsterSection{Actions}
\DndMonsterAction{Multiattack}
The haunt makes two Spectral Wrath attacks.

\DndMonsterAction{Spectral Wrath}
\textsl{Melee Spell Attack:} 8 to hit, reach 0 ft., one target in the haunt's space. \textsl{Hit:} 17 (2d12 + 4) force damage.

\DndMonsterAction{Wave of Despair (Recharge 5\textendash 6)}
The haunt unleashes a wave of psychic pain. Each enemy within 20 feet of the haunt must make a DC 16 Wisdom saving throw, taking 33 (6d10) psychic damage on a failed save, or half as much damage on a successful one.

\DndMonsterSection{Bonus Actions}
\DndMonsterAction{Possess Object}
The haunt can magically manipulate a Large or smaller object within 30 feet of them that isn't being worn or carried by another creature. The haunt can exert fine control on objects under their control, such as playing piano keys, slamming doors, opening windows, or writing with a quill.

As part of this bonus action, the haunt can hurl the object up to 30 feet in any direction or use it as a ranged weapon to attack one creature within 30 feet of the object. The object has a +8 bonus to hit and deals 15 (2d10 + 4) bludgeoning damage on a hit. The GM might rule that a specific object deals piercing or slashing damage based on its form.

\end{DndMonster}



\begin{DndMonster}[float=t]{Hellhound}
\DndMonsterType{Medium fiend (Soldier), neutral evil}
\DndMonsterBasics[
armor-class = {15 (natural armor) (18 with Hardened by Flame)},
hit-points = {\DndDice{9d8 + 36}},
speed = {50 ft.}
]
\DndMonsterAbilityScores[
 str = 14,
 dex = 17,
 con = 19,
 int = 11,
 wis = 8,
 cha = 10
]
\DndMonsterDetails[
saving-throws = {Dex +5, Con +6},
skills = {Perception +3, Stealth +7, Survival +3},
damage-immunities = {fire},
senses = {darkvision 60 ft., passive perception 13},
languages = {understands Common and Infernal but can't speak},
challenge = 3,
]
\DndMonsterAction{Hardened by Flame}
When the hound is subjected to fire damage, they take no damage. Instead, their skin darkens and hardens, and their AC increases to 18 until the end of their next turn.

\DndMonsterAction{Stealthy Hunter}
The hound has advantage on Wisdom (Survival) checks to track other creatures and on Dexterity (Stealth) checks to hide from creatures who are unaware of their presence.

\DndMonsterSection{Actions}
\DndMonsterAction{Hellish Bite}
\textsl{Melee Weapon Attack:} 5 to hit, reach 5 ft., one target. \textsl{Hit:} 6 (1d6 + 3) piercing damage plus 4 (1d8) fire damage.

\DndMonsterAction{Hellfire Breath (Recharge 5\textendash 6)}
The hound exhales flame in a 15-foot cone. Each creature in that area must make a DC 14 Dexterity saving throw, taking 14 (3d6 + 4) fire damage on a failed save, or half as much damage on a successful one. A creature who isn't a hellhound who fails the saving throw is lit on fire for 1 minute (save ends at end of turn), or until the target or another creature who can reach them uses an action to extinguish the flames. A creature who is on fire at the start of their turn takes 7 (2d6) fire damage. If a creature who is already on fire is set on fire again on a subsequent turn, the damage isn't cumulative, but the duration of the fire resets to 1 minute.

\DndMonsterSection{Reactions}
\DndMonsterAction{Tug of War}
When an enemy within 5 feet of the hound hits the hound with a melee attack, the hound can make a Hellish Bite attack against that enemy. If the hound's attack hits, the enemy is also grappled (escape DC 12), and if the triggering attack was a weapon attack, the weapon used in the triggering attack can't be used to make any more attacks while the target is grappled. The grapple ends if the hound attacks a different target with Hellish Bite.

\end{DndMonster}



\begin{DndMonster}[float*=b,width=\textwidth + 8pt]{High Mage Vairae}
\begin{multicols}{2}
\DndMonsterType{Medium undead (Solo), neutral evil}
\DndMonsterBasics[
armor-class = {21 (natural armor)},
hit-points = {\DndDice{42d8 + 126}},
speed = {30 ft., fly 50 ft. \textit{(hover)}}
]
\DndMonsterAbilityScores[
 str = 14,
 dex = 19,
 con = 16,
 int = 23,
 wis = 20,
 cha = 20
]
\DndMonsterDetails[
saving-throws = {Con +11, Int +14, Wis +13, Cha +13},
skills = {Arcana +14, Athletics +10, Deception +13, History +14, Perception +13, Persuasion +13},
damage-resistances = {fire; cold; bludgeoning, piercing, slashing from mundane attacks},
damage-immunities = {necrotic, poison},
condition-immunities = {charmed, dazed, exhaustion, frightened, paralyzed, poisoned, stunned},
senses = {truesight 120 ft., passive perception 23},
languages = {Abyssal, Common, Draconic, Elvish, Infernal, Sylvan, telepathy 120 ft.},
challenge = 25,
]
\DndMonsterAction{Draining Resistance (3/Day)}
If Vairae fails a saving throw, they can choose to succeed instead. When they do, Vairae has disadvantage on spell attacks until the end of their next turn.

\DndMonsterAction{Rejuvenation}
Vairae has the Rejuvenation trait found in the lich stat block.

\DndMonsterAction{Supernatural Resistance}
Vairae has advantage on saving throws against powers, spells, and other supernatural effects.

\DndMonsterAction{Turn Immunity}
Vairae is immune to effects that turn Undead.

\DndMonsterSection{Actions}
\DndMonsterAction{Multiattack}
Vairae uses Glare of Undeath and makes four attacks using Conflagration, Immortal Pain, Magic Drain, or a combination of them.

\DndMonsterAction{Conflagration (4th-Level Spell)}
\textsl{Ranged Spell Attack:} 14 to hit, range 150 ft., one target. \textsl{Hit:} 35 (10d6) fire damage, and if the target takes both an action and a bonus action on their next turn, they take an extra 17 (5d6) fire damage at the end of that turn.

\DndMonsterAction{Immortal Pain (4th-Level Spell)}
\textsl{Melee or Ranged Spell Attack:} 14 to hit, reach 5 ft. or range 90 ft., one creature. \textsl{Hit:} 28 (8d6) psychic damage, and the target is afflicted with pain. While afflicted in this way, the creature's attacks deal half damage. After each attack the afflicted creature makes against an enemy, the afflicted creature can make a DC 22 Constitution saving throw, ending the effect on a success. A cure ailment power, a \textit{lesser restoration} spell, or a similar supernatural effect also ends the effect.

\DndMonsterAction{Magic Drain (4th-Level Spell)}
\textsl{Melee or Ranged Spell Attack:} 14 to hit, reach 5 ft. or range 60 ft., one creature. \textsl{Hit:} 26 (4d12) necrotic damage. If the target is currently affected by any spells not cast by Vairae, the target must succeed on a DC 22 Wisdom saving throw or those effects end for the target.

\DndMonsterAction{Glare of Undeath}
Vairae glares at one creature they can see within 60 feet of them. The target must succeed on a DC 22 Wisdom saving throw or be compelled to hug themself tightly as protection against an unbearably cruel world. A creature hugging themself in this way is restrained for 1 minute (save ends at end of turn). A creature the target can hear and understand can use an action to make a DC 22 Charisma (Persuasion) check, ending the effect on a success.

\DndMonsterAction{Machinations of Bone (2/Day)}
Each enemy within 30 feet of a point Vairae can see within 150 feet of them must make a DC 22 Strength saving throw. On a failed save, a target's innards fly from their body and they take 54 (12d8) piercing damage, fall prone, and can't stand up until they recieve supernatural healing.

\DndMonsterSection{Bonus Actions}
\DndMonsterAction{Necrotic Form (3rd-Level Spell)}
Vairae enters a shadow form and can move up to their speed. While in this form, Vairae is immune to bludgeoning, piercing, and slashing damage and the grappled and restrained conditions, and they can move through a space as narrow as 1 inch without squeezing. Vairae reverts to their true form at the end of their turn.

\DndMonsterSection{Reactions}
\DndMonsterAction{Baleful Teleport (6th-Level Spell)}
When a creature Vairae can see within 120 feet of them targets Vairae with an attack, Vairae utters a word of arcane power. The creature must succeed on a DC 22 Charisma saving throw or be teleported up to 60 feet to an unoccupied space Vairae can see, and if Vairae is no longer a valid target for the triggering attack, the attacker must choose a new target or the attack misses.

\DndMonsterSection{Legendary Actions}
Vairae has three villain actions. They can take each action once during an encounter after an enemy's turn. They can take these actions in any order but can use only one per round.

\begin{DndMonsterLegendaryActions}
\DndMonsterLegendaryAction{Action 1: Cages of Wasting}{A cage of black energy springs forth around each enemy within 60 feet of Vairae. Each target must make a DC 22 Dexterity saving throw. On a failed save, a target takes 27 (5d10) necrotic damage and is restrained in the cage for 1 minute or until Vairae is destroyed (save ends at end of turn). On a successful save, a target takes half as much damage and isn't restrained. If a restrained creature teleports out of their cage, the effect ends for them.
}
\DndMonsterLegendaryAction{Action 2: My Magic Alone}{Vairae unleashes a wave of canceling magic. Each creature within 60 feet of Vairae must make a DC 22 Wisdom saving throw. On a failed save, a target can't benefit from magic items or cast spells of 4th level or higher until the end of Vairae's next turn.
}
\DndMonsterLegendaryAction{Action 3: Necrostorm}{Vairae fires black bolts of magic from their body. Each creature within 60 feet of them must make a DC 22 Dexterity saving throw, taking 78 (12d12) necrotic damage on a failed save, or half as much damage on a successful one. The spirit of a Humanoid killed by this damage rises at the start of Vairae's next turn as a wraith who permanently follows Vairae's verbal commands.
}
\end{DndMonsterLegendaryActions}
\end{multicols}
\end{DndMonster}



\begin{DndMonster}[float=t]{Hobgoblin Brandbearer}
\DndMonsterType{Medium humanoid (hobgoblin, Minion), A}
\DndMonsterBasics[
armor-class = {13 (\textit{hide armor})},
hit-points = {\DndDice{12}},
speed = {30 ft.}
]
\DndMonsterAbilityScores[
 str = 14,
 dex = 12,
 con = 12,
 int = 10,
 wis = 11,
 cha = 15
]
\DndMonsterDetails[
damage-immunities = {fire},
senses = {darkvision 60 ft., passive perception 10},
languages = {Common, Goblin, Infernal},
challenge = 5,
]
\DndMonsterAction{Infernal Ichor}
When a creature within 5 feet of the brandbearer reduces the brandbearer to 0 hit points, the brandbearer's corpse unleashes a spray of burning orange ichor. The creature who reduced the brandbearer to 0 hit points takes 1 fire damage.

\DndMonsterAction{Minion}
If the brandbearer takes damage from an attack or as the result of a failed saving throw, their hit points are reduced to 0. If the brandbearer takes damage from another effect, they die if the damage equals or exceeds their hit point maximum; otherwise they take no damage.

\DndMonsterAction{Too Hot to Handle}
Whenever an enemy takes fire damage from a non-minion creature, the enemy takes an extra 1 fire damage for each brandbearer within 5 feet of them.

\DndMonsterSection{Actions}
\DndMonsterAction{Searing Grasp (Group Attack)}
\textsl{Melee Spell Attack:} 5 to hit, reach 5 ft., one target. \textsl{Hit:} 4 fire damage. If three or more brandbearers joined the attack, the target has disadvantage on saving throws against powers, spells, and other supernatural effects that deal fire damage until the start of the brandbearer's next turn.

\end{DndMonster}


\clearpage

\begin{DndMonster}[float=t]{Hobgoblin Death Captain}
\DndMonsterType{Medium humanoid (hobgoblin, Support), A}
\DndMonsterBasics[
armor-class = {20 (plate armor)},
hit-points = {\DndDice{12d8 + 36}},
speed = {30 ft.}
]
\DndMonsterAbilityScores[
 str = 18,
 dex = 12,
 con = 17,
 int = 15,
 wis = 12,
 cha = 14
]
\DndMonsterDetails[
saving-throws = {Con +6, Wis +4},
skills = {Athletics +7, History +5, Intimidation +5, Perception +4},
damage-resistances = {fire},
senses = {darkvision 60 ft., passive perception 14},
languages = {Common, Goblin, Infernal},
challenge = 6,
]
\DndMonsterAction{Battle Ready}
The captain and each ally within 60 feet of them who can hear them have advantage on initiative rolls and can't be surprised, provided the captain isn't incapacitated.

\DndMonsterAction{Infernal Ichor}
When the captain dies, their corpse unleashes a spray of burning orange ichor. Each creature within 5 feet of the captain takes 5 fire damage.

\DndMonsterSection{Actions}
\DndMonsterAction{Multiattack}
The captain makes two Blightblade or two Eye Fire attacks. They can replace one attack with a use of On My Mark.

\DndMonsterAction{Blightblade}
\textsl{Melee Weapon Attack:} 7 to hit, reach 5 ft., one target. \textsl{Hit:} 8 (1d8 + 4) slashing damage plus 7 (2d6) necrotic damage.

\DndMonsterAction{Eye Fire}
\textsl{Ranged Spell Attack:} 6 to hit, range 60 ft., one target. \textsl{Hit:} 14 (4d6) fire damage.

\DndMonsterAction{On My Mark}
The captain commands an ally within 60 feet of them who can hear them to attack. That ally can make a weapon attack (no action required), dealing an extra 3 (1d6) necrotic damage on a hit.

\DndMonsterSection{Bonus Actions}
\DndMonsterAction{Lead the Charge (1/Day)}
The captain and each ally within 30 feet of them who can see them can move up to their speed without provoking opportunity attacks.

\DndMonsterSection{Reactions}
\DndMonsterAction{Hidden Gift (3/Day)}
When an ally the captain can see within 30 feet of them takes damage, the captain magically turns that ally invisible. The ally remains invisible until the end of their next turn or until they attack, deal damage, or force another creature to make a saving throw.

\end{DndMonster}



\begin{DndMonster}[float=t]{Hobgoblin Firerunner}
\DndMonsterType{Medium humanoid (hobgoblin, Skirmisher), A}
\DndMonsterBasics[
armor-class = {16 (studded leather armor)},
hit-points = {\DndDice{12d8 + 24}},
speed = {40 ft.}
]
\DndMonsterAbilityScores[
 str = 15,
 dex = 18,
 con = 14,
 int = 11,
 wis = 12,
 cha = 10
]
\DndMonsterDetails[
saving-throws = {Dex +6},
skills = {Acrobatics +6},
damage-immunities = {fire},
senses = {darkvision 60 ft., passive perception 11},
languages = {Common, Goblin, Infernal},
challenge = 4,
]
\DndMonsterAction{Fiery Leap}
The firerunner's long jump is up to 30 feet and their high jump is up to 15 feet, with or without a running start. The first time a firerunner lands after a jump on their turn, each creature within 5 feet of them takes 3 (1d6) fire damage.

\DndMonsterAction{Infernal Ichor}
When the firerunner dies, their corpse unleashes a spray of burning orange ichor. Each creature within 5 feet of the firerunner takes 5 fire damage.

\DndMonsterSection{Actions}
\DndMonsterAction{Multiattack}
The firerunner makes two Fire Scimitar attacks.

\DndMonsterAction{Fire Scimitar}
\textsl{Melee Weapon Attack:} 6 to hit, reach 5 ft., one target. \textsl{Hit:} 7 (1d6 + 4) slashing damage plus 7 (2d6) fire damage.

\DndMonsterAction{Blaze Sprint (Recharge 5\textendash 6)}
The firerunner moves up to their speed without provoking opportunity attacks. Each creature they come within 5 feet of during this movement must make a DC 14 Dexterity saving throw, taking 10 (3d6) fire damage on a failed save, or half as much damage on a successful one.

\end{DndMonster}



\begin{DndMonster}[float=t]{Hobgoblin Incendiarist}
\DndMonsterType{Medium humanoid (hobgoblin, Artillery), A}
\DndMonsterBasics[
armor-class = {15 (studded leather armor)},
hit-points = {\DndDice{6d8 + 6}},
speed = {30 ft.}
]
\DndMonsterAbilityScores[
 str = 10,
 dex = 16,
 con = 12,
 int = 10,
 wis = 16,
 cha = 10
]
\DndMonsterDetails[
skills = {Perception +5},
damage-resistances = {fire},
senses = {darkvision 60 ft., passive perception 15},
languages = {Common, Goblin, Infernal},
challenge = 1,
]
\DndMonsterAction{Infernal Ichor}
When the incendiarist dies, their corpse unleashes a spray of burning orange ichor. Each creature within 5 feet of the incendiarist takes 3 fire damage.

\DndMonsterSection{Actions}
\DndMonsterAction{Shortsword}
\textsl{Melee Weapon Attack:} 5 to hit, reach 5 ft., one target. \textsl{Hit:} 6 (1d6 + 3) piercing damage.

\DndMonsterAction{Blazing Crossbow}
\textsl{Ranged Weapon Attack:} 5 to hit, range 100/400 ft., one target. \textsl{Hit:} 8 (1d10 + 3) piercing damage, and if the target is a creature, they are set on fire. While on fire, the target takes 3 (1d6) fire damage at the end of each of their turns, and the target or another creature who can reach them can use an action to put out the flames. A creature who is already on fire suffers no additional effect from being set on fire in this way again.

\DndMonsterAction{Fire Mote (Recharge 6)}
The incendiarist shoots an explosive mote of fire at a creature they can see within 100 feet of them. The target must succeed on a DC 13 Dexterity saving throw or take 7 (2d6) fire damage and, if they are Large or smaller, fall prone. All other creatures within 15 feet of the target must succeed on a DC 13 Dexterity saving throw or take 3 (1d6) fire damage.

\end{DndMonster}



\begin{DndMonster}[float=t]{Hobgoblin Recruit}
\DndMonsterType{Medium humanoid (hobgoblin, Minion), A}
\DndMonsterBasics[
armor-class = {14 (\textit{leather armor})},
hit-points = {\DndDice{7}},
speed = {30 ft.}
]
\DndMonsterAbilityScores[
 str = 14,
 dex = 13,
 con = 12,
 int = 11,
 wis = 12,
 cha = 10
]
\DndMonsterDetails[
damage-immunities = {fire},
senses = {darkvision 60 ft., passive perception 11},
languages = {Common, Goblin, Infernal},
challenge = 1/2,
]
\DndMonsterAction{Infernal Ichor}
When a creature within 5 feet of the recruit reduces the recruit to 0 hit points, the recruit's corpse unleashes a spray of burning orange ichor. The creature who reduced the recruit to 0 hit points takes 1 fire damage.

\DndMonsterAction{Minion}
If the recruit takes damage from an attack or as the result of a failed saving throw, their hit points are reduced to 0. If the recruit takes damage from another effect, they die if the damage equals or exceeds their hit point maximum; otherwise they take no damage.

\DndMonsterAction{Tactical Positioning}
When a non-minion ally attacks a creature, that ally gains a +1 bonus to the attack roll for each recruit within 5 feet of the target (maximum bonus of +5).

\DndMonsterSection{Actions}
\DndMonsterAction{Longsword (Group Attack)}
\textsl{Melee Weapon Attack:} 4 to hit, reach 5 ft., one target. \textsl{Hit:} 1 slashing damage.

\end{DndMonster}



\begin{DndMonster}[float=t]{Hobgoblin Smokebinder}
\DndMonsterType{Medium humanoid (hobgoblin, Ambusher), A}
\DndMonsterBasics[
armor-class = {17 (studded leather armor)},
hit-points = {\DndDice{8d8 + 8}},
speed = {30 ft. (fly 30 ft. in smoke form)}
]
\DndMonsterAbilityScores[
 str = 12,
 dex = 17,
 con = 13,
 int = 14,
 wis = 12,
 cha = 10
]
\DndMonsterDetails[
skills = {Stealth +5},
damage-resistances = {fire; bludgeoning, piercing, slashing from mundane attacks in smoke form},
condition-immunities = {grappled, restrained in smoke form},
senses = {darkvision 60 ft., passive perception 11},
languages = {Common, Goblin, Infernal},
challenge = 2,
]
\DndMonsterAction{Amorphous (Smoke Form Only)}
The smokebinder can move through a space as narrow as 1 inch wide without squeezing.

\DndMonsterAction{Infernal Ichor}
When the smokebinder dies, their corpse unleashes a spray of burning orange ichor. Each creature within 5 feet of the smokebinder takes 3 fire damage.

\DndMonsterSection{Actions}
\DndMonsterAction{Toxic Flame}
\textsl{Melee Spell Attack:} 5 to hit, reach 5 ft., one target. \textsl{Hit:} 10 (3d6) fire damage. If the smokebinder has advantage on the attack roll, the target is poisoned and can't speak until the end of the target's next turn.

\DndMonsterAction{Choking Bolt}
\textsl{Ranged Spell Attack:} 5 to hit, range 60 ft., one target. \textsl{Hit:} 9 (2d8) fire damage. If the smokebinder has advantage on the attack roll, the target can't speak until the end of the target's next turn.

\DndMonsterSection{Bonus Actions}
\DndMonsterAction{Smoke Form (1/Day)}
The smokebinder transforms from flesh into a magical gray smoke for 1 minute. Any equipment they are wearing or carrying is also transformed for the duration.

While in smoke form, the smokebinder gains a flying speed equal to their walking speed, resistance to bludgeoning, piercing, and slashing damage from mundane attacks, and immunity to the grappled and restrained conditions.

The smokebinder can revert to their true form as a bonus action. If the smokebinder becomes incapacitated or dies, they revert to their true form.

\DndMonsterAction{Sneak (Smoke Form Only)}
The smokebinder takes the Hide action.

\end{DndMonster}



\begin{DndMonster}[float=t]{Hobgoblin Tactician}
\DndMonsterType{Medium humanoid (hobgoblin, Retainer), A}
\DndMonsterBasics[
armor-class = {15 (medium armor)},
hit-points = {\DndDice{Seven times their level (number of d8 Hit Dice equal to their level)}},
speed = {30 ft.}
]
\DndMonsterAbilityScores[
 str = 10,
 dex = 10,
 con = 12,
 int = 16,
 wis = 12,
 cha = 10
]
\DndMonsterDetails[
saving-throws = {Str +PB, Dex +PB, Con +PB, Int +PB, Wis +PB, Cha +PB},
skills = {Arcana +3 plus PB, History +3 plus PB},
damage-resistances = {fire},
senses = {darkvision 60 ft., passive perception 11},
languages = {Common, Goblin, Infernal},
]
\DndMonsterSection{Actions}
\DndMonsterAction{Signature Attack (Flame Touch)}
\textsl{Melee or Ranged Spell Attack:} 3 plus PB to hit, reach 5 ft. or range 60 ft., one target. \textsl{Hit:} 5 (1d10) fire damage. A flammable object hit by this spell ignites if it isn't being worn or carried. This spell's damage increases by 1d10 when the tactician reaches 5th level (2d10), 11th level (3d10), and 17th level (4d10).

\DndMonsterAction{3rd Level: Heat Seeker (3/Day)}
As an action, the tactician makes a signature attack but doesn't make an attack roll for it. Instead, the attack automatically hits, and each ally within 30 feet of the target can use their reaction to move up to their speed toward the target.

\DndMonsterAction{5th Level: Beacon (3/Day)}
When the tactician hits with a signature attack, each creature within 5 feet of the target can use a reaction to make a melee weapon attack against the target.

\DndMonsterAction{7th Level: Explosion (1/Day)}
As an action, the tactician hurls a ball of fire at a point they can see within 120 feet of them. The ball explodes in a 20-foot-radius sphere, forcing each enemy in that area to make a DC 10 plus PB Dexterity saving throw. On a failed save, a creature takes PBd12 fire damage and is pushed 20 feet away from the center of the area. On a successful save, a creature takes half as much damage and isn't pushed. The fire spreads around corners and ignites flammable objects in the area that aren't being worn or carried.

\end{DndMonster}



\begin{DndMonster}[float=t]{Hobgoblin Trooper}
\DndMonsterType{Medium humanoid (hobgoblin, Soldier), A}
\DndMonsterBasics[
armor-class = {16 (\textit{ring mail})},
hit-points = {\DndDice{4d8}},
speed = {30 ft.}
]
\DndMonsterAbilityScores[
 str = 14,
 dex = 12,
 con = 11,
 int = 10,
 wis = 10,
 cha = 10
]
\DndMonsterDetails[
skills = {Athletics +4, Perception +2},
damage-resistances = {fire},
senses = {darkvision 60 ft., passive perception 12},
languages = {Common, Goblin, Infernal},
challenge = 1/2,
]
\DndMonsterAction{Infernal Ichor}
When the trooper dies, their corpse unleashes a spray of burning orange ichor. Each creature within 5 feet of the trooper takes 2 fire damage.

\DndMonsterSection{Actions}
\DndMonsterAction{Fire Flail}
\textsl{Melee Weapon Attack:} 4 to hit, reach 5 ft., one target. \textsl{Hit:} 6 (1d8 + 2) bludgeoning damage, and another target within 5 feet of the first takes 2 (1d4) fire damage.

\DndMonsterAction{Brimstone Dagger}
\textsl{Melee or Ranged Weapon Attack:} 4 to hit, reach 5 ft. or range 30/60 ft., one target. \textsl{Hit:} 4 (1d4 + 2) piercing damage plus 2 (1d4) fire damage.

\DndMonsterSection{Bonus Actions}
\DndMonsterAction{Fight Me, Coward}
The trooper hexes a creature within 5 feet of them who isn't already hexed by another trooper. The first time on a turn a hexed creature damages a creature other than the trooper, the hexed creature takes 5 fire damage. The hex lasts until the start of the trooper's next turn.

\end{DndMonster}



\begin{DndMonster}[float=t]{Hobgoblin War Mage}
\DndMonsterType{Medium humanoid (hobgoblin, Controller), A}
\DndMonsterBasics[
armor-class = {15 (\textit{chain shirt})},
hit-points = {\DndDice{10d8 + 20}},
speed = {30 ft.}
]
\DndMonsterAbilityScores[
 str = 10,
 dex = 15,
 con = 15,
 int = 18,
 wis = 14,
 cha = 16
]
\DndMonsterDetails[
saving-throws = {Int +6, Cha +5},
skills = {Arcana +6, History +6},
damage-resistances = {fire},
senses = {darkvision 60 ft., passive perception 12},
languages = {Common, Goblin, Infernal},
challenge = 3,
]
\DndMonsterAction{Infernal Ichor}
When the mage dies, their corpse unleashes a spray of burning orange ichor. Each creature within 5 feet of the mage takes 5 fire damage.

\DndMonsterSection{Actions}
\DndMonsterAction{Infernal Teleport}
\textsl{Melee Spell Attack:} 6 to hit, reach 5 ft., one target. \textsl{Hit:} 11 (2d10) necrotic damage, and if the target is a creature, they must succeed on a DC 14 Charisma saving throw or be teleported up to 15 feet to an unoccupied space the mage can see.

\DndMonsterAction{Rain Hellfire (2nd-Level Spell)}
The mage rains fire down on a point they can see within 60 feet of them. Each creature in a 10-foot-radius sphere centered on that point must make a DC 14 Dexterity saving throw, taking 10 (3d6) fire damage on a failed save, or half as much damage on a successful one. The fire spreads around corners, and it ignites flammable objects in that area that aren't being worn or carried.

\DndMonsterAction{Unhallowed Ground (3rd-Level Spell; Concentration)}
The mage chooses a point on the ground that they can see within 60 feet of them. A 20-foot square of glowing, unholy energy appears on the ground centered on that point and lasts for up to 1 minute. Each enemy in that area is vulnerable to fire damage.

\DndMonsterSection{Reactions}
\DndMonsterAction{Spell Siphon}
When a creature the mage can see within 60 feet of them casts a spell, the mage reflects some of the spell's magical energy back at the caster. The caster must succeed on a DC 14 Dexterity saving throw or take 3 (1d6) force damage per level of the triggering spell (minimum 1d6 force damage).

\end{DndMonster}



\begin{DndMonster}[float=t]{Hulking Brain}
\DndMonsterType{Large aberration (Brute), lawful evil}
\DndMonsterBasics[
armor-class = {16 (natural armor)},
hit-points = {\DndDice{12d10 + 60}},
speed = {30 ft.}
]
\DndMonsterAbilityScores[
 str = 21,
 dex = 10,
 con = 20,
 int = 8,
 wis = 6,
 cha = 12
]
\DndMonsterDetails[
senses = {blindsight 120 ft. (blind beyond this radius), passive perception 8},
languages = {Deep Speech, Undercommon, telepathy 120 ft.},
challenge = 7,
]
\DndMonsterAction{Psionic Immunity}
The hulking brain is unaffected by psionic powers manifested by voiceless talkers unless the brain wishes to be.

\DndMonsterAction{Psychic Scream}
When the brain drops to 0 hit points, each creature within 30 feet of them must make a DC 16 Wisdom saving throw, taking 21 (6d6) psychic damage on a failed save, or half as much damage on a successful one.

\DndMonsterSection{Actions}
\DndMonsterAction{Multiattack}
The brain makes three Bash attacks.

\DndMonsterAction{Bash}
\textsl{Melee Weapon Attack:} 8 to hit, reach 10 ft., one target. \textsl{Hit:} 14 (2d8 + 5) bludgeoning damage, and the target is grappled (escape DC 15). Until the grapple ends, the target is restrained and has disadvantage on saving throws against psionic powers. The brain can grapple up to four creatures at once.

\DndMonsterAction{Squeeze}
Each creature grappled by the brain must make a DC 16 Strength saving throw, taking 22 (4d10) bludgeoning damage on a failed save, or half as much damage on a successful one.

\DndMonsterSection{Bonus Actions}
\DndMonsterAction{Psionic Invigoration (1/Day)}
Each willing voiceless talker within 60 feet of the brain regains the use of one psionic power of the voiceless talker's choice, and the voiceless talker has advantage on attack rolls until the end of the brain's next turn.

\end{DndMonster}



\begin{DndMonster}[float=t]{Human Apprentice Mage}
\DndMonsterType{Medium humanoid (human, Minion), A}
\DndMonsterBasics[
armor-class = {14 (studded leather armor)},
hit-points = {\DndDice{13}},
speed = {30 ft.}
]
\DndMonsterAbilityScores[
 str = 8,
 dex = 14,
 con = 10,
 int = 14,
 wis = 11,
 cha = 10
]
\DndMonsterDetails[
skills = {Arcana +5},
languages = {Common plus any two languages},
challenge = 6,
]
\DndMonsterAction{Empower Patron}
A non-minion creature serves as the mage's patron in exchange for the mage's protective magic. At the start of the patron's turn, the patron gains temporary hit points equal to twice the number of apprentice mages within 60 feet of them who chose them as a patron and can see them.

\DndMonsterAction{Exploit Weakness}
When the mage makes or joins an attack that's made with advantage, the attack deals an extra 1 damage per mage who made or joined the attack.

\DndMonsterAction{Minion}
If the mage takes damage from an attack or as the result of a failed saving throw, their hit points are reduced to 0. If the mage takes damage from another effect, they die if the damage equals or exceeds their hit point maximum; otherwise they take no damage.

\DndMonsterSection{Actions}
\DndMonsterAction{Lightning Strike (Group Attack)}
\textsl{Ranged Spell Attack:} 5 to hit, range 120 ft., one creature. \textsl{Hit:} 2 lightning damage, and if this attack was made by more than one mage, each mage who joined the attack picks one creature within 30 feet of the original target. Each creature picked takes 2 lightning damage for each mage who targeted them; the mage's Exploit Weakness trait does not increase this damage.

\DndMonsterAction{Thunder Crack (Group Attack)}
\textsl{Melee Spell Attack:} 5 to hit, reach 5 ft., one target. \textsl{Hit:} 4 thunder damage, and the target can't make opportunity attacks until the start of their next turn.

\end{DndMonster}



\begin{DndMonster}[float=t]{Human Brawler}
\DndMonsterType{Medium humanoid (human, Brute), A}
\DndMonsterBasics[
armor-class = {14 (studded leather armor)},
hit-points = {\DndDice{7d8 + 21}},
speed = {30 ft.}
]
\DndMonsterAbilityScores[
 str = 17,
 dex = 14,
 con = 16,
 int = 10,
 wis = 10,
 cha = 12
]
\DndMonsterDetails[
skills = {Athletics +5, Intimidation +3},
languages = {Common},
challenge = 2,
]
\DndMonsterAction{Exploit Opening (3/Day)}
When the brawler makes an attack, they have advantage on the attack roll.

\DndMonsterSection{Actions}
\DndMonsterAction{Multiattack}
The brawler makes two attacks using Grab, Haymaker, or both.

\DndMonsterAction{Grab}
\textsl{Melee Weapon Attack:} 5 to hit, reach 5 ft., one Medium or smaller creature. \textsl{Hit:} 5 (1d4 + 3) bludgeoning damage, and the target is grappled (escape DC 13). Until this grapple ends, the brawler can't grab another creature.

\DndMonsterAction{Haymaker}
\textsl{Melee Weapon Attack:} 5 to hit, reach 5 ft., one target. \textsl{Hit:} 5 (1d4 + 3) bludgeoning damage, or 10 (3d4 + 3) bludgeoning damage against a grappled target.

\DndMonsterAction{Throw}
The brawler throws one Medium or smaller creature they are grappling or object they are holding up to 30 feet horizontally. If the thrown target is a creature, they fall prone after this throw. If the thrown target would enter the space of a creature or solid object that is no more than one size smaller than it, the thrown target collides with it and stops in the nearest unoccupied space, taking 3 (1d6) bludgeoning damage for every 10 feet it was thrown. A Large or smaller creature hit by this thrown target must succeed on a DC 13 Dexterity saving throw or take the same amount of damage and fall prone.

\DndMonsterSection{Reactions}
\DndMonsterAction{Meat Shield}
When the brawler is grappling a target and is hit by a ranged attack made by another creature the brawler can see, the brawler gains a +2 bonus to AC against the triggering attack. If this bonus causes the attack to miss the brawler, it hits the grappled target instead.

\end{DndMonster}



\begin{DndMonster}[float=t]{Human Death Acolyte}
\DndMonsterType{Medium humanoid (human, Minion), A}
\DndMonsterBasics[
armor-class = {13 (\textit{leather armor})},
hit-points = {\DndDice{9}},
speed = {30 ft.}
]
\DndMonsterAbilityScores[
 str = 11,
 dex = 14,
 con = 12,
 int = 10,
 wis = 14,
 cha = 14
]
\DndMonsterDetails[
languages = {Common},
challenge = 2,
]
\DndMonsterAction{Death Chant}
When an enemy succeeds on a death saving throw while within 30 feet of three or more acolytes who aren't incapacitated, the success instead counts as a failure, and if the roll was a 20, the enemy doesn't regain hit points.

\DndMonsterAction{Exploit Weakness}
When the acolyte makes or joins an attack that's made with advantage, the attack deals an extra 1 damage per acolyte who made or joined the attack.

\DndMonsterAction{Minion}
If the acolyte takes damage from an attack or as the result of a failed saving throw, their hit points are reduced to 0. If the acolyte takes damage from another effect, they die if the damage equals or exceeds their hit point maximum; otherwise they take no damage.

\DndMonsterSection{Actions}
\DndMonsterAction{Transfer Life (Group Attack)}
\textsl{Melee Spell Attack:} 4 to hit, reach 5 ft., one creature. \textsl{Hit:} 2 necrotic damage, and if this attack was made by more than one acolyte, a creature of the acolytes' choice within 30 feet of them regains 1 hit point per acolyte who joined this group attack.

\DndMonsterAction{Necrotic Bolt (Group Attack)}
\textsl{Ranged Spell Attack:} 4 to hit, range 30 ft., one target. \textsl{Hit:} 2 necrotic damage.

\end{DndMonster}



\begin{DndMonster}[float=t]{Human Death Cultist}
\DndMonsterType{Medium humanoid (human, Support), A}
\DndMonsterBasics[
armor-class = {16 (\textit{chain mail})},
hit-points = {\DndDice{8d8 + 24}},
speed = {30 ft.}
]
\DndMonsterAbilityScores[
 str = 16,
 dex = 10,
 con = 16,
 int = 14,
 wis = 18,
 cha = 12
]
\DndMonsterDetails[
saving-throws = {Wis +6, Cha +3},
skills = {Intimidation +3, Religion +4},
languages = {Common plus any one language},
challenge = 4,
]
\DndMonsterAction{Exploit Opening (3/Day)}
When the cultist makes an attack, they have advantage on the attack roll.

\DndMonsterSection{Actions}
\DndMonsterAction{Scythe}
\textsl{Melee Weapon Attack:} 5 to hit, reach 5 ft., one target. \textsl{Hit:} 8 (1d10 + 3) slashing damage plus 7 (2d6) necrotic damage, and the cultist regains hit points equal to half the necrotic damage dealt.

\DndMonsterAction{Death Bolt}
\textsl{Ranged Spell Attack:} 6 to hit, range 60 ft., one target. \textsl{Hit:} 14 (4d6) necrotic damage, and the target's weapon attacks deal half damage until the start of the cultist's next turn.

\DndMonsterAction{Blackfire Blessing (1/Day)}
The cultist empowers up to 10 non-minion allies within 30 feet of them. For 1 minute or until the cultist dies, each creature's weapons burn with black fire and deal an extra 2 (1d4) necrotic damage on a hit.

\DndMonsterSection{Bonus Actions}
\DndMonsterAction{Rise, My Minions! (1/Day)}
The cultist chooses up to three creatures within 30 feet of them who died within the last minute. These creatures return to life with 1 hit point, but they can't regain hit points, and they die after 1 minute.

\DndMonsterSection{Reactions}
\DndMonsterAction{Life from Death}
When a creature the cultist can see within 30 feet of them fails a death saving throw or dies, the cultist siphons their faltering life energy. The cultist chooses a creature within 30 feet of the cultist who isn't unconscious, and that creature regains 14 (4d6) hit points.

\end{DndMonster}



\begin{DndMonster}[float=t]{Human Guard}
\DndMonsterType{Medium humanoid (human, Minion), A}
\DndMonsterBasics[
armor-class = {16 (\textit{chain shirt})},
hit-points = {\DndDice{5}},
speed = {30 ft.}
]
\DndMonsterAbilityScores[
 str = 12,
 dex = 12,
 con = 12,
 int = 10,
 wis = 10,
 cha = 10
]
\DndMonsterDetails[
skills = {Perception +2},
languages = {Common},
challenge = 1/8,
]
\DndMonsterAction{Exploit Weakness}
When the guard makes or joins an attack that's made with advantage, the attack deals an extra 1 damage per guard who made or joined the attack.

\DndMonsterAction{Minion}
If the guard takes damage from an attack or as the result of a failed saving throw, their hit points are reduced to 0. If the guard takes damage from another effect, they die if the damage equals or exceeds their hit point maximum; otherwise they take no damage.

\DndMonsterAction{Overwhelm}
If a Medium or smaller enemy starts their turn within 5 feet of three or more guards who can see them, until the start of the enemy's next turn, the enemy's speed is reduced by 5 feet for each guard within 5 feet of them. If this reduces the enemy's walking speed to 0, they are restrained until the start of their next turn.

\DndMonsterSection{Actions}
\DndMonsterAction{Spear (Group Attack)}
\textsl{Melee or Ranged Weapon Attack:} 3 to hit, reach 5 ft. or range 20/60 ft., one target. \textsl{Hit:} 1 piercing damage.

\end{DndMonster}



\begin{DndMonster}[float=t]{Human Knave}
\DndMonsterType{Medium humanoid (human, Soldier), A}
\DndMonsterBasics[
armor-class = {16 (\textit{breastplate})},
hit-points = {\DndDice{8d8 + 24}},
speed = {30 ft.}
]
\DndMonsterAbilityScores[
 str = 18,
 dex = 10,
 con = 16,
 int = 10,
 wis = 12,
 cha = 10
]
\DndMonsterDetails[
skills = {Intimidation +4, Perception +3},
languages = {Common},
challenge = 3,
]
\DndMonsterAction{Exploit Opening (3/Day)}
When the knave makes an attack, they have advantage on the attack roll.

\DndMonsterSection{Actions}
\DndMonsterAction{Multiattack}
The knave makes two Morningstar or two Javelin attacks. They can replace one attack with a Shield Bash attack.

\DndMonsterAction{Morningstar}
\textsl{Melee Weapon Attack:} 6 to hit, reach 5 ft., one target. \textsl{Hit:} 8 (1d8 + 4) piercing damage, or 13 (2d8 + 4) piercing damage if the target is prone.

\DndMonsterAction{Shield Bash}
\textsl{Melee Weapon Attack:} 6 to hit, reach 5 ft., one target. \textsl{Hit:} 6 (1d4 + 4) bludgeoning damage, and if the target is a Medium or smaller creature, they must succeed on a DC 14 Strength saving throw or be knocked prone.

\DndMonsterAction{Javelin}
\textsl{Ranged Weapon Attack:} 6 to hit, range 30/120 ft., one target. \textsl{Hit:} 7 (1d6 + 4) piercing damage.

\DndMonsterSection{Bonus Actions}
\DndMonsterAction{Stay Down}
The knave kicks one prone creature within 5 feet of them. The target must succeed on a DC 14 Constitution saving throw or their speed is reduced to 0 until the end of their next turn.

\end{DndMonster}



\begin{DndMonster}[float=t]{Human Mercenary}
\DndMonsterType{Medium humanoid (human, Retainer), A}
\DndMonsterBasics[
armor-class = {15 (medium armor)},
hit-points = {\DndDice{Nine times their level (number of d12 Hit Dice equal to their level)}},
speed = {30 ft.}
]
\DndMonsterAbilityScores[
 str = 16,
 dex = 10,
 con = 14,
 int = 10,
 wis = 10,
 cha = 10
]
\DndMonsterDetails[
saving-throws = {Str +PB, Dex +PB, Con +PB, Int +PB, Wis +PB, Cha +PB},
skills = {Athletics +3 plus PB, Medicine +0 plus PB, Perception +0 plus PB},
languages = {Common},
]
\DndMonsterSection{Actions}
\DndMonsterAction{Signature Attack (Halberd)}
\textsl{Melee Weapon Attack:} 3 plus PB to hit, reach 10 ft., one target. \textsl{Hit:} 1d10 plus PB slashing damage. Beginning at 7th level, the mercenary can make this attack twice, instead of once, when they take the Attack action on their turn.

\DndMonsterAction{3rd Level: Battlefield Medicine (3/Day)}
As an action, the mercenary restores PBd8 hit points to a creature they can see within 5 feet of them.

\DndMonsterAction{5th Level: Exploit Opening (3/Day)}
The mercenary takes the Attack action, making each attack with advantage and dealing an extra PBd4 slashing damage on a hit.

\DndMonsterAction{7th Level: Halberd Master (3/Day)}
When a creature enters a space within the mercenary's reach, the mercenary uses a reaction to make a signature attack with advantage against that creature. If this attack hits, it deals an extra PBd6 slashing damage.

\end{DndMonster}



\begin{DndMonster}[float=t]{Human Raider}
\DndMonsterType{Medium humanoid (human, Skirmisher), A}
\DndMonsterBasics[
armor-class = {14 (padded armor) (12 without shield)},
hit-points = {\DndDice{4d8 + 4}},
speed = {30 ft.}
]
\DndMonsterAbilityScores[
 str = 15,
 dex = 12,
 con = 12,
 int = 10,
 wis = 12,
 cha = 10
]
\DndMonsterDetails[
skills = {Stealth +3},
languages = {Common},
challenge = 1/2,
]
\DndMonsterAction{Charge}
If the raider moves at least 15 feet straight toward a target and then hits the target with a melee attack on the same turn, the target takes an extra 3 (1d6) damage.

\DndMonsterAction{Exploit Opening (3/Day)}
When the raider makes an attack, they have advantage on the attack roll.

\DndMonsterSection{Actions}
\DndMonsterAction{Handaxe}
\textsl{Melee or Ranged Weapon Attack:} 4 to hit, reach 5 ft. or range 20/60 ft., one target. \textsl{Hit:} 5 (1d6 + 2) slashing damage.

\DndMonsterAction{Shield Shove}
\textsl{Melee Weapon Attack:} 4 to hit, reach 5 ft., one target. \textsl{Hit:} 4 (1d4 + 2) bludgeoning damage, and if the target is a Medium or smaller creature, they must succeed on a DC 12 Strength saving throw or be pushed 5 feet away from the raider.

\DndMonsterSection{Reactions}
\DndMonsterAction{Splinter Shield}
When the raider is wielding a shield and is hit by an attack made by a creature they can see, the raider gains a +4 bonus to AC against the triggering attack. If this causes the attack to miss, the raider's shield breaks, and until they get a new shield, their AC is reduced by 2 and they can't use this reaction or make a Shield Shove attack.

\end{DndMonster}



\begin{DndMonster}[float=t]{Human Scoundrel}
\DndMonsterType{Medium humanoid (human, Ambusher), A}
\DndMonsterBasics[
armor-class = {15 (studded leather armor)},
hit-points = {\DndDice{10d8 + 10}},
speed = {30 ft.}
]
\DndMonsterAbilityScores[
 str = 13,
 dex = 16,
 con = 12,
 int = 10,
 wis = 10,
 cha = 14
]
\DndMonsterDetails[
skills = {Deception +4, Sleight Of Hand +5, Stealth +5},
languages = {Common},
challenge = 3,
]
\DndMonsterAction{Ambusher}
In the first round of combat, the scoundrel has advantage on attack rolls against any surprised creature.

\DndMonsterAction{Exploit Opening (3/Day)}
When the scoundrel makes an attack, they have advantage on the attack roll.

\DndMonsterAction{Hit 'Em Where It Hurts}
When the scoundrel has advantage on a weapon attack roll, the attack deals an extra 7 (2d6) damage on a hit.

\DndMonsterSection{Actions}
\DndMonsterAction{Rapier}
\textsl{Melee Weapon Attack:} 5 to hit, reach 5 ft., one target. \textsl{Hit:} 7 (1d8 + 3) piercing damage, and the scoundrel can make a dagger attack against the target with advantage.

\DndMonsterAction{Dagger}
\textsl{Melee or Ranged Weapon Attack:} 5 to hit, reach 5 ft. or range 20/60 ft., one target. \textsl{Hit:} 5 (1d4 + 3) piercing damage.

\end{DndMonster}



\begin{DndMonster}[float=t]{Human Storm Wizard}
\DndMonsterType{Medium humanoid (human, Controller), A}
\DndMonsterBasics[
armor-class = {12 (15 with \textit{mage armor})},
hit-points = {\DndDice{10d8 + 30}},
speed = {30 ft.}
]
\DndMonsterAbilityScores[
 str = 10,
 dex = 14,
 con = 16,
 int = 18,
 wis = 14,
 cha = 10
]
\DndMonsterDetails[
saving-throws = {Int +7, Wis +5},
skills = {Arcana +7, History +7},
languages = {Common plus any two languages},
challenge = 6,
]
\DndMonsterAction{Exploit Opening (3/Day)}
When the wizard makes an attack, they have advantage on the attack roll.

\DndMonsterAction{Spellcasting (Utility)}
In addition to any other spells in this stat block, the wizard can cast the following spells, using Intelligence as the spellcasting ability (spell save DC 15):
\begin{DndMonsterSpells}
  \DndInnateSpellLevel{mage hand\textsuperscript{Casting Times}, \textit{prestidigitation}\textsuperscript{Casting Times}}
  \DndInnateSpellLevel[1]{clairvoyance\textsuperscript{Casting Times}, \textit{mage armor}\textsuperscript{Casting Times}, \textit{see invisibility}\textsuperscript{Casting Times}, \textit{sending}\textsuperscript{Casting Times}}
\end{DndMonsterSpells}
\DndMonsterSection{Actions}
\DndMonsterAction{Arcane Staff}
\textsl{Melee or Ranged Spell Attack:} 7 to hit, reach 5 ft. or range 30 ft., one target. \textsl{Hit:} 21 (6d6) lightning damage.

\DndMonsterAction{Gust of Wind (1/Day; 2nd-Level Spell; Concentration)}
A 10-foot-wide, 60-foot-long line of strong wind gusts from the wizard for 1 minute. Each creature who starts their turn in that area must succeed on a DC 15 Strength saving throw or be pushed 15 feet away from the wizard. Each creature in that area must spend 2 feet of movement for every 1 foot they move toward the wizard.

The gust disperses gas or vapor, and it extinguishes candles, torches, and similar unprotected flames. It has a 50 percent chance to extinguish protected flames like lanterns. The wizard can use a bonus action to change the direction in which the wind blasts from them.

\DndMonsterAction{Lightning Bolt (3/Day; 3rd-Level Spell)}
The wizard fires magical lightning in a 5-foot-wide, 100-foot-long line. Each creature in the line must make a DC 15 Dexterity saving throw, taking 28 (8d6) lightning damage on a failed save, or half as much damage on a successful one. The lightning ignites flammable objects in the area that aren't being worn or carried.

\DndMonsterSection{Reactions}
\DndMonsterAction{Arcane Shield (3/Day)}
When the wizard is hit by an attack, they magically gain a +5 bonus to AC against that attack, potentially causing it to miss. If the attacker is within 10 feet of the wizard, the attacker must succeed on a DC 15 Constitution saving throw or take 18 (4d8) thunder damage and be pushed 10 feet away from the wizard.

\end{DndMonster}



\begin{DndMonster}[float=t]{Human Trickshot}
\DndMonsterType{Medium humanoid (human, Artillery), A}
\DndMonsterBasics[
armor-class = {14},
hit-points = {\DndDice{6d8 + 12}},
speed = {30 ft.}
]
\DndMonsterAbilityScores[
 str = 12,
 dex = 18,
 con = 14,
 int = 10,
 wis = 14,
 cha = 14
]
\DndMonsterDetails[
skills = {Perception +4, Stealth +6},
languages = {Common},
challenge = 2,
]
\DndMonsterAction{Exploit Opening (3/Day)}
When the trickshot makes an attack, they have advantage on the attack roll.

\DndMonsterAction{Point Blank Shooting}
When the trickshot hits a creature within 30 feet of them with a ranged weapon attack, they deal an extra 7 (2d6) piercing damage. Additionally, being within 5 feet of an enemy doesn't impose disadvantage on the trickshot's ranged attack rolls.

\DndMonsterSection{Actions}
\DndMonsterAction{Heavy Crossbow}
\textsl{Ranged Weapon Attack:} 6 to hit, range 100/400 ft., one target. \textsl{Hit:} 9 (1d10 + 4) piercing damage.

\DndMonsterAction{Bayonet}
\textsl{Melee Weapon Attack:} 6 to hit, reach 5 ft., one target. \textsl{Hit:} 7 (1d6 + 4) piercing damage.

\DndMonsterAction{Ricochet Bolt (Recharge 5\textendash 6)}
The trickshot uses a special ricocheting bolt to make one Heavy Crossbow attack against a target within 100 feet of them, then a second Heavy Crossbow attack against a different target within 30 feet of the first target.

\end{DndMonster}


\clearpage

\begin{DndMonster}[float*=b,width=\textwidth + 8pt]{Infernal Chancellor Lazivos}
\begin{multicols}{2}
\DndMonsterType{Medium fiend (devil, Leader), lawful evil}
\DndMonsterBasics[
armor-class = {19 (natural armor)},
hit-points = {\DndDice{28d8 + 112}},
speed = {40 ft., fly 60 ft.}
]
\DndMonsterAbilityScores[
 str = 14,
 dex = 20,
 con = 18,
 int = 19,
 wis = 18,
 cha = 24
]
\DndMonsterDetails[
saving-throws = {Dex +10, Int +9, Wis +9, Cha +12},
skills = {Arcana +9, Deception +12, Insight +9, Perception +9, Persuasion +12, Religion +9},
damage-resistances = {psychic; bludgeoning, piercing, slashing from mundane attacks},
damage-immunities = {fire},
condition-immunities = {charmed, dazed, exhaustion, flanked, frightened, paralyzed, stunned},
senses = {darkvision 120 ft., passive perception 19},
languages = {Common, Infernal},
challenge = 16,
]
\DndMonsterAction{Security for Speed (3/Day)}
When Lazivos fails a saving throw, he can succeed instead. When he does, his speed is halved and he can't take bonus actions or reactions until the end of his next turn.

\DndMonsterAction{Supernatural Resistance}
Lazivos has advantage on saving throws against powers, spells, and other supernatural effects.

\DndMonsterAction{True Name}
If a creature Lazivos can hear within 60 feet of him speaks his true name aloud, Lazivos loses his damage resistances, damage immunities, and Devilish Charm reaction for 24 hours. Lazivos's true name is Minto Abi.

\DndMonsterSection{Actions}
\DndMonsterAction{Multiattack}
Lazivos makes three attacks using Infernal Rapier, Chancellor's Decree, or both.

\DndMonsterAction{Infernal Rapier}
\textsl{Melee Weapon Attack:} 10 to hit, reach 5 ft., one target. \textsl{Hit:} 9 (1d8 + 5) piercing damage plus 9 (2d8) fire damage, and the target is marked with an infernal sigil until the end of their next turn. While marked, a creature can't hide or benefit from being invisible, and attack rolls against them have advantage.

\DndMonsterAction{Chancellor's Decree}
\textsl{Ranged Spell Attack:} 12 to hit, range 120 ft., one target. \textsl{Hit:} 18 (2d10 + 7) psychic damage, and the target must make a DC 20 Wisdom saving throw. On a failed save, the target is charmed by Lazivos until the end of his next turn or until Lazivos or one of his allies harms the target.

\DndMonsterAction{Diabolic Deposition (Recharge 5\textendash 6)}
Lazivos commands the attention of each enemy within 30 feet of him who can hear him. Each target must make a DC 20 Wisdom saving throw. On a failed save, a creature must choose to either take 27 (6d8) fire damage plus 27 (6d8) psychic damage or be dazed for 1 minute (save ends at end of turn). On a successful save, a creature must either choose to take half as much damage or be dazed until the end of their next turn. Succeed or fail, if a creature is already dazed, they must choose the damage

\DndMonsterSection{Bonus Actions}
\DndMonsterAction{As I Say}
Lazivos commands a creature charmed by him to use their reaction, if available, to make a weapon attack against a target of his choice.

\DndMonsterSection{Reactions}
\DndMonsterAction{Devilish Charm}
When Lazivos is targeted by an attack, power, spell, or other supernatural effect by a creature he can see within 60 feet of him, the creature must make a DC 20 Charisma saving throw. On a failed save, the creature is charmed by Lazivos until the start of the creature's next turn, and Lazivos chooses a new target he can see for the triggering effect. The new target must be within the triggering effect's range.

\DndMonsterSection{Legendary Actions}
Lazivos has three villain actions. He can take each action once during an encounter after an enemy's turn. He can take these actions in any order but can use only one per round.

\begin{DndMonsterLegendaryActions}
\DndMonsterLegendaryAction{Action 1: Welcome, Friends}{Lazivos invokes an infernal word of power. Up to three enemies within 60 feet of him who can hear him must make a DC 20 Charisma saving throw. On a failed save, a target is charmed by Lazivos for 1 hour (save ends at end of turn) or until Lazivos or one of his allies harms the target. On a successful save, a target takes 16 (3d10) psychic damage.
}
\DndMonsterLegendaryAction{Action 2: Heed My Commands!}{Lazivos and each ally within 60 feet of him moves up to their speed without provoking opportunity attacks. Each creature charmed by Lazivos moves up to half their speed in a direction he chooses (no action required).
}
\DndMonsterLegendaryAction{Action 3: Deceptive Stratagem}{Lazivos switches places with a willing creature of his choice within 120 feet of him. Immediately afterward, up to five creatures of Lazivos's choice within 120 feet of him can make a melee attack (no action required).
}
\end{DndMonsterLegendaryActions}
\end{multicols}
\end{DndMonster}



\begin{DndMonster}[float*=b,width=\textwidth + 8pt]{Ithu'rath}
\begin{multicols}{2}
\DndMonsterType{Large aberration (Solo), chaotic evil}
\DndMonsterBasics[
armor-class = {19 (natural armor)},
hit-points = {\DndDice{30d10 + 90}},
speed = {30 ft., fly 60 ft. \textit{(hover)}, swim 60 ft.}
]
\DndMonsterAbilityScores[
 str = 23,
 dex = 9,
 con = 16,
 int = 23,
 wis = 14,
 cha = 18
]
\DndMonsterDetails[
saving-throws = {Con +8, Int +11, Wis +7, Cha +9},
skills = {History +11, Insight +12, Intimidation +9, Perception +12},
condition-immunities = {charmed, flanked, frightened, prone},
senses = {blindsight 60 ft., darkvision 120 ft., passive perception 22},
languages = {Common, Deep Speech, Undercommon, telepathy 120 ft.},
challenge = 16,
]
\DndMonsterAction{Amphibious}
Ithu'rath can breathe air and water.

\DndMonsterAction{Immutable Form}
Ithu'rath is immune to any power, spell, or effect that would alter their form.

\DndMonsterAction{Slimy Resistance (3/Day)}
When Ithu'rath fails a saving throw, they can succeed on the save by channeling the power of their protective slime. Until the end of Ithu'rath's next turn, creatures have advantage on the saving throw for Ithu'rath's Devolving Tentacle attack.

\DndMonsterAction{Supernatural Resistance}
Ithu'rath has advantage on saving throws against powers, spells, and other supernatural effects.

\DndMonsterSection{Actions}
\DndMonsterAction{Multiattack}
Ithu'rath makes six Devolving Tentacle attacks and uses Piscine Transformation.

\DndMonsterAction{Devolving Tentacle}
\textsl{Melee Weapon Attack:} 11 to hit, reach 20 ft., one target. \textsl{Hit:} 16 (3d6 + 6) piercing damage. The first time a creature is hit with this attack on a turn, they must succeed on a DC 19 Constitution saving throw or be injected with a degenerative psionic slime (save ends at end of turn). While slimed, each time the creature makes an attack roll or an ability check, they must roll a d4 and subtract the number rolled from the d20 roll. At the end of a slimed creature's turn, the size of the subtracted die increases by one: the d4 becomes a d6, the d6 becomes a d8, and so on, to a maximum of d12. The cure ailment power or \textit{lesser restoration} spell also ends this effect.

\DndMonsterAction{Piscine Transformation (3rd-Order Power)}
Ithu'rath targets a creature with a head, legs, or torso they can see within 60 feet of them, blasting them with a pulse of transformational energy. The target must succeed on a DC 19 Wisdom saving throw or undergo one of the following transformations of the Ithu'rath's choice:


\begin{description}
\item \textbf{Head.} The target's head transforms into the head of a fish, proportionately sized for their body. They can't speak.
\item \textbf{Legs.} The target's legs become fins. Their walking speed is reduced to 10 feet (unless their walking speed is slower), and they gain a swimming speed of 30 feet.
\item \textbf{Torso.} The target's torso becomes the body of a fish with gills. They can only breathe water and can hold their breath for 1 hour. If the target isn't underwater when this transformation takes effect, they begin suffocating.
\end{description}

The transformation lasts until a cure ailment power of 4th order or higher or the \textit{greater restoration} spell affects the target. A target who suffers all three transformations at once is permanently transformed into a hybrid fishlike creature and can only be returned to their original form by a \textit{wish} spell or similar supernatural effects.

\DndMonsterSection{Bonus Actions}
\DndMonsterAction{Jaunt (3rd-Order Power)}
Ithu'rath teleports 30 feet to an unoccupied space they can see.

\DndMonsterSection{Reactions}
\DndMonsterAction{Liquefy (3rd-Order Power)}
When a creature within 30 feet of Ithu'rath hits them with an attack, Ithu'rath can psionically barrage that creature. That creature must succeed on a DC 19 Constitution saving throw or be pushed 5 feet away from Ithu'rath and be restrained until the end of the creature's next turn, as their lower body temporarily melts into an immovable goo.

\DndMonsterSection{Legendary Actions}
Ithu'rath has three villain actions. They can take each action once during an encounter after an enemy's turn. They can take these actions in any order but can use only one per round.

\begin{DndMonsterLegendaryActions}
\DndMonsterLegendaryAction{Action 1: Skeletal Collapse}{Ithu'rath targets one or more creatures with a skeleton within 60 feet of them, attempting to temporarily melt their ligaments. Each target must succeed on a DC 19 Constitution saving throw or take 18 (4d8) necrotic damage and have disadvantage on Strength- and Dexterity-based ability checks, attack rolls, and saving throws for 1 minute (save ends at end of turn).
}
\DndMonsterLegendaryAction{Action 2: Psychic Pulse}{Ithu'rath channels a pulse of telekinetic energy to scour and harry their foes. Each enemy within 60 feet of Ithu'rath must succeed on a DC 19 Strength saving throw or take 18 (4d8) force damage and be moved up to 30 feet in any direction.
}
\DndMonsterLegendaryAction{Action 3: Slime Storm}{Ithu'rath unleashes a cloud of transformative slime. Each creature within 60 feet of Ithu'rath who isn't an olothec or a slime servant must make a DC 19 Constitution saving throw. On a failed save, a target takes 45 (10d8) piercing damage and is blinded until the end of their next turn as their eyes temporarily turn into insectoid feelers. On a successful save, a target takes half as much damage and isn't blinded.
}
\end{DndMonsterLegendaryActions}
\end{multicols}
\end{DndMonster}



\begin{DndMonster}[float*=b,width=\textwidth + 8pt]{Kingfissure Worm}
\begin{multicols}{2}
\DndMonsterType{Gargantuan monstrosity (Solo), U}
\DndMonsterBasics[
armor-class = {16 (natural armor)},
hit-points = {\DndDice{18d20 + 108}},
speed = {60 ft., burrow 50 ft.}
]
\DndMonsterAbilityScores[
 str = 29,
 dex = 10,
 con = 23,
 int = 1,
 wis = 10,
 cha = 5
]
\DndMonsterDetails[
saving-throws = {Con +11},
condition-immunities = {charmed, flanked, frightened, prone},
senses = {blindsight 30 ft., tremorsense 120 ft., passive perception 10},
challenge = 15,
]
\DndMonsterAction{Multiple Tongues}
The worm has three tongues, each of which can be targeted with attacks (AC 20, immunity to psychic damage) while that tongue is grappling a creature. When a tongue takes 35 damage or more in a single turn, that tongue is destroyed, releasing any target it was grappling. The worm can't use a destroyed tongue with their Multiattack or Tongue Grab action until they finish a long rest. Damaging or destroying a tongue deals no damage to the worm.

\DndMonsterAction{Slowed, Not Stopped (3/Day)}
When the worm fails a saving throw, they can choose to succeed instead. When they do, their speed and the reach of their Tongue Grab attack are each reduced by a cumulative 10 feet until they finish a long rest. The worm's speed and reach can't be reduced below 5 feet in this way.

\DndMonsterAction{Tunneler}
The worm can burrow through solid rock at half their burrowing speed and leaves a 10-foot-diameter tunnel in their wake.

\DndMonsterSection{Actions}
\DndMonsterAction{Multiattack}
The worm can make a number of Tongue Grab attacks equal to the number of tongues they currently have (see Multiple Tongues). The worm can replace any number of those attacks with a use of Swallow.

\DndMonsterAction{Tongue Grab}
\textsl{Melee Weapon Attack:} 14 to hit, reach 30 ft., one target. \textsl{Hit:} 14 (4d6) acid damage, and the creature is grappled (escape DC 19). Until this grapple ends, the worm can't use this tongue against another target.

\DndMonsterAction{Maw}
\textsl{Melee Weapon Attack:} 14 to hit, reach 5 ft., one target. \textsl{Hit:} 22 (3d8 + 9) piercing damage. If the target is Huge or smaller, they must succeed on a DC 19 Dexterity saving throw or be swallowed (as if by the worm's Swallow action).

\DndMonsterAction{Swallow}
The worm pulls a Huge or smaller grappled creature into their maw and swallows them. A swallowed creature is no longer grappled, but they are blinded and restrained, have total cover against attacks and other effects outside the worm, and take 28 (8d6) acid damage at the start of each of the worm's turns.

If the worm takes 30 damage or more on a single turn from a creature inside them, the worm must succeed on a DC 21 Constitution check at the end of that turn or regurgitate all swallowed creatures, each of whom fall prone in an unoccupied space within 10 feet of the worm. If the worm dies, a swallowed creature can escape from the corpse by using 20 feet of movement, exiting prone.

\DndMonsterSection{Bonus Actions}
\DndMonsterAction{Tongue Sweep}
The worm violently swings a creature they have grappled. The grappled creature takes 10 (3d6) bludgeoning damage, and each enemy within 10 feet of the grappled creature must succeed on a DC 16 Dexterity saving throw or take 10 (3d6) bludgeoning damage.

\DndMonsterSection{Reactions}
\DndMonsterAction{Tongue Recall}
When one of the worm's tongues takes damage that doesn't immediately destroy it, the worm can release any target that tongue is grappling and recall the tongue back inside the worm's body. That tongue can't be targeted until the worm uses it to grapple a creature again.

\DndMonsterSection{Legendary Actions}
The worm has three villain actions. They can take each action once during an encounter after an enemy's turn. They can take these actions in any order but can use only one per round.

\begin{DndMonsterLegendaryActions}
\DndMonsterLegendaryAction{Action 1: Open the Earth}{The worm writhes, sending tremors through the ground and opening a 20-foot-wide, 100-foot-long, 40-foot-deep fissure in the ground, originating from the worm. Each creature in that area must make a DC 19 Dexterity saving throw. On a failed save, a creature falls into the fissure, taking 14 (4d6) bludgeoning damage from the fall and landing prone. On a successful save, a creature flees to the closest unoccupied space outside the fissure's area.
}
\DndMonsterLegendaryAction{Action 2: Earth Breach}{The worm burrows up to half their speed and then bursts through the surface, breaching the ground in a 10-foot-radius circle. Huge and smaller creatures within that area must succeed on a DC 19 Dexterity saving throw or be swallowed (as if by the worm's Swallow action).
}
\DndMonsterLegendaryAction{Action 3: Better Out Than In}{The worm projectile vomits the contents of their stomach in a 60-foot cone. Each creature in that area must make a DC 19 Dexterity saving throw, taking 21 (6d6) bludgeoning damage plus 21 (6d6) acid damage on a failed save, or half as much damage on a successful one. Additionally, creatures who were swallowed by the worm are regurgitated, taking 21 (6d6) bludgeoning damage and landing prone in an unoccupied space of the worm's choice within that area.
}
\end{DndMonsterLegendaryActions}
\end{multicols}
\end{DndMonster}



\begin{DndMonster}[float*=b,width=\textwidth + 8pt]{Kiona the Dread Lord}
\begin{multicols}{2}
\DndMonsterType{Medium undead (incorporeal, Leader), chaotic evil}
\DndMonsterBasics[
armor-class = {17 (natural armor)},
hit-points = {\DndDice{16d8 + 64}},
speed = {0 ft., fly 60 ft. \textit{(hover)}}
]
\DndMonsterAbilityScores[
 str = 9,
 dex = 18,
 con = 18,
 int = 20,
 wis = 14,
 cha = 13
]
\DndMonsterDetails[
saving-throws = {Int +9, Wis +6},
skills = {Arcana +9, History +9, Perception +6},
damage-resistances = {acid; cold; fire; lightning; thunder; bludgeoning, piercing, slashing from mundane attacks},
damage-immunities = {necrotic, poison},
condition-immunities = {charmed, dazed, exhaustion, grappled, paralyzed, petrified, poisoned, prone, restrained, unconscious},
senses = {blindsight 60 ft., passive perception 16},
languages = {Common, Draconic, Orc},
challenge = 9,
]
\DndMonsterAction{Corrupting Phasing}
Kiona can move through other creatures and objects. Kiona takes 5 (1d10) force damage if she ends her turn inside an object. A creature takes 9 (2d8) necrotic damage the first time Kiona passes through them on a turn.

\DndMonsterAction{Incorporeal Empowerment}
Each incorporeal Undead within 120 feet of Kiona ignores difficult terrain.

\DndMonsterAction{Spiritual Sacrifice (3/Day)}
When Kiona fails a saving throw, she can lose 15 hit points and choose to succeed instead.

\DndMonsterSection{Actions}
\DndMonsterAction{Multiattack}
Kiona makes two Energy Siphon attacks.

\DndMonsterAction{Energy Siphon}
\textsl{Melee or Ranged Spell Attack:} 9 to hit, reach 5 ft. or range 120 ft., one creature. \textsl{Hit:} 15 (3d6 + 5) necrotic damage, and the target can't regain hit points until the start of Kiona's next turn. Kiona gains temporary hit points equal to half the damage dealt by this attack.

\DndMonsterAction{Rise from Death (1/Day)}
Kiona conjures three shadows, two specters, or one wraith under her control. Each summoned creature appears in an unoccupied space within 30 feet of Kiona and acts immediately after her turn on the same initiative count.

\DndMonsterSection{Bonus Actions}
\DndMonsterAction{Path of Pain}
Kiona chooses one other incorporeal Undead within 120 feet of her who can see her. That Undead can move up to their speed.

\DndMonsterSection{Reactions}
\DndMonsterAction{Death's Breath}
When a creature within 5 feet of Kiona hits her with an attack, she breathes a toxin at the creature. The creature must succeed on a DC 17 Constitution saving throw or be poisoned until the end of their next turn.

\DndMonsterSection{Legendary Actions}
Kiona has three villain actions. She can take each action once during an encounter after an enemy's turn. She can take these actions in any order but can use only one per round.

\begin{DndMonsterLegendaryActions}
\DndMonsterLegendaryAction{Action 1: Restrain Them!}{A 20-foot-radius sphere of spectral arms emerges from a point Kiona can see within 120 feet of her. Each enemy in that area must succeed on a DC 17 Dexterity saving throw or be restrained until the end of Kiona's next turn.
}
\DndMonsterLegendaryAction{Action 2: Ruinous March}{Kiona and each ally within 60 feet of her can move up to their speed without provoking opportunity attacks.
}
\DndMonsterLegendaryAction{Action 3: Vitality Drain}{Kiona and each ally within 60 feet of her can make a melee attack (no action required). A creature hit by one of these attacks gains a level of exhaustion.
}
\end{DndMonsterLegendaryActions}
\end{multicols}
\end{DndMonster}



\begin{DndMonster}[float=t]{Kobold Decanus}
\DndMonsterType{Small humanoid (kobold, Retainer), A}
\DndMonsterBasics[
armor-class = {15 (medium armor) (includes shield)},
hit-points = {\DndDice{Eight times their level (number of d10 Hit Dice equal to their level)}},
speed = {30 ft.}
]
\DndMonsterAbilityScores[
 str = 16,
 dex = 10,
 con = 12,
 int = 12,
 wis = 10,
 cha = 10
]
\DndMonsterDetails[
saving-throws = {Str +PB, Dex +PB, Con +PB, Int +PB, Wis +PB, Cha +PB},
skills = {Athletics +3 plus PB, History +1 plus PB, Perception +0 plus PB},
senses = {darkvision 60 ft., passive perception 10 plus PB},
languages = {Common, Draconic},
]
\DndMonsterSection{Actions}
\DndMonsterAction{Signature Attack (Gladius)}
\textsl{Melee Weapon Attack:} 3 plus PB to hit, reach 5 ft., one target. \textsl{Hit:} 1d8 plus PB piercing damage. Beginning at 7th level, the decanus can make this attack twice, instead of once, when they take the Attack action on their turn.

\DndMonsterAction{3rd Level: Reinforce (3/Day)}
When the decanus sees an ally within 30 feet of them being targeted by an attack, the decanus can use a reaction to move up to their speed toward the target. If the decanus ends their movement within 5 feet of the target, the target gains a +PB bonus to their AC until the beginning of the decanus's next turn, including against the triggering attack.

\DndMonsterAction{5th Level: Shield Bash (3/Day)}
The decanus attempts to topple a Medium or smaller creature within 5 feet of them. The target must succeed on a DC 10 plus PB Strength saving throw or be knocked prone. Succeed or fail, the decanus then makes two signature attacks against any creature within reach.

\DndMonsterAction{7th Level: One-Kobold Army (1/Day)}
The decanus moves up to their speed, then makes a number of signature attacks equal to the number of enemies within 5 feet of them (minimum of three attacks). Each enemy within range must be attacked at least once. On a hit, the target must succeed on a DC 10 plus PB Charisma saving throw or be frightened of the decanus for 1 minute (save ends at end of turn).

\end{DndMonster}



\begin{DndMonster}[float=t]{Koptourok}
\DndMonsterType{Medium undead (Controller), chaotic evil}
\DndMonsterBasics[
armor-class = {13},
hit-points = {\DndDice{12d8 + 36}},
speed = {30 ft. (see wings of breath)}
]
\DndMonsterAbilityScores[
 str = 9,
 dex = 16,
 con = 16,
 int = 10,
 wis = 12,
 cha = 8
]
\DndMonsterDetails[
saving-throws = {Dex +6, Con +6},
skills = {Perception +4},
damage-resistances = {necrotic; poison; bludgeoning, piercing, slashing from mundane attacks},
condition-immunities = {charmed, dazed, exhaustion, frightened, paralyzed, petrified, poisoned},
senses = {blindsight 120 ft., passive perception 14},
languages = {understands the languages they knew in life but can't speak},
challenge = 5,
]
\DndMonsterAction{Breathless Aura}
Enemies within 30 feet of the koptourok can't speak, and if they are a creature who needs to breathe, they also have disadvantage on Constitution saving throws.

\DndMonsterAction{Wings of Breath}
The koptourok gains a flying speed equal to their walking speed if they have at least one enemy within 30 feet of them. If the koptourok starts their turn in the air and doesn't have a flying speed, they fall.

\DndMonsterSection{Actions}
\DndMonsterAction{Multiattack}
The koptourok makes two Choking Grasp attacks.

\DndMonsterAction{Choking Grasp}
\textsl{Melee Spell Attack:} 6 to hit, reach 30 ft., one creature. \textsl{Hit:} 14 (4d6) bludgeoning damage, and the target is pulled to an unoccupied space within 5 feet of the koptourok and grappled (escape DC 14). While grappled in this way, the target is restrained and the koptourok's speed isn't halved by the grapple.

\DndMonsterAction{Last Gasp}
The koptourok attempts to steal the breath from each creature they are grappling. Each target must make a DC 14 Constitution saving throw, taking 14 (4d6) bludgeoning damage on a failed save, or half as much damage on a successful one. The koptourok gains temporary hit points equal to 5 times the number of creatures who failed their save.

\DndMonsterSection{Reactions}
\DndMonsterAction{Thunderous Deflation}
When the koptourok drops to 30 hit points or fewer, they unleash a shrieking wail. Each creature within 30 feet of the koptourok must make a DC 14 Constitution saving throw. On a failed save, a creature takes 16 (3d10) thunder damage and is deafened until the end of their next turn. On a successful save, a creature takes half as much damage and isn't deafened. The koptourok then releases all creatures they are grappling and flies up to their speed.

\end{DndMonster}



\begin{DndMonster}[float=t]{Lacuna}
\DndMonsterType{Medium fiend (Skirmisher), chaotic evil}
\DndMonsterBasics[
armor-class = {15 (natural armor)},
hit-points = {\DndDice{26d8 + 78}},
speed = {0 ft.}
]
\DndMonsterAbilityScores[
 str = 5,
 dex = 10,
 con = 17,
 int = 18,
 wis = 16,
 cha = 24
]
\DndMonsterDetails[
saving-throws = {Dex +5, Wis +8, Cha +12},
skills = {Arcana +9, Deception +12, Insight +8, Persuasion +12},
condition-immunities = {charmed, dazed, exhaustion, frightened, grappled, paralyzed, petrified, poisoned, prone, restrained, stunned},
senses = {truesight 60 ft., passive perception 13},
languages = {Abyssal, Common, Infernal},
challenge = 15,
]
\DndMonsterAction{Avoidance}
If the lacuna is subjected to an effect that allows them to make a saving throw to take only half damage, they instead take no damage if they succeed on the saving throw, and only half damage if they fail.

\DndMonsterAction{False Appearance}
While the lacuna remains motionless, they are indistinguishable from a normal statue.

\DndMonsterAction{Fiend Between}
The lacuna doesn't have a walking speed and can't benefit from any bonus to their speed.

\DndMonsterAction{Supernatural Resistance}
The lacuna has advantage on saving throws against powers, spells, and other supernatural effects.

\DndMonsterSection{Actions}
\DndMonsterAction{Multiattack}
The lacuna makes three Grief-Eater attacks and uses Infuse Sadness, if available.

\DndMonsterAction{Grief-Eater}
\textsl{Melee or Ranged Spell Attack:} 12 to hit, reach 5 ft. or range 20 ft., one creature. \textsl{Hit:} 29 (4d10 + 7) psychic damage, and the target's speed is reduced by 10 feet until the end of their next turn.

\DndMonsterAction{Infuse Sadness (1/Day)}
Each enemy the lacuna can see within 30 feet of them must succeed on a DC 20 Wisdom saving throw or be overwhelmed with grief for 1 minute (save ends at end of turn). While a creature is overwhelmed, attacks against them are made with advantage.

\DndMonsterSection{Bonus Actions}
\DndMonsterAction{Slip Between}
The lacuna teleports up to 60 feet to an unoccupied space they can see.

\DndMonsterSection{Reactions}
\DndMonsterAction{Slip Away}
When a creature within 60 feet of the lacuna targets them with an attack, the lacuna teleports up to 60 feet to an unoccupied space they can see. If the lacuna is no longer a valid target for the triggering attack, the attacker must choose a new target or the attack misses.

\end{DndMonster}



\begin{DndMonster}[float*=b,width=\textwidth + 8pt]{Lady Emer}
\begin{multicols}{2}
\DndMonsterType{Medium monstrosity (Solo), lawful evil}
\DndMonsterBasics[
armor-class = {17},
hit-points = {\DndDice{20d8 + 100}},
speed = {40 ft.}
]
\DndMonsterAbilityScores[
 str = 18,
 dex = 24,
 con = 20,
 int = 16,
 wis = 17,
 cha = 18
]
\DndMonsterDetails[
saving-throws = {Str +8, Dex +11, Con +9, Wis +7},
skills = {Athletics +8, Deception +8, Perception +7, Persuasion +8, Stealth +11, Survival +7},
damage-immunities = {poison},
condition-immunities = {poisoned},
senses = {darkvision 60 ft., passive perception 17},
languages = {Common, Elvish},
challenge = 11,
]
\DndMonsterAction{Stone Sacrifice (3/Day)}
When Emer fails a saving throw, she can choose to succeed instead by ending the effects of her Stone Gaze on a creature of her choice within 300 feet of her.

\DndMonsterSection{Actions}
\DndMonsterAction{Multiattack}
Emer uses Pinning Shot or makes three attacks using Snake Bite, Longbow, or both. She also uses Stone Gaze or Envenomed Stone, if available.

\DndMonsterAction{Snake Bite}
\textsl{Melee Weapon Attack:} 11 to hit, reach 10 ft., one creature. \textsl{Hit:} 17 (5d6) poison damage, and the target must succeed on a DC 17 Constitution saving throw or be poisoned until the end of Emer's next turn. While poisoned in this way, the target can't take reactions.

\DndMonsterAction{Longbow}
\textsl{Ranged Weapon Attack:} 11 to hit, range 150/600 ft., one target. \textsl{Hit:} 11 (1d8 + 7) piercing damage plus 14 (4d6) poison damage.

\DndMonsterAction{Pinning Shot}
Emer makes three Longbow attacks against a creature touching the ground within 150 feet of her. If at least two of those attacks hit, the target is restrained by arrows pinning their limbs to the ground. A creature restrained in this way or another creature who can reach them can use an action to pull out the arrows and end the condition.

\DndMonsterAction{Stone Gaze}
Emer fires beams of energy from her eyes at up to three creatures she can see within 60 feet of her. Each target must succeed on a DC 17 Constitution saving throw or begin turning to stone (save ends at end of turn). While turning to stone, a creature's speed is reduced by 10 feet and they have disadvantage on Dexterity checks and saving throws. If a creature who is turning to stone is targeted by this action again and fails their save, their speed is reduced by another 10 feet and they are dazed until the effect ends.

If a creature who is turning to stone and dazed in this way is targeted by this action again and fails their save, they are petrified and can no longer make saving throws to end the effects of Stone Gaze. The petrification lasts until the creature is restored by Emer's Stone Sacrifice trait, a cure ailment power of 4th order or higher, a \textit{greater restoration} spell, or a similar supernatural effect.

\DndMonsterAction{Envenomed Stone (Recharge 5\textendash 6)}
Emer utters an ancient hex. Each creature within 60 feet of her who is being turned to stone by her Stone Gaze must make a DC 17 Constitution saving throw, taking 44 (8d10) poison damage on a failed save, or half as much damage on a successful one.

\DndMonsterSection{Bonus Actions}
\DndMonsterAction{Mesmerizing Eyes}
The glowing eyes of Emer's snakes gaze at one creature Emer can see within 60 feet of her. If the target can see Emer, they must succeed on a DC 17 Wisdom saving throw or be charmed until the start of the target's next turn. A target charmed in this way must immediately use their reaction, if available, to move up to their speed in a direction of Emer's choice.

\DndMonsterSection{Reactions}
\DndMonsterAction{Blinding Mucus}
When a creature Emer can see within 5 feet of her hits her with a melee attack, Emer spits venom at their eyes. The target must succeed on a DC 17 Dexterity saving throw or take 7 (2d6) acid damage and be blinded until the start of their next turn.

\DndMonsterSection{Legendary Actions}
Emer has three villain actions. She can take each action once during an encounter after an enemy's turn. She can take these actions in any order but can use only one per round.

\begin{DndMonsterLegendaryActions}
\DndMonsterLegendaryAction{Action 1: I See You!}{Emer uses Stone Gaze against each enemy she can see within 60 feet of her.
}
\DndMonsterLegendaryAction{Action 2: Medusa's Evolution}{Emer grows wings from her back, gaining a flying speed of 40 feet, then she moves up to her speed without provoking opportunity attacks. After 1 minute, the wings crumble to dust and her flying speed is lost.
}
\DndMonsterLegendaryAction{Action 3: Stone Puppets}{Emer mentally manipulates the stone in each creature who is being turned to stone by her Stone Gaze. Emer and each of those creatures move up to their speed then make a weapon attack against a creature of Emer's choice (no action required).
}
\end{DndMonsterLegendaryActions}
\end{multicols}
\end{DndMonster}



\begin{DndMonster}[float=t]{Lightbender}
\DndMonsterType{Large monstrosity (Ambusher), U}
\DndMonsterBasics[
armor-class = {15 (natural armor)},
hit-points = {\DndDice{12d10 + 36}},
speed = {50 ft.}
]
\DndMonsterAbilityScores[
 str = 18,
 dex = 14,
 con = 16,
 int = 6,
 wis = 12,
 cha = 8
]
\DndMonsterDetails[
skills = {Perception +4, Stealth +8},
senses = {darkvision 60 ft., passive perception 14},
challenge = 5,
]
\DndMonsterAction{Avoidance}
If the lightbender is subjected to an effect that allows them to make a saving throw to take only half damage, they instead take no damage if they succeed on the saving throw, and only half damage if they fail.

\DndMonsterAction{Pounce}
If the lightbender moves at least 20 feet straight toward a creature and then hits them with a Claw attack on the same turn, that target must succeed on a DC 15 Strength saving throw or be knocked prone. If the target is prone, the lightbender can make one Bite attack against them as a bonus action.

\DndMonsterSection{Actions}
\DndMonsterAction{Multiattack}
The lightbender makes one Bite attack and two Claw attacks, or uses their Tail Whip twice.

\DndMonsterAction{Bite}
\textsl{Melee Weapon Attack:} 7 to hit, reach 5 ft., one target. \textsl{Hit:} 8 (1d8 + 4) piercing damage.

\DndMonsterAction{Claw}
\textsl{Melee Weapon Attack:} 7 to hit, reach 5 ft., one target. \textsl{Hit:} 7 (1d6 + 4) slashing damage.

\DndMonsterAction{Tail Whip}
\textsl{Melee Weapon Attack:} 7 to hit, reach 15 ft., one target. \textsl{Hit:} 7 (1d6 + 4) bludgeoning damage plus 7 (2d6) radiant damage.

\DndMonsterAction{Hypnotic Mane (1/Day)}
The lightbender discharges the light absorbed in their mane into a brilliant, mesmerizing display. Each creature within 15 feet of the lightbender must succeed on a DC 14 Wisdom saving throw or be charmed by the lightbender for 1 minute. While charmed in this way, a creature is incapacitated and has a speed of 0. If a creature charmed in this way takes damage or if someone else uses an action to shake the creature out of their stupor, the condition ends for that creature.

\DndMonsterSection{Reactions}
\DndMonsterAction{Afterimage}
When the lightbender is hit by an attack, they can reveal that the attacker is attacking a past visual imprint of the lightbender. The lightbender appears in an unoccupied space they can see within 30 feet of their imprint, the attack misses, then the imprint disappears. The lightbender can't use this reaction if the attacker relies on senses other than sight, such as blindsight, or if they can perceive illusions as false, as with truesight.

\end{DndMonster}



\begin{DndMonster}[float=t]{Lightbender Companion}
\DndMonsterType{Large monstrosity (Companion), U}
\DndMonsterBasics[
armor-class = {13 plus PB (natural armor)},
hit-points = {\DndDice{7 + seven times caregiver's level (number of d8 Hit Dice equal to their caregiver's level)}},
speed = {50 ft.}
]
\DndMonsterAbilityScores[
 str = 16,
 dex = 14,
 con = 14,
 int = 6,
 wis = 12,
 cha = 8
]
\DndMonsterDetails[
saving-throws = {Str +3 plus PB, Dex +2 plus PB},
skills = {Perception +1 plus PB, Stealth +2 plus PB},
senses = {darkvision 60 ft., passive perception 11 plus PB},
]
\DndMonsterSection{Actions}
\DndMonsterAction{Signature Attack (Bite)}
\textsl{Melee Weapon Attack:} 3 plus PB to hit, reach 5 ft., one target. \textsl{Hit:} 1d6 plus PB piercing damage.

\DndMonsterAction{1st Level: Tail Whip (2 Ferocity)}
The lightbender makes a signature attack. On a hit, the attack deals an extra 0d1 + PB radiant damage to the target, and a different creature the lightbender chooses within 15 feet of them takes PB radiant damage.

\DndMonsterAction{3rd Level: Silent Pounce (5 Ferocity)}
The lightbender teleports 30 feet to an unoccupied space they can see. Before or after teleporting, the lightbender can make a signature attack. If the attack hits, the target is knocked prone.

\DndMonsterAction{5th Level: Hypnotic Mane (8 Ferocity)}
Each creature within 10 feet of the lightbender must succeed on a DC 10 plus PB Wisdom saving throw or be charmed by the lightbender until the end of the lightbender's next turn. While charmed in this way, a creature is incapacitated and has a speed of 0. If a creature charmed in this way takes damage or if someone else uses an action to shake the creature out of their stupor, the condition ends on that creature.

\DndMonsterSection{Reactions}
\DndMonsterAction{Shared Afterimage (Recharges after a Long Rest)}
When the lightbender and their caregiver are within 30 feet of each other and one of them is hit by an attack, the lightbender reveals that both they and their caregiver are past visual imprints. The lightbender and the caregiver each appear in an unoccupied space they can see within 30 feet of their imprints, the attack misses, then the imprints disappear. The lightbender can't use this reaction if the attacker relies on senses other than sight, such as blindsight, or if they can perceive illusions as false, as with truesight.

\end{DndMonster}



\begin{DndMonster}[float=t]{Lightthief}
\DndMonsterType{Tiny fiend (Ambusher), chaotic evil}
\DndMonsterBasics[
armor-class = {15 (natural armor)},
hit-points = {\DndDice{11d4 + 11}},
speed = {20 ft., fly 30 ft.}
]
\DndMonsterAbilityScores[
 str = 6,
 dex = 17,
 con = 12,
 int = 6,
 wis = 16,
 cha = 5
]
\DndMonsterDetails[
skills = {Perception +5, Stealth +5},
senses = {blindsight 60 ft., passive perception 15},
languages = {understands Abyssal, Common, and Infernal but can't speak},
challenge = 2,
]
\DndMonsterSection{Actions}
\DndMonsterAction{Necrotic Drain}
\textsl{Melee Spell Attack:} 5 to hit, reach 5 ft., one target. \textsl{Hit:} 10 (2d6 + 3) necrotic damage, or 13 (3d6 + 3) necrotic damage if the target can't see the lightthief.

\DndMonsterAction{Into the Light (Recharge 6)}
The lightthief teleports inside a light source they can see within 15 feet of them that is emitting bright or dim light. While inside the light source, the lightthief has total cover against attacks and other effects outside the light source. If the light source is extinguished by any means other than the lightthief's Steal Light action, the lightthief is forced out, appears in an unoccupied space of their choice within 5 feet of the light source, and is paralyzed until the end of their next turn.

\DndMonsterAction{Steal Light (Inside Light Source Only)}
The lightthief can cause the light they are inside to erupt in a burst of brilliant radiance. Each creature within 15 feet of the light source must make a DC 13 Constitution saving throw. On a failed save, a creature takes 14 (4d6) radiant damage and loses any darkvision they have for 1 minute. On a successful save, a creature takes half as much damage and doesn't lose their darkvision. After this, the light source extinguishes and the lightthief appears in an unoccupied space of their choice within 5 feet of the light source.

\end{DndMonster}



\begin{DndMonster}[float=t]{Lizardfolk Hunter}
\DndMonsterType{Medium humanoid (lizardfolk, Retainer), A}
\DndMonsterBasics[
armor-class = {15 (medium armor)},
hit-points = {\DndDice{Eight times their level (number of d10 Hit Dice equal to their level)}},
speed = {30 ft., swim 30 ft.}
]
\DndMonsterAbilityScores[
 str = 16,
 dex = 10,
 con = 10,
 int = 10,
 wis = 14,
 cha = 10
]
\DndMonsterDetails[
saving-throws = {Str +PB, Dex +PB, Con +PB, Int +PB, Wis +PB, Cha +PB},
skills = {Athletics +3 plus PB, Nature +0 plus PB, Perception +2 plus PB},
languages = {Common, Draconic},
]
\DndMonsterSection{Actions}
\DndMonsterAction{Signature Attack (Greatsword)}
\textsl{Melee Attack?} 3 plus PB to hit, reach 5 ft., one target. \textsl{Hit:} 2d6 plus PB slashing damage. Beginning at 7th level, the hunter can make this attack twice, instead of once, when they take the Attack action on their turn.

\DndMonsterAction{3rd Level: Wrenching Bite (3/Day)}
As an action, the hunter makes a signature attack against a Large or smaller creature. On a hit, the target is grappled (escape DC 10 plus PB). Until the grapple ends, the target takes PB piercing damage at the end of each of their turns, and the hunter can't use their Wrenching Bite against another target.

\DndMonsterAction{5th Level: Thrashing Tail (3/Day)}
As an action, the hunter makes a signature attack against a Large or smaller creature. On a hit, the target must succeed on a DC 10 plus PB Dexterity saving throw or fall prone. While the hunter is within 5 feet of this prone target, the target can't stand up unless they use an action to make a DC 10 plus PB Strength (Athletics) check. On a success, the target can stand up.

\DndMonsterAction{7th Level: Hypnotic Sleep (1/Day)}
As a bonus action, the hunter drones lullingly. Each creature within 5 feet of the hunter who can hear them must succeed on a DC 10 plus PB Constitution saving throw or fall unconscious for 1 minute. An unconscious target wakes up if they take damage, the hunter becomes incapacitated, the hunter moves more than 5 feet away from the target, or another creature uses an action to rouse the target.

\end{DndMonster}



\begin{DndMonster}[float*=b,width=\textwidth + 8pt]{Lord Syuul}
\begin{multicols}{2}
\DndMonsterType{Medium aberration (Solo), lawful evil}
\DndMonsterBasics[
armor-class = {18 (natural armor)},
hit-points = {\DndDice{26d8 + 104}},
speed = {30 ft., fly 30 ft.}
]
\DndMonsterAbilityScores[
 str = 15,
 dex = 16,
 con = 18,
 int = 22,
 wis = 20,
 cha = 21
]
\DndMonsterDetails[
saving-throws = {Con +9, Int +11, Wis +10},
skills = {Arcana +16, Deception +10, Insight +10, Perception +10, Persuasion +10, Stealth +8},
damage-resistances = {bludgeoning},
condition-immunities = {charmed, frightened},
senses = {darkvision 120 ft., passive perception 20},
languages = {Deep Speech, Undercommon, telepathy 120 ft.},
challenge = 15,
]
\DndMonsterAction{Draw In}
If Lord Syuul fails a saving throw while grappling at least one enemy, he can choose to release all grappled creatures and succeed on the saving throw instead.

\DndMonsterSection{Actions}
\DndMonsterAction{Multiattack}
Lord Syuul makes two Tentacle or two Psionic Repeater attacks then manifests a power or uses Memory Transfer.

\DndMonsterAction{Tentacle}
\textsl{Melee Weapon Attack:} 11 to hit, reach 15 ft., one target. \textsl{Hit:} 28 (4d10 + 6) psychic damage, and if the target is Large or smaller, they are grappled (escape DC 19).

\DndMonsterAction{Psionic Repeater}
\textsl{Ranged Weapon Attack:} 11 to hit, range 150/600 ft., one target. \textsl{Hit:} 22 (3d10 + 6) force damage.

\DndMonsterAction{Brain Overload (1/Day; 6th-Order Power)}
Lord Syuul creates a sudden surge of energy in the mind of a creature he can see within 120 feet of him. The target must make a DC 19 Constitution saving throw, taking 77 (14d10) psychic damage on a failed save, or half as much damage on a successful one. If the target has a brain and is reduced to 0 hit points, their brain explodes, and if they can't survive without their brain, they die.

\DndMonsterAction{Flay (6th-Order Power)}
Lord Syuul shoots a 15-foot cone of psionic energy from his eyes. Each creature in that area must make a DC 19 Intelligence saving throw, taking 35 (10d6) psychic damage on a failed save, or half as much damage on a successful one.

\DndMonsterAction{Memory Transfer}
Lord Syuul psionically plunders the mind of each creature he is grappling. Each target must succeed on a DC 19 Intelligence saving throw or take 44 (8d10) psychic damage. If a target fails their saving throw and doesn't escape the grapple by the end of their next turn, they must choose one of these effects:


\begin{itemize}
\item The target's proficiency bonus drops to 0, they can't form new thoughts or speak, and they have disadvantage on all ability checks, attack rolls, and saving throws. Only the cure ailment power, a \textit{lesser restoration} spell, or a similar supernatural effect can end this effect. Additionally, until Lord Syuul finishes a long rest, he gains a cumulative +5 bonus to damage rolls.
\item The target is charmed by Lord Syuul for 1 hour. While the target is charmed, Lord Syuul can issue the target telepathic commands (no action required), which the target does their best to obey. Each time the target takes damage, they can make a DC 19 Wisdom saving throw, ending the effect on a success.
\end{itemize}

\DndMonsterAction{Guise (3rd-Order Power)}
Lord Syuul projects a psionic image over his body, transforming his appearance for 1 hour into that of a Medium creature he has seen. When he manifests this power, he can also change the appearance of any equipment he carries.

The changes wrought by this power fail to hold up to physical inspection. A creature can use an action to inspect Lord Syuul's appearance and make a DC 19 Intelligence (Investigation) check, noticing the image is a projection on a success.

\DndMonsterSection{Bonus Actions}
\DndMonsterAction{Grappling Jaunt}
Lord Syuul and each creature grappled by him teleport up to 30 feet to an unoccupied space he can see.

\DndMonsterSection{Reactions}
\DndMonsterAction{Brain Drain}
When a creature grappled by Lord Syuul makes a saving throw against Memory Transfer, Lord Syuul momentarily weakens the creature and the creature has disadvantage on the saving throw.

\DndMonsterSection{Legendary Actions}
Lord Syuul has three villain actions. He can take each action once during an encounter after an enemy's turn. He can take these actions in any order but can use only one per round.

\begin{DndMonsterLegendaryActions}
\DndMonsterLegendaryAction{Action 1: Mindblind}{Each enemy Lord Syuul can see within 60 feet of him must make a DC 19 Wisdom saving throw. On a failed save, a creature can't see creatures other than Lord Syuul for 1 minute (save ends at end of turn).
}
\DndMonsterLegendaryAction{Action 2: Was I Ever Here?}{Lord Syuul becomes invisible, teleports 60 feet to an unoccupied space he can see, and can take the Hide action. At the same time, an illusory psionic image of Lord Syuul appears in the space he left, gesturing, speaking, and behaving as he chooses (no action required). A creature who touches the image for the first time on a turn or makes a melee attack against it must succeed on a DC 19 Constitution saving throw or take 21 (6d6) psychic damage because things can pass through it. Lord Syuul's invisibility ends and his image disappears at the end of his next turn.
}
\DndMonsterLegendaryAction{Action 3: Mindshatter}{Each enemy Lord Syuul can see within 60 feet of him must make a DC 19 Wisdom saving throw. On a failed save, a target takes 55 (10d10) psychic damage and is dazed for 1 minute (save ends at end of turn). On a successful save, a target takes half as much damage and isn't dazed.
}
\end{DndMonsterLegendaryActions}
\end{multicols}
\end{DndMonster}



\begin{DndMonster}[float=t]{Mimic}
\DndMonsterType{Medium monstrosity (shapechanger, Ambusher), N}
\DndMonsterBasics[
armor-class = {11},
hit-points = {\DndDice{10d8 + 20}},
speed = {15 ft., climb 15 ft.}
]
\DndMonsterAbilityScores[
 str = 17,
 dex = 12,
 con = 15,
 int = 5,
 wis = 12,
 cha = 10
]
\DndMonsterDetails[
skills = {Deception +2, Stealth +5},
damage-immunities = {acid; bludgeoning damage from falling},
condition-immunities = {prone},
senses = {darkvision 60 ft., passive perception 11},
challenge = 2,
]
\DndMonsterAction{False Appearance (Object Form Only)}
While the mimic remains motionless, they are indistinguishable from an ordinary object.

\DndMonsterAction{Mimicry}
The mimic can imitate any sounds they have heard, including voices. A creature who hears the sounds can tell they are imitations with a successful DC 12 Wisdom (Insight) check.

\DndMonsterSection{Actions}
\DndMonsterAction{Pseudopod}
\textsl{Melee Weapon Attack:} 5 to hit, reach 10 ft., one target. \textsl{Hit:} 7 (1d8 + 3) bludgeoning damage. If the target is a Large or smaller creature, they are grappled (escape DC 13).

\DndMonsterAction{Gnashing Maw}
\textsl{Melee Weapon Attack:} 5 to hit, reach 5 ft., one target. \textsl{Hit:} 7 (1d8 + 3) piercing damage plus 3 (1d6) acid damage. If the target is grappled by the mimic, they take an extra 3 (1d6) acid damage.

\DndMonsterAction{Shapechanger}
The mimic polymorphs into a Small, Medium, or Large object or back into their true amorphous form. Other than size, the mimic's statistics are the same in each form. Anything they are wearing or carrying isn't transformed. They revert to their true form if they die.

\DndMonsterAction{Envelop}
While grappling a creature, the mimic uses their Shapechanger action to polymorph into an object that is the same size or bigger than the grappled creature (if the mimic isn't already in such a form). The creature is then enveloped by the mimic. An enveloped creature is no longer grappled, but they can't breathe, are blinded and restrained, have total cover against attacks and other effects outside the mimic, and take 10 (3d6) acid damage at the start of each of the mimic's turns. When the mimic moves, the enveloped creature moves with them. The mimic can envelop one creature at a time.

An enveloped creature can use their action to try to escape by making a DC 13 Strength check. On a successful check, the creature is no longer enveloped and is shunted to an unoccupied space of their choice within 5 feet of the mimic. The creature automatically escapes if the mimic uses Shapechanger or Panic Shift, becomes incapacitated, or dies.

\DndMonsterSection{Reactions}
\DndMonsterAction{Panic Shift (1/Day)}
When the mimic takes damage, they release any grappled or enveloped creatures and then polymorph into a Small dense object with spikes, such as an iron ball or helm. While in this form, the mimic's walking and climbing speeds are doubled, they have resistance to bludgeoning, piercing, and slashing damage from mundane attacks, they can only take the Dash, Disengage, or Shapechanger action, and a creature who touches the mimic or hits them with a melee attack while within 5 feet of them takes 10 (3d6) piercing damage. This effect ends when the mimic dies or uses their Shapechanger action to polymorph into another form.

\end{DndMonster}



\begin{DndMonster}[float=t]{Mindkiller}
\DndMonsterType{Small aberration (Ambusher), lawful evil}
\DndMonsterBasics[
armor-class = {12},
hit-points = {\DndDice{10d6 + 10}},
speed = {10 ft., fly 30 ft.}
]
\DndMonsterAbilityScores[
 str = 14,
 dex = 15,
 con = 12,
 int = 16,
 wis = 15,
 cha = 14
]
\DndMonsterDetails[
skills = {Deception +4, Stealth +4},
senses = {darkvision 120 ft., passive perception 12},
languages = {Deep Speech, Undercommon, telepathy 120 ft.},
challenge = 2,
]
\DndMonsterAction{Amorphous}
The mindkiller can move through a space as narrow as 1 inch wide without squeezing.

\DndMonsterAction{Psionic Immunity}
The mindkiller is unaffected by psionic powers manifested by voiceless talkers unless the mindkiller wishes to be.

\DndMonsterSection{Actions}
\DndMonsterAction{Claws}
\textsl{Melee Weapon Attack:} 4 to hit, reach 5 ft., one target. \textsl{Hit:} 7 (2d4 + 2) slashing damage plus 7 (2d6) psychic damage, and if the target is a Medium or smaller creature, they are grappled (escape DC 12).

\DndMonsterAction{Concealing Strike}
\textsl{Ranged Psionic Attack:} 5 to hit, range 30 ft., one creature. \textsl{Hit:} 13 (3d6 + 3) psychic damage, and the mindkiller is invisible to the target until the end of the mindkiller's next turn.

\DndMonsterAction{Mindwipe}
A Humanoid grappled by the mindkiller must succeed on a DC 12 Strength saving throw or the mindkiller enters the Humanoid's body. While inside a Humanoid, the mindkiller has total cover against attacks and other effects originating outside the Humanoid, and the only action the mindkiller can take is to leave the body, exiting in an unoccupied space within 5 feet of the body.

When a Humanoid ends their turn with the mindkiller inside of them, they must succeed on a DC 13 Constitution saving throw or take 10 necrotic damage. If the Humanoid is reduced to 0 hit points, they die and the mindkiller takes over the body, which regains hit points equal to the Humanoid's hit point maximum. The mindkiller retains their Intelligence, Wisdom, and Charisma scores, their understanding of Deep Speech and Undercommon, and their telepathy. They otherwise adopt the target's statistics and can take the actions the creature could take. They know everything the creature knew, including spells, class features, traits, and languages. If the body is reduced to 0 hit points after the mindkiller takes control, the mindkiller must leave it.

A creature wielding a sharp tool or weapon within reach of a Humanoid host or body with a mindkiller inside can use an action to attempt to remove the mindkiller, making an attack roll against the Humanoid's AC if the host is unwilling. On a hit, the creature deals 11 (2d10) slashing damage to the host and must make a DC 15 Wisdom (Medicine) check. On a successful check, the creature cuts the mindkiller out of the host. On a faiIed check, if this slashing damage reduced the host to 0 hit points, the mindkiller kills the host; otherwise, there is no effect.

\end{DndMonster}



\begin{DndMonster}[float=t]{Mindkiller Whelp}
\DndMonsterType{Small aberration (Minion), lawful evil}
\DndMonsterBasics[
armor-class = {12},
hit-points = {\DndDice{9}},
speed = {10 ft., fly 30 ft.}
]
\DndMonsterAbilityScores[
 str = 10,
 dex = 14,
 con = 11,
 int = 14,
 wis = 12,
 cha = 12
]
\DndMonsterDetails[
senses = {darkvision 120 ft., passive perception 11},
languages = {Deep Speech, Undercommon, telepathy 120 ft.},
challenge = 2,
]
\DndMonsterAction{Amorphous}
The whelp can move through a space as narrow as 1 inch wide without squeezing.

\DndMonsterAction{Minion}
If the whelp takes damage from an attack or as the result of a failed saving throw, their hit points are reduced to 0. If the whelp takes damage from another effect, they die if the damage equals or exceeds their hit point maximum; otherwise they take no damage.

\DndMonsterAction{Psionic Immunity}
The whelp is unaffected by psionic powers manifested by voiceless talkers unless the whelp wishes to be.

\DndMonsterAction{Resistance Drain}
When an enemy within 5 feet of three or more whelps makes a saving throw against a power or psionic effect, the enemy takes a penalty to the save equal to the number of whelps within 5 feet of them.

\DndMonsterSection{Actions}
\DndMonsterAction{Claws (Group Attack)}
\textsl{Melee Weapon Attack:} 4 to hit, reach 5 ft., one target. \textsl{Hit:} 2 slashing damage.

\end{DndMonster}



\begin{DndMonster}[float=t]{Minotaur Devastator}
\DndMonsterType{Large monstrosity (minotaur, Retainer), A}
\DndMonsterBasics[
armor-class = {13 (light armor)},
hit-points = {\DndDice{Eight times their level (number of d10 Hit Dice equal to their level)}},
speed = {30 ft.}
]
\DndMonsterAbilityScores[
 str = 16,
 dex = 10,
 con = 14,
 int = 10,
 wis = 10,
 cha = 10
]
\DndMonsterDetails[
saving-throws = {Str +PB, Dex +PB, Con +PB, Int +PB, Wis +PB, Cha +PB},
skills = {Athletics +3 plus PB, Survival +0 plus PB},
senses = {darkvision 60 ft., passive perception 10},
languages = {Common},
]
\DndMonsterSection{Actions}
\DndMonsterAction{Signature Attack (Goring Horns)}
\textsl{Melee Weapon Attack:} 3 plus PB to hit, reach 5 ft., one target. \textsl{Hit:} 1d10 plus PB piercing damage. Beginning at 7th level, the devastator can make this attack twice, instead of once, when they take the Attack action on their turn.

\DndMonsterAction{3rd Level: Furious Charge (3/Day)}
As an action, the devastator moves up to their speed. At the end of this movement, they can make a signature attack. If that attack hits a Large or smaller creature, the target is pushed up to 10 feet away from the devastator, and the devastator and their allies have advantage on attack rolls against the target until the end of the devastator's next turn.

\DndMonsterAction{5th Level: Fearsome Gore (3/Day)}
As an action, the devastator makes a signature attack against a creature within 5 feet of them. On a hit, the attack deals an extra PBd10 piercing damage, and the target must succeed on a DC 10 plus PB Wisdom saving throw or be frightened of the devastator until the end of the devastator's next turn.

\DndMonsterAction{7th Level: Retaliatory Charge (1/Day)}
When a creature the devastator can see within 40 feet of them deals damage to the devastator, the devastator can use a reaction to move up to their speed toward the creature. If the devastator ends this movement within 5 feet of the creature, they can make two signature attacks against the creature with a +PB bonus to those attack and damage rolls.

\end{DndMonster}



\begin{DndMonster}[float=t]{Mohler}
\DndMonsterType{Small beast (Controller), U}
\DndMonsterBasics[
armor-class = {12},
hit-points = {\DndDice{2d6 + 2}},
speed = {20 ft., burrow 40 ft.}
]
\DndMonsterAbilityScores[
 str = 7,
 dex = 14,
 con = 12,
 int = 2,
 wis = 12,
 cha = 5
]
\DndMonsterDetails[
senses = {blindsight 20 ft., tremorsense 40 ft., passive perception 11},
challenge = 1/4,
]
\DndMonsterSection{Actions}
\DndMonsterAction{Claw}
\textsl{Melee Weapon Attack:} 4 to hit, reach 5 ft., one target. \textsl{Hit:} 5 (1d6 + 2) slashing damage.

\DndMonsterAction{Earth Bump}
As the mohler burrows under a Large or smaller creature on the ground within 5 feet of them, the mohler shifts the ground. The creature above them must succeed on a DC 12 Dexterity saving throw or be knocked prone.

\DndMonsterAction{Ground Grinder}
The mohler burrows up to their burrowing speed. Ground within 5 feet of where the mohler burrows becomes difficult terrain.

\DndMonsterSection{Reactions}
\DndMonsterAction{Zip}
If the mohler fails a Dexterity saving throw while burrowing, they succeed instead.

\end{DndMonster}



\begin{DndMonster}[float=t]{Mohler Companion}
\DndMonsterType{Small beast (Companion), U}
\DndMonsterBasics[
armor-class = {13 plus PB (natural armor)},
hit-points = {\DndDice{7 + seven times caregiver's level (number of d8 Hit Dice equal to their caregiver's level)}},
speed = {20 ft., burrow 40 ft.}
]
\DndMonsterAbilityScores[
 str = 7,
 dex = 16,
 con = 14,
 int = 2,
 wis = 12,
 cha = 5
]
\DndMonsterDetails[
saving-throws = {Dex +3 plus PB, Con +2 plus PB},
skills = {Acrobatics +3 plus PB, Stealth +3 plus PB},
senses = {blindsight 20 ft., tremorsense 40 ft., passive perception 11},
]
\DndMonsterSection{Actions}
\DndMonsterAction{Signature Attack (Claw)}
\textsl{Melee Weapon Attack:} 3 plus PB to hit, reach 5 ft., one target. \textsl{Hit:} 1d6 plus PB slashing damage.

\DndMonsterAction{1st Level: Earth Bump (2 Ferocity)}
The mohler burrows up to their speed and chooses one Large or smaller creature they come within 5 feet of during the move. The target must succeed on a DC 10 plus PB Dexterity saving throw or be knocked prone.

\DndMonsterAction{3rd Level: Sinkhole (5 Ferocity)}
The mohler burrows up to their speed and targets a creature they can sense on the surface within 10 feet of them. The target must succeed on a DC 10 plus PB Strength saving throw or be restrained by the ground until the start of the mohler's next turn.

\DndMonsterAction{5th Level: Terranova (8 Ferocity)}
Each creature on the ground within 10 feet of the mohler must succeed on a DC 10 plus PB Strength saving throw or take 1d6 plus PB bludgeoning damage and be knocked prone. A creature who fails this saving throw by 5 or more is restrained by the ground (save ends at end of turn). A creature can use their action to free themself or another creature within their reach, ending the effect.

\DndMonsterSection{Reactions}
\DndMonsterAction{Quick Pit (Recharges after a Long Rest)}
If the mohler's caregiver fails a Dexterity saving throw while they are within 10 feet of the mohler and the mohler isn't incapacitated or restrained, the mohler pulls the caregiver out of danger. The caregiver succeeds on the saving throw instead.

\end{DndMonster}



\begin{DndMonster}[float=t]{Mummy}
\DndMonsterType{Medium undead (Brute), lawful evil}
\DndMonsterBasics[
armor-class = {12 (natural armor)},
hit-points = {\DndDice{8d8 + 24}},
speed = {30 ft.}
]
\DndMonsterAbilityScores[
 str = 16,
 dex = 8,
 con = 16,
 int = 10,
 wis = 14,
 cha = 10
]
\DndMonsterDetails[
skills = {Arcana +2, History +2, Religion +2},
damage-immunities = {necrotic, poison},
condition-immunities = {blinded, charmed, deafened, exhaustion, frightened, poisoned},
senses = {blindsight 60 ft., passive perception 12},
languages = {the languages they knew in life},
challenge = 3,
]
\DndMonsterAction{Flammable}
Whenever the mummy takes fire damage, they take an extra 5 (1d10) fire damage.

\DndMonsterAction{Mummy Dust}
Whenever the mummy takes piercing or slashing damage, each creature within 5 feet of them takes 4 (1d8) poison damage.

\DndMonsterSection{Actions}
\DndMonsterAction{Multiattack}
The mummy makes two Desiccating Slam attacks and uses Guardian Curse.

\DndMonsterAction{Desiccating Slam}
\textsl{Melee Weapon Attack:} 5 to hit, reach 5 ft., one target. \textsl{Hit:} 10 (2d6 + 3) bludgeoning damage. Until the start of the mummy's next turn, the target's speed is reduced by 10 feet and the target can't regain hit points.

\DndMonsterAction{Guardian Curse}
The mummy calls down a scourge on one creature they can see within 60 feet of them. The target must succeed on a DC 13 Wisdom saving throw or be cursed. While cursed in this way, whenever the target takes damage, they take an extra 3 (1d6) necrotic damage that can't be reduced in any way. The curse lasts until removed by a cure ailment power, a \textit{remove curse} spell, or a similar supernatural effect. A target who succeeds on their saving throw is immune to the Guardian Curse of all mummies for the next 24 hours.

\end{DndMonster}


\clearpage

\begin{DndMonster}[float=t]{Night Hag}
\DndMonsterType{Medium fey (Controller), neutral evil}
\DndMonsterBasics[
armor-class = {17 (natural armor)},
hit-points = {\DndDice{13d8 + 39}},
speed = {30 ft., fly 40 ft.}
]
\DndMonsterAbilityScores[
 str = 15,
 dex = 18,
 con = 16,
 int = 16,
 wis = 17,
 cha = 19
]
\DndMonsterDetails[
skills = {Deception +10, Insight +6, Perception +6, Stealth +7 (+17 with pass without trace)},
damage-resistances = {fire; necrotic; bludgeoning, piercing, slashing from mundane attacks},
condition-immunities = {charmed},
senses = {darkvision 120 ft., passive perception 16},
languages = {Common, Infernal, Sylvan},
challenge = 5,
]
\DndMonsterAction{Supernatural Resistance}
The hag has advantage on saving throws against powers, spells, and other supernatural effects.

\DndMonsterAction{Spellcasting (Utility)}
In addition to any other spells in this stat block, the hag can cast the following spells, using Charisma as the spellcasting ability (spell save DC 15):
\begin{DndMonsterSpells}
  \DndInnateSpellLevel{alter self\textsuperscript{Casting Times}, \textit{detect magic}\textsuperscript{Casting Times}, \textit{pass without trace}\textsuperscript{Casting Times}}
  \DndInnateSpellLevel[]{scrying\textsuperscript{Casting Times}}
  \DndInnateSpellLevel[3]{dream\textsuperscript{Casting Times}, \textit{legend lore}\textsuperscript{Casting Times}, \textit{mirage arcane}\textsuperscript{Casting Times}}
\end{DndMonsterSpells}
\DndMonsterSection{Actions}
\DndMonsterAction{Multiattack}
The hag makes two Dream Claw attacks.

\DndMonsterAction{Dream Claw}
\textsl{Melee Weapon Attack:} 7 to hit, reach 5 ft., one target. \textsl{Hit:} 9 (1d10 + 4) slashing damage plus 3 (1d6) psychic damage. If a target is hit by this attack more than once on a turn, they must make a DC 15 Constitution saving throw. On a failed save, the target falls unconscious for 1 minute (save ends at end of turn), until the target takes damage, or until another creature who can reach the target uses their action to wake them. Creatures who can't be supernaturally put to sleep automatically succeed on this saving throw.

\DndMonsterAction{Vicious Visions (1/Day; 6th-Level Spell)}
The hag releases magic dream dust in a 60-foot cone. Each creature must succeed on a DC 15 Wisdom saving throw or be charmed for 1 minute or until they deal damage to another creature. While charmed, the creature sees all allies as enemies and all enemies as allies.

\DndMonsterAction{Tiny Time}
The hag magically shrinks down to a height of 1 inch, becoming Tiny along with any items they wear or carry. While Tiny, the hag has advantage on Dexterity (Stealth) checks made to hide, and they can use a bonus action to enter the head of an unconscious Humanoid they can reach or to grow back to the hag's normal size. The hag's game statistics otherwise remain the same.

While the hag is inside a Humanoid's head, that Humanoid is considered haunted by the hag, and the hag has total cover against attacks and other effects outside the head. The hag can speak telepathically to the haunted Humanoid, but the Humanoid is unaware of the source of this telepathic speech. The haunted Humanoid can't gain any benefit from rest unless the hag allows it, and when the Humanoid attempts to sleep, the hag can fill the Humanoid's mind with nightmarish visions.

The hag is expelled from the head and grows back to their normal size if the haunted Humanoid is affected by a \textit{dispel evil and good} spell or similar supernatural effects, or if the hag chooses to exit the head as a bonus action.

\DndMonsterSection{Bonus Actions}
\DndMonsterAction{Sleepwalk}
The hag moves one unconscious creature they can see within 60 feet of them up to 30 feet horizontally.

\DndMonsterSection{Reactions}
\DndMonsterAction{Aren't You Tired?}
When a creature within 5 feet of the hag hits them with an attack, the hag spits sleep dust at the attacker. The attacker must make a DC 15 Constitution saving throw. On a failed save, until the end of the hag's next turn, the attacker feels tremendously sleepy, they can't take reactions, and they have disadvantage on attack rolls and saving throws. Creatures who can't be supernaturally put to sleep are immune to this effect.

\end{DndMonster}



\begin{DndMonster}[float=t]{Oracle of Storms}
\DndMonsterType{Huge elemental (air, water, Artillery), A}
\DndMonsterBasics[
armor-class = {17 (natural armor)},
hit-points = {\DndDice{18d12 + 72}},
speed = {40 ft., fly 40 ft. \textit{(hover)}}
]
\DndMonsterAbilityScores[
 str = 17,
 dex = 21,
 con = 18,
 int = 12,
 wis = 18,
 cha = 17
]
\DndMonsterDetails[
saving-throws = {Dex +10, Cha +8},
skills = {History +11, Nature +11, Perception +9, Performance +8},
damage-immunities = {lightning, poison, thunder},
condition-immunities = {dazed, grappled, paralyzed, petrified, poisoned, prone, restrained},
senses = {blindsight 60 ft., passive perception 19},
languages = {Aquan, Auran, Common},
challenge = 16,
]
\DndMonsterAction{Turbulent Escalation}
When the oracle starts their turn outdoors while not incapacitated, they can stoke a thunderstorm (no action required). For 1 minute or until the oracle uses Turbulent Escalation again, the weather in a 1-mile radius centered on the oracle changes as follows:


\begin{itemize}
\item If the sky is cloudless, it becomes cloudy.
\item If the sky is cloudy, it begins to rain or snow (oracle's choice).
\item If it's raining or snowing, the weather becomes stormy.
\end{itemize}

\DndMonsterAction{Spellcasting (Utility)}
In addition to any other spells in this stat block, the oracle can cast the following spells, using Wisdom as the spellcasting ability (spell save DC 17):
\begin{DndMonsterSpells}
  \DndInnateSpellLevel[3]{commune\textsuperscript{Casting Times}, \textit{commune with nature}\textsuperscript{Casting Times}, \textit{legend lore}\textsuperscript{Casting Times}, \textit{scrying}\textsuperscript{Casting Times}}
\end{DndMonsterSpells}
\DndMonsterSection{Actions}
\DndMonsterAction{Multiattack}
The oracle uses Tempest Bolt twice. They can replace one use with a Cyclone Staff attack.

\DndMonsterAction{Cyclone Staff}
\textsl{Melee Weapon Attack:} 10 to hit, reach 10 ft., one target. \textsl{Hit:} 23 (4d8 + 5) bludgeoning damage plus 22 (4d10) thunder damage, and the oracle moves the target up to 10 feet in any direction.

\DndMonsterAction{Tempest Bolt}
A bolt of lightning streaks from the oracle to strike one creature they can see within 90 feet of them, or within 300 feet of them if the target is outdoors and the weather is stormy. The target must make a DC 18 Dexterity saving throw, taking 31 (7d8) lightning damage on a failed save, or half as much damage on a successful one. If the target is outdoors and the weather is stormy, they have disadvantage on the saving throw.

\DndMonsterSection{Bonus Actions}
\DndMonsterAction{Convocation of Storms (2/Day)}
The oracle imbues the power of storm in themself or another Elemental they can see within 30 feet of them, granting one of the following effects of that Elemental's choice:


\begin{description}
\item \textbf{Lightning Within.} The Elemental becomes charged with lightning. The next time the Elemental hits a target with an attack within the next minute, the attack deals an extra 33 (6d10) lightning damage.
\item \textbf{Thunder Burst.} A concussive shock wave erupts from the Elemental toward four creatures they can see within 60 feet of them. Each target must succeed on a DC 17 Constitution saving throw or take 22 (4d10) thunder damage, be pushed up to 10 feet away from the Elemental, and be deafened until the end of the target's next turn.
\end{description}

\end{DndMonster}



\begin{DndMonster}[float=t]{Orc Blacksmith}
\DndMonsterType{Medium humanoid (orc, Retainer), A}
\DndMonsterBasics[
armor-class = {15 (medium armor)},
hit-points = {\DndDice{Eight times their level (number of d10 Hit Dice equal to their level)}},
speed = {35 ft.}
]
\DndMonsterAbilityScores[
 str = 16,
 dex = 10,
 con = 12,
 int = 10,
 wis = 12,
 cha = 12
]
\DndMonsterDetails[
saving-throws = {Str +PB, Dex +PB, Con +PB, Int +PB, Wis +PB, Cha +PB},
skills = {Athletics +3 plus PB, Perception +1 plus PB},
senses = {darkvision 60 ft., passive perception 11 plus PB},
languages = {Common, Orc},
]
\DndMonsterSection{Actions}
\DndMonsterAction{Signature Attack (Warhammer)}
\textsl{Melee Weapon Attack:} 3 plus PB to hit, reach 5 ft., one target. \textsl{Hit:} 1d10 plus PB bludgeoning damage. Beginning at 7th level, the blacksmith can make this attack twice, instead of once, when they take the Attack action on their turn.

\DndMonsterAction{3rd Level: Relentless Rush (3/Day)}
When the blacksmith isn't incapacitated and they take damage but aren't killed outright, they can make an attack against an enemy (no action required) before the hit point reduction is resolved. If the attack hits, the blacksmith regains twice their PB hit points.

\DndMonsterAction{5th Level: Reinforce Armor (3/Day)}
As a bonus action, the blacksmith touches a piece of armor and fortifies it. For the next 10 minutes, a creature wearing the armor gains a bonus to their AC equal to half the blacksmith's PB.

\DndMonsterAction{7th Level: Fortify Weapon (3/Day)}
As a bonus action, the blacksmith touches a weapon and polishes it. For 1 minute, the weapon is magical and deals an extra 0d1 + PB damage.

\end{DndMonster}



\begin{DndMonster}[float=t]{Orc Blitzer}
\DndMonsterType{Medium humanoid (orc, Minion), A}
\DndMonsterBasics[
armor-class = {13 (studded leather armor)},
hit-points = {\DndDice{7}},
speed = {35 ft.}
]
\DndMonsterAbilityScores[
 str = 15,
 dex = 12,
 con = 12,
 int = 10,
 wis = 11,
 cha = 11
]
\DndMonsterDetails[
senses = {darkvision 60 ft., passive perception 10},
languages = {Common, Orc},
challenge = 1/2,
]
\DndMonsterAction{Gnashing Horde}
If an enemy the blitzer can see starts their turn within 5 feet of three or more blitzers, the enemy takes 1 piercing damage for each blitzer within 5 feet of them.

\DndMonsterAction{Minion}
If the blitzer takes damage from an attack or as the result of a failed saving throw, their hit points are reduced to 0. If the blitzer takes damage from another effect, they die if the damage equals or exceeds their hit point maximum; otherwise they take no damage.

\DndMonsterAction{Relentless Minion}
When the blitzer is reduced to 0 hit points while they aren't incapacitated, they can deal 1 piercing damage to an enemy within 5 feet of them.

\DndMonsterSection{Actions}
\DndMonsterAction{Spear (Group Attack)}
\textsl{Melee or Ranged Weapon Attack:} 4 to hit, reach 5 ft. or range 20/60 ft., one target. \textsl{Hit:} 1 piercing damage.

\end{DndMonster}



\begin{DndMonster}[float=t]{Orc Bloodrunner}
\DndMonsterType{Medium humanoid (orc, Skirmisher), A}
\DndMonsterBasics[
armor-class = {15 (shield)},
hit-points = {\DndDice{2d8 + 6}},
speed = {60 ft.}
]
\DndMonsterAbilityScores[
 str = 14,
 dex = 16,
 con = 16,
 int = 11,
 wis = 9,
 cha = 10
]
\DndMonsterDetails[
saving-throws = {Dex +5},
skills = {Acrobatics +5, Athletics +4},
senses = {darkvision 60 ft., passive perception 9},
languages = {Common, Orc},
challenge = 1/2,
]
\DndMonsterAction{Relentless (1/Turn)}
When the bloodrunner isn't incapacitated and they are reduced to 0 hit points but not killed outright, they can make an attack against an enemy (no action required) before the hit point reduction is resolved. If the attack hits and its damage reduces the target to 0 hit points, the bloodrunner drops to 1 hit point instead of 0 hit points.

\DndMonsterAction{Unimpeded}
The bloodrunner can occupy a prone creature's space. The first time the bloodrunner enters a prone creature's space on the bloodrunner's turn, the prone creature takes 3 (1d6) bludgeoning damage. When a prone creature in the bloodrunner's space stands up, the creature enters the nearest unoccupied space of their choice.

\DndMonsterSection{Actions}
\DndMonsterAction{Spiked Shield}
\textsl{Melee Weapon Attack:} 5 to hit, reach 5 ft., one target. \textsl{Hit:} 6 (1d6 + 3) piercing damage. If the bloodrunner moves at least 10 feet before hitting a Large or smaller creature with this attack, the bloodrunner can choose to knock the target prone or push them up to 10 feet away.

\DndMonsterAction{Javelin}
\textsl{Melee or Ranged Weapon Attack:} 4 to hit, reach 5 ft. or range 30/120 ft., one target. \textsl{Hit:} 5 (1d6 + 2) piercing damage.

\DndMonsterSection{Reactions}
\DndMonsterAction{Juke}
When the bloodrunner makes a Dexterity saving throw, they can gain advantage on the save. To do so, the bloodrunner can't be prone, grappled, or restrained.

\end{DndMonster}



\begin{DndMonster}[float=t]{Orc Conduit}
\DndMonsterType{Medium humanoid (orc, Artillery), A}
\DndMonsterBasics[
armor-class = {13 (studded leather armor)},
hit-points = {\DndDice{4d8 + 12}},
speed = {35 ft.}
]
\DndMonsterAbilityScores[
 str = 14,
 dex = 12,
 con = 16,
 int = 10,
 wis = 12,
 cha = 16
]
\DndMonsterDetails[
skills = {Perception +3},
damage-resistances = {determined by the conduit's affinity trait},
senses = {darkvision 60 ft., passive perception 13},
languages = {Common, Orc},
challenge = 1,
]
\DndMonsterAction{Affinity}
The conduit has a lifelong affinity for one of the following damage types: cold, fire, or lightning. The chosen type determines the conduit's damage resistance and the damage of their Elemental Discharge and Power Burst actions.

\DndMonsterAction{Relentless (1/Turn)}
When the conduit isn't incapacitated and they are reduced to 0 hit points but not killed outright, they can make an attack against an enemy (no action required) before the hit point reduction is resolved. If the attack hits and its damage reduces the target to 0 hit points, the conduit drops to 1 hit point instead of 0 hit points.

\DndMonsterSection{Actions}
\DndMonsterAction{Elemental Discharge (Cantrip)}
\textsl{Melee or Ranged Spell Attack:} 5 to hit, reach 5 ft. or range 120 ft., one target. \textsl{Hit:} 9 (1d12 + 3) damage of the type determined by the conduit's Affinity trait.

\DndMonsterAction{Power Burst (2/Day; 1st-Level Spell)}
The conduit unleashes explosive energy in a 15-foot cone. The ground in that area becomes difficult terrain, and each creature in that area must make a DC 13 Dexterity saving throw. On a failed save, a creature takes 10 (3d6) damage of the type determined by the conduit's Affinity trait, and they are pushed 10 feet away from the conduit. On a successful save, a creature takes half as much damage and isn't pushed.

\DndMonsterSection{Bonus Actions}
\DndMonsterAction{Rush}
The conduit takes the Dash action.

\end{DndMonster}



\begin{DndMonster}[float=t]{Orc Fury}
\DndMonsterType{Medium humanoid (orc, Brute), A}
\DndMonsterBasics[
armor-class = {13 (studded leather armor)},
hit-points = {\DndDice{3d8 + 9}},
speed = {35 ft.}
]
\DndMonsterAbilityScores[
 str = 17,
 dex = 12,
 con = 16,
 int = 9,
 wis = 11,
 cha = 12
]
\DndMonsterDetails[
saving-throws = {Str +5},
skills = {Athletics +5},
senses = {darkvision 60 ft., passive perception 10},
languages = {Common, Orc},
challenge = 1/2,
]
\DndMonsterAction{Relentless (1/Turn)}
When the fury isn't incapacitated and they are reduced to 0 hit points but not killed outright, they can make an attack against an enemy (no action required) before the hit point reduction is resolved. If the attack hits and its damage reduces the target to 0 hit points, the fury drops to 1 hit point instead of 0 hit points.

\DndMonsterSection{Actions}
\DndMonsterAction{Haymaker Greataxe}
\textsl{Melee Weapon Attack:} 5 to hit, reach 5 ft., one target. \textsl{Hit:} 9 (1d12 + 3) slashing damage, or 12 (1d12 + 6) slashing damage if the target is prone. If the target is a Medium or smaller creature, they must succeed on a DC 13 Strength saving throw or be knocked prone.

\DndMonsterAction{Javelin}
\textsl{Melee or Ranged Weapon Attack:} 5 to hit, reach 5 ft. or range 30/120 ft., one target. \textsl{Hit:} 6 (1d6 + 3) piercing damage.

\DndMonsterSection{Bonus Actions}
\DndMonsterAction{Healing Rally (1/Day)}
The fury regains 5 hit points.

\end{DndMonster}



\begin{DndMonster}[float=t]{Orc Garroter}
\DndMonsterType{Medium humanoid (orc, Ambusher), A}
\DndMonsterBasics[
armor-class = {14 (studded leather armor)},
hit-points = {\DndDice{4d8 + 12}},
speed = {35 ft.}
]
\DndMonsterAbilityScores[
 str = 15,
 dex = 14,
 con = 16,
 int = 10,
 wis = 12,
 cha = 9
]
\DndMonsterDetails[
saving-throws = {Dex +4},
skills = {Athletics +4, Stealth +6},
senses = {darkvision 60 ft., passive perception 11},
languages = {Common, Orc},
challenge = 1,
]
\DndMonsterAction{Relentless (1/Turn)}
When the garroter isn't incapacitated and they are reduced to 0 hit points but not killed outright, they can make an attack against an enemy (no action required) before the hit point reduction is resolved. If the attack hits and its damage reduces the target to 0 hit points, the garroter drops to 1 hit point instead of 0 hit points.

\DndMonsterSection{Actions}
\DndMonsterAction{Strangle}
\textsl{Melee Weapon Attack:} 4 to hit, reach 5 ft., one target. \textsl{Hit:} 5 (1d6 + 2) bludgeoning damage, and the target is grappled (escape DC 12). Until the grapple ends, the target can't speak intelligibly or cast spells with verbal components. If the attack was made with advantage, the target must make a DC 12 Constitution saving throw. On a failed save, the target falls unconscious for 1 minute (save ends at start of turn) or until they take damage. A creature can use an action to shake another creature awake who falls unconscious in this way.

\DndMonsterAction{Dagger}
\textsl{Melee or Ranged Weapon Attack:} 4 to hit, reach 5 ft. or range 20/60 ft., one target. \textsl{Hit:} 4 (1d4 + 2) piercing damage. If the attack was made with advantage, the target takes an extra 7 (2d6) piercing damage.

\DndMonsterSection{Bonus Actions}
\DndMonsterAction{Cloak (1/Day)}
The garroter magically turns invisible for 1 minute. Any equipment the garroter wears or carries is invisible with them. While invisible, the garroter's movements make no sound. If the garroter hits a creature with an attack, casts a spell, or takes damage, the effect ends.

\end{DndMonster}



\begin{DndMonster}[float=t]{Orc Godcaller}
\DndMonsterType{Medium humanoid (orc, Support), A}
\DndMonsterBasics[
armor-class = {14 (studded leather armor)},
hit-points = {\DndDice{11d8 + 33}},
speed = {35 ft.}
]
\DndMonsterAbilityScores[
 str = 14,
 dex = 14,
 con = 16,
 int = 12,
 wis = 13,
 cha = 19
]
\DndMonsterDetails[
saving-throws = {Wis +3},
skills = {Arcana +3, Insight +3, Performance +8},
condition-immunities = {charmed},
senses = {darkvision 60 ft., passive perception 11},
languages = {Common, Orc},
challenge = 4,
]
\DndMonsterAction{Relentless (1/Turn)}
When the godcaller isn't incapacitated and they are reduced to 0 hit points but not killed outright, they can make an attack against an enemy (no action required) before the hit point reduction is resolved. If the attack hits and its damage reduces the target to 0 hit points, the godcaller drops to 1 hit point instead of 0 hit points.

\DndMonsterSection{Actions}
\DndMonsterAction{Power Chord (1st-Level Spell)}
\textsl{Melee or Ranged Spell Attack:} 6 to hit, reach 5 ft. or range 30 ft., one creature who can hear the godcaller. \textsl{Hit:} 18 (4d6 + 4) thunder damage.

\DndMonsterAction{Cadenza}
The godcaller chooses another creature within 30 feet of them. If the target can hear the godcaller, the target can use their reaction to move up to their speed and make an attack.

\DndMonsterAction{Song of the Gods (2nd-Level Spell)}
The godcaller and each ally within 30 feet of them who can hear them has advantage on attack rolls until the start of the godcaller's next turn. This effect ends early if the godcaller takes any damage.

\DndMonsterSection{Bonus Actions}
\DndMonsterAction{Rallying Ostinato (1/Day)}
The godcaller and up to three allies within 60 feet of them who can hear them regain 20 hit points, and these creatures ignore difficult terrain for 1 minute.

\end{DndMonster}



\begin{DndMonster}[float=t]{Orc Rampart}
\DndMonsterType{Medium humanoid (orc, Soldier), A}
\DndMonsterBasics[
armor-class = {18 (\textit{chain mail})},
hit-points = {\DndDice{5d8 + 20}},
speed = {35 ft.}
]
\DndMonsterAbilityScores[
 str = 18,
 dex = 12,
 con = 18,
 int = 10,
 wis = 11,
 cha = 12
]
\DndMonsterDetails[
skills = {Athletics +6, Perception +2},
senses = {darkvision 60 ft., passive perception 12},
languages = {Common, Orc},
challenge = 2,
]
\DndMonsterAction{Relentless (1/Turn)}
When the rampart isn't incapacitated and they are reduced to 0 hit points but not killed outright, they can make an attack against an enemy (no action required) before the hit point reduction is resolved. If the attack hits and its damage reduces the target to 0 hit points, the rampart drops to 1 hit point instead of 0 hit points.

\DndMonsterSection{Actions}
\DndMonsterAction{Multiattack}
The rampart makes two Spear attacks. If the rampart targets the same creature with both attacks, the target has disadvantage on attack rolls against creatures other than the rampart until the start of the rampart's next turn.

\DndMonsterAction{Spear}
\textsl{Melee or Ranged Weapon Attack:} 6 to hit, reach 5 ft. or range 20/60 ft., one target. \textsl{Hit:} 7 (1d6 + 4) piercing damage, or 8 (1d8 + 4) piercing damage if used with two hands to make a melee attack.

\DndMonsterSection{Reactions}
\DndMonsterAction{No!}
When an enemy within 5 feet of the rampart targets another creature with an attack, the attacker must target the rampart instead.

\end{DndMonster}



\begin{DndMonster}[float=t]{Orc Terranova}
\DndMonsterType{Medium humanoid (orc, Controller), A}
\DndMonsterBasics[
armor-class = {15 (\textit{breastplate})},
hit-points = {\DndDice{5d8 + 15}},
speed = {35 ft.}
]
\DndMonsterAbilityScores[
 str = 14,
 dex = 12,
 con = 16,
 int = 12,
 wis = 9,
 cha = 16
]
\DndMonsterDetails[
skills = {Athletics +4},
condition-immunities = {prone},
senses = {darkvision 60 ft., tremorsense 30 ft., passive perception 9},
languages = {Common, Orc},
challenge = 2,
]
\DndMonsterAction{Relentless (1/Turn)}
When the terranova isn't incapacitated and they are reduced to 0 hit points but not killed outright, they can make an attack against an enemy (no action required) before the hit point reduction is resolved. If the attack hits and its damage reduces the target to 0 hit points, the terranova drops to 1 hit point instead of 0 hit points.

\DndMonsterAction{Tremorspeak}
The terranova can sense anything a willing creature within 30 feet of them can sense with tremorsense or Tremorspeak.

\DndMonsterAction{Unhindered}
The terranova ignores difficult terrain.

\DndMonsterSection{Actions}
\DndMonsterAction{Earth Pillar (1st-Level Spell)}
\textsl{Melee or Ranged Spell Attack:} 5 to hit, reach 5 ft. or range 30 ft., one target touching the ground. \textsl{Hit:} 12 (2d8 + 3) bludgeoning damage, and the target is knocked prone.

\DndMonsterAction{Quake}
The terranova chooses up to five creatures on the ground within 15 feet of them and stomps the ground. Each target must succeed on a DC 13 Strength saving throw or be knocked prone.

\DndMonsterAction{Sinkhole}
The terranova chooses up to two creatures on the ground within 30 feet of them. Each target must succeed on a DC 13 Strength saving throw or be restrained. A creature can use their action to free themself or another creature within their reach, ending the effect on that creature.

\DndMonsterSection{Bonus Actions}
\DndMonsterAction{Rush}
The terranova takes the Dash action.

\DndMonsterAction{Unearth Mohlers (1/Day)}
The terranova stomps, and 1d3 + 1 \textbf{mohler}s pop out of the ground in unoccupied spaces within 30 feet of them. The mohlers are friendly to the terranova and act immediately after the terranova's turn on the same initiative count.

\end{DndMonster}



\begin{DndMonster}[float=t]{Owlbear}
\DndMonsterType{Large monstrosity (Brute), U}
\DndMonsterBasics[
armor-class = {13 (natural armor)},
hit-points = {\DndDice{11d10 + 22}},
speed = {40 ft.}
]
\DndMonsterAbilityScores[
 str = 21,
 dex = 12,
 con = 15,
 int = 5,
 wis = 12,
 cha = 6
]
\DndMonsterDetails[
skills = {Perception +3},
condition-immunities = {frightened},
senses = {darkvision 60 ft., passive perception 13},
challenge = 3,
]
\DndMonsterAction{Deadly Leap}
The owlbear's long jump is up to 30 feet and their high jump is 15 feet, with or without a running start. If the owlbear leaps at least 20 feet toward a target and then hits them with a Bite or Claw attack, the target must succeed on a DC 15 Strength saving throw or be knocked prone.

\DndMonsterAction{Keen Sight and Smell}
The owlbear has advantage on Wisdom (Perception) checks that rely on sight or smell.

\DndMonsterSection{Actions}
\DndMonsterAction{Multiattack}
The owlbear makes one Bite and one Claw attack.

\DndMonsterAction{Bite}
\textsl{Melee Weapon Attack:} 7 to hit, reach 10 ft., one target. \textsl{Hit:} 11 (1d12 + 5) piercing damage, and the owlbear can move the target 5 feet in any direction.

\DndMonsterAction{Claw}
\textsl{Melee Weapon Attack:} 7 to hit, reach 5 ft., one target. \textsl{Hit:} 12 (2d6 + 5) slashing damage.

\DndMonsterAction{Bear Hug (Recharge 4\textendash 6)}
The owlbear attempts to grab and crush a creature they can see within 5 feet of them. The target must make a DC 15 Dexterity saving throw. On a failed save, the target takes 22 (4d10) bludgeoning damage and is grappled (escape DC 15). On a successful save, the target takes half as much damage and is not grappled. Until this grapple ends, the target is restrained and takes 5 (1d10) bludgeoning damage at the start of each of their turns. The grapple ends if the owlbear uses Bear Hug on another target or makes a Claw attack against another target.

\DndMonsterSection{Reactions}
\DndMonsterAction{Hulking Rush}
When the owlbear takes damage, they can move up to half their speed without provoking opportunity attacks.

\end{DndMonster}



\begin{DndMonster}[float=t]{Pitling}
\DndMonsterType{Tiny fiend (demon, category 1, Minion), chaotic evil}
\DndMonsterBasics[
armor-class = {14 (natural armor)},
hit-points = {\DndDice{11}},
speed = {30 ft., fly 30 ft.}
]
\DndMonsterAbilityScores[
 str = 11,
 dex = 14,
 con = 12,
 int = 3,
 wis = 6,
 cha = 6
]
\DndMonsterDetails[
damage-immunities = {poison},
condition-immunities = {poisoned},
senses = {darkvision 60 ft., soulsight 10 ft., passive perception 8},
languages = {Abyssal},
challenge = 4,
]
\DndMonsterAction{Horrid Stench}
If an enemy starts their turn within 10 feet of three or more pitlings, the enemy is poisoned until the start of their next turn.

\DndMonsterAction{Minion}
If the pitling takes damage from an attack or as the result of a failed saving throw, their hit points are reduced to 0. If the pitling takes damage from another effect, they die if the damage equals or exceeds their hit point maximum; otherwise they take no damage.

\DndMonsterAction{Soul Feast}
When the pitling reduces a creature who isn't a Construct or an Undead to 0 hit points or deals damage to a dying creature, the creature must make a DC 11 Wisdom saving throw. On a failed save, the pitling consumes the creature's soul, and the pitling becomes a Category 2 demon of the GM's choice with 1 soul. A creature whose soul is consumed in this way immediately dies, and they can't be restored to life by any means short of a \textit{wish} spell. If the creature was reduced to 0 hit points by a group attack, the GM picks one pitling who joined the attack to consume the soul.

\DndMonsterSection{Actions}
\DndMonsterAction{Spit (Group Attack)}
\textsl{Melee or Ranged Weapon Attack:} 4 to hit, reach 5 ft. or range 15/30 ft., one target. \textsl{Hit:} 4 poison damage.

\end{DndMonster}



\begin{DndMonster}[float*=b,width=\textwidth + 8pt]{Qazyldrath}
\begin{multicols}{2}
\DndMonsterType{Gargantuan dragon (Solo), chaotic evil}
\DndMonsterBasics[
armor-class = {20 (natural armor)},
hit-points = {\DndDice{23d20 + 161}},
speed = {40 ft., fly 80 ft., swim 40 ft.}
]
\DndMonsterAbilityScores[
 str = 27,
 dex = 12,
 con = 24,
 int = 14,
 wis = 16,
 cha = 18
]
\DndMonsterDetails[
saving-throws = {Con +14, Wis +10, Cha +11},
skills = {Arcana +9, History +9, Perception +10, Stealth +8},
damage-immunities = {acid, necrotic},
condition-immunities = {charmed, dazed, frightened, paralyzed, petrified, poisoned, stunned},
senses = {blindsight 120 ft., truesight 60 ft., passive perception 20},
languages = {Common, Draconic},
challenge = 21,
]
\DndMonsterAction{Amphibious}
Qazyldrath can breathe air and water.

\DndMonsterAction{Consume Shadow (3/Day)}
When Qazyldrath fails a saving throw, they can succeed instead. When they do, any areas of magical darkness in the lair are dispelled, and they can't use Enshroud or Thwart Healing until the end of their next turn.

\DndMonsterAction{Shadow Strength}
Acid and necrotic damage dealt by Qazyldrath ignore damage resistance.

\DndMonsterSection{Actions}
\DndMonsterAction{Multiattack}
Qazyldrath makes one Bite attack and two Claw attacks.

\DndMonsterAction{Bite}
\textsl{Melee Weapon Attack:} 15 to hit, reach 15 ft., one target. \textsl{Hit:} 19 (2d10 + 8) piercing damage plus 9 (2d8) necrotic damage. If the target is a light source or is wearing or holding a light source, that light is extinguished.

\DndMonsterAction{Claw}
\textsl{Melee Weapon Attack:} 15 to hit, reach 10 ft., one target. \textsl{Hit:} 15 (2d6 + 8) slashing damage, and Qazyldrath can move the target up to 15 feet horizontally.

\DndMonsterAction{Vacuous Breath (Recharge 5\textendash 6)}
Qazyldrath exhales dark energy in a line that is 90 feet long and 20 feet wide. Each creature in that area must make a DC 22 Dexterity saving throw, taking 24 (7d6) acid damage plus 24 (7d6) necrotic damage on a failed save, or half as much damage on a successful one.

Additionally, that area is filled with magical darkness for 1 minute. A creature with darkvision can't see through this darkness, and no light except a \textit{daylight} spell or a light-creating spell of 5th level or higher can illuminate it.

\DndMonsterSection{Bonus Actions}
\DndMonsterAction{Enshroud}
Shadows cling to one creature Qazyldrath can see within 120 feet of them. The target must succeed on a DC 19 Dexterity saving throw or be blinded and vulnerable to necrotic damage until the start of Qazyldrath's next turn.

\DndMonsterSection{Reactions}
\DndMonsterAction{Thwart Healing}
When a creature within 60 feet of Qazyldrath regains hit points, Qazyldrath forces them to make a DC 19 Constitution saving throw. On a failed save, the creature regains half the number of hit points instead.

\DndMonsterSection{Legendary Actions}
Qazyldrath has three villain actions. They can take each action once during an encounter after an enemy's turn. They can take these actions in any order but can use only one per round.

\begin{DndMonsterLegendaryActions}
\DndMonsterLegendaryAction{Action 1: Burning Globs}{Qazyldrath spits acid globules at each enemy they can see within 90 feet of them. Each target must succeed on a DC 19 Dexterity saving throw or have a globule attached to their body. A creature attached to a globule takes 14 (4d6) acid damage at the start of their turns. A creature can use an action to remove a globule from themself or a creature they can reach.
}
\DndMonsterLegendaryAction{Action 2: Shadow Form}{Qazyldrath becomes semiincorporeal, gaining resistance to bludgeoning, piercing, and slashing damage until the end of their next turn. Qazyldrath then teleports up to 120 feet to an unoccupied space that they can see.
}
\DndMonsterLegendaryAction{Action 3: Sinking Gloom}{Qazyldrath summons sticky black acid centered on a point they can see on the ground within 120 feet of them, creating a 5-foot-deep, 20-foot-radius pool. The pool is difficult terrain and lasts for 1 minute. An enemy who starts their turn in the pool takes 24 (7d6) acid damage and must succeed on a DC 19 Strength saving throw or be restrained until the start of their next turn. An enemy who starts their turn flying within 30 feet of the pool must succeed on a DC 19 Dexterity saving throw or be pulled down by living acid, landing in the nearest unoccupied space of their choice within the pool, take 24 (7d6) acid damage, and be restrained until the start of their next turn.
}
\end{DndMonsterLegendaryActions}
\end{multicols}
\end{DndMonster}



\begin{DndMonster}[float*=b,width=\textwidth + 8pt]{Queen Bargnot}
\begin{multicols}{2}
\DndMonsterType{Small humanoid (goblin, Leader), neutral evil}
\DndMonsterBasics[
armor-class = {17 (studded leather armor)},
hit-points = {\DndDice{12d6 + 12}},
speed = {30 ft., climb 20 ft.}
]
\DndMonsterAbilityScores[
 str = 10,
 dex = 17,
 con = 13,
 int = 14,
 wis = 12,
 cha = 13
]
\DndMonsterDetails[
saving-throws = {Dex +5, Wis +3},
skills = {Insight +3, Intimidation +3, Perception +3, Stealth +5},
senses = {darkvision 60 ft., passive perception 13},
languages = {Common, Goblin},
challenge = 3,
]
\DndMonsterAction{Crafty}
Queen Bargnot doesn't provoke opportunity attacks when she moves out of an enemy's reach.

\DndMonsterAction{Take My Pain (3/Day)}
When Queen Bargnot fails a saving throw against a spell or other supernatural effect, she can choose a willing creature within 30 feet of her. Queen Bargnot succeeds on the saving throw, the creature is targeted with the same spell or effect as if they were in her space, and they fail their saving throw automatically.

\DndMonsterSection{Actions}
\DndMonsterAction{Multiattack}
Queen Bargnot makes three Shortsword attacks or two Shortbow attacks.

\DndMonsterAction{Shortsword}
\textsl{Melee Weapon Attack:} 5 to hit, reach 5 ft., one target. \textsl{Hit:} 6 (1d6 + 3) piercing damage.

\DndMonsterAction{Shortbow}
\textsl{Ranged Weapon Attack:} 5 to hit, range 80/320 ft., one target. \textsl{Hit:} 6 (1d6 + 3) piercing damage.

\DndMonsterSection{Bonus Actions}
\DndMonsterAction{Get In Here}
Queen Bargnot shouts for aid and 1d4 \textbf{goblin lackey}s appear in unoccupied spaces within 60 feet of her.

\DndMonsterSection{Reactions}
\DndMonsterAction{No Dying!}
When an ally Queen Bargnot can see within 30 feet of her is reduced to 0 hit points, they are reduced to 1 hit point instead.

\DndMonsterSection{Legendary Actions}
Queen Bargnot has three villain actions. She can take each action once during an encounter after an enemy's turn.

She can take these actions in any order but can use only one per round.

\begin{DndMonsterLegendaryActions}
\DndMonsterLegendaryAction{Action 1: What Are You Waiting For?!}{Each ally within 60 feet of Queen Bargnot who can hear her can move up to their speed or make a melee weapon attack (no action required).
}
\DndMonsterLegendaryAction{Action 2: Focus Fire}{Queen Bargnot chooses an enemy she can see with 60 feet of her. Queen Bargnot and each ally within 60 feet of her who can hear her can move up to their speed toward the target.
}
\DndMonsterLegendaryAction{Action 3: Kill!}{Each ally within 60 feet of Queen Bargnot who can hear her can make a weapon attack with advantage (no action required). If the attack hits, it deals an extra 3 (1d6) damage.
}
\end{DndMonsterLegendaryActions}
\end{multicols}
\end{DndMonster}



\begin{DndMonster}[float=t]{Rot Angel}
\DndMonsterType{Medium celestial (Skirmisher), lawful neutral}
\DndMonsterBasics[
armor-class = {15 (natural armor)},
hit-points = {\DndDice{12d8 + 36}},
speed = {0 ft., fly 50 ft. \textit{(hover)}}
]
\DndMonsterAbilityScores[
 str = 6,
 dex = 16,
 con = 17,
 int = 8,
 wis = 14,
 cha = 19
]
\DndMonsterDetails[
saving-throws = {Con +5, Cha +6},
skills = {Perception +6},
damage-resistances = {necrotic, radiant},
condition-immunities = {charmed, flanked, frightened},
senses = {blindsight 10 ft., darkvision 60 ft., passive perception 16},
languages = {understands Celestial but can't speak it, telepathy 120 ft.},
challenge = 4,
]
\DndMonsterAction{Flyby}
The angel doesn't provoke opportunity attacks when they fly out of an enemy's reach.

\DndMonsterSection{Actions}
\DndMonsterAction{Decaying Touch}
\textsl{Melee Weapon Attack:} 5 to hit, reach 5 ft., one target. \textsl{Hit:} 18 (4d8) necrotic damage, or if the target is a corpse or Undead, the damage increases to 36 (8d8).

Additionally, for the next 24 hours, the target is marked for death. A creature who drops to 0 hit points while marked for death immediately fails one death saving throw. This effect isn't cumulative.

\DndMonsterSection{Bonus Actions}
\DndMonsterAction{Hypnotic Beam}
One of the angel's eyes shoots a magic ray at a creature the angel can see within 60 feet of them. The target must succeed on a DC 14 Wisdom saving throw or use their reaction, if available, to move as far as their speed allows toward the angel, avoiding obviously dangerous ground.

\DndMonsterSection{Reactions}
\DndMonsterAction{Radiant Puff}
When the angel is hit with an attack, they release a cloud of spores. Each creature within 10 feet of the angel must succeed on a DC 14 Constitution saving throw or be blinded until the end of the rot angel's next turn.

\end{DndMonster}



\begin{DndMonster}[float=t]{Ruinant}
\DndMonsterType{Medium fiend (demon, category 2, Skirmisher), chaotic evil}
\DndMonsterBasics[
armor-class = {15 (natural armor)},
hit-points = {\DndDice{14d8 + 42}},
speed = {60 ft.}
]
\DndMonsterAbilityScores[
 str = 15,
 dex = 18,
 con = 16,
 int = 14,
 wis = 18,
 cha = 16
]
\DndMonsterDetails[
saving-throws = {Wis +7, Cha +6},
skills = {Deception +6, Perception +7},
damage-resistances = {necrotic},
senses = {darkvision 120 ft., soulsight 30 ft., passive perception 17},
languages = {Abyssal, Common, telepathy 120 ft.},
challenge = 6,
]
\DndMonsterAction{Lethe}
When the ruinant's soul count is 0, they have advantage on attack rolls, disadvantage on saving throws, and their Intelligence score becomes 3 (-4). Additionally, the ruinant must use their movement on each of their turns to move as close as possible to the nearest creature they can sense with their soulsight, and then if they are able, they must use their action to attack and attempt to kill that creature. The ruinant can't act with any other purpose until they add 1 to their soul count.

\DndMonsterAction{Soul Devourer}
When the ruinant reduces a creature who isn't a Construct or an Undead to 0 hit points or deals damage to a dying creature, the creature must make a DC 11 Wisdom saving throw. On a failed save, the ruinant consumes the creature's soul and adds 1 to the ruinant's soul count. A creature whose soul is consumed in this way immediately dies, and they can't be restored to life by any means short of a \textit{wish} spell.

\DndMonsterSection{Actions}
\DndMonsterAction{Multiattack}
The ruinant makes three Bloodletting Claws attacks.

\DndMonsterAction{Bloodletting Claws}
\textsl{Melee Weapon Attack:} 7 to hit, reach 5 ft., one creature. \textsl{Hit:} 7 (1d6 + 4) piercing damage plus 7 (2d6) necrotic damage, and the target can't take reactions this turn.

\DndMonsterAction{Salt Wounds (Costs 1 Soul)}
The ruinant chooses up to three creatures they can see within 60 feet of them who don't have all their hit points. Each target must make a DC 15 Constitution saving throw, taking 16 (3d10) necrotic damage on a failed save, or half as much damage on a successful one.

\DndMonsterSection{Reactions}
\DndMonsterAction{Corrupt Healing (Costs 1 Soul)}
When a creature within 60 feet of the ruinant regains hit points from a power, a spell, or a similar supernatural effect, the ruinant corrupts the effect. The target regains no hit points, and the target and each of the ruinant's enemies within 5 feet of the ruinant must succeed on a DC 15 Constitution saving throw or take necrotic damage equal to half the number of hit points the effect would have restored.

\end{DndMonster}



\begin{DndMonster}[float=t]{Scyza}
\DndMonsterType{Gargantuan monstrosity (Brute), U}
\DndMonsterBasics[
armor-class = {14 (natural armor)},
hit-points = {\DndDice{8d20 + 24}},
speed = {40 ft.}
]
\DndMonsterAbilityScores[
 str = 21,
 dex = 9,
 con = 17,
 int = 5,
 wis = 11,
 cha = 7
]
\DndMonsterDetails[
saving-throws = {Str +7, Con +5},
condition-immunities = {frightened},
senses = {darkvision 120 ft., passive perception 10},
challenge = 4,
]
\DndMonsterAction{Siege Monster}
The scyza deals double damage to objects and structures.

\DndMonsterAction{War Harness}
While wearing a war harness, the scyza can carry up to thirty-six Medium or smaller creatures.

\DndMonsterSection{Actions}
\DndMonsterAction{Multiattack}
The scyza makes two Claw attacks.

\DndMonsterAction{Claw}
\textsl{Melee Weapon Attack:} 7 to hit, reach 5 ft., one target. \textsl{Hit:} 12 (2d6 + 5) slashing damage, and the target is knocked prone.

\DndMonsterAction{Tail Whip}
\textsl{Melee Weapon Attack:} 7 to hit, reach 20 ft., one target. \textsl{Hit:} 14 (2d8 + 5) bludgeoning damage. If the target is on top of the scyza, the target falls off and lands prone.

\DndMonsterAction{Head Crest}
\textsl{Melee Weapon Attack:} 7 to hit, reach 10 ft., one target. \textsl{Hit:} 23 (4d8 + 5) bludgeoning damage, and if the target is a creature who can hear, they are dazed until the end of their next turn by the crest's deep ringing.

\DndMonsterAction{Claw Twister (Recharge 6)}
The scyza claws at the ground in a flurry. Each creature of the scyza's choice within 5 feet of them must make a DC 15 Dexterity saving throw. On a failed save, a creature takes 13 (3d8) slashing damage, is knocked prone, and is blinded by a cloud of dust until the start of the scyza's next turn. On a successful save, a creature takes half as much damage and isn't knocked prone or blinded.

\end{DndMonster}



\begin{DndMonster}[float=t]{Sea Hag Apprentice}
\DndMonsterType{Medium fey (hag, Retainer), A}
\DndMonsterBasics[
armor-class = {13 (light armor)},
hit-points = {\DndDice{Six times their level (number of d6 Hit Dice equal to their level)}},
speed = {30 ft., swim 30 ft.}
]
\DndMonsterAbilityScores[
 str = 10,
 dex = 12,
 con = 10,
 int = 10,
 wis = 12,
 cha = 16
]
\DndMonsterDetails[
saving-throws = {Str +PB, Dex +PB, Con +PB, Int +PB, Wis +PB, Cha +PB},
skills = {Deception +3 plus PB, Intimidation +3 plus PB},
damage-resistances = {cold, lightning},
senses = {darkvision 60 ft., passive perception 11},
languages = {Aquan, Common, Sylvan},
]
\DndMonsterAction{Amphibious}
The apprentice can breathe air and water.

\DndMonsterSection{Actions}
\DndMonsterAction{Signature Attack (Zap)}
\textsl{Melee or Ranged Spell Attack:} 3 plus PB to hit, reach 5 ft. or range 60 ft., one target. \textsl{Hit:} 4 (1d8) lightning damage. This spell's damage increases by 1d8 when the apprentice reaches 5th level (2d8), 11th level (3d8), and 17th level (4d8).

\DndMonsterAction{3rd Level: Wash Away (3/Day)}
As a bonus action, the apprentice rides a crashing wave. They make a signature attack against an enemy within 5 feet of them and move up to their walking speed without provoking opportunity attacks. If the attack hits, it deals an extra 0d1 + PB bludgeoning damage.

\DndMonsterAction{5th Level: Ride the Lightning (3/Day)}
As an action, the apprentice supercharges their foe's nervous system. The apprentice makes a signature attack, and if it deals damage to a creature, the target must make a DC 10 plus PB Constitution saving throw. On a failed save, the target is incapacitated and their speed is halved until the end of their next turn.

\DndMonsterAction{7th Level: Drown (1/Day)}
As an action, the apprentice chooses three creatures the apprentice can see within 60 feet of them. Each target feels as though their lungs are filling with water and must make a DC 10 plus PB Constitution saving throw. A creature who can breathe water or who doesn't need air automatically succeeds on this saving throw. On a failed save, a target takes PBd8 force damage, and for 1 minute, they can't speak, their speed is halved, and they are dazed (save ends at end of turn). On a successful save, the target takes half as much damage and suffers no other effect. Another creature within 5 feet of a target can use their action to make a DC 10 plus PB Strength (Athletics) or Intelligence (Medicine) check, ending the effect on a success.

\end{DndMonster}



\begin{DndMonster}[float=t]{Shade}
\DndMonsterType{Medium undead (incorporeal, Minion), chaotic evil}
\DndMonsterBasics[
armor-class = {11},
hit-points = {\DndDice{8}},
speed = {0 ft., fly 50 ft. \textit{(hover)}}
]
\DndMonsterAbilityScores[
 str = 1,
 dex = 13,
 con = 10,
 int = 10,
 wis = 10,
 cha = 14
]
\DndMonsterDetails[
damage-immunities = {necrotic, poison},
condition-immunities = {charmed, dazed, exhaustion, grappled, paralyzed, petrified, poisoned, prone, restrained, unconscious},
senses = {darkvision 60 ft., passive perception 10},
languages = {the languages they knew in life},
challenge = 1,
]
\DndMonsterAction{Incorporeal Movement}
The shade can move through other creatures and objects as if they were difficult terrain. The shade is destroyed if they end their turn inside an object.

\DndMonsterAction{Minion}
If the shade takes damage from an attack or as the result of a failed saving throw, their hit points are reduced to 0. If the shade takes damage from another effect, they die if the damage equals or exceeds their hit point maximum; otherwise they take no damage.

\DndMonsterAction{Terrifying}
If an enemy starts their turn within 5 feet of three or more shades, the enemy must succeed on a Wisdom saving throw or be frightened of the shades until the start of their next turn. The DC for this saving throw equals 10 plus the number of shades within 5 feet of the enemy. On a successful save, the enemy is immune to the Terrifying trait of all shades for the next 24 hours.

\DndMonsterSection{Actions}
\DndMonsterAction{Life Drain (Group Attack)}
\textsl{Melee Spell Attack:} 4 to hit, reach 5 ft., one creature. \textsl{Hit:} 1 necrotic damage.

\end{DndMonster}


\clearpage

\begin{DndMonster}[float=t]{Shadow}
\DndMonsterType{Medium undead (incorporeal, Ambusher), chaotic evil}
\DndMonsterBasics[
armor-class = {12},
hit-points = {\DndDice{3d8}},
speed = {40 ft.}
]
\DndMonsterAbilityScores[
 str = 6,
 dex = 14,
 con = 11,
 int = 10,
 wis = 10,
 cha = 14
]
\DndMonsterDetails[
skills = {Stealth +4},
damage-resistances = {acid; lightning; necrotic; thunder; bludgeoning, piercing, slashing from mundane attacks},
damage-immunities = {cold, poison},
condition-immunities = {dazed, exhaustion, frightened, grappled, paralyzed, petrified, poisoned, prone, restrained},
senses = {darkvision 60 ft., passive perception 10},
languages = {the languages they knew in life},
challenge = 1/2,
]
\DndMonsterAction{Chilling Phasing}
The shadow can move through other creatures and objects as if they were difficult terrain. The shadow takes 5 (1d10) force damage if they end their turn inside an object. A creature takes 1 cold damage the first time a shadow passes through them on a turn.

\DndMonsterAction{Dark Stalker}
While in dim light or darkness, the shadow has advantage on Dexterity (Stealth) checks.

\DndMonsterSection{Actions}
\DndMonsterAction{Chilling Grasp}
\textsl{Melee Spell Attack:} 4 to hit, reach 5 ft., one creature. \textsl{Hit:} 6 (1d8 + 2) cold damage, and the target's Dexterity score is reduced by 2. The target dies if this reduces its Dexterity to 0. Otherwise, the reduction lasts until the target finishes a short or long rest.

\DndMonsterAction{Freezing Dark (1/Day)}
The shadow pours magical darkness out of their hand in a 15-foot cone. A creature with darkvision who isn't an incorporeal Undead can't see through this darkness, and light can't illuminate it. Each creature in this darkness when it first appears must succeed on a DC 12 Constitution saving throw or take 7 (2d6) cold damage. The darkness remains in that area for 1 minute or until the shadow is destroyed.

\DndMonsterSection{Bonus Actions}
\DndMonsterAction{Shadow Bind}
The shadow targets one creature they can see within 15 feet of them who is in bright or dim light. The target must succeed on a DC 12 Strength saving throw or be grappled by their own shadow (escape DC 12). The grapple ends if the target is no longer in bright or dim light, or if the shadow who used this bonus action is destroyed.

\end{DndMonster}



\begin{DndMonster}[float*=b,width=\textwidth + 8pt]{Shtriga Nonna}
\begin{multicols}{2}
\DndMonsterType{Medium fey (Solo), neutral evil}
\DndMonsterBasics[
armor-class = {17 (natural armor)},
hit-points = {\DndDice{14d8 + 56}},
speed = {30 ft., fly 30 ft.}
]
\DndMonsterAbilityScores[
 str = 18,
 dex = 16,
 con = 18,
 int = 14,
 wis = 15,
 cha = 19
]
\DndMonsterDetails[
saving-throws = {Wis +5},
skills = {Deception +7, Perception +5, Stealth +6, Survival +5},
damage-resistances = {thunder; bludgeoning, piercing, slashing from mundane attacks},
damage-immunities = {cold, fire},
condition-immunities = {charmed, frightened},
senses = {darkvision 60 ft., passive perception 15},
languages = {Common, Giant, Infernal, Sylvan},
challenge = 6,
]
\DndMonsterAction{Feline Resilience (3/Day)}
When Shtriga Nonna fails a saving throw, she can immediately turn into a cat of her size (Medium in her normal form, Large while under the effect of Grow) and succeed instead. While in cat form, Shtriga Nonna retains her game statistics and ability to speak, except she can't cast spells or use Soul Steal. She reverts to her previous form at the end of her next turn.

\DndMonsterAction{Supernatural Resistance}
Shtriga Nonna has advantage on saving throws against powers, spells, and other supernatural effects.

\DndMonsterAction{Spellcasting (Utility)}
In addition to any other spells in this stat block, Shtriga Nonna can cast the following spells, using Charisma as the spellcasting ability (spell save DC 15):
\begin{DndMonsterSpells}
  \DndInnateSpellLevel{alter self\textsuperscript{Casting Times}, \textit{detect thoughts}\textsuperscript{Casting Times}, \textit{thaumaturgy}\textsuperscript{Casting Times}}
  \DndInnateSpellLevel[]{scrying\textsuperscript{Casting Times}}
  \DndInnateSpellLevel[3]{fabricate\textsuperscript{Casting Times}, \textit{legend lore}\textsuperscript{Casting Times}}
\end{DndMonsterSpells}
\DndMonsterSection{Actions}
\DndMonsterAction{Multiattack}
Shtriga Nonna makes two Corrosive Claw attacks. She can replace one attack with a use of Soul Steal, if available.

\DndMonsterAction{Corrosive Claw}
\textsl{Melee Weapon Attack:} 7 to hit, reach 5 ft. (10 ft. with Grow), one target. \textsl{Hit:} 11 (2d6 + 4) slashing damage, or 14 (3d6 + 4) slashing damage while under the effect of Grow. If the target is wearing mundane metal armor, they must succeed on a DC 15 Dexterity throw or their armor takes a permanent and cumulative -1 penalty to the AC it offers. If the armor is reduced to an AC of 10, it is destroyed.

\DndMonsterAction{Soul Steal (Recharge 5\textendash 6)}
Shtriga Nonna inhales deeply in a 30-foot cone. Each creature in that area must make a DC 15 Constitution saving throw, taking 16 (3d10) necrotic damage on a failed save, or half as much damage on a successful one. Shtriga Nonna regains 10 hit points for each creature who fails this saving throw.

\DndMonsterSection{Bonus Actions}
\DndMonsterAction{Grow (2/Day; 3rd-Level Spell)}
For 1 minute, Shtriga Nonna magically increases in size, along with anything she is wearing or carrying. While enlarged, she is Large, her reach becomes 10 feet, she deals an extra 3 (1d6) damage on Strength-based weapon attacks (included in the attacks), and she makes Strength checks and Strength saving throws with advantage. If Shtriga Nonna lacks the room to become Large, she attains the maximum size possible in the space available.

\DndMonsterAction{Gobble You Up}
Shtriga Nonna feeds on the body of an unconscious creature within 5 feet of her. The target takes 10 (3d6) piercing damage, and Shtriga Nonna recharges her Soul Steal action.

\DndMonsterSection{Reactions}
\DndMonsterAction{Turned Upside Down}
When a creature hits Shtriga Nonna with a melee attack, she magically reverses gravity around her. Each creature within 5 feet of her must succeed on a DC 15 Strength saving throw or rise vertically 10 feet into the air. A creature levitating in this way remains suspended until the end of their next turn, then falls unless they have an ability that keeps them aloft.

\DndMonsterSection{Legendary Actions}
Shtriga Nonna has three villain actions. She can take each action once during an encounter after an enemy's turn. She can take these actions in any order but can use only one per round.

\begin{DndMonsterLegendaryActions}
\DndMonsterLegendaryAction{Action 1: Snackies for Sweeties}{The air within 30 feet of Shtriga Nonna magically fills with the scent of rotten food. Each other creature in that area must succeed on a DC 15 Wisdom saving throw or be poisoned for 1 minute (save ends at end of turn). Creatures are affected even if they hold their breath or don't need to breathe.
}
\DndMonsterLegendaryAction{Action 2: Feline Felicity}{Shtriga Nonna grows four cat legs in a burst of magical energy. Each creature within 5 feet of her must make a DC 15 Dexterity saving throw. On a failed save, a creature takes 10 (4d4) force damage and is pushed 10 feet away from her. On a successful save, a creature takes half as much damage and isn't pushed. Shtriga Nonna can then move up to twice her speed, after which the extra cat legs disappear.
}
\DndMonsterLegendaryAction{Action 3: Open the Oven}{Shtriga Nonna opens a portal to her massive oven in the palms of her hands, creating a blast of heat and noise in a 30-foot cone. Each creature in that area must make a DC 15 Dexterity saving throw, taking 13 (3d8) fire damage plus 13 (3d8) thunder damage on a failed save, or half as much damage on a successful one.
}
\end{DndMonsterLegendaryActions}
\end{multicols}
\end{DndMonster}



\begin{DndMonster}[float=t]{Slaughter Demon}
\DndMonsterType{Huge fiend (demon, category 3, Brute), chaotic evil}
\DndMonsterBasics[
armor-class = {15 (natural armor)},
hit-points = {\DndDice{12d12 + 60}},
speed = {40 ft.}
]
\DndMonsterAbilityScores[
 str = 22,
 dex = 10,
 con = 20,
 int = 8,
 wis = 12,
 cha = 10
]
\DndMonsterDetails[
saving-throws = {Str +9, Dex +3, Con +8},
skills = {Athletics +9, Perception +7},
damage-resistances = {fire},
condition-immunities = {flanked},
senses = {darkvision 60 ft., soulsight 30 ft., passive perception 17},
languages = {Abyssal, Goblin, Infernal},
challenge = 8,
]
\DndMonsterAction{Lethe}
When the demon's soul count is 0, they have advantage on attack rolls, disadvantage on saving throws, and their Intelligence score becomes 3 (-4). Additionally, the demon must use their movement on each of their turns to move as close as possible to the nearest creature they can sense with their soulsight, and then if they are able, they must use their action to attack and attempt to kill that creature. The demon can't act with any other purpose until they add 1 to their soul count.

\DndMonsterAction{Soul Devourer}
When the demon reduces a creature who isn't a Construct or an Undead to 0 hit points or deals damage to a dying creature, the creature must make a DC 11 Wisdom saving throw. On a failed save, the demon consumes the creature's soul and adds 1 to the demon's soul count. A creature whose soul is consumed in this way immediately dies, and they can't be restored to life by any means short of a \textit{wish} spell.

\DndMonsterSection{Actions}
\DndMonsterAction{Multiattack}
The demon makes four attacks using Pincer, Sword, Javelin, or a combination of them. They can replace one of these attacks with a Tail Spike attack. The demon can burn 1 soul to make a fifth attack using Pincer, Sword, or Javelin.

\DndMonsterAction{Pincer}
\textsl{Melee Weapon Attack:} 9 to hit, reach 5 ft., one target. \textsl{Hit:} 13 (2d6 + 6) bludgeoning damage, and the target is grappled (escape DC 17). While grappled, the target is restrained. The demon has six pincers, each of which can grapple a target or wield a weapon.

\DndMonsterAction{Sword}
\textsl{Melee Weapon Attack:} 9 to hit, reach 10 ft., one target. \textsl{Hit:} 15 (2d8 + 6) slashing damage.

\DndMonsterAction{Javelin}
\textsl{Melee or Ranged Weapon Attack:} 9 to hit, reach 10 ft. or range 60/240 ft., one target. \textsl{Hit:} 13 (2d6 + 6) piercing damage.

\DndMonsterAction{Tail Spike}
\textsl{Melee Weapon Attack:} 9 to hit, reach 15 ft., one target. \textsl{Hit:} 11 (1d10 + 6) piercing damage, and the target must succeed on a DC 16 Constitution saving throw or be poisoned for 1 minute (save ends at end of turn).

\DndMonsterSection{Bonus Actions}
\DndMonsterAction{Dive Deep (Costs 1 Soul)}
The demon takes the Disengage action and gains a burrowing speed equal to their walking speed until the end of their next turn.

\end{DndMonster}



\begin{DndMonster}[float=t]{Specter}
\DndMonsterType{Medium undead (incorporeal, Skirmisher), chaotic evil}
\DndMonsterBasics[
armor-class = {11},
hit-points = {\DndDice{3d8 + 6}},
speed = {0 ft., fly 50 ft. \textit{(hover)}}
]
\DndMonsterAbilityScores[
 str = 1,
 dex = 13,
 con = 14,
 int = 10,
 wis = 10,
 cha = 15
]
\DndMonsterDetails[
damage-resistances = {acid; cold; fire; lightning; thunder; bludgeoning, piercing, slashing from mundane attacks},
damage-immunities = {necrotic, poison},
condition-immunities = {charmed, dazed, exhaustion, grappled, paralyzed, petrified, poisoned, prone, restrained, unconscious},
senses = {darkvision 60 ft., passive perception 10},
languages = {the languages they knew in life},
challenge = 1,
]
\DndMonsterAction{Corrupting Phasing}
The specter can move through other creatures and objects as if they were difficult terrain. The specter takes 5 (1d10) force damage if they end their turn inside an object. A creature takes 2 (1d4) necrotic damage the first time a specter passes through them on a turn.

\DndMonsterSection{Actions}
\DndMonsterAction{Decaying Touch}
\textsl{Melee Spell Attack:} 4 to hit, reach 5 ft., one creature. \textsl{Hit:} 7 (2d6) necrotic damage, and the target must succeed on a DC 12 Constitution saving throw or spend 1 Hit Die without any benefit. If the target has no Hit Dice to spend, they drop to 0 hit points instead.

If a Humanoid dies within 1 minute of being hit by this attack, their spirit rises as a specter under the GM's control in an unoccupied space nearest to where that Humanoid died.

\DndMonsterSection{Bonus Actions}
\DndMonsterAction{Hidden Movement (Recharge 5\textendash 6)}
The specter turns invisible and then moves up to their speed. At the end of this movement, the invisibility ends.

\end{DndMonster}



\begin{DndMonster}[float=t]{Sunlight Nexus}
\DndMonsterType{Large elemental (air, fire, Support), A}
\DndMonsterBasics[
armor-class = {18 (natural armor)},
hit-points = {\DndDice{16d10 + 64}},
speed = {40 ft., fly 90 ft.}
]
\DndMonsterAbilityScores[
 str = 20,
 dex = 15,
 con = 19,
 int = 14,
 wis = 16,
 cha = 20
]
\DndMonsterDetails[
saving-throws = {Wis +7, Cha +9},
skills = {Insight +7, Intimidation +9, Perception +7, Religion +6},
damage-resistances = {fire, radiant},
damage-immunities = {poison},
condition-immunities = {charmed, exhaustion, frightened, poisoned},
senses = {darkvision 60 ft., passive perception 17},
languages = {Common, Primordial},
challenge = 12,
]
\DndMonsterAction{Corona}
The nexus sheds bright sunlight in a 120-foot radius and dim light for another 120 feet.

\DndMonsterSection{Actions}
\DndMonsterAction{Multiattack}
The nexus makes three Solar Arc attacks. They can replace one attack with a use of Heavenly Collapse or Rejuvenating Flare, if available.

\DndMonsterAction{Solar Arc}
\textsl{Melee or Ranged Spell Attack:} 9 to hit, reach 5 ft. or range 30 ft., one target. \textsl{Hit:} 18 (2d12 + 5) radiant damage.

\DndMonsterAction{Heavenly Collapse (2/Day)}
The nexus shines harsh light on an enemy they can see within 60 feet of them. Each ally of the nexus within 5 feet of the target regains 20 hit points, and the target must succeed on a DC 17 Constitution saving throw or lose any damage resistances they have and be blinded for 1 minute (save ends at end of turn).

\DndMonsterAction{Rejuvenating Flare (1/Day)}
One ally the nexus can see within 30 feet of them is bathed in rejuvenating light. The ally regains 27 (5d10) hit points, and if they are charmed, dazed, frightened, or poisoned, those conditions end for them.

\DndMonsterSection{Bonus Actions}
\DndMonsterAction{Convocation of Sunlight (2/Day)}
The nexus imbues the power of sunlight in themself or another Elemental they can see within 30 feet of them, granting one of the following effects of that Elemental's choice:


\begin{description}
\item \textbf{Healing Radiance.} The Elemental regains 30 hit points. Until the end of the nexus's next turn, the Elemental refracts shimmering light, and creatures within 10 feet of them who aren't Elementals are blinded.
\item \textbf{Protective Light.} Four brilliant motes orbit the Elemental for 1 minute. When the Elemental deals damage to a creature with an attack, the Elemental can extinguish one mote (no action required) to deal an extra 14 (4d6) radiant damage to that creature.
\end{description}

\end{DndMonster}



\begin{DndMonster}[float=t]{Swamp Dryad}
\DndMonsterType{Large fey (Controller), chaotic neutral}
\DndMonsterBasics[
armor-class = {17 (natural armor)},
hit-points = {\DndDice{18d10 + 72}},
speed = {30 ft.}
]
\DndMonsterAbilityScores[
 str = 18,
 dex = 12,
 con = 18,
 int = 12,
 wis = 21,
 cha = 10
]
\DndMonsterDetails[
saving-throws = {Str +9, Con +9, Wis +10},
skills = {Acrobatics +6, Animal Handling +10, Arcana +6, Medicine +10, Survival +10},
damage-resistances = {acid; poison; bludgeoning, piercing, slashing from mundane attacks},
condition-immunities = {poisoned},
senses = {darkvision 120 ft., passive perception 15},
languages = {Sylvan, any languages the dryad spoke before their binding},
challenge = 13,
]
\DndMonsterAction{Supernatural Resistance}
The dryad has advantage on saving throws against supernatural effects.

\DndMonsterAction{Speak with Beasts and Plants}
The dryad can communicate with Beasts and Plants as if they shared a language.

\DndMonsterAction{Tree Stride (1/Turn)}
The dryad can use 10 feet of movement to step into one living tree within their reach and emerge from a second living tree within 60 feet of the first tree, appearing in an unoccupied space within 5 feet of the second tree. Both trees must be Large or larger.

\DndMonsterSection{Actions}
\DndMonsterAction{Multiattack}
The dryad makes two Slam attacks and uses Entangling Roots, if available.

\DndMonsterAction{Slam}
\textsl{Melee Weapon Attack:} 9 to hit, reach 10 ft., one target. \textsl{Hit:} 13 (2d8 + 4) piercing damage, and the target must make a DC 18 Constitution saving throw. On a failed save, the target takes 14 (4d6) poison damage and is poisoned for 1 minute. On a successful save, the target takes half as much damage and isn't poisoned.

\DndMonsterAction{Entangling Roots (Recharge 5\textendash 6)}
The dryad chooses a point they can see within 60 feet of them, causing tree roots to pull free of the ground and attack in a 15-foot radius centered on that point. If the dryad is within 5 feet of their tree, the radius of this effect doubles to 30 feet. Each creature in that area must succeed on a DC 18 Strength saving throw or take 22 (4d10) bludgeoning damage and be grappled (escape DC 14).

While grappled in this way, a creature is restrained and their flesh begins to harden and turn to wood.

At the end of a creature's third consecutive turn being grappled in this way, they become petrified and turn into a hardwood statue. The petrified condition ends for that creature if the dryad uses an action to end it, or if the dryad dies.

\DndMonsterAction{Fey Intimidation}
The dryad targets up to five creatures they can see within 30 feet of them. Each target who can see the dryad must succeed on a DC 18 Wisdom saving throw or be frightened of the dryad (save ends at end of turn). If a target succeeds on their saving throw or the effect ends for them, the target is immune to the swamp dryad's Fey Intimidation for the next 24 hours.

\DndMonsterSection{Bonus Actions}
\DndMonsterAction{Swamp Medicine (3/Day)}
The dryad touches a creature, who regains 18 (3d8 + 5) hit points. Alternatively, the creature removes one disease or neutralizes one poison afflicting them.

\end{DndMonster}



\begin{DndMonster}[float=t]{Swarm of Skitterlings}
\DndMonsterType{Medium beasts (Soldier), U}
\DndMonsterBasics[
armor-class = {13},
hit-points = {\DndDice{8d8}},
speed = {30 ft., fly 60 ft.}
]
\DndMonsterAbilityScores[
 str = 10,
 dex = 16,
 con = 10,
 int = 3,
 wis = 12,
 cha = 5
]
\DndMonsterDetails[
damage-resistances = {bludgeoning, piercing, slashing},
condition-immunities = {charmed, dazed, frightened, grappled, paralyzed, petrified, prone, restrained, stunned},
senses = {darkvision 60 ft., passive perception 11},
challenge = 2,
]
\DndMonsterAction{In Your Face}
While the swarm occupies an enemy's space, that creature has disadvantage on attack rolls made against any target other than the swarm and takes 3 (1d6) slashing damage if they attack a creature other than the swarm.

\DndMonsterAction{Swarm}
The swarm can occupy another creature's space and vice versa, and the swarm can move through any opening large enough for a Tiny skitterling. The swarm can't regain hit points or gain temporary hit points.

\DndMonsterSection{Actions}
\DndMonsterAction{Claws}
\textsl{Melee Weapon Attack:} 5 to hit, reach 0 ft., one target in the swarm's space. \textsl{Hit:} 14 (4d6) slashing damage, or 7 (2d6) slashing damage if the swarm has half of their hit points or fewer. The target must succeed on a DC 10 Constitution saving throw or be blinded until the start of their next turn. Creatures who are immune to poison damage or the poisoned condition succeed on this saving throw automatically.

\end{DndMonster}



\begin{DndMonster}[float=t]{Swarm of Spiders}
\DndMonsterType{Medium beasts (Controller), U}
\DndMonsterBasics[
armor-class = {12},
hit-points = {\DndDice{4d8}},
speed = {20 ft., climb 20 ft.}
]
\DndMonsterAbilityScores[
 str = 3,
 dex = 14,
 con = 10,
 int = 2,
 wis = 9,
 cha = 2
]
\DndMonsterDetails[
damage-resistances = {bludgeoning, piercing, slashing},
condition-immunities = {charmed, dazed, flanked, frightened, grappled, paralyzed, petrified, prone, restrained, stunned},
senses = {blindsight 10 ft., darkvision 30 ft., passive perception 9},
challenge = 1/2,
]
\DndMonsterAction{Spider Climb}
The swarm can climb difficult surfaces, including upside down on ceilings, without needing to make an ability check.

\DndMonsterAction{Swarm}
The swarm can occupy another creature's space and vice versa, and the swarm can move through any opening large enough for a Tiny spider. The swarm can't regain hit points or gain temporary hit points.

\DndMonsterAction{Web Walker}
The swarm ignores movement restrictions caused by webbing.

\DndMonsterSection{Actions}
\DndMonsterAction{Bite}
\textsl{Melee Weapon Attack:} 4 to hit, reach 0 ft., one target in the swarm's space. \textsl{Hit:} 5 (2d4) piercing damage plus 4 (1d8) poison damage, or 2 (1d4) piercing damage plus 3 (1d6) poison damage if the swarm has half their hit points or fewer.

\DndMonsterAction{Too Many Legs}
The swarm mobs a creature in their space and attaches to them. When the swarm takes damage other than psychic damage, the attached creature also takes half as much damage. The attached creature has disadvantage on attacks against creatures other than the swarm, and attacks against the attached creature have advantage.

While attached, the swarm can't attack another target. The attached creature or another creature who can reach them can shake off the swarm and detach them as an action. The swarm can detach by spending 5 feet of their movement.

\end{DndMonster}



\begin{DndMonster}[float=t]{Thornblood}
\DndMonsterType{Gargantuan plant (Controller), N}
\DndMonsterBasics[
armor-class = {16},
hit-points = {\DndDice{24d20 + 72}},
speed = {30 ft., climb 30 ft.}
]
\DndMonsterAbilityScores[
 str = 22,
 dex = 23,
 con = 17,
 int = 11,
 wis = 16,
 cha = 8
]
\DndMonsterDetails[
saving-throws = {Con +9, Wis +9},
damage-resistances = {cold; lightning; bludgeoning, piercing, slashing from mundane attacks},
damage-immunities = {poison, psychic},
condition-immunities = {charmed, dazed, flanked, frightened, grappled, poisoned, prone, stunned},
senses = {blindsight 120 ft., passive perception 13},
languages = {Sylvan},
challenge = 20,
]
\DndMonsterAction{Corpse Retreat (1/Day)}
The thornblood is a mass of thick, thorny vines that grow from a Medium humanoid corpse. When the thornblood is reduced to 0 hit points, they are reduced to 1 hit point instead and retreat into the corpse that grew them.

The thornblood becomes immune to all damage until the end of their next turn. Additionally, until the thornblood finishes a long rest, their size becomes Medium, their walking and climbing speeds become 60 feet, and they can't make Throttle attacks or benefit from the Overgrowth trait.

\DndMonsterAction{Overgrowth}
The thornblood can move their vines through a space as narrow as 1 inch wide without squeezing. The thornblood can occupy another creature's space and vice versa.

\DndMonsterAction{Supernatural Resistance}
The thornblood has advantage on saving throws against powers, spells, and other supernatural effects.

\DndMonsterSection{Actions}
\DndMonsterAction{Multiattack}
The thornblood makes four attacks using Vinerip, Throttle, or both.

\DndMonsterAction{Vinerip}
\textsl{Melee Weapon Attack:} 12 to hit, reach 30 ft., one target. \textsl{Hit:} 20 (4d6 + 6) bludgeoning damage plus 11 (2d10) piercing damage, and the target is pulled up to 30 feet toward the thornblood.

\DndMonsterAction{Throttle}
\textsl{Melee Weapon Attack:} 12 to hit, reach 0 ft., one target in the thornblood's space. \textsl{Hit:} 20 (4d6 + 6) piercing damage, and the tar get is entangled in thorny vines. While entangled, the target is restrained and moves with the thornblood, can't see outside their space, and takes 20 (4d6 + 6) piercing damage at the start of each of their turns. The target or another creature within 5 feet of them can use their action to take 20 (4d6 + 6) slashing damage and make a DC 20 Strength check, freeing the target on a success.

\DndMonsterAction{From Life, Lightning (2/Day)}
The thornblood must have at least one creature entangled to use this action. Each creature entangled by the thornblood must make a DC 20 Constitution saving throw, taking 55 (10d10) necrotic damage on a failed save, or half as much damage on a successful one. Then the thornblood shoots a bolt of lightning at each enemy within 60 feet of them who isn't entangled by them. Each target must make a DC 20 Dexterity saving throw, taking 55 (10d10) lightning damage on a failed save, or half as much damage on a successful one.

\DndMonsterSection{Bonus Actions}
\DndMonsterAction{You're Mine}
The thornblood moves an entangled creature to a space occupied by the thornblood that isn't also occupied by another creature or object.

\DndMonsterSection{Reactions}
\DndMonsterAction{Hostile Roots}
When an enemy ends their turn in the thornblood's space, the thornblood attempts to pull them down. The target must make a DC 17 Dexterity saving throw or fall prone.

\end{DndMonster}



\begin{DndMonster}[float=t]{Time Raider Freebooter}
\DndMonsterType{Medium humanoid (time raider, Retainer), A}
\DndMonsterBasics[
armor-class = {15 (medium armor)},
hit-points = {\DndDice{Seven times their level (number of d8 Hit Dice equal to their level)}},
speed = {30 ft.}
]
\DndMonsterAbilityScores[
 str = 16,
 dex = 10,
 con = 10,
 int = 12,
 wis = 12,
 cha = 10
]
\DndMonsterDetails[
saving-throws = {Str +PB, Dex +PB, Con +PB, Int +PB, Wis +PB, Cha +PB},
skills = {Athletics +3 plus PB, Perception +1 plus PB},
damage-resistances = {psychic},
condition-immunities = {blinded, charmed},
senses = {darkvision 120 ft., passive perception 11 plus PB},
languages = {Common, Kuran'zoi},
]
\DndMonsterSection{Actions}
\DndMonsterAction{Signature Attack (Cutlass)}
\textsl{Melee Weapon Attack:} 3 plus PB to hit, reach 5 ft., one target. \textsl{Hit:} 1d10 plus PB slashing damage. Beginning at 7th level, the freebooter can make this attack twice, instead of once, when they take the Attack action on their turn.

\DndMonsterAction{7th Level: Wraithstep (1/Day)}
As an action, the freebooter becomes invisible for 1 minute. While invisible, the freebooter has resistance to bludgeoning, piercing, and slashing damage from mundane attacks, and the freebooter can move through other creatures and objects as if they were difficult terrain. The freebooter takes 5 (1d10) force damage if they end their turn inside an object.

\DndMonsterSection{Bonus Actions}
\DndMonsterAction{3rd Level: Blip Strike (3/Day)}
As a bonus action, the freebooter teleports up to 30 feet to an unoccupied space they can see. If this space is within 5 feet of an enemy the freebooter can see, the freebooter can make a signature attack against them with advantage as part of this bonus action.

\DndMonsterAction{5th Level: Kinetic Clamp (3/Day)}
As a bonus action, the freebooter tries to telekinetically grip one creature they can see within 60 feet of them. The target must make a DC 10 plus PB Strength saving throw. On a failed save, the target takes PBd6 force damage and is restrained until the start of the freebooter's next turn. On a successful save, the target takes half as much damage and is not restrained.

\end{DndMonster}



\begin{DndMonster}[float=t]{Time Raider Nemesis}
\DndMonsterType{Medium humanoid (time raider, Skirmisher), A}
\DndMonsterBasics[
armor-class = {17 (precog reflexes)},
hit-points = {\DndDice{20d8 + 40}},
speed = {30 ft.}
]
\DndMonsterAbilityScores[
 str = 18,
 dex = 16,
 con = 14,
 int = 18,
 wis = 12,
 cha = 10
]
\DndMonsterDetails[
saving-throws = {Str +7, Con +5, Int +7},
skills = {Arcana +7, Athletics +7, Intimidation +6, Perception +4, Stealth +6},
damage-resistances = {psychic},
condition-immunities = {blinded, charmed},
senses = {darkvision 120 ft., passive perception 14},
languages = {Common, Kuran'zoi, telepathy 60 ft.},
challenge = 8,
]
\DndMonsterAction{Precog Reflexes}
The nemesis adds their Intelligence modifier (+4) to initiative rolls and to their AC (included in Armor Class).

\DndMonsterAction{Psychic Scar}
The nemesis is immune to any effect that would sense their emotions, read their thoughts, reveal they are lying, or reveal their alignment or location.

\DndMonsterSection{Actions}
\DndMonsterAction{Multiattack}
The nemesis makes two Golden Scythe attacks. They can replace one attack with a use of Kinetic Crush or Phase, if available.

\DndMonsterAction{Golden Scythe}
\textsl{Melee Weapon Attack:} 7 to hit, reach 10 ft., one target. \textsl{Hit:} 11 (2d6 + 4) slashing damage plus 13 (3d8) necrotic damage. If this attack reduces a creature who is in an astral body (as with the \textit{astral projection} spell) to 0 hit points, that creature's soul returns to their material body, and then their material body is also reduced to 0 hit points.

\DndMonsterAction{Kinetic Crush (4th-Order Power)}
The nemesis attempts to telekinetically grip one creature they can see within 60 feet of them. The target must make a DC 15 Strength saving throw. On a failed save, the target takes 27 (6d8) force damage, and their speed becomes 0 until the start of the nemesis's next turn. On a successful save, a target takes half as much damage and their speed isn't reduced.

\DndMonsterAction{Phase (1/Day; 4th-Order Power; Concentration)}
For 1 minute, the nemesis can move through other creatures and objects as if they were difficult terrain. The nemesis takes 5 (1d10) force damage if they end their turn inside an object. If this power ends while the nemesis is inside an object, they are shunted out of the object into the nearest unoccupied space of their choice.

\DndMonsterSection{Bonus Actions}
\DndMonsterAction{Jaunt (3/Day; 3rd-Order Power)}
The nemesis teleports up to 30 feet to an unoccupied space they can see.

\end{DndMonster}



\begin{DndMonster}[float=t]{Tormenauk}
\DndMonsterType{Large fiend (demon, category 3, Brute), chaotic evil}
\DndMonsterBasics[
armor-class = {17 (natural armor)},
hit-points = {\DndDice{13d10 + 65}},
speed = {30 ft.}
]
\DndMonsterAbilityScores[
 str = 20,
 dex = 10,
 con = 20,
 int = 15,
 wis = 15,
 cha = 18
]
\DndMonsterDetails[
saving-throws = {Con +8, Wis +5, Cha +7},
skills = {Intimidation +7, Perception +5},
senses = {darkvision 120 ft., soulsight 30 ft., passive perception 15},
languages = {Abyssal, Common, telepathy 120 ft.},
challenge = 8,
]
\DndMonsterAction{Lethe}
When the tormenauk's soul count is 0, they have advantage on attack rolls, disadvantage on saving throws, and their Intelligence score becomes 3 (-4). Additionally, the tormenauk must use their movement on each of their turns to move as close as possible to the nearest creature they can sense with their soulsight, and then if they are able, they must use their action to attack and attempt to kill that creature. The tormenauk can't act with any other purpose until they add 1 to their soul count.

\DndMonsterAction{Soul Devourer}
When the tormenauk reduces a creature who isn't a Construct or an Undead to 0 hit points or deals damage to a dying creature, the creature must make a DC 11 Wisdom saving throw. On a failed save, the tormenauk consumes the creature's soul and adds 1 to the tormenauk's soul count. A creature whose soul is consumed in this way immediately dies, and they can't be restored to life by any means short of a \textit{wish} spell.

\DndMonsterSection{Actions}
\DndMonsterAction{Multiattack}
The tormenauk makes two Slam attacks and one Many Maws attack.

\DndMonsterAction{Slam}
\textsl{Melee Weapon Attack:} 8 to hit, reach 5 ft., one target. \textsl{Hit:} 14 (2d8 + 5) bludgeoning damage. If the target is Large or smaller, they are grappled (escape DC 15), and they are restrained until the grapple ends. The tormenauk can grapple up to two targets at a time.

\DndMonsterAction{Many Maws}
\textsl{Melee Weapon Attack:} 8 to hit, reach 5 ft., one target grappled by the tormenauk. \textsl{Hit:} 21 (3d10 + 5) piercing damage.

\DndMonsterAction{Agony Wail (Costs 1 Soul)}
The tormenauk screeches, broadcasting a consumed soul's lifelong pain. Each creature within 30 feet of the tormenauk who can hear them must make a DC 15 Wisdom saving throw, taking 35 (10d6) psychic damage on a failed save, or half as much damage on a successful one. Creatures who knew the owner of the burned soul have disadvantage on the saving throw.

\DndMonsterSection{Reactions}
\DndMonsterAction{Share Agony (Costs 1 Soul)}
When the tormenauk takes damage, they choose a creature they can see within 30 feet of them to share the pain. The target must make a DC 15 Wisdom saving throw, taking psychic damage equal to the triggering damage taken by the tormenauk on a failed save, or half as much damage on a successful one.

\end{DndMonster}



\begin{DndMonster}[float=t]{Troll}
\DndMonsterType{Large giant (Brute), chaotic evil}
\DndMonsterBasics[
armor-class = {15 (natural armor)},
hit-points = {\DndDice{9d10 + 45}},
speed = {40 ft.}
]
\DndMonsterAbilityScores[
 str = 18,
 dex = 13,
 con = 20,
 int = 7,
 wis = 9,
 cha = 7
]
\DndMonsterDetails[
saving-throws = {Con +8},
skills = {Athletics +7, Perception +5},
senses = {darkvision 60 ft., passive perception 15},
languages = {Giant},
challenge = 5,
]
\DndMonsterAction{Relentless Hunger}
When the troll is reduced to 0 hit points by any damage other than acid or fire damage, they don't die or fall unconscious, and can continue moving and taking actions as usual. The troll only dies if they end their turn with 0 hit points, if acid or fire damage reduces them to 0 hit points, or if they take acid or fire damage while they have 0 hit points.

\DndMonsterAction{Sluggish}
For 1 hour after sleep or another period of unconsciousness, the troll's speed is halved and they can't use Multiattack.

\DndMonsterSection{Actions}
\DndMonsterAction{Multiattack}
The troll makes one Bite attack and two Claw attacks.

\DndMonsterAction{Bite}
\textsl{Melee Weapon Attack:} 7 to hit, reach 5 ft., one creature. \textsl{Hit:} 14 (3d6 + 4) piercing damage. If the target is not a Construct, an Elemental, or a Plant, the troll regains hit points equal to the damage dealt. If this attack reduces the target to 0 hit points, the troll regains twice as many hit points as they otherwise would with this attack.

\DndMonsterAction{Claw}
\textsl{Melee Weapon Attack:} 7 to hit, reach 10 ft., one target. \textsl{Hit:} 7 (1d6 + 4) slashing damage.

\DndMonsterAction{Crash Through}
The troll moves up to their speed in a straight line, smashing through mundane obstacles in their path. The troll can enter the spaces of Large or smaller creatures during this movement, forcing those creatures to make a DC 15 Strength saving throw. On a failed save, a creature takes 14 (4d6) bludgeoning damage and falls prone. On a successful save, a creature takes half as much damage and doesn't fall prone.

During this move, mundane objects that aren't worn or carried by a creature take 14 (4d6) damage when the troll enters their space. If the troll moves into the space of a Large or larger object and this damage doesn't destroy it, the troll's movement stops and they are stunned until the end of their next turn.

\DndMonsterSection{Reactions}
\DndMonsterAction{Spiteful Retort}
When the troll is reduced to 0 hit points and doesn't die, they can make a Bite attack against a creature within 5 feet of them.

\end{DndMonster}



\begin{DndMonster}[float=t]{Troll Whelp}
\DndMonsterType{Medium giant (Minion), chaotic evil}
\DndMonsterBasics[
armor-class = {15 (natural armor)},
hit-points = {\DndDice{12}},
speed = {40 ft.}
]
\DndMonsterAbilityScores[
 str = 18,
 dex = 13,
 con = 20,
 int = 7,
 wis = 9,
 cha = 7
]
\DndMonsterDetails[
saving-throws = {Con +8},
skills = {Athletics +7, Perception +2},
senses = {darkvision 60 ft., passive perception 12},
languages = {Giant},
challenge = 5,
]
\DndMonsterAction{Minion}
If the whelp takes damage from an attack or as the result of a failed saving throw, their hit points are reduced to 0. If the whelp takes damage from another effect, they die if the damage equals or exceeds their hit point maximum; otherwise they take no damage.

\DndMonsterAction{Relentless Hunger}
When the whelp is reduced to 0 hit points by any damage other than acid or fire damage, they don't die or fall unconscious, and can continue moving and taking actions as usual. The whelp only dies if they end their turn with 0 hit points, if acid or fire damage reduces them to 0 hit points, or if they take acid or fire damage while they have 0 hit points.

\DndMonsterAction{Ripping Claws}
Spaces within reach of one or more whelps are difficult terrain for Medium or smaller enemies. When an enemy leaves the reach of one or more whelps who can see them, the enemy takes 2 slashing damage.

\DndMonsterSection{Actions}
\DndMonsterAction{Bite (Group Attack)}
\textsl{Melee Weapon Attack:} 7 to hit, reach 5 ft., one target. \textsl{Hit:} 4 piercing damage, and if the target is not a Construct, an Elemental, or a Plant, each whelp who joined this attack regains all lost hit points.

\end{DndMonster}



\begin{DndMonster}[float*=b,width=\textwidth + 8pt]{Vampire}
\begin{multicols}{2}
\DndMonsterType{Medium undead (Skirmisher), neutral evil}
\DndMonsterBasics[
armor-class = {18 (natural armor)},
hit-points = {\DndDice{24d8 + 96}},
speed = {30 ft., fly 30 ft.}
]
\DndMonsterAbilityScores[
 str = 20,
 dex = 20,
 con = 18,
 int = 18,
 wis = 16,
 cha = 20
]
\DndMonsterDetails[
saving-throws = {Dex +10, Wis +8, Cha +10},
skills = {Deception +10, History +9, Insight +8, Intimidation +10, Perception +8, Persuasion +10, Stealth +10},
damage-resistances = {necrotic; bludgeoning, piercing, slashing from mundane attacks},
condition-immunities = {charmed, frightened},
senses = {blindsight 60 ft., darkvision 120 ft., passive perception 18},
languages = {the languages they knew in life},
challenge = 13,
]
\DndMonsterAction{Radiant Aversion}
Each time the vampire takes radiant damage, they take an extra 10 radiant damage.

\DndMonsterAction{Resting Place}
When the vampire drops to 0 hit points and is not in sunlight, running water, or their resting place, the vampire teleports to their resting place. While in their resting place, the vampire is stable. After spending 1 hour in their resting place with 0 hit points, they regain 1 hit point.

\DndMonsterAction{Turn Resistance}
The vampire has advantage on saving throws against any effect that turns undead.

\DndMonsterAction{Spellcasting (Utility)}
In addition to any other spells in this stat block, the vampire can cast the following spells, using Charisma as the spellcasting ability (spell save DC 18):
\begin{DndMonsterSpells}
  \DndInnateSpellLevel{charm person\textsuperscript{Casting Times}, \textit{detect thoughts}\textsuperscript{Casting Times}, \textit{disguise self}\textsuperscript{Casting Times}, \textit{sending}\textsuperscript{Casting Times}}
  \DndInnateSpellLevel[1]{clairvoyance\textsuperscript{Casting Times}, \textit{geas}\textsuperscript{Casting Times}}
\end{DndMonsterSpells}
\DndMonsterSection{Actions}
\DndMonsterAction{Multiattack}
The vampire makes two Claw attacks and, if possible, a Bite attack.

\DndMonsterAction{Claw}
\textsl{Melee Weapon Attack:} 10 to hit, reach 5 ft., one target. \textsl{Hit:} 12 (2d6 + 5) slashing damage plus 10 (3d6) necrotic damage. If the target is Large or smaller, they are grappled (escape DC 18). While grappled in this way, the target is restrained. Moving while grappling a Medium or smaller creature doesn't halve the vampire's speed. A vampire can have only one target grappled in this way at a time.

\DndMonsterAction{Bite}
\textsl{Melee Weapon Attack:} 10 to hit, reach 5 ft., one creature who is grappled by the vampire, incapacitated, or restrained. \textsl{Hit:} 8 (1d6 + 5) piercing damage plus 10 (3d6) necrotic damage. The target's hit point maximum is reduced by a number equal to the necrotic damage taken, and the vampire regains hit points equal to that number. The reduction lasts until reversed by a cure ailment power of 4th order or higher, a \textit{greater restoration} spell, or a similar supernatural effect. The target dies if this attack reduces their hit point maximum to 0. A Humanoid slain in this way rises the following midnight as a vampire spawn under the vampire's control.

\DndMonsterAction{Exsanguinating Mist (Recharge 6)}
The vampire, along with anything they are wearing or carrying, turns into a cloud of blood-sucking mist. When they do, they release any creature they were grappling, and if the vampire was grappled or restrained, that effect ends for them. Then the vampire flies up to twice their speed to an unoccupied space. During this move, the vampire ignores difficult terrain, doesn't provoke opportunity attacks, and can move through creatures and objects. Each creature the vampire passes through must make a DC 18 Constitution saving throw, taking 35 (10d6) necrotic damage on a failed save, or half as much damage on a successful one. At the end of this movement, the vampire transforms back into their previous form.

\DndMonsterSection{Bonus Actions}
\DndMonsterAction{Beguile}
The vampire targets one creature they can see within 30 feet of them and who can see them. The target must make a DC 18 Wisdom saving throw. On a failed save, the target uses their reaction, if available, to move up to their speed to a location chosen by the vampire, then either makes a melee attack against a target of the vampire's choice or falls prone (vampire's choice).

Creatures who can't be charmed automatically succeed on this saving throw, and creatures who have advantage on saving throws against being charmed make this saving throw with advantage.

\DndMonsterAction{Inhuman Agility}
The vampire moves up to their speed without provoking opportunity attacks.

\DndMonsterSection{Reactions}
\DndMonsterAction{Run, My Child}
When an ally the vampire can see takes damage, the vampire can command the ally to move up to their speed without provoking opportunity attacks.

\end{multicols}
\end{DndMonster}



\begin{DndMonster}[float=t]{Vampire Spawn}
\DndMonsterType{Medium undead (Skirmisher), neutral evil}
\DndMonsterBasics[
armor-class = {14},
hit-points = {\DndDice{15d8 + 45}},
speed = {30 ft., climb 30 ft.}
]
\DndMonsterAbilityScores[
 str = 18,
 dex = 18,
 con = 16,
 int = 12,
 wis = 11,
 cha = 16
]
\DndMonsterDetails[
saving-throws = {Wis +3, Cha +6},
skills = {Acrobatics +7, Intimidation +6, Perception +3, Persuasion +6, Stealth +7},
damage-resistances = {necrotic},
senses = {darkvision 60 ft., passive perception 13},
languages = {the languages they knew in life},
challenge = 5,
]
\DndMonsterAction{Radiant Aversion}
Each time the spawn takes radiant damage, they take an extra 10 radiant damage.

\DndMonsterSection{Actions}
\DndMonsterAction{Multiattack}
The spawn makes two Claw attacks. They can replace one attack with a Bite attack.

\DndMonsterAction{Claw}
\textsl{Melee Weapon Attack:} 7 to hit, reach 5 ft., one target. \textsl{Hit:} 11 (2d6 + 4) slashing damage, and if the target is Large or smaller, they are grappled (escape DC 15). Moving while grappling a Medium or smaller creature doesn't halve the spawn's speed. The spawn can have only one target grappled in this way at a time.

\DndMonsterAction{Bite}
\textsl{Melee Weapon Attack:} 7 to hit, reach 5 ft., one creature who is grappled by the spawn, incapacitated, or restrained. \textsl{Hit:} 7 (1d6 + 4) piercing damage plus 7 (2d6) necrotic damage. The target's hit point maximum is reduced by a number equal to the necrotic damage taken, and the spawn regains hit points equal to that number. The reduction lasts until the target finishes a long rest or is targeted by a cure ailment power of 4th order or higher, a \textit{greater restoration} spell, or a similar supernatural effect. The target dies if this attack reduces their hit point maximum to 0.

\DndMonsterSection{Bonus Actions}
\DndMonsterAction{Inhuman Speed}
The spawn moves up to their speed.

\end{DndMonster}



\begin{DndMonster}[float=t]{Voiceless Talker}
\DndMonsterType{Medium aberration (Controller), lawful evil}
\DndMonsterBasics[
armor-class = {16 (natural armor)},
hit-points = {\DndDice{14d8 + 28}},
speed = {30 ft., fly 30 ft.}
]
\DndMonsterAbilityScores[
 str = 10,
 dex = 16,
 con = 14,
 int = 19,
 wis = 15,
 cha = 14
]
\DndMonsterDetails[
saving-throws = {Con +5, Int +7, Wis +5},
skills = {Arcana +7, Deception +5, Insight +5, Perception +5, Persuasion +5, Stealth +6},
damage-resistances = {bludgeoning},
senses = {darkvision 120 ft., passive perception 15},
languages = {Deep Speech, Undercommon, telepathy 120 ft.},
challenge = 7,
]
\DndMonsterSection{Actions}
\DndMonsterAction{Multiattack}
The voiceless talker manifests a power and makes one Tentacle or Psionic Pistol attack.

\DndMonsterAction{Tentacle}
\textsl{Melee Weapon Attack:} 7 to hit, reach 10 ft., one target. \textsl{Hit:} 15 (2d10 + 4) psychic damage, and if the target is Large or smaller, they are grappled (escape DC 15).

\DndMonsterAction{Psionic Pistol}
\textsl{Ranged Weapon Attack:} 7 to hit, range 80/320 ft., one target. \textsl{Hit:} 15 (2d10 + 4) force damage.

\DndMonsterAction{Memory Thief (4th-Order Power)}
The voiceless talker psionically plunders the mind of a creature they can see within 30 feet of them. The target must make a DC 15 Intelligence saving throw. On a failed save, the target takes 22 (4d10) psychic damage, and until they finish a long rest or die, their proficiency bonus is cumulatively lowered by 1 and the voiceless talker gains a cumulative +2 bonus to damage rolls. On a successful save, a target takes half as much damage and doesn't have their proficiency bonus reduced.

A creature whose proficiency bonus drops to 0 can't form new thoughts or speak, and they have disadvantage on ability checks, attack rolls, and saving throws.

\DndMonsterAction{Flay (3/Day; 5th-Order Power)}
The voiceless talker shoots forth a 15-foot cone of pure psionic energy. Each creature in the area must make a DC 15 Intelligence saving throw, taking 28 (8d6) psychic damage on a failed save, or half as much damage on a successful one.

\DndMonsterAction{Guise (3rd-Order Power)}
The voiceless talker projects a psionic image over their body, transforming their appearance for 1 hour into that of a Medium creature they have seen. When they manifest this power, they can also change the appearance of any equipment they carry for the duration.

The changes wrought by this power fail to hold up to physical inspection. A creature can use an action to inspect the voiceless talker's appearance and make a DC 15 Intelligence (Investigation) check, noticing the image is a projection on a success.

\DndMonsterSection{Reactions}
\DndMonsterAction{Brain Drain}
When a creature grappled by the voiceless talker makes a saving throw against one of the voiceless talker's powers, the voiceless talker momentarily weakens the creature and the creature has disadvantage on the saving throw.

\end{DndMonster}



\begin{DndMonster}[float=t]{Voiceless Talker Artillerist}
\DndMonsterType{Medium aberration (Artillery), lawful evil}
\DndMonsterBasics[
armor-class = {16 (natural armor)},
hit-points = {\DndDice{20d8 + 60}},
speed = {30 ft., fly 30 ft.}
]
\DndMonsterAbilityScores[
 str = 12,
 dex = 16,
 con = 16,
 int = 21,
 wis = 17,
 cha = 16
]
\DndMonsterDetails[
saving-throws = {Con +7, Int +9, Wis +7},
skills = {Arcana +9, Deception +7, Insight +7, Perception +7, Persuasion +7, Stealth +7},
damage-resistances = {bludgeoning},
senses = {darkvision 120 ft., passive perception 17},
languages = {Deep Speech, Undercommon, telepathy 120 ft.},
challenge = 10,
]
\DndMonsterAction{Phasing Rifle (1/Turn)}
When the artillerist hits a target with their Psionic Rifle attack, the artillerist can attempt to teleport the target (no action required). The target must succeed on a DC 17 Wisdom saving throw or be teleported to an unoccupied space the artillerist can see within 60 feet of the artillerist.

\DndMonsterSection{Actions}
\DndMonsterAction{Multiattack}
The artillerist makes two Psionic Rifle attacks, or they manifest a power and make one Phasing Tentacles attack.

\DndMonsterAction{Psionic Rifle}
\textsl{Ranged Weapon Attack:} 9 to hit, range 150/600 ft., one target. \textsl{Hit:} 21 (3d10 + 5) force damage.

\DndMonsterAction{Phasing Tentacles}
\textsl{Melee Weapon Attack:} 9 to hit, reach 10 ft., one target. \textsl{Hit:} 21 (3d10 + 5) psychic damage, and if the target is Large or smaller, the artillerist can choose one of the following effects: the target is grappled (escape DC 17) or teleported up to 15 feet to an unoccupied space the artillerist can see.

\DndMonsterAction{*Memory Thief (5th-Order Power)}
The artillerist psionically plunders the mind of a creature they can see within 30 feet of them. The target must make a DC 17 Intelligence saving throw. On a failed save, the target takes 33 (6d10) psychic damage, and until they finish a long rest or die, their proficiency bonus is lowered by 1 and the artillerist gains a cumulative +2 bonus to damage rolls. On a successful save, a target takes half as much damage and doesn't have their proficiency bonus reduced.

A creature whose proficiency bonus drops to 0 can't form new thoughts or speak, and they have disadvantage on ability checks, attack rolls, and saving throws.

\DndMonsterAction{Guise (3rd-Order Power)}
The artillerist projects a psionic image over their body, transforming their appearance for 1 hour into that of a Medium creature they have seen. When they manifest this power, they can also change the appearance of any equipment they carry for the duration.

The changes wrought by this power fail to hold up to physical inspection. A creature can use an action to inspect the artillerist's appearance and make a DC 17 Intelligence (Investigation) check, noticing the image is a projection on a success.

\DndMonsterSection{Bonus Actions}
\DndMonsterAction{*Vanish for One (3/Day; 3rd-Order Power; Concentration)}
The artillerist cloaks themself from the sight of a creature the artillerist can see within 60 feet of them. The creature must succeed on a DC 17 Wisdom saving throw or the artillerist becomes invisible to the creature for 1 minute (save ends at end of turn). This effect ends early if the artillerist attacks the creature, deals damage to them, or creates an effect that forces them to make a saving throw.

\end{DndMonster}



\begin{DndMonster}[float=t]{War Spider}
\DndMonsterType{Huge beast (Brute), U}
\DndMonsterBasics[
armor-class = {15 (natural armor)},
hit-points = {\DndDice{7d12 + 28}},
speed = {40 ft., climb 40 ft.}
]
\DndMonsterAbilityScores[
 str = 18,
 dex = 14,
 con = 18,
 int = 2,
 wis = 11,
 cha = 4
]
\DndMonsterDetails[
senses = {blindsight 10 ft., darkvision 60 ft., passive perception 10},
challenge = 3,
]
\DndMonsterAction{Baby Burst}
When the spider dies, a \textbf{swarm of spiders} bursts forth from their remains and acts on the same initiative count as the spider did.

\DndMonsterAction{Rider Launcher}
An allied rider who jumps off the spider has a long jump of 30 feet and a high jump of 15 feet, with or without a running start. If an allied rider jumps off the spider, the first melee weapon attack the ally makes on the same turn has advantage.

\DndMonsterAction{Spider Climb}
The spider can climb difficult surfaces, including upside down on ceilings, without needing to make an ability check.

\DndMonsterAction{Web Walker}
The spider ignores movement restrictions caused by webbing.

\DndMonsterAction{Wide Back}
Two Small or Tiny creatures can share a space while riding the spider, allowing up to eighteen Small or seventy-two Tiny creatures to ride the spider at once.

\DndMonsterSection{Actions}
\DndMonsterAction{Bite}
\textsl{Melee Weapon Attack:} 6 to hit, reach 5 ft., one target. \textsl{Hit:} 11 (2d6 + 4) piercing damage plus 7 (2d6) poison damage.

\DndMonsterAction{Bladed Leg}
\textsl{Melee Weapon Attack:} 6 to hit, reach 10 ft., one target. \textsl{Hit:} 11 (2d6 + 4) slashing damage.

\DndMonsterAction{Web Spray}
The spider sprays webbing from their abdomen in a 15-foot cube. Each creature in that area must succeed on a DC 14 Dexterity saving throw or be restrained by the webbing. A creature can use their action to make a DC 14 Strength or Dexterity check, freeing themself or another creature they can reach on a success.

\DndMonsterAction{Trample (Recharge 6)}
The spider can move up to their speed and move through the spaces of other creatures as if they were difficult terrain. The spider can make one Bladed Leg attack against each creature whose space they move into during the move. A creature hit by this attack can't take reactions until the start of their next turn.

\end{DndMonster}



\begin{DndMonster}[float=t]{Water Spark}
\DndMonsterType{Medium elemental (water, Minion), A}
\DndMonsterBasics[
armor-class = {19 (natural armor)},
hit-points = {\DndDice{27}},
speed = {30 ft., swim 50 ft., fly 30 ft.}
]
\DndMonsterAbilityScores[
 str = 20,
 dex = 18,
 con = 16,
 int = 9,
 wis = 16,
 cha = 11
]
\DndMonsterDetails[
senses = {darkvision 60 ft., passive perception 13},
languages = {Aquan, Common},
challenge = 20,
]
\DndMonsterAction{Minion}
If the spark takes damage from an attack or as the result of a failed saving throw, their hit points are reduced to 0. If the spark takes damage from another effect, they die if the damage equals or exceeds their hit point maximum; otherwise they take no damage.

\DndMonsterAction{Whirlpool Suction}
If an enemy starts their turn within 5 feet of three or more water sparks, the enemy must succeed on a Strength saving throw or have their speed reduced to 0. The DC for this saving throw equals 14 plus the number of water sparks within 5 feet of the enemy.

\DndMonsterSection{Actions}
\DndMonsterAction{Whelming Wave (Group Attack)}
\textsl{Melee or Ranged Spell Attack:} 11 to hit, reach 5 ft. or range 30 ft., one target. \textsl{Hit:} 10 bludgeoning damage, and if the attack was made by more than one spark, the target must succeed on a Dexterity saving throw or be knocked prone and be unable to stand up until the end of the sparks' next turn. The DC for this saving throw equals 14 plus the number of sparks who joined the attack.

\end{DndMonster}


\clearpage

\begin{DndMonster}[float=t]{Wight}
\DndMonsterType{Medium undead (Soldier), neutral evil}
\DndMonsterBasics[
armor-class = {14 (natural armor)},
hit-points = {\DndDice{10d8 + 20}},
speed = {30 ft.}
]
\DndMonsterAbilityScores[
 str = 16,
 dex = 13,
 con = 15,
 int = 12,
 wis = 13,
 cha = 16
]
\DndMonsterDetails[
skills = {Intimidation +5, Perception +3},
damage-resistances = {necrotic},
damage-immunities = {poison},
condition-immunities = {exhaustion, poisoned},
senses = {darkvision 60 ft., passive perception 13},
languages = {the languages they knew in life},
challenge = 3,
]
\DndMonsterSection{Actions}
\DndMonsterAction{Multiattack}
The wight makes one Longsword attack and one Vitality Theft attack, or they make two Heavy Crossbow attacks.

\DndMonsterAction{Longsword}
\textsl{Melee Weapon Attack:} 5 to hit, reach 5 ft., one target. \textsl{Hit:} 7 (1d8 + 3) slashing damage.

\DndMonsterAction{Heavy Crossbow}
\textsl{Ranged Weapon Attack:} 3 to hit, range 100/400 ft., one target. \textsl{Hit:} 6 (1d10 + 1) piercing damage.

\DndMonsterAction{Vitality Theft}
\textsl{Melee Spell Attack:} 5 to hit, reach 5 ft., one creature. \textsl{Hit:} 9 (2d8) necrotic damage, and the wight regains a number of hit points equal to half the damage dealt. Until the start of the wight's next turn, the target has disadvantage on attack rolls made against creatures other than the wight.

\DndMonsterSection{Reactions}
\DndMonsterAction{Decaying Guard}
When an ally the wight can see within 5 feet of them is hit by an attack, the wight forces the attacker to reroll the attack and to instead target the wight or a willing ally within 5 feet of the original target.

\end{DndMonster}



\begin{DndMonster}[float=t]{Wight Deathguard}
\DndMonsterType{Medium undead (undead, Retainer), A}
\DndMonsterBasics[
armor-class = {15 (medium armor)},
hit-points = {\DndDice{Eight times their level (number of d10 Hit Dice equal to their level)}},
speed = {30 ft.}
]
\DndMonsterAbilityScores[
 str = 16,
 dex = 10,
 con = 12,
 int = 10,
 wis = 10,
 cha = 12
]
\DndMonsterDetails[
saving-throws = {Str +PB, Dex +PB, Con +PB, Int +PB, Wis +PB, Cha +PB},
skills = {Intimidation +3 plus PB, Perception +0 plus PB},
damage-resistances = {necrotic},
damage-immunities = {poison},
condition-immunities = {exhaustion, poisoned},
senses = {darkvision 60 ft., passive perception 10 plus PB},
languages = {the languages they knew in life},
]
\DndMonsterSection{Actions}
\DndMonsterAction{Signature Attack (Longsword)}
\textsl{Melee Attack?} 3 plus PB to hit, reach 5 ft., one target. \textsl{Hit:} 1d8 plus PB slashing damage. Beginning at 7th level, the deathguard can make this attack twice, instead of once, when they take the Attack action on their turn.

\DndMonsterAction{3rd Level: Blood for Blood (3/Day)}
As a reaction to an ally the deathguard can see within 5 feet of them being hit with an attack, the deathguard redirects the attack to themself, potentially causing the attack to miss. If the attacker is within 5 feet of the deathguard, the deathguard can make a signature attack against the attacker.

\DndMonsterAction{5th Level: Soul Thief (3/Day)}
As an action, the deathguard shrouds themself in dark energy. The deathguard regains PBd6 hit points, and each enemy within 5 feet of the deathguard must make a DC 10 plus PB Wisdom saving throw. On a failed save, a target takes PBd6 necrotic damage. On a successful save, a target takes half as much damage.

\DndMonsterAction{7th Level: This Way! (3/Day)}
As an action, the deathguard moves up to their speed without provoking opportunity attacks. Each ally the deathguard passes within 5 feet of during the move can use a reaction to move up to their speed in the same direction as the deathguard without provoking opportunity attacks. Until the start of the deathguard's next turn, attacks against these allies have disadvantage.

\end{DndMonster}



\begin{DndMonster}[float=t]{Wobalas}
\DndMonsterType{Medium fiend (demon, category 4, Artillery), chaotic evil}
\DndMonsterBasics[
armor-class = {17 (natural armor)},
hit-points = {\DndDice{24d8 + 48}},
speed = {30 ft.}
]
\DndMonsterAbilityScores[
 str = 14,
 dex = 22,
 con = 14,
 int = 16,
 wis = 18,
 cha = 18
]
\DndMonsterDetails[
saving-throws = {Dex +10, Wis +8, Cha +8},
skills = {Deception +8, Intimidation +8, Perception +8},
senses = {darkvision 120 ft., soulsight 30 ft., passive perception 18},
languages = {Abyssal, Common, telepathy 120 ft.},
challenge = 10,
]
\DndMonsterAction{Lethe}
When the wobalas's soul count is 0, they have advantage on attack rolls, disadvantage on saving throws, and their Intelligence score becomes 3 (-4). Additionally, the wobalas must use their movement on each of their turns to move as close as possible to the nearest creature they can sense with their soulsight, and then if they are able, they must use their action to attack and attempt to kill that creature. The wobalas can't act with any other purpose until they add 1 to their soul count.

\DndMonsterAction{Soul Devourer}
When the wobalas reduces a creature who isn't a Construct or an Undead to 0 hit points or deals damage to a dying creature, the creature must make a DC 11 Wisdom saving throw. On a failed save, the wobalas consumes the creature's soul and adds 1 to the wobalas's soul count. A creature whose soul is consumed in this way immediately dies, and they can't be restored to life by any means short of a \textit{wish} spell.

\DndMonsterSection{Actions}
\DndMonsterAction{Multiattack}
The wobalas makes three Despair Bolt attacks.

\DndMonsterAction{Despair Bolt}
\textsl{Ranged Weapon Attack:} 10 to hit, range 150/600 ft., one creature. \textsl{Hit:} 16 (3d6 + 6) piercing damage plus 10 (3d6) psychic damage. The wobalas can burn 1 soul to make the target frightened of them for 1 minute. If the target is immune to being frightened, the wobalas regains the burned soul.

\DndMonsterAction{Banishing Touch}
\textsl{Melee Spell Attack:} 8 to hit, reach 5 ft., one creature. \textsl{Hit:} 22 (4d10) psychic damage, and the target is teleported up to 30 feet to an unoccupied space the wobalas can see.

\DndMonsterSection{Bonus Actions}
\DndMonsterAction{Flee, Mortal (Costs 1 Soul)}
The wobalas chooses a creature they can see within 60 feet of them. The target must succeed on a DC 16 Charisma saving throw or use their reaction, if available, to move their speed away from the wobalas by the most direct route possible, with no regard for their own safety. Creatures who can't be frightened succeed on this saving throw automatically.

\end{DndMonster}



\begin{DndMonster}[float=t]{Wraith}
\DndMonsterType{Medium undead (incorporeal, Controller), chaotic evil}
\DndMonsterBasics[
armor-class = {12},
hit-points = {\DndDice{9d8 + 18}},
speed = {0 ft., fly 60 ft. \textit{(hover)}}
]
\DndMonsterAbilityScores[
 str = 6,
 dex = 14,
 con = 15,
 int = 12,
 wis = 14,
 cha = 17
]
\DndMonsterDetails[
damage-resistances = {acid; cold; fire; lightning; necrotic; thunder; bludgeoning, piercing, slashing from mundane attacks},
damage-immunities = {poison, psychic},
condition-immunities = {charmed, dazed, exhaustion, grappled, paralyzed, petrified, poisoned, prone, restrained, unconscious},
senses = {darkvision 60 ft., passive perception 12},
languages = {the languages they knew in life},
challenge = 5,
]
\DndMonsterAction{Agonizing Phasing}
The wraith can move through other creatures and objects as if they were difficult terrain. The wraith takes 5 (1d10) force damage if they end their turn inside an object. A creature takes 5 (1d10) psychic damage the first time a wraith passes through them on a turn.

\DndMonsterAction{Throes of Oblivion}
When the wraith dies, they collapse inward, creating a burst of painful psychic energy. Each creature within 20 feet of the wraith must make a DC 14 Constitution saving throw, taking 18 (4d8) psychic damage on a failed save, or half as much damage on a successful one.

\DndMonsterSection{Actions}
\DndMonsterAction{Agonizing Touch}
\textsl{Melee Spell Attack:} 6 to hit, reach 5 ft., one creature. \textsl{Hit:} 14 (4d6) psychic damage, and the target must succeed on a DC 14 Wisdom saving throw or be dazed (save ends at end of turn).

Each time a target dazed in this way fails a saving throw to end the condition, their hit point maximum is halved. This effect is cumulative, and the reduction lasts until the target finishes a short or long rest.

If a Humanoid dies within 1 minute of being hit by this attack, their spirit rises as a specter under the wraith's control in an unoccupied space nearest to where that Humanoid died.

\DndMonsterAction{Psychic Enervation}
\textsl{Ranged Spell Attack:} 6 to hit, range 60 ft., one creature. \textsl{Hit:} 10 (3d6) psychic damage, and the target must succeed on a DC 14 Wisdom saving throw or fall prone and be frightened of the wraith for 1 minute (save ends at end of turn). While frightened in this way, the target can't stand up.

\DndMonsterSection{Reactions}
\DndMonsterAction{Denied Vitality}
When a creature who the wraith can see within 30 feet of them regains hit points, the wraith can attempt to sap the life energy restoring the creature. The creature must succeed on a DC 14 Constitution saving throw or only regain half the number of hit points they otherwise would regain.

\end{DndMonster}



\begin{DndMonster}[float=t]{Wyvern}
\DndMonsterType{Large dragon (Controller), U}
\DndMonsterBasics[
armor-class = {12 (natural armor)},
hit-points = {\DndDice{14d10 + 42}},
speed = {20 ft., fly 60 ft.}
]
\DndMonsterAbilityScores[
 str = 17,
 dex = 12,
 con = 16,
 int = 4,
 wis = 14,
 cha = 6
]
\DndMonsterDetails[
skills = {Athletics +6, Perception +5},
damage-immunities = {acid},
senses = {darkvision 60 ft., passive perception 15},
challenge = 6,
]
\DndMonsterAction{Stubborn Rage}
The wyvern is immune to the frightened and stunned conditions while they have fewer than half their hit points or are within 60 feet of another wyvern they can hear or see.

\DndMonsterSection{Actions}
\DndMonsterAction{Multiattack}
The wyvern makes one Stinger attack and one Bite attack. They can replace one attack with a use of Tail Sweep.

\DndMonsterAction{Stinger}
\textsl{Melee Weapon Attack:} 6 to hit, reach 15 ft., one target. \textsl{Hit:} 13 (3d6 + 3) piercing damage. The target must succeed on a DC 14 Constitution saving throw or take 7 (2d6) acid damage at the start of each of their turns for 1 minute (save ends at end of turn). A cure ailment power or a \textit{lesser restoration} spell ends this effect on a target.

\DndMonsterAction{Bite}
\textsl{Melee Weapon Attack:} 6 to hit, reach 10 ft., one target. \textsl{Hit:} 16 (2d12 + 3) piercing damage. If the target is Large or smaller, they are grappled (escape DC 13). Until this grapple ends, the target is restrained and the wyvern can't bite another target.

\DndMonsterAction{Tail Sweep}
The wyvern swipes their tail in a 15-foot cone. Each creature in that area must succeed on a DC 14 Strength saving throw or take 14 (2d10 + 3) bludgeoning damage and be pushed up to 10 feet away from the wyvern.

\DndMonsterSection{Bonus Actions}
\DndMonsterAction{Seek}
The wyvern takes the Search action.

\end{DndMonster}



\begin{DndMonster}[float*=b,width=\textwidth + 8pt]{Xaantikorijek}
\begin{multicols}{2}
\DndMonsterType{Gargantuan dragon (Solo), N}
\DndMonsterBasics[
armor-class = {20 (natural armor)},
hit-points = {\DndDice{28d20 + 224}},
speed = {40 ft., burrow 40 ft., fly 80 ft.}
]
\DndMonsterAbilityScores[
 str = 28,
 dex = 10,
 con = 27,
 int = 24,
 wis = 16,
 cha = 20
]
\DndMonsterDetails[
saving-throws = {Con +15, Wis +10, Cha +12},
skills = {Arcana +21, History +21, Nature +21, Perception +10, Religion +21},
damage-immunities = {lightning, thunder},
condition-immunities = {charmed, dazed, frightened, paralyzed, petrified, stunned},
senses = {blindsight 120 ft., truesight 60 ft., passive perception 20},
languages = {all},
challenge = 23,
]
\DndMonsterAction{Static Shield (3/Day)}
When Xaantikorijek fails a saving throw, he can succeed instead. When he does, his sight, blindsight, and truesight are reduced to a range of 30 feet until the start of his next turn.

\DndMonsterAction{Storm of the Gods}
Lightning and thunder damage dealt by Xaantikorijek ignore damage resistance.

\DndMonsterSection{Actions}
\DndMonsterAction{Multiattack}
Xaantikorijek makes one Bite attack and two Claw attacks.

\DndMonsterAction{Bite}
\textsl{Melee Weapon Attack:} 16 to hit, reach 15 ft., one target. \textsl{Hit:} 20 (2d10 + 9) piercing damage plus 5 (1d10) lightning damage, and Xaantikorijek can move the target up to 20 feet in any direction.

\DndMonsterAction{Claw}
\textsl{Melee Weapon Attack:} 16 to hit, reach 10 ft., one target. \textsl{Hit:} 16 (2d6 + 9) slashing damage plus 7 (2d6) thunder damage, and the target is knocked prone.

\DndMonsterAction{Fulminating Breath (Recharge 5\textendash 6)}
Xaantikorijek exhales lightning in a line that is 120 feet long and 30 feet wide. Each creature in that area must make a DC 23 Dexterity saving throw. On a failed save, a creature takes 66 (12d10) lightning damage and has disadvantage on saving throws until the end of Xaantikorijek's next turn. On a successful save, a creature takes half as much damage and suffers no other effect.

\DndMonsterSection{Bonus Actions}
\DndMonsterAction{Lightning Discharge}
Xaantikorijek makes one Bite attack. On a hit, Xaantikorijek's teeth spark with energy, dealing an extra 11 (2d10) lightning damage to the target and to one other enemy within 60 feet of the target.

\DndMonsterSection{Reactions}
\DndMonsterAction{Voice of Negation}
When a creature Xaantikorijek can see within 60 feet of him uses an action or bonus action to cast a spell, Xaantikorijek speaks an arcane word, forcing the creature to make a DC 23 saving throw using their spellcasting ability. On a failed save, the target takes 16 (3d10) thunder damage and is unable to cast spells until the end of their turn, but the spell slot is not expended and the creature's action is not lost.

\DndMonsterSection{Legendary Actions}
Xaantikorijek has three villain actions. He can take each action once during an encounter after an enemy's turn. He can take these actions in any order but can use only one per round.

\begin{DndMonsterLegendaryActions}
\DndMonsterLegendaryAction{Action 1: Voice of Reverence}{Xaantikorijek speaks an ancient word of authority into the mind of each creature of his choice within 60 feet of him. Each target must succeed on a DC 20 Wisdom saving throw or fall prone. A creature who falls prone in this way takes 22 (4d10) thunder damage if they stand up before the end of Xaantikorijek's next turn.
}
\DndMonsterLegendaryAction{Action 2: Fulguration}{Xaantikorijek sheds his physical form to become a being of lightning until the end of his next turn. While in this form, Xaantikorijek gains the following benefits:


\begin{itemize}
\item He is resistant to all damage.
\item He can move through a space as narrow as 1 inch wide without squeezing.
\item His movement doesn't provoke opportunity attacks.
\item He can move through any creature's space, and when he moves into a creature's space for the first time on a turn, that creature takes 16 (3d10) lightning damage.
\end{itemize}
}
\DndMonsterLegendaryAction{Action 3: Voice of the Ages}{Xaantikorijek recites an ancient lament that reverberates in nearby souls. Each enemy within 60 feet of Xaantikorijek who can hear him must make a DC 20 Wisdom saving throw. On a failed save, a creature takes 44 (8d10) psychic damage, falls prone, and is unable to stand up for 1 minute as they weep uncontrollably (save ends at end of turn). On a successful save, a creature takes half as much damage and suffers no other effect.
}
\end{DndMonsterLegendaryActions}
\end{multicols}
\end{DndMonster}



\begin{DndMonster}[float*=b,width=\textwidth + 8pt]{Xorannox}
\begin{multicols}{2}
\DndMonsterType{Large aberration (Solo), lawful evil}
\DndMonsterBasics[
armor-class = {18 (natural armor)},
hit-points = {\DndDice{24d10 + 96}},
speed = {0 ft., fly 20 ft. \textit{(hover)}}
]
\DndMonsterAbilityScores[
 str = 10,
 dex = 16,
 con = 18,
 int = 19,
 wis = 14,
 cha = 18
]
\DndMonsterDetails[
saving-throws = {Int +9, Wis +7, Cha +9},
skills = {Arcana +9, Deception +9, Insight +7, Intimidation +9, Perception +7},
condition-immunities = {charmed, flanked, frightened, prone},
senses = {darkvision 120 ft., passive perception 17},
languages = {Common, Deep Speech, Undercommon},
challenge = 14,
]
\DndMonsterAction{Painful Resistance (3/Day)}
Eight eyes float around Xorannox and create his Eye Psionics. If Xorannox fails a saving throw, he can destroy a random floating eye (chosen by rolling a d8) and succeed instead. If an eye is destroyed, a new eye pops of out of Xorannox's face to replace it 24 hours later.

\DndMonsterSection{Actions}
\DndMonsterAction{Bite}
\textsl{Melee Weapon Attack:} 5 to hit, reach 5 ft., one target. \textsl{Hit:} 14 (4d6) piercing damage.

\DndMonsterAction{Eye Psionics}
Xorannox creates three of the following psionic eye effects. Unless otherwise stated, each eye targets a creature who he can see within 120 feet of him. He can't use the same effect twice on a turn.


\begin{description}
\item[1. Charm Beam.]
The targeted creature must succeed on a DC 17 Wisdom saving throw or be charmed by Xorannox for 1 hour, or until he harms the creature or one of their allies.

\item[2. Compulsion Beam.]
The targeted creature must succeed on a DC 17 Intelligence saving throw or use their reaction, if available, to move up to their speed toward the closest ally they can see and make one weapon attack against them. Creatures who can't be charmed are unaffected.

\item[3. Memory Beam.]
The targeted creature must make a DC 17 Intelligence saving throw. On a failed save, if the creature has unexpended spell slots, they expend a spell slot of their highest remaining level, with no effect. If they don't have unexpended spell slots, they instead lose proficiency with a weapon they're holding or carrying for 1 minute (save ends at end of turn).

\item[4. Toxic Vapors.]
Xorannox chooses a point he can see within 120 feet of him. Each creature within 10 feet of the point must succeed on a DC 17 Constitution saving throw or be poisoned for 1 minute (save ends at end of turn).

\item[5. Telekinetic Field.]
Xorannox chooses a point he can see within 120 feet of him. Each enemy within 10 feet of the point must succeed on a DC 17 Strength saving throw or be moved up to 30 feet in any direction and take 7 (2d6) force damage at the end of the move. Additionally, Xorannox can target any object within 10 feet of the point that weighs 300 pounds or less and isn't being worn or carried, moving it up to 30 feet in any direction. Xorannox can also exert fine control on objects with this effect, such as manipulating a simple tool or opening a door or a container.

\item[6. Lightning Bolt.]
Xorannox shoots a 5-foot-wide, 60-foot-long line of lightning. Each creature in the line must succeed on a DC 17 Dexterity saving throw, taking 36 (8d8) lightning damage on a failed save, or half as much damage on a successful one.

\item[7. Explosion.]
Xorannox chooses a point he can see within 120 feet of him. Each creature within 10 feet of the point must make a DC 17 Dexterity saving throw, taking 36 (8d8) fire damage on a failed save, or half as much damage on a successful one.

\item[8. Necrosis Beam.]
The targeted creature must succeed on a DC 17 Constitution saving throw or immediately take 55 (10d10) necrotic damage and lose 7 (2d6) hit points at the start of each of their turns. If this effect reduces a target's hit points to 0, they die. The effect ends if Xorannox dies or can't see the target. This psionic effect can't be used again while this effect persists.

\end{description}

\DndMonsterSection{Bonus Actions}
\DndMonsterAction{The Great Eye (Recharge 6)}
Xorannox's central eye turns solid black and projects a 150-foot cone of energy. If Xorannox or any creature in that area is affected by a spell, the spell's effects immediately end for that creature.

\DndMonsterSection{Reactions}
\DndMonsterAction{Cower!}
When a creature hits Xorannox with an attack, he shoots a psionic fear beam at them. The creature must succeed on a DC 17 Wisdom saving throw or take 21 (6d6) psychic damage and be frightened of Xorannox for 1 minute (save ends at end of turn).

\DndMonsterSection{Legendary Actions}
Xorannox has three villain actions. He can take each action once during an encounter after an enemy's turn. He can take these actions in any order but can use only one per round.

\begin{DndMonsterLegendaryActions}
\DndMonsterLegendaryAction{Action 1: Disruption Beams}{Xorannox shoots a psionic disruption beam at three creatures he can see within 120 feet of him. Each target must make a DC 17 Wisdom saving throw. On a failed save, whenever the target makes more than one weapon attack on a turn or casts a spell that isn't a cantrip, they take 14 (4d6) psychic damage. This effect lasts for 1 minute (save ends at end of turn).
}
\DndMonsterLegendaryAction{Action 2: Disappearing Act}{Xorannox turns invisible until the end of his next turn and teleports up to 120 feet to an unoccupied space he can see.
}
\DndMonsterLegendaryAction{Action 3: Megabeam}{Xorannox unleashes the effects of all of his remaining psionic eyes at once in a 20-foot-wide, 120-foot-long line. Each creature in that area must make a save against two random eye effects from Xorannox's Eye
}
\DndMonsterLegendaryAction{Psionics (reroll duplicates for each target)}{If the effect usually affects an area, it instead only affects individual creatures.
}
\end{DndMonsterLegendaryActions}
\end{multicols}
\end{DndMonster}



\begin{DndMonster}[float*=b,width=\textwidth + 8pt]{Yserthrax}
\begin{multicols}{2}
\DndMonsterType{Gargantuan dragon (Solo), neutral evil}
\DndMonsterBasics[
armor-class = {19 (natural armor)},
hit-points = {\DndDice{24d20 + 168}},
speed = {40 ft., fly 80 ft., swim 40 ft.}
]
\DndMonsterAbilityScores[
 str = 28,
 dex = 10,
 con = 25,
 int = 22,
 wis = 17,
 cha = 18
]
\DndMonsterDetails[
saving-throws = {Con +14, Wis +10, Cha +11},
skills = {Arcana +13, Deception +11, Insight +17, Perception +17, Persuasion +11, Stealth +7},
damage-immunities = {poison, psychic},
condition-immunities = {blinded, charmed, dazed, frightened, paralyzed, petrified, poisoned, stunned},
senses = {blindsight 120 ft., truesight 60 ft., passive perception 27},
languages = {Common, Deep Speech, Draconic},
challenge = 22,
]
\DndMonsterAction{Amphibious}
Yserthrax can breathe air and water.

\DndMonsterAction{Corruption}
Poison damage dealt by Yserthrax ignores damage resistance and treats damage immunity as damage resistance. Whenever Yserthrax deals poison damage to a creature who isn't a Construct or an Undead, that creature must succeed on a DC 19 Constitution saving throw or be corrupted (save ends at end of turn). The corruption lasts until cured by a cure ailment power, \textit{lesser restoration} spell, or similar supernatural effect.

While corrupted, a creature takes 14 (4d6) psychic damage at the start of each of their turns as their mind and body are ravaged by reality-warping anomalies. If a creature dies while corrupted, their body melts into a gruesome puddle, becoming a \textbf{gibbering mouther} under Yserthrax's control.

\DndMonsterAction{Withering Resistance (3/Day)}
When Yserthrax fails a saving throw, she can take 16 (3d10) necrotic damage and succeed instead.

\DndMonsterSection{Actions}
\DndMonsterAction{Multiattack}
Yserthrax makes one Bite attack and two Claw attacks.

\DndMonsterAction{Bite}
\textsl{Melee Weapon Attack:} 16 to hit, reach 15 ft., one target. \textsl{Hit:} 20 (2d10 + 9) piercing damage plus 11 (2d10) poison damage.

\DndMonsterAction{Claw}
\textsl{Melee Weapon Attack:} 16 to hit, reach 10 ft., one target. \textsl{Hit:} 16 (2d6 + 9) slashing damage, and Yserthrax can move the target up to 15 feet horizontally.

\DndMonsterAction{Corrupting Breath (Recharge 5\textendash 6)}
Yserthrax exhales poison in a 90-foot cone. Each creature in that area must make a DC 22 Constitution saving throw. On a failed save, a target takes 56 (16d6) poison damage and is poisoned until the end of their next turn. On a successful save, a target takes half as much damage and isn't poisoned.

\DndMonsterSection{Bonus Actions}
\DndMonsterAction{Ulcerate}
Yserthrax magically aggravates a creature she can see within 60 feet of her who is corrupted by her Corruption trait. The target must succeed on a DC 19 Constitution saving throw or be dazed until the end of their next turn.

\DndMonsterSection{Reactions}
\DndMonsterAction{Eldritch Disruption}
When a creature Yserthrax can see within 60 feet of her hits her with an attack or succeeds on a saving throw, Yserthrax conjures a reality-warping anomaly around them. The creature must choose to either take 10 (3d6) psychic damage or to reroll the attack or a saving throw and use the new result.

\DndMonsterSection{Legendary Actions}
Yserthrax has three villain actions. She can take each action once during an encounter after an enemy's turn. She can take these actions in any order but can use only one per round.

\begin{DndMonsterLegendaryActions}
\DndMonsterLegendaryAction{Action 1: Summon Monoliths}{Yserthrax conjures three monoliths of otherworldly green crystal, which grow out of the ground centered on three points she can see within 500 feet of her. Each monolith is a 10-foot-radius, 60-foot-high cylinder with AC 19, 100 hit points, and immunity to poison and psychic damage. The monoliths don't block Yserthrax's vision. Each creature within 10 feet of a monolith other than Yserthrax is vulnerable to psychic damage. When a monolith is destroyed, each creature within 30 feet of it must make a DC 19 Wisdom saving throw, taking 21 (6d6) psychic damage on a failed save, or half as much damage on a successful one.
}
\DndMonsterLegendaryAction{Action 2: Spatial Alteration}{Yserthrax bats her wings erratically, releasing foul energy that warps space around her. Each creature within 60 feet of her must succeed on a DC 19 Charisma saving throw or be teleported up to 120 feet to an unoccupied space Yserthrax can see.
}
\DndMonsterLegendaryAction{Action 3: Elder Scream}{Yserthrax unleashes an eldritch scream, piercing the mind of each enemy she can see within 120 feet of her. Each target must make a DC 19 Wisdom saving throw, taking 77 (14d10) psychic damage on a failed save, or half as much damage on a successful one.
}
\end{DndMonsterLegendaryActions}
\end{multicols}
\end{DndMonster}



\begin{DndMonster}[float*=b,width=\textwidth + 8pt]{Zenith Aastrika}
\begin{multicols}{2}
\DndMonsterType{Huge giant (Leader), lawful neutral}
\DndMonsterBasics[
armor-class = {18 (natural armor)},
hit-points = {\DndDice{18d12 + 108}},
speed = {50 ft.}
]
\DndMonsterAbilityScores[
 str = 25,
 dex = 16,
 con = 23,
 int = 10,
 wis = 23,
 cha = 14
]
\DndMonsterDetails[
saving-throws = {Str +11, Dex +7, Con +10, Wis +10},
skills = {Athletics +11, Perception +10, Persuasion +6},
damage-immunities = {fire},
condition-immunities = {charmed, frightened},
languages = {Giant, Primordial},
challenge = 12,
]
\DndMonsterAction{Battle Tranquility (3/Day)}
When Aastrika fails a saving throw, she can choose to succeed instead. When she does so, she can't use Sweltering Heat until the end of her next turn.

\DndMonsterAction{Molten Flesh}
The first time a creature other than a fire giant touches Aastrika or hits her with a melee attack on a turn, that creature takes 7 (2d6) fire damage.

\DndMonsterSection{Actions}
\DndMonsterAction{Multiattack}
Aastrika takes the Disengage action and makes up to three attacks using Unarmed Strike, Hurl Flame, or both.

\DndMonsterAction{Unarmed Strike}
\textsl{Melee Weapon Attack:} 11 to hit, reach 10 ft., one target. \textsl{Hit:} 21 (4d6 + 7) bludgeoning damage plus 7 (2d6) fire damage. If the target is a creature, they are outlined in red light until the start of Aastrika's next turn, and for the duration, attack rolls against them have advantage.

\DndMonsterAction{Hurl Flame}
\textsl{Ranged Spell Attack:} 10 to hit, range 180 ft., one target. \textsl{Hit:} 20 (4d6 + 6) fire damage.

\DndMonsterAction{Lava Pillar (Recharge 5\textendash 6)}
Aastrika causes a 20-foot-radius, 60-foot-high cylinder of lava to erupt from the ground at a point she can see within 120 feet of her. Each creature in that area must make a DC 18 Dexterity saving throw, taking 38 (11d6) fire damage on a failed save, or half as much damage on a successful one.

\DndMonsterSection{Bonus Actions}
\DndMonsterAction{Sweltering Heat}
Aastrika's inner flame pulses harshly. Each creature within 20 feet of her who doesn't have resistance or immunity to fire damage must succeed on a DC 18 Constitution saving throw or gain a level of exhaustion.

\DndMonsterSection{Reactions}
\DndMonsterAction{Power of Flame}
When another fire giant Aastrika can see within 60 feet of her misses with a weapon attack, Aastrika causes flame to trail from the attacker. The attacker rolls damage for the attack as if they had hit, and the target takes half that amount as fire damage instead of taking the full amount of normal damage.

\DndMonsterSection{Legendary Actions}
Aastrika has three villain actions. She can take each action once during an encounter after an enemy's turn. She can take these actions in any order but can use only one per round.

\begin{DndMonsterLegendaryActions}
\DndMonsterLegendaryAction{Action 1: Forward!}{Aastrika and up to three creatures she can see can move up to their speed without provoking opportunity attacks and make one melee attack each.
}
\DndMonsterLegendaryAction{Action 2: The Manifold Self}{Aastrika disappears in a flash of light and teleports, reappearing in an unoccupied space she can see within 30 feet of her. At the same time, five \textbf{fire giant trooper}s appear in unoccupied spaces within 30 feet of the space she teleported from. The troopers act immediately after this villain action is taken and on the same initiative count in subsequent rounds.
}
\DndMonsterLegendaryAction{Action 3: This. Ends. Now}{Aastrika and each willing fire giant within 60 feet of her unleashes a burst of heat. Each creature within 10 feet of Aastrika or one of these fire giants must make a DC 18 Dexterity saving throw, taking 55 (10d10) fire damage on a failed save, or half as much damage on a successful one.
}
\end{DndMonsterLegendaryActions}
\end{multicols}
\end{DndMonster}


\end{document}
